\documentstyle{article}

\setlength{\topmargin}{0pt}
\setlength{\oddsidemargin}{0pt}
\setlength{\evensidemargin}{0pt}
\setlength{\textwidth}{6.5in}
\setlength{\textheight}{8.5in}
\setlength{\headheight}{0pt}

%here are some mathematical symbols
\def\emptyset{{\hbox{\rm \O}}}
\def\set#1{\{#1\}}
\newcommand{\type}[1]{\mbox{\tt #1}}
\newcommand{\str}{?$*$}
\def\R{\rlap{\rm I}\,{\rm R}}
\def\powset{{\cal P}}
\def\glb{{\mbox{\rm glb}}}\def\lub{{\mbox{\rm lub}}}

\newcommand{\nm}{non\-mon\-o\-ton\-ic}
\newcommand{\Nm}{Non\-mon\-o\-ton\-ic}

\newcommand{\prolog}{{\sc prolog}}\newcommand{\Prolog}{{\sc Prolog}}
\newcommand{\atms}{{\sc atms}}\newcommand{\Atms}{{\sc Atms}}

%this is where the title page begins

\def\draftline{}

\def\authors{Matthew L. Ginsberg}

\newcommand{\nodblspace}{\baselineskip=\normalbaselineskip}
\def\group{\nodblspace Department of Computer Science\\
Stanford University\\
Stanford, California 94305\\
\bigskip
Updated \today
}

\def\titlebottom{\vspace{1.55in}
{\large \authors}\\
\vspace{1.55in}
\group
\end{center}
\clearpage}

\def\xtitletop{\vspace*{1.35in}\begin{center}}
\def\xtitl#1#2{\xtitletop
{\Large\bf #2}\\
\titlebottom}

\newcommand{\xtitle}[2]
{\def\tphrase{#2}\begin{titlepage}\draftline\xtitl{#1}{#2}\end{titlepage}}

\newcommand{\mvl}{{\sc mvl}}
\newcommand{\Mvl}{{\sc Mvl}}
\newcommand{\lisp}{{\sc lisp}}
\newcommand{\Lisp}{{\sc Lisp}}
\newcommand{\dg}{{\sc dag}}
\newcommand{\Dg}{{\sc Dag}}
\newcommand{\T}{\type{T}}
\newcommand{\NIL}{\type{NIL}}
\newcommand{\apa}{{\tt allow-\-procedural-\-attachment}}
\newcommand{\at}{{\tt *anytime-\-trace*}}
\newcommand{\mi}{{\tt mvl-\-invocation}}
\newcommand{\wbli}{{\tt *warn-\-on-\-bad-\-lisp-\-invocations*}}
\newcommand{\somo}{{\tt succeed-\-on-\-modal-\-ops}}
\newcommand{\et}{{\tt *equivalence-\-translation*}}
\newcommand{\itype}[1]{\type{#1}\index{#1 **}}

\makeindex
%\renewcommand{\index}[1]{}

\begin{document}
\xtitle{}{User's Guide to the MVL System}

\tableofcontents
\clearpage

\section{Introduction}
\subsection{What you got}

The \mvl\ system is currently distributed via anonymous ftp from
t.stanford.edu.  A .DVI version of this file, the \mvl\ manual,
can be found in the \type{papers} directory on this machine.

The \type{mvl} directory contains the source code for the \mvl\ system
itself, which is written in standard Common Lisp.  It should be
possible to port the system to any Common Lisp without substantial
difficulty.  In addition, the \type{papers} directory contains .DVI
versions of the following papers that relate to the functioning of the
system:
 \begin{enumerate}
 \item ``Multivalued logics: A uniform approach to inference in
artificial intelligence.''  This paper gives a formal description of
multivalued logics and the \mvl\ system.
 \item ``Bilattices and modal operators.''  This paper describes the
formal connection between the \mvl\ system and modal operators.
 \end{enumerate}

The \mvl\ sources are public domain.  You are free to use, modify, and
distribute them as you see fit, provided that such use or distribution
is not for commercial purposes.  No warranties are made with regard to
the performance of the system, however, and no commitment of support
has been made.

\subsection{Overview}

This user's guide has been written to help you install and use the \mvl\
system.  It describes the basic structure of the system, and shows you
how to take advantage of the flexibility of multivalued inference.

Essentially, \mvl\ is just like any other logic programming system.
There is a database consisting of those sentences about which the
system knows something, and there are facilities for determining
whether or not some query follows from the information the database
contains.  Viewed in this fashion, the \mvl\ system is not much
different from a conventional logic programming language such as
\prolog.

What makes \mvl\ unique is that each sentence in the database is
labelled with a {\em truth value\/}\index{truth value}.  Thus,
instead of simply assuming a statement to be true and dumping it into
the database, we will {\em label\/} the statement as true when we
insert it.  Of course, we can use other labels as well: true by
default, true by virtue of some assumption, or even false or false by
default.  \Mvl\ considers these truth values when determining how to
respond to a \prolog-like query.

This manual is divided into two parts.  In the first part (Sections
\ref{s.fol1} -- \ref{s.support}), we describe the use of the system
where the truth values assigned to sentences are simply ``true'' or
``false.''  Used in this fashion, \mvl\ is not much different from
\prolog.

The second part of the manual shows you how to use logics involving
more complex truth values.  We examine the logics that have been
predefined as part of the \mvl\ system, and describe the various ways
in which you can define your own logics.  Several tools for this
purpose have been included as part of the \mvl\ system.  In addition,
we describe \mvl's facilities for handling modal operators, which
are operators that accept logical sentences as arguments.

\subsection{Why MVL is unique}

As mentioned in the previous section, what distinguishes \mvl\ from
other declarative programming languages is its support for multiple
truth values.  There are other differences, however, and these will
be discussed in detail in subsequent sections.  Here's a quick summary:
 \begin{enumerate}
 \item As mentioned, the principal feature offered by \mvl\ is its
support of multiple truth values.
 \item The basic theorem prover used by the system includes facilities
to detect loops and control recursion encountered in proof search.
The first-order theorem prover is discussed in Section \ref{s.fo},
although we will have nothing further to say about recursion control.
 \item The system can handle negative subgoals containing free variables
as discussed in Section \ref{s.effort}.
 \item The \mvl\ unifier includes an occurcheck and full support for
sequence variables.  This is discussed in Section \ref{s.rep}.
 \end{enumerate}

\section{Installation and testing}

The \mvl\ system is currently developed and maintained on a
Sparcstation running Allegro Common Lisp and is tested only
infrequently on other platforms.  Although the system has been written
in standard Common Lisp to maximize ease of portability and no
problems have been encountered when moving the system to a variety of
platforms, performance on other systems is not guaranteed.

\Mvl\ currently supports both ANSI Common Lisp and Common Lisp as
described in Steele's 1984 book.  In order to determine whether or not
the ANSI standard is conformed to, the system checks to see if
\type{in-package} is a macro or a function.  If the former, it is
assumed that the entire ANSI standard is supported; if the latter, it
is assumed that the language is as in Steele and various extensions
are defined.  It is expected that some future release of \mvl\ will
support the ANSI standard only.

\subsection{Sun-4}

In order to load the system on a Sparcstation running Allegro Common
Lisp, do the following:
 \begin{enumerate}

 \item (ANSI) Set up a logical pathname translation so that logical
pathnames beginning with \type{MVL} point to the direction where \mvl\
is kept.  (non-ANSI) Edit the file \type{mvldcl.lisp} so that the
\type{:default-pathname} field in the \mvl\ system definition points
to the directory where \mvl\ will be kept.  This directory should also
contain the following files:
 \begin{enumerate}\setlength{\itemsep}{0pt}\raggedright
 \item {\tt f1.mvl, recursion.mvl, def-test.mvl, circ-test-1.mvl,
circ-test-2.mvl, circ-test-3.mvl, circ-test-4.mvl, completed-db-1.mvl,
completed-db-2.mvl, modal-test.mvl, actions.mvl, shooting.mvl}
 \item {\tt mvl.test, bc.test, recursion.test, atms.test,
defaults.test, circumscription.test, completed-db.test, modal.test,
fc.test, temporal.test}
 \end{enumerate}
 \item Compile and load the file \type{mvldcl.lisp}.
 \item Execute the function
\type{(mvl:compile-mvl)}\index{compile-mvl} to compile and load the
\mvl\ system.  Using this function in the future will only compile 
those files that have been changed since last compiled.
 \item Invoke the function \type{(use-package 'mvl)} to access the
functions and symbols of the \mvl\ system (which is kept in its own
package; see Section \ref{s.package}).
 \end{enumerate}

Now invoke the function {\tt (test-mvl)}, which should return 0
(i.e., 0 errors encountered).  If it doesn't, the problem is probably
that \mvl\ is unable to locate the files mentioned in (1) above.

\subsection{Symbolics 3600 series}

The \mvl\ system will run on Symbolics Genera 7.2 or later.  It is
possible to run the system using Genera 7.0 or 7.1, although a bug in
the implementation of the Common Lisp \type{\&allow-other-keys}
keyword may cause the system to dump to the cold load stream if
certain predicates are procedurally attached.  (See Section
\ref{s.attach}.)

You should begin by editing the file \type{mvldcl.lisp} so that the
default pathname in the \mvl\ system declaration points to the
directory in which the files are actually kept.  You can then load
the \mvl\ system using either the Load System command or the Lisp
function \type{(mvl:load-mvl)}\index{load-mvl} that is defined in the
file \type{mvldcl.lisp}.  Now access the \mvl\ package using
\type{(use-package 'mvl)} and execute the function \itype{test-mvl}.
This function should return 0 (i.e., 0 errors encountered).  If it
doesn't, the problem is probably that \mvl\ is unable to find the
directory its test files are on (they're included in the distribution
tape).

\subsection{Other machines}

The \mvl\ system has also been run and tested on a wide variety of
other machines.  In order to load the system, follow the instructions
provided for Sun-4's.

\section{The MVL database}
\label{s.fol1}
\subsection{Representation}
\subsubsection{Structure of the database}
\label{s.rep}

As already remarked, the \mvl\ database consists of sentences labelled
with truth values.  In this section, we will describe the form of the
sentences themselves; in the next section, we comment briefly on the
truth values used to label them.

Sentences in the database are represented as \lisp\ s-expressions
using the prenex normal form \index{prenex normal form} of the
standard logical connectives \type{NOT}, \type{OR} and \type{AND}.
Thus the statement that Fred is tall and Harry is not would be
written as:
 \[\type{(AND (TALL FRED) (NOT (TALL HARRY)))}\]

Any \lisp\ atom can be used as a constant term, with the exception of
variables (see below) and the following reserved connectives:
 \[\type{NOT AND OR => <= <=> IF IFF}\]
The connective \type{<=} is used in backward-chaining rules;
 \[\type{(<= CONCLUSION PREMISE-1 ... PREMISE-N)}\]
is logically equivalent to
 \[\type{(IF (AND PREMISE-1 ... PREMISE-N) CONCLUSION)}\]
as is the forward-chaining version
 \[\type{(=> PREMISE-1 ... PREMISE-N CONCLUSION)}\]
 Although logically equivalent, these expressions behave differently
in practice.  The version
 \[\type{(IF (AND PREMISE-1 ... PREMISE-N) CONCLUSION)}\]
 in which the connective \type{IF} is used represents a pair of
rules, depending on the value of the parameter \itype{*if-translation*}.
If the value is \type{bc} (the default), then two backward-chaining
rules are produced, so that
 \[\type{(IF (A) (B))}\]
is equivalent to
 \[\type{(AND (<= (B) (A)) (<= (NOT (A)) (NOT (B))))}\]
 the conjunction of the original rule and its contraposition.
If the value of \type{*if-translation*} is \type{fc}, then
forward-chaining rules are produced and the translation is to
 \[\type{(AND (=> (A) (B)) (=> (NOT (B)) (NOT (A))))}\]
 Finally, if \type{*if-translation*} has the value \type{mix}, the
pair of noncontraposed rules
 \[\type{(AND (<= (B) (A)) (=> (A) (B)))}\]
 is produced.

Equivalences are represented using either \type{IFF} or \type{<=>}.
Thus
 \[\type{(IFF (FOO FRED) (GOO FRED))}\]
is equivalent to
 \[\type{(AND (IF (FOO FRED) (GOO FRED)) (IF (GOO FRED) (FOO FRED)))}\]
 Translations involving \type{<=>} depend on the value of the
parameter \et\index{*equivalence-translation*}.  As above, if the
value is \type{bc} (the default), then two backward-chaining rules
are produced, so that
 \[\type{(<=> (FOO FRED) (GOO FRED))}\]
translates to
 \[\type{(AND (<= (FOO FRED) (GOO FRED)) (<= (GOO FRED) (FOO FRED)))}\]
 If \et\ has the value \type{fc}, two forward-chaining rules are produced
so that the translation is to
 \[\type{(AND (=> (FOO FRED) (GOO FRED)) (=> (GOO FRED) (FOO FRED)))}\]
 Finally, if the value is \type{mix}, one forward-chaining rule and
one backward-chaining rule are produced:
 \[\type{(AND (=> (FOO FRED) (GOO FRED)) (<= (FOO FRED) (GOO FRED)))}\]

Variables \index{variables} are represented using \lisp\ atoms
beginning with the character \type{?} and are handled as in \prolog,
so that free variables in database statements are implicitly
universally quantified.  Again as in \prolog, existential
quantification is handled via Skolemization; thus the
backward-chaining sentence ``Every rich person is happy'' and the
sentence ``There is a rich person'' would be written as
 \begin{equation}
\type{(<= (HAPPY ?P) (PERSON ?P) (RICH ?P))}
\label{e.if}
\end{equation}
and
 \[\type{(AND (PERSON P0017) (RICH P0017))}\]
respectively.

The variables in database sentences are standardized apart from the
other variables used by the system.  This is done by using uninterned
variables in the actual representation of a database fact, so that
the internal representation of (\ref{e.if}) would be
 \[\type{(<= (HAPPY \#:?P) (PERSON \#:?P) (RICH \#:?P))}\]
 The original variables are generally reinstated when \mvl\ is asked
to display a database proposition.

\subsubsection{Sequence variables}

In addition to the simple variables discussed in the previous
section, \mvl\ supports the use of ``sequence
variables.''\index{sequence variables} These variables begin with the
characters \type{\str}, and will match not a single atom, but an
arbitrary (possibly empty) sequence of atoms.

Thus, for example,
 \[\type{(FOO A B)}\]
 will unify with
 \[\type{(FOO \str X)},\]
 returning the binding list
 \[\type{((\str X . (A B)))}\]
 in which \type{?*X} has been bound to the list \type{(A B)}.  This
binding list will presumably be displayed as
 \[\type{((\str X A B))}.\]
 Similarly, the most general unifier of $\type{(FOO)}$ and
$\type{(FOO \str X)}$ is given by $\type{((\str X))}$.

The presence of sequence variables may cause the most general unifier
of two patterns to be nonunique.  For example, the result of unifying
 \[\type{(FOO A B)}\]
 with
 \[\type{(FOO \str A \str B)}\]
 can be any one of the three answers
 \begin{eqnarray*}
& \type{((\str A A B) (\str B))} & \\
& \type{((\str A A) (\str B B))} & \\
& \type{((\str A) (\str B A B))} &
 \end{eqnarray*}

For this reason, the unification functions return not a single answer,
but a {\em list\/} of possible unifiers.

\subsubsection{Functions to manipulate database sentences}
\label{s.manipulate}

The following functions and global parameters are used to manipulate
sentences in the database:
 \begin{description}
 \item [\itype{varp(x)}] Returns \T\ if $x$ is a variable, and
\NIL\ otherwise.
 \item [\type{varp?(x)}]\index{varp-query} Returns \T\ if $x$ is a
standard (i.e., non-sequence) variable, and \NIL\ otherwise.
 \item [\type{varp$*$(x)}]\index{varp-star} Returns \T\ if $x$ is a
sequence variable, and \NIL\ otherwise.
 \item [\itype{vartype(x)}] Returns ? if $x$ is a standard variable,
$*$ if $x$ is a sequence variable, and \NIL\ is $x$ is not a
variable.
 \item [\type{new-?var()}]\index{new-query-var} Returns a new
temporary variable.
 \item [\type{new-$*$var()}]\index{new-star-var} Returns a new
temporary sequence variable.
 \item [\itype{normal-form(p)}] Returns an equivalent form of $p$
suitable for insertion into the database (i.e., with the connectives
\type{IFF}, \type{IF} and \type{<=>} reduced to \type{<=} and \type{=>}).
 \item [\itype{*if-translation*}] Controls the fashion in which
\type{normal-form} deals with the connective \type{if}.
 \item [\itype{*equivalence-translation*}] Controls the fashion in which
\type{normal-form} deals with the connective \type{<=>}.
 \item [\itype{negate(p)}] Returns the negation of $p$.
 \item [\itype{standardize-operators(p)}] Returns an equivalent form
of $p$ with the connectives \type{<=>}, \type{IF} and \type{IFF}
removed.  (Note: This function does not use the control variables
\type{*if-translation*} and \et, since it is used when proving
sentences, not when inserting them into the database.)
 \item [\itype{cnf(p)}] Returns $p$ in conjunctive normal form, as a
list of conjuncts, where each conjunct is a list of disjuncts.  Thus
 \[\type{(CNF '(AND A B))}\]
returns
 \[\type{((A) (B))}\]
\item [\itype{dnf(p)}] Returns $p$ in disjunctive normal form, as a list
of disjuncts, where each disjunct is a list of conjuncts.  Thus
 \[\type{(DNF '(AND A B))}\]
returns
 \[\type{((A B))}\]
\end{description}

\subsection{Truth values}

The truth values used to label the sentences in the \mvl\ database
are elements of a bilattice that can be selected by the user with
the function \type{load-logic}.  We will defer a discussion of these
truth values to the second portion of this manual, which discusses
the use of bilattices that label sentences as other than simply true
or false.  At this point, we merely remark that if it is our intention
to assert the sentence \type{(not (foo))}, this is actually accomplished
by asserting that the negated sentence \type{(foo)} has truth value
\type{false}.

\subsection{Database form and internals}
\label{s.database}

The database is maintained using a discrimination net, as suggested
by Charniak et.~al in their book on {\sc lisp} programming.  Any
particular sentence in the database, say \type{(tall tom)}, is
associated with a \lisp\ atom such as \type{P27} that is its
``proposition number.''  The propositions are numbered consecutively,
beginning with \type{P1}.  At any given point in time, the function
\type{(props)} can be invoked to return a list of all of the
propositions in the database.

Very few of the functions that are used to manipulate the internal form
of the database are likely to be needed by a user of the \mvl\ system.  The
following are the most useful:
 \begin{description}
 \item [\itype{props()}] This function returns a list of all
propositions about which any information has been recorded.
 \item [\itype{denotes(a)}] Returns the proposition associated with a
particular \lisp\ atom.
 \item [\itype{truth-value(a)}] Returns the truth value associated
with a particular \lisp\ atom.  Thus, if the result of asserting
\type{(not (short tom))} into the database was the atom \type{P38},
\type{(denotes 'p38)} would evaluate to \type{(short tom)}, and
\type{(truth-value 'p38)} would evaluate to \type{false}. 
 \item [\itype{index(p)}] Updates the discrimination net as
needed to insert the proposition $p$ into the database.
 \item [\itype{unindex(p)}] Removes from the discrimination net
information relating to the proposition $p$.
 \end{description}

\section{Quantification}

As we have remarked, quantification is treated in \mvl\ in a fashion
similar to that used by \prolog\ systems:  Free variables appearing
in the database are assumed to be universally quantified; free variables
appearing in queries are existentially quantified.  Skolemization is
used to handle existential quantification in database statements or
universal quantification in queries.

The use of this approach means that the answer to an existentially
quantified query will often be of a form that indicates how the
variables should be bound in the result.  In \mvl\, these binding
lists \index{binding lists} are represented as expressions of the
form
 \[\type{((var-1 . bdg-1) ... (var-n . bdg-n))}\]
 where \type{var-1} through \type{var-n} are variables and
\type{bdg-1} through \type{bdg-n} are their respective bindings.
If no variable is bound to a value, the empty binding list \NIL\
is used.

As an example, suppose that the database contains the single sentence
 \[\type{(LIKES COCO ?X)}\]
 indicating that Coco likes everything.  Now the query
 \[\type{(LIKES ?Y MATT)}\]
 will be interpreted as, ``Find something that likes Matt.''  The answer
returned will be
 \[\type{((?Y . COCO))}\]
since Coco likes Matt.  Since the queries
 \[\type{(LIKES COCO MATT)}\]
and
 \[\type{(LIKES COCO ?Z)}\]
 both succeed without binding a variable, the answer \NIL\ will be
returned for each.\footnote{In order to distinguish the empty binding
list \NIL\ from a failed database query, all queries return a binding
list as well as a truth value.  This is discussed in Section
\ref{s.answer}.}

\subsection{Functions used to manipulate binding lists}

A variety of functions within \mvl\ are used to construct and manipulate
binding lists.  Here are some of the more useful ones:
 \begin{description}
 \item [\itype{groundp(p)}] Returns \T\ if the proposition $p$ is
ground (i.e., contains no variables) and \NIL\ otherwise.
 \item [\itype{vars-in(p)}] Returns a list of the variables
appearing in $p$.
 \item [\itype{get-bdg(var bdg-list)}] Returns the value assigned to
the variable \type{var} by the binding list \type{bdg-list}.  If the
variable does not appear in the binding list, \NIL\ is returned.
 \item [\itype{delete-bdg(var bdg-list)}] Destructively modifies
\type{bdg-list}, removing any binding for the variable \type{var}.
 \item [\itype{plug(p bdg-list)}] Returns the result of instantiating
the proposition $p$ using the given binding list.  This function is
nondestructive.
 \item [\itype{samep(p1 p2)}] Returns \T\ if the propositions $p_1$
and $p_2$ are identical (up to variable renaming), and \NIL\
otherwise.  If the optional \type{return-answer?} argument is
supplied and the propositions are the same, a binding list is
returned instead of \T.
 \item [\itype{unifyp(p1 p2)}] Returns a list of the most general
unifiers of $p_1$ and $p_2$, or \NIL\ if the propositions fail to
unify.
 \item [\itype{instp(p1 p2)}] If $p_2$ is an instance of $p_1$,
return a list of their most general unifiers; only variables in
$p_1$ are bound.  Returns \NIL\ if $p_2$ is not an instance of $p_1$.
 \item [\itype{matchp(pattern atom)}] Returns a list of elements of
the form \type{(b . d)} where $b$ is a most general binding list such
that \type{(plug pattern b)} is an instance of the proposition
labelled with \type{atom} and $d$ is a binding list for the variables
in the database proposition.  \NIL\ is returned if no such binding
list exists.  Thus if \type{P1} is \type{(TALL \#:?2)}, and \type{P2}
is \type{(TALL TOM)}, then
 \[\type{(matchp (TALL TOM) P1)}\]
returns a list consisting of the single entry 
 \[\type{(NIL . ((\#:?2 . TOM)))}\]
 The \type{car} of this entry is \NIL, indicating that no variables
in the query need to be bound to match \type{P1}; the \type{cdr} is
\type{((\#:?2 . TOM))} since the variable \type{\#:?2} in \type{P1}
needs to be bound to \type{TOM} to match the query.
 \[\type{(matchp (TALL ?X) P2)}\]
returns \type{((((?X . TOM)) . NIL))}.
 \end{description}

Additional functions are described in Section \ref{s.binding-dag}.

\subsection{Negative subgoals containing free variables}
\label{s.effort}

Consider the following \prolog\ fragment:
\begin{verbatim}
     f(x) :- not(g(x)).
     g(A).
\end{verbatim}
 The interpretation we will assign to the implication
\type{f(x) :- not(g(x))} will be that it is universally quantified
over $x$, and therefore corresponds to the implication
 \[\forall x.[\neg g(x) \supset f(x)]\]
 Given this interpretation, what should the response be if we provide
the system with the query \type{f(x)}?

In light of the negation as failure rule of \prolog\ (which has analogs
in the \mvl\ system), we would expect the query to succeed for all values
of $x$ {\em except\/} $x=A$.

Most \prolog\ systems cannot respond to queries such as this one, since
a nonground negative subgoal is encountered during the proof process.
One approach to these queries, suggested by David Chan, is to do just
what we have suggested in the previous paragraph, computing those binding
lists for which the query fails.

Unfortunately, there are many queries that succeed in general, but
for which it is impractical to compute all of the bindings for which
the query fails.  This may be because there are infinitely many such
failures, or because the time needed to compute them explicitly is
prohibitive.  Situations such as this arise frequently in
applications of declarative programming techniques to planning
problems; here, we would like simply to return an indication that the
query ``succeeds in general.''  In the above example, we might
report that the query had succeeded without binding the variable $x$;
even though $x=A$ is an exception, the report of success is true in
general.

Both approaches are supported by \mvl, and the global variable
\type{*completed-database*} is used to distinguish between them.  If
this variable has the value \T, all exceptions are computed; if it
has the value \NIL, the other approach is taken.  The default value
is \NIL:
 \begin{description}
 \item [\itype{*completed-database*}] If \T, find all exceptions
to default proofs.  If \NIL\ (the default), accept defaults proofs
with possible exceptions as complete.
 \end{description}

\section{Modifying the MVL database}

\subsection{Keyboard input}

The simplest functions to modify the \mvl\ database simply accept
propositions from the user and either add or remove them from the
knowledge base.  The primitive function for adding to the database is
\itype{stash}; the function \itype{unstash} is used for removal.

The function \type{stash} accepts a proposition and up to three
keywords:
 \begin{description}
 \item [\itype{value}]  This keyword is used to supply a specific truth
value to be assigned to the proposition.  If no value is supplied, the
default value is used (this value is supplied when the logic is chosen;
see Section \ref{s.bilattice}).
 \item [\itype{increment-p}] If the proposition in question already
appears in the database, the user has a choice: The given truth value
can be used to replace the one currently appearing in the database,
or it can be combined with the existing value using the combination
function supplied with the logic being used.  The truth value is
replaced if \type{increment-p} is \NIL, and is updated if it is \T\
(the default).
 \item [\apa\index{allow-procedural-attachment}] This keyword,
which defaults to \T, controls whether or not it is possible to
procedurally attach the action of stashing the supplied proposition.
See Section \ref{s.attach}.
 \end{description}

Somewhat more useful than \type{stash} is the function \itype{state}.
This function is identical to \type{stash}, except it translates the
supplied sentence into a form suitable for inference by using the
function \itype{normal-form} described in Section \ref{s.manipulate}.
In addition, \type{state} cannot be procedurally attached, so the
keyword \apa\ is not allowed.

For removing sentences from the database, the functions \type{unstash}
and \itype{unstate} are used.

The function \type{unstash} accepts a proposition to be removed from
the database, and up to three keywords: \type{value},
\itype{decrement-p} and \apa.  The last of these is treated as for
\type{stash}.  In general, if no keywords are used, the effect is to
remove the proposition from the database.

If \type{decrement-p} is \NIL, then \type{unstash} behaves like
\type{stash}, except the default value used for unstashing is not
necessarily the same as the default value used for stashing.

If \type{decrement-p} is \T, then the new truth value assigned to
the proposition is given by
 \begin{equation}
o \cdot \neg v
\label{e.dot-star}
\end{equation}
 where $o$ is the old value and $v$ is the value supplied via the
\type{value} keyword.  The $\cdot$ and $\neg$ operations are as
described in Section \ref{s.bilattice}.

There is one additional function for removing sentences from the
database, \itype{empty}.  This function accepts no arguments, but
does take a single keyword, \type{value}.  The function calls
\type{unstash} for every sentence in the database.

All of the functions described in this section work by calling the
primitive function \itype{new-val}.  This function accepts as
arguments a proposition name and a truth value, and sets the truth
value of the proposition to the given value.  If the truth value is
\type{unknown}, indicating that nothing is known about the
proposition, it is removed from the database and the discrimination
net.

\subsection{Input from files}

It is also possible to enter a database from a disk file.  These files
should in general have the extension \type{.MVL}, and should be
located either on the currently active disk area or on the \mvl\
directory.  The files should consist simply of the sentences to be
added to the database.

In addition to consisting of such sentences, the files may contain the
keyword \type{:value}, followed by the value to be assigned to each
sentence, or the keyword \type{:function}, followed by a function to
be applied to each sentence in order to {\em compute\/} the value to
be assigned to it.  The keyword \type{:empty} empties the database
before loading the file.  Examples of these constructs can be found in
the various \type{.MVL} files included on the distribution tape.

The sentences are loaded with \itype{mvl-load}, which accepts a file
name and an optional keyword \type{empty}.  If the keyword is
supplied, the database is emptied before the file is read; in any
event, \type{state} is then applied to each sentence in the given
file.  The sentences in the file can be removed from the database
with the function \itype{mvl-unload}, which accepts a file name as its
argument.  These functions each call \type{mvl-file}, which accepts
a string and returns a path to the \type{.MVL} file associated to it.

Related to \type{mvl-load} is \itype{mvl-save}, which also accepts a
file name as an argument, but instead of loading a declarative
database from the file, writes the current database {\em to\/} the
file.

\subsection{The editor interface} \index{editor interface}
\subsubsection{Sun EMACS editor}

In addition to the file commands \type{mvl-load} and
\type{mvl-unload}, there is an interface to the {\sc emacs} editor on
Sparcstations.  For a file of mode \mvl, certain keys are bound as
follows (each of these commands also recognizes the
\type{:value} and \type{:function} keywords):

 \begin{description}
 \item [\type{meta-S}] Applies \type{state} to each s-expression in the
current region.
 \item [\type{meta-U}] Applies \type{unstate} to each s-expression in the
current region.
 \item [\itype{meta-s}] Applies \type{state} to each s-expression in the
current defun (i.e., s-expression beginning in column 1).
 \item [\itype{meta-u}] Applies \type{unstate} to each s-expression in the
current defun.
 \end{description}

In order to use these commands, certain modifications will need to be made
to your \type{.emacs} file.  The file \type{dot.emacs} is an example of
what to do.

\subsubsection{Symbolics ZMACS editor}

In addition to the file commands \type{mvl-load} and
\type{mvl-unload}, there is an interface to the {\sc zmacs} editor on
Symbolics 3600 series machines.  For a file of mode \mvl, certain
keys are bound as follows (each of these commands also recognizes the
\type{:value} and \type{:function} keywords):

 \begin{description}
 \item [\itype{super-s}] Applies \type{state} to each s-expression in the
current region.
 \item [\itype{super-u}] Applies \type{unstate} to each s-expression in the
current region.
 \item [\itype{super-f}] Applies \type{fc} to each s-expression in the
current region.  See Section \ref{s.fc}.
 \item [\itype{super-e}] Applies \type{erase} to each s-expression in the
current region.  See Section \ref{s.fc}.
 \end{description}

\subsection{Summary of database modification functions}

\begin{description}
 \item [\itype{stash(p)}] Inserts $p$ into the database.  Allowed
keywords: \type{value}, \type{increment-p} and \apa.
 \item [\itype{state(p)}] Inserts $p$ into the database after
converting it to normal form.  Keywords: \type{value},
\type{increment-p}.
 \item [\itype{unstash(p)}] Removes $p$ from the database.  Keywords:
\type{value}, \type{decrement-p}, \apa.
 \item [\itype{unstate(p)}] Removes $p$ from the database after
converting it to normal form.  Keywords: \type{value},
\type{decrement-p}.
 \item [\itype{empty}] Unstashes every sentence in the database.  Keyword:
\type{value}.
 \item [\itype{new-val(a v)}] Sets the truth value of the sentence
associated with the atom $a$ to $v$.  If $v$ is the truth value
\type{unknown}, the sentence is removed from the database.
 \item [\itype{mvl-file(string)}] Returns the path for the \mvl\ file
associated with the given name.  The file name may include an
extension, one may be supplied as an optional argument, or
\type{.MVL} will be used by default.  There is an optional keyword
\type{no-error} indicating that even if the file cannot be found, the
path should be returned.
 \item [\itype{mvl-load(file)}] Loads the sentences in the given file into
the \mvl\ database.
 \item [\itype{mvl-unload(file)}] Removes the sentences in the
given file from the \mvl\ database.
 \item [\itype{mvl-save(file)}] Writes the \mvl\ database out to the
given file.  This function also accepts the keywords \type{cutoffs}
and \type{props}, which are treated as for \type{contents} (see
Section \ref{s.contents}).  Finally, the keyword \type{if-exists} is
passed to the \type{with-open-file} macro.
 \end{description}

\section{Querying the MVL database}

\subsection{MVL responses to queries}
\label{s.answer}

Because of the fact that a particular fact may be stored in the
database with any of a variety of truth values, \mvl\ responds to a
query by returning not simply a binding list (as most \prolog\
interpreters do, for example), but by returning a binding list {\em
and a truth value}. 

In actuality, what is returned is an instance of the \type{answer}
structure\index{answer structure}.  This structure contains the
following fields:
 \begin{description}
 \item[\type{binding}] Bindings for some of the variables in the query.
 \item[\type{value}] A truth value.  The bound version of the query has
this truth value.
 \end{description}

The definition of this structure leads Common Lisp to define the
access functions \type{answer-binding} and \type{answer-value};
in addition, a function \type{equal-answer} is defined by \mvl\ to
test two answers for equality:
 \begin{description}
 \item[\itype{answer-binding(a)}] Returns the binding field of the
answer \type a.
 \item[\itype{answer-value(a)}] Returns the value field of the answer
\type a.
 \item[\itype{equal-answer(a1 a2)}] Returns \T\ if the two answers
\type{a1} and \type{a2} are equal, and \NIL\ otherwise.  The
functions used are \type{binding-equal} and \type{mvl-eq}, the
functions normally used to test binding lists and truth values for
equality.
 \end{description}

Thus, in response to the query
 \[\type{(flies ?x)}\]
 the response of an answer with binding \type{((?x . Tweety))} and
truth value \type{dt} would indicate that
 \[\type{(flies Tweety)}\] 
 has truth value \type{dt}.  In other words, Tweety's flying is a
conclusion that is true by default.  (See Section \ref{s.j} for a
description of default truth values in the \mvl\ system.)

As a side effect, this allows the user to distinguish a response of
\NIL\ (indicating that the query has failed) from a response in which
the binding is \NIL\ (indicating that the query has succeeded without
binding any variables).

\subsection{Searching for specific sentences}

There are two functions that are most commonly used to retrieve facts
that have been stored in the database, \itype{lookup} and
\itype{lookups}.  Each of these functions accepts as arguments a
proposition to search for and up to four keywords, and returns two
values.

The keywords and their meanings are as follows:
 \begin{description}
 \item [\itype{cutoffs}] This keyword is used to give conditions on
the sentences being considered as possible matches to the query.  Any
sentence that is considered as a match must have an associated truth
value \type{v} such that \type{(test v cutoffs)} evaluates to \T\ (see
Section \ref{s.cutoffs}).  The default value of this keyword is defined
by the bilattice in use (see Section \ref{s.bilattice}).
 \item [\itype{inv}] This keyword defaults to \NIL.  If it is \T,
however, any truth values found should be negated before they are
returned.
 \item [\apa] This is similar to the keyword for \type{stash}.
 \item [\itype{succeed-on-modal-ops}] If \T, then any attempt to look
up a proposition involving a modal operator will immediately succeed;
if \NIL\ (the default), then the result returned will be the result
of applying the modal operator to the truth values of the
propositions that are its arguments.  See Section \ref{s.modal}.
 \end{description}

The two functions \type{lookup} and \type{lookups} are identical
except that \type{lookup} looks for only a single match for the
supplied proposition, while \type{lookups} searches for all such
matches.  Thus while \type{lookup} returns two values, \type{lookups}
returns two {\em lists\/} of values.  Each entry on each of these
lists is analogous to the corresponding value returned by
\type{lookup}.  The values returned by \type{(lookup q)} are as
follows:
 \begin{enumerate}
 \item First, an \type{answer} is returned.  If this answer contains
the binding list \type b and the truth value \type v, then applying
the binding list to the query produces a result with the given truth
value.  In other words, \type{(plug q b)} is an instantiation of a
sentence appearing explicitly in the database with truth value \type
v.
 \item Second, one of the following three values is returned:
 \begin{enumerate}
 \item If the sentence was found by procedural attachment, the atom
\type{attach} is returned.
 \item If the sentence was found by applying a modal operator, the
atom \type{modal} is returned.
 \item Under other conditions, the instantiated version of the
proposition that resulted in the match is returned.
 \end{enumerate}
 \end{enumerate}

As already mentioned, \type{lookups} returns two lists of values.

Here are summaries of \type{lookup} and \type{lookups}:

 \begin{description}
 \item [\itype{lookup(q)}] Searches the database for a sentence
matching the query $q$.  Allowed keywords: \type{cutoffs}, \type{inv},
\apa, \somo.  Returns two values: an answer and the source of the
match.  Returns \NIL\ if the query cannot be found in the database.
 \item [\itype{lookups(q)}] Similar to \type{lookup}, but looks
for {\em all\/} sentences matching the query.  Returns two lists of
values.
 \end{description}

\subsection{Searching the entire database}
\label{s.contents}

In addition to searching for a particular sentence, it is also
possible to display large portions of the entire \mvl\ database.  The
basic functions for doing this are \type{contents}, which displays
the entire database, and \type{prfacts}, which accepts a \lisp\ atom
as an argument and displays all of those facts concerning that atom.

Each of these two functions also accepts a variety of keywords:
 \begin{description}
 \item [\itype{cutoffs}] As for \type{lookup}, it is possible to
supply a test that must be passed by the truth value of any database
sentence before that sentence is displayed.
 \end{description}
 In addition, \type{contents} accepts a keyword \type{props}, which
defaults to \type{(props)} (see Section \ref{s.database}) and is a
list of the sentences to be displayed.

Related to \type{prfacts} is \itype{unstash-facts-with-atom}, which
accepts an atom as an argument and removes from the database any
fact concerning that atom.  \type{Unstash-facts-with-atom} accepts
the same keywords as \type{unstash} does.

Here, then, are the functions described in this section:

\begin{description}
 \item [\itype{contents}] Displays the contents of the database.  Keywords:
\type{cutoffs}, \type{props}.
 \item [\itype{prfacts(atom)}] Displays the database facts concerning the
given atom.  Keyword: \type{cutoffs}.
 \item [\itype{unstash-facts-with-atom(atom)}] Unstashes from the
database any facts concerning the given atom.  Keywords:
\type{value}, \type{decrement-p}, \apa.
 \end{description}

\subsubsection{Command line interface}

\paragraph{Sun}
In addition to the above functions, on Sparcstations, the key
$\langle$meta$\rangle$-* may be typed to the command line processor
running from the editor.  The result is that any information in the \mvl\
database concerning the currently open predicate will be displayed to the
user.

\paragraph{Symbolics}
On Symbolics systems, the key $\langle$super-shift$\rangle$-A may be typed
to the command line processor.  The result is that any information in the
\mvl\ database concerning the currently open predicate will be displayed
to the user.  This feature is the \mvl\ analog to the
$\langle$ctrl-shift$\rangle$-A key, which displays information about the
arguments to the currently open function.

\section{The first-order prover}
\label{s.fo}

As discussed in the technical papers distributed with \mvl, inference
in \mvl\ proceeds using a conventional first-order theorem prover; in
this section, we will describe this prover and the interface to it.
Be warned that the information in this section is unlikely to be used
by other than the most sophisticated users of \mvl; most readers will
want to proceed directly to Section \ref{s.inf}.

In general, there will be a variety of invocations to the first-order
prover active at any particular time.  Each such invocation is
associated with a task that the multivalued prover needs to complete,
and each such task is a \lisp\ object that is an instance of the
\type{consider} structure.\footnote{In fact, each task is an instance
of a structure that {\em includes\/} the \type{consider} structure.
See the file \type{bc.lisp} for further details.}

The multiple invocations of the first-order prover are run in
parallel by \mvl.  To each such invocation is associated an instance
of the structure \type{proof} that is used to record temporary
information about the progress of the proof effort.  We will discuss
first the first-order prover itself, then the construction and
termination of the invocations of it, and finally the structure of
the multivalued proving tasks.

\subsection{Nature}
\label{s.simple}

The prover itself is a simple natural-deduction style complete
theorem prover.  Only inferences against the initial clauses of
database statements are allowed; the result of this is that after
applying \type{state} to a sentence of the form
 \[\type{(<= CONCLUSION PREMISE-1 ... PREMISE-N)}\]
 \type{PREMISE-1} will always be tried first when the prover attempts
to prove \type{CONCLUSION}.  In a similar way, stating the
forward-chaining rule
 \[\type{(=> PREMISE-1 ... PREMISE-N CONCLUSION)}\]
 will result in a situation where only \type{PREMISE-1} will trigger
forward chaining.

Control of the proof search is handled by a function that evaluates
the ``merit'' of each of the ongoing proof attempts, and then
attempts to select the best one.  The evaluation function is stored
in the value of the global constant \itype{*node-evaluation-fn*},
which defaults to \itype{best-first-fn}, a function that attempts to
determine the value of a proof by looking at the truth value being
accumulated.  An alternative to \type{best-first-fn} is the function
\itype{depth-fn}, which evaluates a node simply by its depth in the
proof tree.

Three global parameters are used to control the proof search:
 \begin{description}
 \item [\itype{*final-limit*}] is the limit point beyond which the
search is abandoned, so that any node that evaluates to a value worse
than this will be pruned.  The default value of this parameter is
\NIL, indicating that no nodes are to be pruned.
 \item [\itype{*initial-limit*}] is used to indicate the type of search.
If \NIL\ (the default), then the values assigned to the nodes are not
used to order the search, and depth-first search is used.
 \item [\itype{*limit-delta*}] is used if \type{*initial-limit*} is
non-\NIL.  In this case, the nodes are ordered by increasing
 \[\frac{\type{value(node)}-\type{*initial-limit*}}{\type{*limit-delta*}}\]
 A typical value for *initial-limit* is 0, and for *limit-delta* is
1.  The value of this parameter is irrelevant if
\type{*initial-limit*} is \NIL.
 \end{description}

These three parameters may be reset by the user:
 \begin{description}
 \item [\itype{set-search-parameters}] This function is the primitive
used to modify the parameters used to control the first-order theorem
prover's search.  It accepts as keywords \type{initial-limit},
\type{limit-delta}, and \type{final-limit}.
 \item [\itype{activate-breadth-first}] Activate breadth-first search
by setting \type{*node-evaluation-fn*} to \type{depth-fn}.  Allowed
keywords: \type{initial-limit}, \type{limit-delta}, and
\type{final-limit}.
 \item [\itype{activate-depth-first}] Activate depth-first search by
setting \type{*initial-limit*} to \NIL.  Allowed keyword:
\type{final-limit}.
 \item [\itype{activate-best-first}] Activate best-first search by
setting \type{*node-evaluation-fn*} to \type{best-first-fn}.  Allowed
keywords: \type{initial-limit}, \type{limit-delta}, and
\type{final-limit}.
 \end{description}

\subsection{Invoking the prover}

The prover is invoked using the function \type{init-prover}.  This
function accepts as arguments a proposition and a pointer to the task
responsible for this proof attempt.

When \type{init-prover} is invoked, an instance of \itype{proof}
is created and returned that contains auxiliary information about the
proof attempt.  Here are the fields in the structure:
 \begin{description}
 \item [\itype{active-nodes}] A list of the nodes that are still
under active consideration.
 \item [\itype{answers}] A list of answers that have been found.
 \item [\itype{indexp} and \itype{indexn}] These entries are used to
index the various subtasks generated by the prover.
 \item [\itype{initial-limit}, \itype{final-limit}, \itype{limit-delta}
and \itype{evaluation-fn}] These entries are the values of the search
control parameters used in this proof attempt.
 \item [\itype{task}] A pointer back to the task responsible
for this invocation of the first-order prover.
 \item [\itype{vars}] The variables that appear in the proposition
being proven.
 \end{description}

Each node is an instance of the structure \itype{goal}, which is
described in the file \type{first-order.lisp}.

To continue a proof process, the function \type{cont-prover} is used.
This function accepts as an argument an instance of \type{proof} that
contains the auxiliary information for the proof process being
resumed.  The value returned by \type{cont-prover} is a structure
describing the results of the proof process itself;
\type{cont-prover} returns as soon as the prover finds a single
solution to its query.  The value returned is an instance of the
\type{bcanswer} structure\index{\type{bcanswer} structure}; the
fields in this structure are as follows:
 \begin{description}
 \item[\type{answer}] This field contains a binding list and truth
value that the system has determined can be assigned to the given
instantiation of the original query.
 \item[\type{just}] This field contains a list of those database
sentences that were used in the proof.
 \end{description}

If the proof attempt fails, \NIL\ is returned.

Here is a summary of the functions described in this section:

\begin{description}
 \item [\itype{init-prover(p effort)}] This function
initializes and returns an instance of \type{proof} containing
auxiliary information about the proof process.
 \item [\itype{cont-prover(proof)}] Resume the task associated to the
auxiliary information in the given instance of the \type{proof}
structure.  This function returns an instance of \type{bcanswer}
corresponding to the result obtained by the prover, and returns as
soon as a solution to the original query is found.
 \end{description}

\subsection{The multivalued proof tasks}
\label{s.consider}

The individual proof attempts generated by a multivalued query are
instances of structures that include the structure
\type{consider}\index{consider structure}.  The exact details of this
and related structures are unlikely to be of much interest to a user
of the \mvl\ system, but we include them here in the interests of
completeness.  These are the fields in the
\type{consider} structure:
 \begin{description}
 \item [\itype{prop}] The proposition being investigated by this task.
 \item [\itype{parent}] The task for which the current task is a
child.
 \item [\itype{children}] Children of the current task.
 \item [\itype{truth-value}] The current best estimate of the truth value
that should be assigned to the proposition being considered.
 \item [\itype{relevant}] A list of answers that, if returned by this
task, are known to be relevant to the original query.
 \item [\itype{irrelevant}] A list of answers know {\em not\/} to be
relevant to the original query.
 \item [\itype{paths}] The \mvl\ system goes to a great deal of effort
to ensure that it does not consider proofs that will not affect the
eventual response to a user's query; intermediate information used in
this computation is stored in this field.  Further details can be found
in the file \type{bc.lisp} and the files to which it points.
 \end{description}

When the prover must select a task to work on next, it does so via the
following function:
 \begin{description}
 \item [\itype{choose-task}] Selects a task to work on next.  The
current implementation of \mvl\ simply selects any active task (see
Section \ref{s.bc}).
 \end{description}

\section{Inference}
\label{s.inf}

\subsection{Backward chaining}
\label{s.bc}

The information in Section \ref{s.fo} is more technical than most users
of \mvl\ are likely to need.  What is of much greater interest are those
database query tools that are like \type{lookup} and \type{lookups} but,
instead of simply looking for a fact in the database, attempt to derive
it as a consequence of other database information.  These functions use
those of the previous section as described in the accompanying technical
papers.

The functions accept three sorts of cutoff information, instead of the
simple information provided to \type{lookup} and \type{lookups}.  The
first, provided via the \type{cutoffs} keyword, is used to construct a
test to decide whether an answer computed by the prover is good enough
to be returned when the proof effort is complete.  The second,
provided via the \type{succeeds} keyword, indicates that an answer
should be returned as soon as it can be shown that the final answer
will pass the given test.  The final cutoff, supplied via the
\type{terminate} keyword, indicates that the answer should be returned
as soon as there is any indication that the final answer {\em might\/}
pass the given test.

As an example, suppose that we were trying to find a plan to achieve some
goal.  We might well be prepared to terminate our search as soon as we found
a plan that was {\em guaranteed\/} to achieve the goal; if we completed our
search {\em without\/} finding such a plan, we might then want to return
all plans that we {\em expected\/} would achieve the goal.  It is to
facilitate such possibilities that the backward-chaining functions expect
three sets of cutoff information.  The default value for \type{cutoffs}
is the standard one associated to the active bilattice (see Section
\ref{s.bilattice}); the default choice for \type{succeeds} corresponds
to complete information for the singular versions of the queries and
guarantees failure for the plural versions (since an answer might be
about to be discovered for a different instantiation of the query).
The default value for \type{terminate} always fails.

Here, then are the various backward-chaining functions.  They all
return two values: an answer (or list of answers for the plural
versions \type{bcs} and \type{consider}) and an instance of the
\type{analysis} structure describing tasks remaining to be considered
by the prover.  In addition, they all accept an optional
\itype{analysis} keyword that, if supplied, should be an instance of
the \type{analysis} structure that is to be continued instead of a
new one being created.  This is to allow the user to continue proof
efforts that have returned for some reason.
 \begin{description}
 \item [\itype{bc(q)}] Attempts to prove the query $q$ from information
in the database, looking only for a single answer.
 \item [\itype{bcs(q)}] Attempts to prove the query $q$ from information
in the database, looking for all possible answers.
 \end{description}

All of these functions work by maintaining an instance of the
\type{analysis} structure\index{\type{analysis} structure} that
includes a variety of information about the status of the current
proof attempt.  Here are the fields in this structure:
 \begin{description}
 \item [\type{prop}] The proposition being considered by the prover.
 \item [\type{cutoffs}] The cutoff information that must be satisfied
by any result before it is to be returned.
 \item [\type{succeeds}] The cutoff information that will allow the
prover to terminate early if an answer is guaranteed to pass the test.
 \item [\type{terminate}] The cutoff information that will allow the
prover to terminate early if an answer might pass the test.
 \item [\type{active-tasks}] A list of the proof tasks (see Section
\ref{s.consider}) that still need to be investigated.
 \item [\type{main-task}] The task that was created upon the 
original invocation of the prover.
 \item [\type{returned-answers}] A list of answers that have already
been returned to the user; this is to keep \mvl\ from compulsively
re-returning these answers if the \type{analysis} is reinvoked.
 \end{description}

The top-level functions using this structure are as follows:
 \begin{description}
 \item [\itype{init-analysis(prop)}] The \mvl\ primitive upon which
the proving functions are based, this function accepts a proposition
and creates and returns an instance of the \type{analysis} structure.
 \item [\itype{cont-analysis(analysis)}] This function continues the
investigation of a specific analysis, and can be used to resume a
proof effort that has been curtailed for some reason.  It returns two
values: an instance of the \type{answer} structure giving the results
of the proof attempts, and a modified version of the \type{analysis}
structure that can be used to continue the proof process if desired.
 \end{description}

The functions \type{init-analysis} and the user functions such as
\type{bc} all accept the following keywords:
 \begin{description}
 \item [\type{cutoffs}] Cutoff test as described above.
 \item [\type{succeeds}] Success test as described above.
 \item [\type{terminate}] Termination test as described above.
 \item [\type{initial-value}] The truth value assigned to the query in
the absence of database information.
 \item [\type{analysis}]  If supplied, use this analysis instead of
creating a new one.
 \end{description}

\subsection{Forward chaining}
\label{s.fc}

Forward chaining in \mvl\ is not as simple as it is in many other
logic programming languages because of the \nm\ nature of multivalued
inference.  It follows from this that inserting a new sentence into
the database may cause a previous conclusion to be retracted, and
this may in turn cause other sentences to become unsupported.

There are two functions that use this information to modify the
database.  The more common is called \itype{fc} and accepts as an
argument a proposition to be added to the database.  The proposition
is added and provided that it was not already in the database when the
forward chainer was invoked, its consequences are also added to the
database and forward chained upon.  The \type{fc} function accepts
the following keywords:
 \begin{description}
 \item [\itype{value}] This keyword is similar to that used by
\type{stash} and \type{state}.
 \item [\itype{increment-p}] This keyword is also similar to that used
by \type{stash} and \type{state}.  The default value is \NIL.
 \item [\itype{fc-increment-p}] When the information content of some
sentence in the database has {\em fallen}, should this drop be
propagated to other sentences that appear to depend on the sentence
in question?  The default value is \T.
 \item [\itype{cutoffs}] Only forward chain if the sentence being
inserted into the database has a truth value that passes the test
corresponding to \type{cutoffs}.  (See Section \ref{s.cutoffs}.)  The
default value used is the one associated with the active bilattice.
 \item [\itype{include-modal}] If \T\ (the default), then the forward
chainer also looks for rules that fire when the truth value of a
modal expression involving $p$ changes (where $p$ is the sentence
being forward chained upon).  See Section \ref{s.modal}.  The reason
this keyword is included is that searching for rules of this form is
computationally expensive.
 \end{description}

Slightly less common is \itype{erase}, the function used to remove a
proposition from the database.  This function accepts a database
sentence as an argument, and also the keywords \itype{value},
\itype{decrement-p} and \itype{cutoffs}, which are similar to the
corresponding keywords for \type{fc}.

Summarizing:
 \begin{description}
 \item [\itype{fc(p)}] Inserts the sentence $p$ into the database and
forward chains.  Allowed keywords: \type{value}, \type{increment-p},
\type{fc-increment-p}, \type{cutoffs} and \type{include-modal}.
 \item [\itype{erase(p)}] Removes $p$ from the database and forward
chains.  Allowed keywords: \type{value}, \type{decrement-p} and
\type{cutoffs}.
 \end{description}

\subsection{Debugging tools}

Debugging facilities are limited in \mvl.  The facilities that exist
describe the prover's activities during backward-chaining,
forward-chaining, or while the first-order theorem prover is looking
for a conventional proof of some proposition.

\subsubsection{Backward chaining}

Diagnostics are displayed during backward chaining using the two
global variables \itype{*bc-trace*} and \at\index{*anytime-trace*}.

Whenever a task is encountered that involves a proposition that is an
instance of some element of \type{*bc-trace*}, information concerning
the task will be displayed to the user.  The variable \at\ is
similar, but only reports when an {\em inference\/} task is
encountered.  An inference task is one that is trying to prove
something, as opposed to one analyzing a modal operator; this
distinction is described in greater detail in the file {\type
bc.lisp}.  In practice, the distinction between \type{*bc-trace*} and
\at\ is that the former should be used for debugging purposes, while
the latter should be used if you are interested in watching the
prover approximate and then refine its eventual answer to some
declarative query.

Although the user can manipulate \type{*bc-trace*} and \at\ directly,
there are four functions written to help do so:
 \begin{description}
 \item [\itype{bc-trace(pat)}] Include \type{pat} on
\type{*bc-trace*}, so that any instance of \type{pat} will be
traced.
 \item [\itype{bc-untrace(pat)}] Remove any instance of \type{pat}
from \type{*bc-trace*}.
 \item [\itype{anytime-trace(pat)}] Include \type{pat} on \at, so
that any instance of \type{pat} will be traced.
 \item [\itype{anytime-untrace(pat)}] Remove any instance of
\type{pat} from \at.
 \end{description}
 For all of these functions, the argument is optional.  If omitted,
\type{(?*)} (of which everything is an instance) is used.

\subsubsection{Forward chaining}

Tracing the activity of the forward chainer is essentially identical
to tracing the behavior of the backward chainer; the global variable
\type{*fc-trace*} is used and is manipulated using the functions:
 \begin{description}
 \item [\itype{fc-trace(pat)}] Include \type{pat} on
\type{*fc-trace*}, so that any instance of \type{pat} will be
traced.
 \item [\itype{fc-untrace(pat)}] Remove any instance of \type{pat}
from \type{*fc-trace*}.
 \end{description}
 Once again, the argument is optional.

\subsubsection{The first-order prover}

Tracing the activity of the first-order theorem prover itself is
somewhat different.  The facility for doing this is similar to the
Common Lisp \type{trace} facility, but instead of reporting when a
function is invoked, it reports when an attempt is made to prove an
instance of a particular predicate, and also when such an instance is
found successfully in the database.

The functions used to control this are \itype{bc-watch} and
\itype{bc-unwatch}.  They are similar to \type{trace} and
\type{untrace}:

 \begin{description}
 \item [\itype{bc-watch(preds)}] Accepts as arguments any number of
predicates that should be watched as the prover proceeds.  Each
predicate can also be a list of the form \type{(predicate depth)}
where \type{depth} is the number of successive ancestors to be
displayed when the prover encounters an instance of the predicate.
(This can be used to understand {\em why\/} a particular clause is
relevant to the proof being attempted.)  \type{bc-watch} returns a
list of the predicates currently being {\sc watch}ed.
 \item [\itype{bc-unwatch(preds)}] Removes the given predicates from
the list of {\sc watch}ed predicates.  If no arguments are supplied,
all predicates are un{\sc watch}ed.
 \end{description}

There are two additional features involved in tracing the activities
of the first-order prover.  The first is used whenever the global
parameter \itype{*save-nodes*} is \T\ (which it is by default).  When
the first-order prover displays information regarding a particular
proof node, the node is labelled with a number.  Subsequent nodes
will then use this label if they need to refer to the node in
question (e.g., as a parent).

If \type{*save-nodes*} is \T, then \type{*break-node*} can be set
to an integer; when a node is displayed that is labelled with this
integer, a \type{(break)} will be executed.  This is to allow the
user to interrupt the first-order prover if he needs to examine the
proof stack or some other \mvl\ internal.

 \begin{description}
 \item [\itype{*save-nodes*}] If \T\ (the default), nodes are
labelled with consecutive integers when they are displayed by the
first-order prover.
 \item [\itype{*break-node*}] The number of a node at which the
proof process should be interrupted.
 \end{description}

\section{Procedural attachment}
\label{s.attach}

In some cases, it may be desirable to store a fact in the database
{\em without\/} using the functions supplied in \mvl.  Suppose, for
example, that we have an $n+1$-ary relation $R$ that corresponds to
an $n$-ary function $f$, so that $R(x_1,\dots,x_n,y)$ holds just in
case $f(x_1,\dots,x_n)=y$.  Now, if we assert $R(c_1,\dots,c_n,d)$
into the database, where $d$ and the $c_i$ are constants, we should
probably also {\em remove\/} any sentences of the form
$R(c_1,\dots,c_n,d')$, where $d'\neq d$.  This sort of a situation
might arise if we were monitoring the temperature in a room, and
represented the observed temperature using a database sentence such
as
 \begin{equation}
\type{(TEMPERATURE ROOM 72)}
\label{e.temp}
\end{equation}

We handle this by informing \mvl\ that to stash a sentence of the
form \type{(TEMPERATURE ...)}, it should use not the usual function
\type{stash}, but some other user-defined function, say
\type{unique-stash}.  This new function would first \type{unstash}
any sentence that agreed with the given one except for the value of
the final term, and would then invoke \type{stash} on the original
sentence, but with \apa\ set to \NIL\ so that \type{unique-stash} was
not invoked a second time.

To modify the function used by \mvl\ to stash sentences such as this,
the user should simply place the desired function (in this case,
\type{unique-stash}) on the \type{stash} property of the relation in
question (in this case, \type{TEMPERATURE}).\footnote{As described in
Section \ref{s.invocation}, the \type{stash} property of a relation
can also be an association list, if multiple invocations of the
system are running in parallel.} The system will then automatically
invoke \type{unique-stash} instead of \type{stash} when a fact such
as (\ref{e.temp}) is asserted.

The functions within \mvl\ that can be modified in this fashion are
\type{stash}, \type{unstash}, \type{lookup} and \type{lookups}.  In
addition, the procedural attachments for \type{lookup} and
\type{lookups} are linked, so that if \type{lookup} has a procedural
attachment and \type{lookups} does not, a call to \type{lookups} will
invoke the procedural attachment for \type{lookup} and then pluralize
the result.  Similar remarks hold if \type{lookups} and not
\type{lookup} has a procedural attachment.

The procedural attachment is effected using the following functions:
 \begin{description}
 \item [\itype{invoke(fn p)}] This function invokes \type{fn} on the
sentence $p$.  In addition, any keyword arguments are passed to
\type{fn}, but with \type{\&allow-other-keys} so that the keywords
will be ignored if \type{fn} is not expecting them.
 \item [\itype{lookup-call(fn p)}] This function is similar to
\type{invoke}, but modifies the values returned to make them agree
with those expected to be returned by \type{lookup}.  Unless the
function \type{fn} supplies them, the truth value is assumed to be
\type{true}.  The second value is \type{attach}.
 \item [\itype{lookups-call(fn p)}] The plural version of
\type{lookup-call}.
 \end{description}

\subsection{Facilities to create procedural attachments}

Some facilities exist to create procedural attachments within the
\mvl\ system.  The following two macros are used for this purpose:
 \begin{description}
 \item [\itype{lisp-predicate(pred)}] This macro is used to indicate
that the given predicate can be evaluated by a \lisp\ call.  It will
also cause \mvl\ to signal an error if an attempt is made to stash
or unstash an instance of the function.
 \item [\itype{lisp-defined(fn)}] This macro is used to indicate that
the given function can be evaluated by a \lisp\ call.  It will also
cause \mvl\ to signal an error if an attempt is made to stash or
unstash an instance of the function.  \type{Lisp-defined} also
accepts an optional argument, which is an inverter function.  The
overall result of attempting to lookup an instance of the function
can be described as follows:
 \begin{enumerate}
 \item If none of the arguments to the function (i.e., all but the
last of the arguments of the relation) is a variable, then the function
is called and the result is unified with the final argument of the
relation.
 \item If the inverter function has been supplied and is equal to \T,
then the function is called and the result is unified with the final
argument of the relation even if some of the arguments are variables.
 \item If some of the arguments are variables but the result is not a
variable, then the inverter function is invoked if one has been
supplied.  The arguments passed to the inverter function are the
result (i.e., the last argument of the relation), and a list of all
of the arguments (i.e., all but the last argument of the relation).
If there is no inverter function, then \NIL\ is returned.  In
addition, if \wbli\ \index{*warn-on-bad-lisp-invocations*} is not
\NIL, the user is warned.  (The parameter \wbli\ defaults to \NIL.)
 \item If both some of the arguments and the result are variables, then
\NIL\ is returned.  In addition, if \wbli\ is not \NIL, the user is
warned.
 \end{enumerate}
 \end{description}

The current version of \mvl\ invokes \type{lisp-predicate} on the
following \lisp\ predicates: \type{numberp}, $<$, $>$, \type{equal},
\type{null}, \type{atom}, \type{listp}, \type{varp}, \type{varp?},
\type{varp$*$}, \type{groundp}, and $=$.  The procedural attachment
to $=$ attempts to unify its arguments.

The macro \type{lisp-defined} is invoked on the following:
\type{car}, \type{cdr}, \type{cadr}, \type{cons}, \type{list},
\type{list*}, \type{revappend}, \type{append}, \type{reverse},
\type{length}, \type{assoc}, \type{delete}, \type{+}, \type{*},
\type{-}, \type{/}, \type{floor}, \type{ceiling}, \type{truncate},
\type{round}, \type{float}, \type{mod}, \type{rem}, \type{max},
\type{min}, \type{sin}, \type{cos}, \type{atan}, \type{sqrt}, 
and \type{eval}.

\subsection{Other procedurally attached relations}

In addition to the \lisp\ functions and predicates described in the
previous section, the following relations have been defined and
are procedurally attached to suitable \lisp\ functions:
 \begin{description}
 \item [\itype{(bagof}\type{ ?x ?p ?b)}] This relation holds for $x$,
$p$ and $b$ if and only if $b$ is a list of all $x$ satisfying $p$.
 \item [\itype{(done exp-1 ... exp-n)}] This predicate, which is
useful for stashing only, causes \lisp\ to evaluate each of the subsequent
expressions.
 \item [\itype{(value symbol val)}] Stashing this predicate sets the
given symbol to the value \type{val}.  Unstashing it makes the symbol
unbound if its current value unifies with \type{val}.
 \item [\itype{(property symbol value indicator)}] Stashing this
predicate sets the symbol's \type{indicator} property to \type{value}
for the current invocation of \mvl.  Unstashing it removes the
\type{indicator} property for the current invocation of \mvl\ if the
previous value of this property unifies with \type{value}.  (See the
definitions of \type{get-mvl} and \type{rem-mvl} in Section
\ref{s.invocation}.)
 \item [\itype{(known p)}]  Stashing this predicate simply stashes $p$;
unstashing it causes an error.  Looking it up looks up $p$.
 \item [\itype{(unknown p)}]  Stashing this predicate unstashes $p$;
unstashing it causes an error.  Looking it up fails if $p$
can be found in the database, and succeeds otherwise.
 \item [\itype{(member item list)}] Tries to unify \type{item} with
the various elements in \type{list}.
 \end{description}

\section{Testing the MVL system}
\label{s.foln}

\subsection{As a diagnostic tool}

There are two functions available to help you to test both \mvl\ and
any applications you write using it.  The first, \type{test-mvl}, is
one you should already have run; if \mvl\ has been installed successfully,
it should return 0.

\type{Test-mvl} works by invoking the generic testing function
\type{mvl-test-file} on the file \type{mvl.test} in the \mvl\
directory.\footnote{\type{Test-mvl} works in its own invocation of
\mvl, so that the existing database is not modified.  This is not true
of \type{mvl-test-file}, however.} This file also contains
documentation regarding the details of the various tests being
executed.

In general, \type{mvl-test-file} accepts as input a file name, and
looks for a file with an extension \type{.TEST} on either the \mvl\
or the current directory.  It then works through the file by
evaluating the first s-expression it contains, and comparing the
value returned to the second s-expression.  If these two
s-expressions match, the next s-expression is evaluated, and so on.
If there is a mismatch, \type{mvl-test-file} reports the error;
eventually, it returns the total number of errors found in the entire
file.  When doing the comparison, \type{**} acts as a wildcard.
Thus, if the second s-expression is simply \type{**}, it will always
match; if it is \type{(** ** **)}, it will match whenever the first
expression returns a list of three values, and so on.

\begin{description}
 \item [\itype{test-mvl}] Tests the \mvl\ system by invoking the function
\type{mvl-test-file} on the file \type{MVL.TEST}.
 \item [\itype{mvl-test-file(string)}] Can be used to evaluate
various forms by the \mvl\ system and compare the results to
predefined values.  Both the forms and the expected answers are in
the \type{.TEST} file whose name is supplied to the function.
Returns the number of errors found.
 \end{description}

\subsection{As a tutorial}

The information used by the function \type{test-mvl} can also be used
as a tutorial in the \mvl\ system.  Take a look at the file
\type{mvl.test} that this function uses; you will see that the first
two lines are
 \begin{verbatim}
     (notime (mvl-test-file "bc" "backward chainer")
     0
\end{verbatim}
 This means that the function \type{mvl-test-file} will be invoked on
the file \type{bc.test}, and that this function should return 0
(i.e., 0 errors).  So let's have a look at \type{bc.test}:
 \begin{verbatim}
     (load-logic #\f)                    ;first-order logic
     **
\end{verbatim}
 Here, the system loads the bilattice corresponding to first-order
logic and is not interested in the value returned by the loading
operation.  Next, a declarative database corresponding to a full
adder is loaded (with the timer turned off).  The sentences in this
database can be found in the file \type{f1.mvl} and are also
described in Section 15 and Appendix B.2 of the paper ``Multivalued
logics: A uniform approach to inference in artificial intelligence.''

Next, we see in \type{mvl.test} the two lines
 \begin{verbatim}
     (get-bdg '?x (answer-binding (bc '(val (out 1 f1) 1 ?x))))
     OFF                   ;the first output (the sum line) should be off
\end{verbatim}
 We invoke the backward chainer on the sentence
 \[\type{(val (out 1 f1) 1 ?x))}\]
 to see if we can prove that the first output of the full adder takes
some value \type{?x}.  Then we get the binding of this variable to
see what the value actually is -- \type{OFF}, in this case.  The
remainder of the various test files can be analyzed similarly; many
of the examples in these files can also be found in the various
technical papers included in the \mvl\ distribution package.

\section{System support}
\label{s.support}

\subsection{Functions that compile and load the MVL system}

The following three functions can be used to manipulate the \mvl\
system as a whole:
 \begin{description}
 \item [\itype{load-mvl}] Loads the \mvl\ system.
 \item [\itype{compile-mvl}] Compiles and loads the \mvl\ system.  On
an Allegro system, the keyword \type{recompile} forces a recompile of
all of the files and the keyword \type{reload} is ignored (since the
files will all be reloaded anyway).  On a Symbolics system, the
keyword \type{reload} forces all of the files to be reloaded and the
keyword \type{recompile} forces them all to be recompiled.
 \item [\itype{release-mvl}] (Symbolics only)  Compiles and loads the
\mvl\ system, and declares the new version to be \type{released}.  The
keyword \type{no-recompile} can be used to skip the compilation step.
 \end{description}

\subsection{The MVL package}
\label{s.package}

In order to avoid name conflicts with existing code, the \mvl\ system
resides in its own package, \type{MVL}\index{MVL package}.  All
of the symbols and functions appearing in the appendix are exported
from this package and can therefore be accessed by invoking the
function \type{use-package(mvl)}.  When a new bilattice is loaded,
the modal operators it defines are also exported. 

To use the \mvl\ system, invoke the function \type{(use-package 'mvl)}
to make these symbols accessible.\footnote{As an alternative,
\type{(in-package 'mvl)} wil put you in the MVL package itself.}

\subsection{Running several copies of MVL in parallel}
\label{s.invocation}

Some users may wish to run several copies of the \mvl\ system
simultaneously, perhaps associating them with distinct {\sc lisp}
processes or with individual agents.  (The use of multiple copies of
\mvl\ to represent distinct components of a single agent's knowledge
goes against the grain of the multivalued approach and is not
recommended.)  In order to support this facility, the following
functions and macros have been defined:
 \begin{description}
 \item [\itype{existing-mvl-invocation}] This function returns a
single value that is an instance of the structure \mi.  It is useful
primarily for saving the current copy of \mvl\ so that it can be
restored later.
 \item [\itype{create-mvl-invocation}] This function sets up a new
copy of \mvl.  It accepts an optional argument that is an instance of
\mi; this can be used to reset the state of \mvl\ to a previous one.
If the optional argument is omitted, default values are used.  The
value returned is the \mvl\ invocation that was active before the
function was called, so that another call to
\type{create-mvl-invocation} can be used to restore the original
environment.
 \item [\itype{with-mvl-invocation}] This is a macro similar to the
Common Lisp macro \type{with-open-file}.  It accepts an optional
argument that is an instance of \mi\ and is followed by a body of
forms (presumably including a \type{load-logic}) that are evaluated
in a locally-bound copy of \mvl.
 \end{description}

These functions work by locally or globally rebinding the values of
the various parameters used by the \mvl\ system
(\type{*active-bilattice*}, \type{true} and so on).  The only
difficulty with this is that procedural attachments for a relation
$R$ use the property list of $R$, and this cannot be rebound.

In order to cater to this difficulty, there is an additional global
variable \itype{*mvl-invocation*} that labels the currently active
copy of \mvl.  To determine what procedure is attached to the
\type{stash} operation for $R$, the \type{stash} property of $R$ is
examined.  There are now the following possibilities:
 \begin{enumerate}
 \item If the value is \NIL, \type{stash} is not procedurally attached,
so the usual \mvl\ function is called.
 \item If the value is an atom, then it is assumed that this atom is
the name of a function that is the procedural attachment for \type{stash}
independent of the copy of \mvl\ that is currently active.
 \item If the value is a list, it is treated as an association list
and a cons whose car is \type{*mvl-invocation*} is searched
for.
 \begin{enumerate}
 \item If such a cons is found, its cdr is used as the procedural
attachment to \type{stash}. 
 \item If no such cons is found, but there is a cons whose car is \T,
then the cdr of that cons is used.
 \item If no cons has either \type{*mvl-invocation*} or \T\ as its
car, then \type{stash} is assumed to be unattached.
 \end{enumerate}
 \end{enumerate}
 The other functions described in Section \ref{s.attach} are handled
similarly.

The association lists described above are manipulated using the following
functions:
 \begin{description}
 \item [\itype{get-mvl (symbol indicator)}]  Like the Common Lisp
\type{get}, but if the value is a list, treats it as an a-list and
searches for an item whose car matches \type{*mvl-invocation*}.
The value taken can be changed with \type{setf}.
 \item [\itype{rem-mvl (symbol indicator)}] Like the Common Lisp
\type{remprop}, but if the value is a list, treats it as an a-list
and removes the item whose car matches \type{*mvl-invocation*}.
 \end{description}

\section{Selecting a bilattice}
\label{s.bilattice}

\subsection{Bilattice structure}

Of course, most of the power of the \mvl\ system comes not from the
features that were described in Sections \ref{s.fol1}--\ref{s.support}
of this manual, but from the fact that the user can change
representation language conveniently within the system.  He does this
by changing the {\em bilattice\/} that he is using.

A bilattice is the mathematical structure from which the truth values
used by the system are selected, and the formal ideas underlying this
approach are discussed at length in the accompanying technical
papers.  In the \mvl\ system, a bilattice is described as an instance
of the \type{bilattice} structure.  This structure has 27 fields; of
these, 17 are used to provide the formal information required by the
system, and the others are used to simplify the user interface.  Here
are 13 of the fields used to describe the bilattice itself:

\begin{description}
 \item [\itype{true}] A \lisp\ object that is the truth value assigned
to a sentence that is unconditionally true, such as the tautology
\type{(OR (P) (NOT (P)))}.
 \item [\itype{false}] A \lisp\ object that is the truth value assigned
to a sentence that is unconditionally false, such as the negation of a
tautology.
 \item [\itype{unknown}] A \lisp\ object that is the truth value assigned
to a sentence about which nothing is known.
 \item [\itype{bottom}] A \lisp\ object that is the truth value assigned
to a sentence that is both true and false.  If any sentence is assigned
this truth value (or can be shown to have it), it indicates the presence
of a contradiction in the database.
 \item [\itype{eq}] A \lisp\ function that takes two truth values as
arguments and returns \T\ if they are equal and \NIL\ if they are
not.  This function is needed because it may be possible to represent
a single element of the bilattice as a \lisp\ object in a variety of
fashions.
 \item [\itype{and}] A \lisp\ function that takes two truth values as
arguments and returns the truth value corresponding to their
conjunction.  If the two sentences $p$ and $q$ are labelled with
truth values $v_1$ and $v_2$ respectively, then the conjunction of
the truth values is the truth value that can be assigned to $p\wedge q$.
 \item [\itype{or}] A \lisp\ function that takes two truth values as
arguments and returns the truth value corresponding to their
disjunction.
 \item [\itype{not}] A \lisp\ function that takes a truth value as its
argument and returns the truth value corresponding to its negation.
 \item [\itype{plus}] A \lisp\ function that takes two truth values as
arguments and returns the truth value corresponding to their ``sum.''
If there is reason to believe that a sentence $p$ should be labelled
with truth value $v_1$, and independent reason to believe that it
should be labelled with truth value $v_2$, then the actual label is
given by the sum $v_1+v_2$.  This operation therefore corresponds to
combination of evidence.
 \item [\itype{dot}] A \lisp\ function that is the dual of the
\type{plus} operation.  It takes two truth values $v_1$ and $v_2$ as
arguments and returns the greatest truth value that could
subsequently be strengthened to a conclusion of either $v_1$ or
$v_2$.  This is a ``consensus'' operation.
 \item [\itype{kle}] A \lisp\ function that takes two truth values $x$
and $y$ and determines whether $x+y=y$.  Although this can also
be done using the bilattice $+$ and $=$ operations, this field can be
provided to give a more efficient implementation of this operation.
This field is optional.
 \item [\itype{tle}] A \lisp\ function that takes two truth values $x$
and $y$ and determines whether $x\vee y=y$.  This field is also
optional.
 \item [\itype{modal-ops}] A list of the modal operators associated
with the given bilattice.  See Section \ref{s.modal}.
 \end{description}

Three of the remaining formal slots are concerned with the incorporation
of binding information into truth values.  In some cases, when a
proposition is retrieved from the database, it is only an instance of
the database proposition that is used, and the truth value may have to
be modified to reflect that:
 \begin{description}
 \item [\itype{plug}]  This \lisp\ function accepts as arguments
a truth value and a binding list, and returns a potentially modified
truth value.  The default value is a function that leaves the truth
value unchanged.
 \end{description}

There is also a function that accepts a truth value and returns the
variables appearing within it:
 \begin{description}
 \item [\itype{vars}]  This \lisp\ function accepts a truth value
as its argument and returns a list of the variables appearing in it.
The default value is a function that returns \NIL\ independent of its
argument.
 \end{description}

The function \itype{absorb} is needed for subtle technical reasons
when working with bilattices with non-\NIL\ plugging functions;
further information on this can be found in the file \type{load.lisp}.
All three of these functions are \NIL\ if there is no interaction
between the bilattice truth values and the variables that they
contain.

The final formal field has to do with search control:
 \begin{description}
 \item [\itype{simplicity}] This \lisp\ function accepts a truth
value (which will frequently be the value obtained by conjoining the
premises in a particular proof attempt) and returns a figure
indicating the ``merit'' that should be assigned to the proof attempt
in question.  The value returned should be a nonempty list of
integers and is viewed lexicographically, with high values being
preferred to low ones.  The default value is a function that returns
the list \type{(0)} independent of the truth value supplied.
 \end{description}

The remaining ten fields in the \type{bilattice} structure are as
follows:
 \begin{description}
 \item [\itype{dws}] This \lisp\ function, ``dot with star,'' accepts
as arguments two truth values, $v_1$ and $v_2$.  The truth value returned
is $v_1$, modified so as to remove any contribution that could be
attributed to a sentence with truth value $v_2$.
 \item [\itype{stash-val}] A function that, given a sentence $p$,
returns the default truth value to be passed to \type{stash}.
 \item [\itype{unstash-val}] A function that, given a sentence $p$,
returns the default truth value to be passed to \type{unstash}.
 \item [\itype{cutoff-val}] The default cutoff point that must be
exceeded by any sentence if it is to be returned by \type{lookup}.
See Section \ref{s.cutoffs}.
 \item [\itype{print-fn}] This function is used to pretty-print
truth values in the associated bilattice.  Its three arguments
are a truth value, the bilattice in question, and a stream.
 \item [\itype{long-desc}] A long description of the bilattice, used
for input and output only.
 \item [\itype{short-desc}] A short description of the bilattice, also
used for input and output only.
 \item [\itype{char}] A one-character identifier for the bilattice.
This character can be used to change the bilattice using the function
\type{load-logic} described in Section \ref{s.change}.
 \item [\itype{type}] An overall description of the mathematical form
of the bilattice.  Perhaps it is a product of two other bilattices,
or formed from them in some other fashion.  This defaults to \type{basic}
if no other value is supplied; various methods for constructing new
bilattices from old ones are described in Section \ref{s.construct}. 
 \item [\itype{components}]  A list of the components used to construct
a bilattice whose \type{type} is not \type{basic}.
 \end{description}

\subsection{The active bilattice}

At any particular point in time, there will be a single bilattice
that has been selected by the user to be used in the application he
is developing.  This bilattice is stored as the value of
\type{*active-bilattice*}; in addition, many of the slots of the
active bilattice are directly available to the user.  Here they
are:

\begin{description}
 \item [\itype{*active-bilattice*}] The currently active bilattice.
 \item [\itype{true}] The truth value \type{true} in the active bilattice.
 \item [\itype{false}] The truth value \type{false} in the active bilattice.
 \item [\itype{unknown}] The truth value \type{unknown} in the active
bilattice.
 \item [\itype{bottom}] The truth value \type{bottom} in the active bilattice.
 \item [\itype{std-cutoffs}] Cutoff information constructed from the
value of \type{cutoff-val} in the active bilattice.  See Section
\ref{s.cutoffs}.
 \item [\itype{stash-value}] The function \type{stash-val} in the
active bilattice.
 \item [\itype{unstash-value}] The function \type{unstash-val} in the
active bilattice.
 \item [\itype{mvl-plus}] The function \type{plus} in the active bilattice.
 \item [\itype{mvl-dot}] The function \type{dot} in the active bilattice.
 \item [\itype{mvl-and}] The function \type{and} in the active bilattice.
 \item [\itype{mvl-or}] The function \type{or} in the active bilattice.
 \item [\itype{mvl-not}] The function \type{not} in the active bilattice.
 \item [\itype{mvl-eq}] The function \type{eq} in the active bilattice.
 \item [\itype{mvl-simplicity}] The function \type{simplicity} in the active
bilattice.
 \item [\itype{dot-with-star}] The function \type{dws} in the active 
bilattice.
 \item [\itype{mvl-vars}] The function \type{vars} in the active bilattice.
 \item [\itype{mvl-plug}] The function \type{plug} in the active bilattice.
 \item [\itype{mvl-print}] This function accepts a single argument, a truth
value, and modifies it in a way that will cause it to be pretty-printed.
The value returned by mvl-print should not be used for any purpose other
than printing.
 \end{description}

All of the functions above accept an optional argument that can be
used to evaluate their result using some bilattice other than the
active one.

\subsubsection{Functions derived from the active bilattice}

There are a variety of functions associated with the functions that
are defined when a new bilattice is made active.  In order to
simplify the description of these functions, we will denote
\type{mvl-and} by $\wedge$, \type{mvl-or} by $\vee$, \type{mvl-not}
by $\neg$, \type{mvl-plus} by $+$, and \type{mvl-dot} by $\cdot$.
The first auxiliary function defined is \type{dot-with-not}:
 \begin{description}
 \item [\itype{dot-with-not(x)}] Computes $x\cdot \neg x$.
 \end{description}

The other functions are constructed from the partial orders defined
using the bilattice operations.  We will write $x\leq_t y$ just in
case $x\wedge y = x$, for example, and $x\leq_k y$ just in case
$x\cdot y = x$.  (All of this notation follows that in the paper,
``Multivalued logics: A uniform approach to inference in artificial
intelligence''.)

Given these definitions, the following \lisp\ functions are defined:
 \begin{description}
 \item [\itype{k-le(x y)}] Returns \T\ if $x \leq_k y$, and \NIL\ otherwise.
 \item [\itype{k-ge(x y)}] Returns \T\ if $x \geq_k y$, and \NIL\ otherwise.
 \item [\itype{t-le(x y)}] Returns \T\ if $x \leq_t y$, and \NIL\ otherwise.
 \item [\itype{t-ge(x y)}] Returns \T\ if $x \geq_t y$, and \NIL\ otherwise.
 \item [\itype{k-lt(x y)}] Returns \T\ if $x <_k y$, and \NIL\ otherwise.
 \item [\itype{k-gt(x y)}] Returns \T\ if $x >_k y$, and \NIL\ otherwise.
 \item [\itype{t-lt(x y)}] Returns \T\ if $x <_t y$, and \NIL\ otherwise.
 \item [\itype{t-gt(x y)}] Returns \T\ if $x >_t y$, and \NIL\ otherwise.
 \item [\itype{k-not-le(x y)}] Returns \T\ if $x \not\leq_k y$, and
\NIL\ otherwise.
 \item [\itype{k-not-ge(x y)}] Returns \T\ if $x \not\geq_k y$, and
\NIL\ otherwise.
 \item [\itype{t-not-le(x y)}] Returns \T\ if $x \not\leq_t y$, and
\NIL\ otherwise.
 \item [\itype{t-not-ge(x y)}] Returns \T\ if $x \not\geq_t y$, and
\NIL\ otherwise.
 \item [\itype{k-not-lt(x y)}] Returns \T\ if $x \not<_k y$, and \NIL\
otherwise.
 \item [\itype{k-not-gt(x y)}] Returns \T\ if $x \not>_k y$, and \NIL\
otherwise.
 \item [\itype{t-not-lt(x y)}] Returns \T\ if $x \not<_t y$, and \NIL\
otherwise.
 \item [\itype{t-not-gt(x y)}] Returns \T\ if $x \not>_t y$, and \NIL\
otherwise.
 \end{description}

All of these functions accept an optional argument that is the
bilattice to be used in the comparison.  This argument defaults to
the currently active bilattice.

\subsubsection{Cutoff information}
\label{s.cutoffs}

We will refer to each of the 16 functions defined at the end of the
previous section (e.g., excluding \type{dot-with-not}) as a {\em tag}
\index{tag}.  Thus \type{k-le} is a tag, and so on.

We will now define a {\em simple test} to be a dotted pair
 \[\type{(v . tag)}\]
 where \type{v} is a truth value and \type{tag} is a tag.  We will
say that a truth value \type{x} {\em satisfies\/} the simple test
if and only if the function \type{tag}, when applied to the arguments
$x$ and $v$ (in that order), returns \T.

As an example, consider the simple test
 \[\type{(true . t-ge)}\]
 This test will be satisfied by any truth value that is $\geq_t
\type{true}$.  Since \type{true} itself is the only truth value
satisfying this test, the test will be satisfied for \type{true} and
no other truth values.  This is the test that is frequently used to
decide to terminate a proof effort.\footnote{Tests in the \mvl\ system
are in fact \lisp\ structures with \type{value} and \type{tag} fields;
we are using dotted lists in this manual only to make the presentation
clearer.}

Alternatively, the test used to determine whether or not to return
an answer after the theorem prover runs to completion is often given
by
 \[\type{(unknown . k-gt)}\]
 This test will succeed for any truth value about which more is known
that \type{unknown} -- in other words, for any truth value other than
\type{unknown} itself.

In general, the simple tests are combined into {\em cutoffs}.  A cutoff
is simply a list of lists of simple tests, and is interpreted as being
in conjunctive normal form.  Thus the cutoff corresponding to the simple
test $s$ is given simply by
 \[\type{((s))}\]

A cutoff corresponding to the conjunction of tests $s_1,\dots,s_n$ is given
by
 \[\type{((s-1) ... (s-n))}\]
 while that corresponding to their {\em disjunction\/} is given by
 \[\type{((s-1 ... s-n))}\]

In general, we can build up conjunctions of disjunctions of simple
tests in as complex a manner as we wish.  The cutoff \NIL\
corresponds to the empty conjunction and therefore always succeeds;
the cutoff \type{(NIL)} contains an empty disjunction and therefore
always fails.  For convenience, these two cutoffs are stored in the
global constants \itype{*always-succeeds*} and
\itype{*never-succeeds*} respectively.

Given a truth value $x$ and a cutoff $c$, the expression \type{(test x c)}
evaluates to \T\ if $x$ passes the test, and to \NIL\ if it fails it.
There is also a function that negates cutoffs:
 \begin{description}
 \item [\itype{test(value cutoffs \&optional bilattice)}] Returns \T\
if the truth value passes the test associated to \type{cutoffs}, and
\NIL\ otherwise.  The test is performed for the supplied bilattice,
or for the currently active one if the optional argument is omitted.
 \item [\itype{cutoffs-not(cutoffs \&optional bilattice)}] Returns a
new cutoffs $c$ such that
 \[\type{(test x c)} = \type{(test (mvl-not x) cutoffs)}\]
 for all truth values $x$.  In other words, $x$ passes $c$ if and
only if $\neg x$ passed the cutoffs supplied originally.
 \end{description}

\subsection{Selecting a new bilattice}
\label{s.change}

To change the bilattice being used, simply invoke the function
\type{load-logic}.  There are three ways to specify the new bilattice:
 \begin{enumerate}
 \item Do nothing.  You will then be asked to select the bilattice
from the list of bilattices known to the system.
 \item Supply the one-character identifier of the new bilattice as an
argument to \type{load-logic}.  The new bilattice must already be
known to the system.
 \item Supply the new bilattice itself as an argument to
\type{load-logic}.
 \end{enumerate}
 \type{Load-logic} empties the database before the bilattice is
changed.

The bilattices known to the system are kept in the global variable
\type{*bilattices*}.  To add a new bilattice to \type{*bilattices*},
invoke the function \type{bilattice} with the bilattice as an
argument.  This function will warn you if you are overwriting another
bilattice that has the same one-character designator.

Summarizing:
 \begin{description}
 \item [\itype{*bilattices*}] A list of the bilattices known to the system.
 \item [\itype{bilattice(b)}] Add the bilattice $b$ to \type{*bilattices*}.
 \item [\itype{load-logic}] Empty the database and make a new bilattice
active.  Accepts as an optional argument either a character designating
a bilattice, or the new bilattice itself.  If no argument is supplied,
the user is prompted for the new bilattice.
 \end{description}

\section{Bilattices available in the MVL system}

Here are the bilattices available in the \mvl\ system as distributed.
Each is discussed more fully in a subsequent section:
 \begin{description}
 \item [\itype{*first-order-bilattice*}] The bilattice corresponding to
first-order logic.
 \item [\itype{*atms-bilattice*}] The bilattice corresponding to an \atms,
where the justification information is maintained in disjunctive normal
form as suggested by de Kleer.
 \item [\itype{*cnf-atms-bilattice*}] An alternative version of the \atms\
bilattice where the justification information is maintained in conjunctive
normal form.
 \item [\itype{*default-bilattice*}] A bilattice that can be used for simple
default reasoning.
 \item [\itype{*hierarchy-bilattice*}] A bilattice similar to the
default bilattice but where the defaults are labelled with
integer-valued strengths.
 \item [\itype{*first-order-atms-bilattice*}] This bilattice corresponds
to a fully first-order \atms, as opposed to the more conventional one
dealing only with propositional logic.
 \item [\itype{*circumscription-bilattice*}] The bilattice used to compute
the results of the circumscription axiom.
 \item [\itype{*probability-bilattice*}] An extended version of the \atms\
bilattice that maintains probabilistic information as well.
 \item [\itype{*time-bilattice*}] A compound bilattice that can be used
for reasoning about events occurring against a temporal background.
 \end{description}

\subsection{First-order logic} \index{first-order logic}

The simplest bilattice is that corresponding to first-order logic, and
contains only the four truth values \type{true} (displayed as \type
T), \type{false} (displayed as \type F), \type{unknown} (\type U) and
\type{bottom} (\type{T/F}).  This bilattice is identified using the
character {\bf f}, and is active when \mvl\ is first loaded.

\subsection{Assumption-based truth maintenance}
\index{assumption-based truth maintenance}

Except for the \nm\ and probabilistic \atms\ bilattices, all of the
bilattices corresponding to \atms's have truth values that are dotted
pairs \type{(j . k)}, where $j$ is a justification for the sentence
being labelled, and $k$ is a justification for its negation.  The
bilattices differ in that the justifications are represented
differently.

\subsubsection{The ATMS bilattice}

For the simplest \atms\ bilattice, the justifications are simply
disjunctive normal form expressions.  Thus the justification
 \begin{equation}
\type{((P23 P37) (P1))}
\label{e.just}
\end{equation}
 is interpreted to mean that the sentence in question is a
consequence either of \type{P1}, or of \type{P23} and \type{P37} in
conjunction.

This bilattice is identified by the character {\bf a}.

\subsubsection{ATMS's in conjunctive normal form}

Alternatively, we can represent the justifications in conjunctive
normal form, as opposed to disjunctive normal form.  This representation
is useful, for example, in diagnosing multiple faults, since we want to
be able to explain an observed failure {\em either\/} in terms of the
collective failures of one set of components, or the collective failures
of another set, and so on.

This bilattice is identified by the character {\bf b}.

\subsubsection{First-order ATMS's}
\label{s.fo-atms}

It is also possible to extend these ideas slightly to obtain a fully
first-order \atms.  In order to do this, we need simply modify
justifications such as that appearing in (\ref{e.just}) so that the
primitive elements of the justification are not atoms referring to
propositions, but pairs
 \[\type{(atom . binding-list)}\]
 where the atom refers to a proposition and the binding list
indicates precisely what instantiation of that proposition is
actually needed to justify the sentence in question.

As an example, suppose that an assumption in our database is \type{P1},
given by
 \begin{equation}
\type{(FRIENDLY ?X)}
\label{e.assume}
\end{equation}
 and indicating that everything is friendly.  If we also have the rule
of inference
 \[\type{(<= (HAPPY ?X) (FRIENDLY ?X))}\]
 indicating that anything friendly is happy, and this rule is unconditionally
true (i.e., not an assumption), then the truth value assigned to the
query \type{(HAPPY HARRY)} will be
 \[\type{((((P1 . ((?X . HARRY))))))}\]
 indicating that Harry is happy by virtue of (\ref{e.assume}) applied
to $x=\mbox{Harry}$, and that there is no justification for the
conclusion that Harry is not happy.  (Recall that all of the \atms\
bilattices combine justification for and against each sentence into a
single truth value.)

This bilattice is identified by the character {\bf o}.

\subsubsection{Probabilistic ATMS's}
\label{s.prob}

Truth values in the probabilistic \atms\ bilattice consist of dotted
pairs, where the \type{car} of a truth value is simply an element of
the \atms\ bilattice and the \type{cdr} is the probability that the
given assumption is valid.

When entering an assumption, the user will be queried for the
probability in question; the probabilities of compound sentences are
then computed automatically by the system.  The algorithm used is
that described in Grnarov, Kleinrock and Gerla in the 1979
Proceedings of the Computer Networking Symposium.

The probabilistic \atms\ bilattice is identified by the character
{\bf p}.
  
\subsection{Default reasoning} \index{default reasoning}
\label{s.j}

The default bilattice consists of seven truth values.  In addition to
the usual elements \type{true}, \type{false}, \type{unknown} and
\type{bottom}, there are elements \type{dt}, \type{df} and \type{*}.

The truth value \type{dt} labels a conclusion that is true by default
(such as Tweety flying in the canonical example), while \type{df} labels
a conclusion that is {\em false\/} by default.  The final value,
\type{*}, labels conclusions that are both true {\em and\/} false
by default, as in the well-known Nixon-Republican-Quaker example.

This bilattice is identified by the character {\bf d}.

\subsubsection{Prioritized default reasoning}

Another version of the default bilattice allows for priorities among
the defaults.  In this bilattice, truth values are of the form
 \[\type{(n . x)}\]
 where $n$ is a nonnegative integer and $x$ is an element of the
first-order bilattice.  Truth values with $n=0$ are displayed as for
the first-order bilattice; with $n=1$ as for the default bilattice.
For $n>1$, we might see \type{D(3)T} or something like it.  This
particular truth value is the same as
 \[\type{(3 . true)}\]
 Truth values with small $n$ override those with larger $n$.
\type{Unknown} in this bilattice is given by the truth value \NIL.
The one-character identifier for the hierarchical default bilattice is
{\bf h}.

\subsection{Circumscription} \index{circumscription}
\label{s.circ}

The default bilattices are not able to deal with all of the
situations in which multiple extensions arise because single truth
values such as \type{*} label sentences that are true in some
extensions and false in others without providing any information
about precisely which extensions are which.  In order to handle this
difficulty, we need to work with a more complex bilattice.

One way to view the default bilattice is as a copy of the first-order
bilattice in which the point \type{unknown} is interpreted as, ``not
known with certainty,'' and is in fact replaced with another copy of
the first-order bilattice that is used to record uncertain
information (i.e., default truth, default falsity, or the presence of
a default contradiction).

The circumscriptive bilattice is formed similarly, except that
instead of simply recording default truth, falsity, and so on, we
record a {\em justification\/} for default conclusions.  This is done
by replacing the point \type{unknown} of the first-order bilattice
not with another copy of the first-order bilattice, but with a copy
of the first-order \atms\ bilattice.  (This idea is in fact a general
way to construct new bilattices from old ones, and is called a {\em
surgery}.  It is discussed in Section \ref{s.surgery}.)  The
designator for the circumscriptive bilattice is {\bf c}.
\type{Circum} and \type{circums} use this bilattice to accept a
sentence as an argument and attempt to prove it from the
circumscription axiom; they resemble \type{bc} and \type{bcs}.

The truth values in the circumscription bilattice corresponding to
certain information are of the form \type{(T . x)}, where $x$ is an
element of the first-order bilattice.  Those corresponding to uncertain
information are of the form \type{(NIL . y)}, where $y$ is an element
of the first-order \atms\ bilattice.

Finally, the circumscriptive theorem prover will output diagnostics
concerning its progress if the global variable \type{*circum-trace*}
has a non-\NIL\ value.

Here is a summary of the circumscriptive facilities:
 \begin{description} 
 \item [\itype{*circum-trace*}] If not \NIL, the circumscriptive proof
functions will output diagnostics describing their progress.
 \item [\itype{circum(p)}] Like \type{bc}, but looks for instances of
$p$ that follow from the circumscription axiom.  Allowed keywords:
\type{cutoffs} and \type{succeeds}.
 \item [\itype{circums(p)}] Like \type{bcs}, but looks for instances of
$p$ that follow from the circumscription axiom.  Allowed keywords:
\type{cutoffs} and \type{succeeds}.
 \end{description}

\subsection{Temporal reasoning}\index{temporal reasoning}

\type{*Time-bilattice*} is a bilattice of the truth values that might
be assigned to propositions that have a time-dependent nature, such as
whether or not a particular switch is on.  The truth values represent
the entire histories of the propositions, so that a single truth value
is a complete function from a set of time points (taken to be the
integers) into a set of ``base'' truth values (take to be elements of
the nonmonotonic {\sc atms} bilattice). This bilattice is discussed at
greater length in the paper, ``Computational considerations in
reasoning about action,'' and a sample use of this bilattice can be
found by examining the file \type{temporal.test} that is included with
\mvl.

\section{Facilities to create new bilattices}
\label{s.construct}

There are a variety of ways in which it is possible to construct
new bilattices from old ones, or from old ones in conjunction
with instances of some other structure.  In this section, we
describe the extent to which these facilities exist within the
\mvl\ system.

\subsection{Products}
\label{s.product}

Perhaps the simplest way in which a new bilattice can be constructed
is as the product of two existing bilattices.  If $A$ and $B$ are
bilattices, then a truth value in the product bilattice $A\times B$
will simply be a pair of truth values \type{(x . y)}, where $x$ is a
truth value in $A$ and $y$ is one in $B$.  All of the bilattice
operations on the product bilattice are computed pointwise using the
corresponding operations on the component bilattices $A$ and $B$.

The product construction is performed using the following function:
 \begin{description}
 \item [\itype{product-bilattice(b1 b2)}] Returns the bilattice
that is the product of the bilattices $b_1$ and $b_2$.
 \end{description}

\subsection{Functional attachments}

Another way in which to construct a bilattice is to start with a
bilattice $B$ and a function $f$ that maps $B$ into a fixed set $S$.
The elements of the new bilattice are simply pairs of the form
 \[\type{(x .~(f x))}\]
 where $x$ is an element of $B$.  The function $f$ here is simply
tagging along for the ride, and has no impact on the bilattice
operations.  It can be used to provide additional information about
the truth values, and is used to compute probabilistic information in
the bilattice {\bf p} described in Section \ref{s.prob}.

The relevant function is the following one:
 \begin{description}
 \item [\itype{append-bilattice(b f)}] Returns the bilattice obtained
by appending the result of applying the function $f$ to the points of
the original bilattice $b$.
 \end{description}

\subsection{Surgeries}
\label{s.surgery}

A somewhat more interesting construction is the one we described in
Section \ref{s.circ} as a {\em surgery}.  Here, we take one bilattice
$b_1$ and replace the point corresponding to \type{unknown} in this
bilattice with another bilattice $b_2$.  The truth values in the
resulting bilattice are either of the form \type{(T . x)} for a point
$x$ in the ``upper'' bilattice $b_1$, or \type{(NIL . y)} for a point
$y$ in the lower bilattice.  Truth values from $b_1$ take priority over
values from $b_2$.

To perform this construction generally, we have: 
 \begin{description}
 \item [\itype{splice-bilattice(upper lower)}] Returns the bilattice
obtained by replacing the point \type{unknown} in the bilattice
\type{upper} with a copy of the bilattice \type{lower}.
 \end{description}

As an example, the default bilattice {\bf d} of Section \ref{s.j} is
given by
 \[\type{(splice-bilattice *first-order-bilattice* *first-order-bilattice*)}\]
 The circumscriptive bilattice of Section \ref{s.circ} is given by
 \[\type{(splice-bilattice *first-order-bilattice*
*first-order-atms-bilattice*)}\]

\subsection{Hierarchies}

It is possible to iterate the surgery construction of the previous
section, producing a bilattice that is an infinite descending
``chain'' of copies of a basic bilattice $b$.  This is done by
the function \type{hierarchy-bilattice}:
 \begin{description}
 \item [\itype{hierarchy-bilattice(b)}] Returns the bilattice constructed
as an infinite descending chain of copies of $b$.
 \end{description}

Truth values in this bilattice have the form \type{(n . x)}, where
$x$ is an element of $b$, and $n$ is an integer giving the hierarchy
level.  The top hierarchy level of the bilattice has $n=0$; the
second level is $n=1$, and so on.  The point \type{unknown} is given
by \NIL.

This function is used to construct the hierarchical version of the
default bilattice.

\subsection{Lattices}

Bilattices can also be constructed from lattices.  A lattice is a set
equipped with a {\em single\/} partial order and greatest lower bound
and least upper bound operations, which we denote by $+$ and $\cdot$
respectively.

An element of the associated bilattice is a {\em pair\/} of elements
of the lattice, and the bilattice operations are given by:
 \begin{eqnarray*}
\neg(a,b) & = & (b,a) \\
(a,b) + (c,d) & = & (a+c,b+d) \\
(a,b) \cdot (c,d) & = & (a\cdot c,b\cdot d) \\
(a,b) \vee (c,d) & = & (a+c, b\cdot d) \\
(a,b) \wedge (c,d) & = & (a\cdot c,b+d)
 \end{eqnarray*}

Roughly speaking, the first component of the bilattice element labels
the ``truth'' of the sentence, and the second labels its falsity.  This
is how the \atms\ bilattices are constructed.

The function used is \type{lattice-to-bilattice}:

 \begin{description}
 \item [\itype{lattice-to-bilattice(l)}] Returns the bilattice constructed
from the lattice $l$.
 \end{description}

A lattice is a structure rather like a bilattice, but with fewer slots
(because it only has a single partial order).  The fields in an instance
of the \itype{lattice} structure are:

\begin{description}
 \item [\type{max}] A \lisp\ object that is the maximal element of
the lattice.
 \item [\type{min}] A \lisp\ object that is the minimal element of
the lattice.
 \item [\type{eq}] A \lisp\ function that takes two lattice elements
as arguments and returns \T\ if they are equal and \NIL\ if they are
not.
 \item [\type{and}] A \lisp\ function that takes two lattice elements
as arguments and returns their greatest lower bound.
 \item [\type{or}] A \lisp\ function that takes two lattice elements
as arguments and returns their least upper bound.
 \item [\type{le}] This optional field is a \lisp\ function that
accepts two arguments $x$ and $y$ and determines whether or not
$x\leq y$ under the partial order corresponding to the lattice.
 \item [\type{vars}] This \lisp\ function accepts a lattice element
as an argument and returns a list of the variables that it contains.
The default value returns \NIL\ independent of its argument.
 \item [\type{plug}] This \lisp\ function accepts as arguments a
lattice value and a binding list, and returns a potentially modified
lattice value.  The default value is a function that leaves the lattice
value unchanged.
 \item [\type{simplicity}] The lattice analog of the bilattice
\type{simplicity} slot.  When constructing a bilattice from a
lattice, the \type{simplicity} of a truth value $(a,b)$ is the result
of applying the simplicity function from the lattice to the lattice
element $a$.  The value $b$ is ignored in the simplicity computation.
 \item [\type{dws}] This \lisp\ function, ``dot with star,'' accepts
as arguments two lattice values, $v_1$ and $v_2$.  The lattice
element returned is $v_1$, modified so as to remove any contribution
that could be attributed to the lattice element $v_2$.
 \item [\type{stash-val}] The lattice element entering into
the bilattice function \type{stash-val}.
 \item [\type{long-desc}] A long description of the lattice, used
for input and output only.
 \item [\type{short-desc}] A short description of the lattice, also
used for input and output only.
 \item [\type{char}] A one-character identifier for the lattice.
 \end{description}

\subsection{Directed acyclic graphs}
\label{s.dag}

The final fashion in which \mvl\ supports the construction of new
bilattices is an extension of the product construction of Section
\ref{s.product}.  There, we formed the product of two bilattices
$A\times B$.  If the bilattices were identical, we would have formed
the product bilattice $B\times B$, or $B^2$.  Another way to view
this bilattice is as the collection of functions from the two point
set $\set{1,2}$ into the original bilattice $B$.  Since such a
function $f$ can be described as the pair $(f(1),f(2))$, there is no
difference between this description and the previous one.

Given this analogy, it is not hard to see that we can extend the
construction to the bilattice $B^S$ for an arbitrary set $S$, where
elements of this bilattice are functions from the set $S$ into $B$,
and all of the bilattice functions ($\wedge$, $\cdot$ and so on) are
once again computed pointwise.

The difficulty with this idea is that for large sets $S$, it is not
clear how we can represent a function from $S$ into $B$ in a
convenient way.  It turns out, however, that if $S$ has the structure
of a directed acyclic graph, or \dg, this problem can be overcome.

\subsubsection{Defining the DAG}

As a preliminary, we need some way to define the \dg\ itself.  In
\mvl, a \dg\ \index{dag} is an instance of the \type{dag} structure;
this structure has the following fields:
 \begin{description}
 \item [\itype{root}] The root of the \dg.
 \item [\itype{eq}] A function that accepts two elements of the \dg\
and returns \T\ if they are the same element and \NIL\ otherwise.
 \item [\itype{leq}] A function that accepts two elements of the \dg\
and returns \T\ if the first is farther from the root of the \dg\
than the second.  If not supplied by the user, this function can by
computed from the \type{dot} operation on the \dg.
 \item [\itype{dot}] A function that accepts two elements of the \dg\
and returns a list of their first common descendants, or \NIL\ if no
such descendant exists.
 \item [\itype{long}] A long description of the \dg.
 \item [\itype{short}] A short description of the \dg.
 \item [\itype{inherit}] A function describing the fashion in which
truth values are inherited within the \dg.  This function $f$ should take
four arguments: a bilattice $b$, two elements of the \dg\ $x$ and $y$,
and an element $z$ of the bilattice.  If $z$ is the element of the
bilattice associated to the \dg\ element $x$, then $f(b,x,y,z)$ will
be the value associated to the \dg\ element $y$ in the absence of
information to the contrary.  By default, this function ignores
its first three arguments and assigns $y$ the same value as $x$ has.
 \item [\type{vars}] This \lisp\ function accepts a \dg\ element
as an argument and returns a list of the variables that it contains.
The default value returns \NIL\ independent of its argument.
 \item [\type{plug}] This \lisp\ function accepts as arguments a
\dg\ value and a binding list, and returns a potentially modified
\dg\ value.  The default value is a function that leaves the \dg\
value unchanged.
 \item [\type{select}] In some applications (none of which is distributed
with \mvl), it appears that truth values should be inherited from only
{\em some\/} of the elements of the \dg\ which might potentially
impact the value of the function at some specific point.  The
selection function is given a point where the truth value is being
computed and a list of points that apparently contribute; it returns a
list of points that {\em actually\/} contribute.  By default, all
points contribute; most users of \mvl\ are very unlikely to need this
facility.
 \end{description}

\subsubsection{Constructing the bilattice}

Having selected a \dg\ $D$ and a bilattice $B$, we can now construct
the bilattice $B^D$.  Specific functions from $D$ to $B$ are described
as structures with the following slots:
 \begin{description}
 \item [\type{dag}] The domain of the function, $D$.
 \item [\type{bilattice}] The range of the function, $B$.
 \item [\type{list}] A list of elements of the form \type{(p . v)},
where \type p is a point where the value changes in a fashion {\em
other\/} than the one predicted by the inheritance function associated
with $D$ and \type v is the value taken there.
 \end{description}
 We will refer to these functions as {\em \dg-functions}.

There are a variety of functions available to manipulate \dg-functions,
and these are used to construct the bilattice $B^D$.  In some cases,
however, the user may want to access these functions directly.  Here
are some of them:
 \begin{description}
 \item [\itype{get-val(d f)}] Returns the value of the
function $f$ at the point $d$ of the \dg.
 \item [\itype{make-tv(dag bilattice d b)}] Returns a \dg-function that
takes the value $b$ at the point $d$.  If $d$ is not the root of the
\dg, the value at the root of the \dg\ is generally the point
\type{unknown} in the bilattice, but the user may supply the keyword
\type{root-fn} if he wishes, in which case the value at the root of
the \dg\ is the result of applying \type{root-fn} to the bilattice.
 \item [\itype{make-root(dag bilattice b)}] Returns a \dg-function that
takes the value $b$ at the root of the \dg.
 \item [\itype{dag-change-tv(f b-f)}] The argument
\type{b-f} should be a function from the bilattice to itself; a
modified version of the original \dg-function $f$ is returned in which
every point in the \dg-function has had the function \type{b-f}
applied to it.  There is an optional keyword \type{simplify}; if
supplied, the computed function is simplified before being returned.
(See \type{simplify}.)
 \item [\itype{dag-change-dag(f d-f)}] The argument
\type{d-f} should be a function from the \dg\ to itself; a modified
version of the original \dg-function $f$ is returned in which every
point at which the \dg-function takes a new value has had the function
\type{d-f} applied to it.  The bilattice argument is expected to be
the bilattice of functional truth values as opposed to the bilattice
that is the target of the functions.  There is an optional keyword
\type{simplify}; if supplied, the computed function is simplified
before being returned.  (See \type{simplify}.)
 \item [\itype{dag-list-of-answers(f)}] Given a \dg-function $f$, assumed
to be from the \dg\ of binding lists into the currently active bilattice, this
function constructs a list of answers corresponding to the points where
the \dg-function takes values.
 \item [\itype{dag-accumulate-fn(dag bilattice list)}] Given a \dg, a
bilattice, and a list of elements of the form
 \[\type{(dag-elt . bilattice-elt)}\]
a \dg-function is constructed that takes the given values at the
given elements of the \dg.

A keyword to \type{dag-accumulate-fn} is \type{modifier}.  This
function is used if the \type{d-list} contains the same \dg\ element
more than once, in which case the modifier function is called with
the new and old values (in that order) and the result is used as the
value of the \dg-function at that particular point.  The default
action is to simply overwrite the old value with the new one.  (In
other words, the default modifier function simply returns its first
argument while ignoring the second.)
 \item [\itype{dag-prune(f pred)}] Given a \dg-function
$f$ and a predicate on the bilattice, this function prunes from $f$
all points $d$ where $f(d)$ fails to satisfy the predicate.  There is
an optional keyword \type{modify-fn} that, if supplied, should be a
function $m$ from the bilattice to itself.  In this case, the points
where $m(f(d))$ fails to satisfy the predicate are pruned.
 \item [\itype{combine-fns(f1 f2 fn comp)}]
This function takes the two \dg-functions $f_1$ and $f_2$ and
combines them using the function \type{fn}, which should be a
function that takes two elements of the bilattice and returns a
single one.  The comparison function \type{comp} is used to determine
if one of the two functions can be returned directly; if, for example
\type{(comp f1 f2)} holds, than  we simply return $f_1$, and similarly
for $f_2$.  As an example, if the combining function is the bilattice
operation $\cdot$, the comparison function is \type{k-le}.
 \item [\itype{combine-1(f params f1 ... fn)}] This is an
$n$-ary version of \type{combine-fns}, where $f$ is an $n$-ary
bilattice function and there are $n$ \dg-functions supplied as well.
The additional argument \type{params} indicates which of the $n$ arguments
to this function are not necessarily \dg-functions, but are instead parameters
expected by the overall function $f$.  This additional argument is a list
of \T\ (if the $n$th argument is a parameter) or \NIL\ (if it isn't).
If the parameter list ends early, all of the remaining entries are
assumed to be \NIL.
 \item [\itype{add-dag(d b f)}] Given the \dg-function
$f$, add the value $b$ at the point $d$.  This function works by side
effect only.  There is a keyword \type{modifier} that, if supplied,
should be a function that takes two elements of the bilattice and
returns a single one.  Given such a function $m$, the \dg-function is
changed so that the new value at $d$ is $m(b,f(d))$, where $f(d)$ is
the old value.  (This is similar to the same keyword in
\type{dag-accumulate-fn}.)  The keyword \type{simplify} is also available.
 \item [\itype{simplify(f)}] In some cases,
manipulating the \dg-function $f$ may cause it to contain
redundancies, where the value at some point is presented explicitly
even though this value is in fact the one that would normally be
inherited from its parents.  This function returns a modification of
$f$ where these irrelevant entries have been removed.
 \item [\itype{fn-and(f1 f2)}] This function conjoins
the values of $f_1$ and $f_2$ using the bilattice $\wedge$ operation.
The new \dg-function is returned.
 \item [\itype{fn-or(f1 f2)}] This function disjoins
the values of $f_1$ and $f_2$ using the bilattice $\vee$ operation.
The new \dg-function is returned.
 \item [\itype{fn-plus(f1 f2)}] This function adds
the values of $f_1$ and $f_2$ using the bilattice $+$ operation.
The new \dg-function is returned.
 \item [\itype{fn-dot(f1 f2)}] This function combines
the values of $f_1$ and $f_2$ using the bilattice $\cdot$ operation.
The new \dg-function is returned.
 \item [\itype{fn-dws(f1 f2)}] This function combines
the values of $f_1$ and $f_2$ using the bilattice operation
\type{dot-with-star}.  The new \dg-function is returned.
 \item [\itype{fn-not(f1)}] This function negates the
values of $f_1$ using the bilattice $\neg$ operation.  The new
\dg-function is returned.
 \item [\itype{fn-eq(f1 f2)}] This function compares
the values of $f_1$ and $f_2$ using the bilattice comparison operation.
Either \T\ or \NIL\ is returned.
 \item [\itype{fn-comp(f1 f2 fn)}] This functions compares
the values of $f_1$ and $f_2$ using the bilattice comparison function
\type{fn}.  Either \T\ or \NIL\ is returned.
 \end{description}

Of course, the point of all of this is to enable us to define the
following function:
 \begin{description}
 \item [\itype{dag-bilattice(dag bilattice)}] Given a bilattice $B$
and a dag $D$, returns the bilattice $B^D$.
 \end{description}

The following functions are also available:
 \begin{description}
 \item [\itype{lattice-to-dag(lattice)}] Construct a \dg\ from the
given lattice.
 \item [\itype{lattice-to-dag-to-bilattice(lattice bilattice)}]  First
construct a \dg\ $D$ from the given lattice.  Then use the bilattice
$B$ to construct the bilattice $B^D$.
 \end{description}

\subsubsection{The DAG of theories}
\label{s.theory}

Two specific \dg s are provided for use in the \mvl\ system.  The first
uses a \dg\ corresponding to the existence of multiple theories within
a single database; the second uses a \dg\ to handle binding information.

The theory \dg\ consists of a hierarchy of theories.  In a single
database, for example, we might want to have rules that applied to
celestial bodies, rules applying to stars, rules applying only to red
giants or only to our sun, and so on.  The truth value of a rule
applicable to all celestial bodies should be inherited if we are
considering our sun (but it should be possible to overwrite it with
a different truth value if we wish), but a rule applicable to red
giants should not be inherited in this fashion.  So we see that our
\dg\ construction is well-suited for this sort of an application.

To activate the theory mechanism, simply invoke the \lisp\ function
\type{activate-theory-mechanism}.  This function replaces the active
bilattice, say $B$, with the bilattice constructed of functions from
the theory \dg\ into $B$.  The original bilattice can still be
accessed, since it is the second \type{component} of the newly formed
bilattice.  The function \type{theory-active-p} can be used to
determine whether or not the theory mechanism has been activated.

The \dg\ of theories is the value of the global variable
\type{theory-dag}, and has as its root the theory whose name is the
value of the constant \type{*global*}.  To indicate that a theory
\type{big} includes another theory \type{little} as a subtheory,
invoke the function 
 \[\type{(includes big little)}\]
 Thus, for example, one would introduce a new theory called
\type{theory} by writing
 \[\type{(includes *global* 'theory)}\]
 The effects of \type{includes} can be reversed by invoking the
function \type{unincludes}.

Note that since the theory \dg\ is used by the functions that
manipulate the bilattice $B^D$, it is important that this \dg\ be
maintained as new theories are introduced.  If \type{includes} is not
used to ``locate'' these new theories in the overall theory \dg\
itself, the theories will be inaccessible to many of the \mvl\
functions.

At any point in time, any new fact will be inserted into the database
as true in a particular theory; the theory into which it is inserted
is the value of \type{*active-theory*}.  The value of this variable
should not be changed directly, but should instead be modified by
invoking the function \type{activate} on the theory into which
subsequent facts should be written.  The variable
\type{true-in-theory} contains the truth value of a sentence that is
true in the currently active writing theory.

The theory \type{*active-theory*} is also the theory that is active
for reading at any point in time.  Of course, any parents of the
active theory are also active by virtue of the structure of the theory
\dg\ itself; if it is desired to read from multiple theories
simultaneously, a new theory should be created that is a child of all
of them.

Here is a summary of the functions described so far:
 \begin{description}
 \item [\itype{activate-theory-mechanism}] Activates the theory
mechanism.  Accepts as an optional argument either a one-character
bilattice descriptor, or a bilattice.
 \item [\itype{theory-active-p}] Returns \T\ or \NIL\ depending on
whether or not the theory mechanism has been activated.
 \item [\itype{theory-dag}] The \dg\ of theories.
 \item [\itype{*global*}] The root of the \dg\ of theories.
 \item [\itype{includes(big little)}] This function is used to state that
the theory \type{little} is a subtheory of the theory \type{big}, and
should inherit from it.
 \item [\itype{unincludes(big little)}] This function is used to state
that the theory \type{little} is no longer a subtheory of the theory
\type{big}, and should not inherit from it.
 \item [\itype{*active-theory*}] The theory that is currently active
for reading and writing.
 \item [\itype{activate(th)}] Activate the theory \type{th} for reading
and writing.
 \item [\itype{true-in-theory}]  The truth value of a sentence that is
true in the currently-active writing theory.
 \end{description}

Finally, five functions have been added to simplify the user
interface somewhat.  The function {\tt theory-\-contents} displays
the contents of a given theory, and \type{empty-theory} removes all
of the propositions in this theory.  To modify the contents of a
specific theory, the following functions can be used:
 \begin{description}
 \item [\itype{state-in-theory(p)}] Adds the proposition $p$ to the
currently active theory.  Allowed keywords: \type{theory} can
be used to specify a different theory to which $p$ should be added;
\type{value} and \type{increment-p} have the same meaning as for
\type{state}.  (See Section \ref{s.database}; if these keywords are
not used, the function \type{state} can be used equally well, since
the default stash value with the theory mechanism activated also
inserts the proposition into the currently active theory.)
 \item [\itype{unstate-from-theory(p)}] Removes the proposition $p$
from the currently active theory.  Allowed keywords:
\type{theory} can be used to specify a different theory from which
$p$ should be removed; \type{value} and \type{decrement-p} have the
same meaning as for \type{unstate}.  (See Section \ref{s.database}.)
 \item [\itype{erase-from-theory(p)}] Removes the proposition $p$
from the currently active theory, and forward chains on the result.
Allowed keywords: \type{theory} can be used to specify the theory from
which $p$ should be removed; \type{value}, \type{decrement-p} and
\type{cutoffs} have the same meaning as for \type{erase}.  (See
Section \ref{s.fc}.)
 \item [\itype{theory-contents()}] Displays the contents of the
currently active theory, which can be changed using the
keyword \type{theory}.  The keywords \type{cutoffs} and \type{props}
are allowed and are treated as for \type{contents}.  (See Section
\ref{s.contents}.)  There is an additional keyword,
\itype{this-theory-only}, which defaults to \T\ and, if supplied,
indicates that only sentences whose truth values {\em change\/} in the
given theory are to be displayed.  (This avoids displaying all
sentences in the global theory, for example.)
 \item [\itype{empty-theory()}] Empties the currently active
theory, although this theory may be changed using the keyword
\type{theory}.  The keyword \type{value} has a meaning similar to its
meaning as a keyword to \type{erase}.
 \end{description}

\subsubsection{The DAG of binding lists}
\label{s.binding-dag}

The other \dg\ supplied with the \mvl\ system is used for binding
lists since these, too, admit a natural inheritance structure.  The
root of the \dg\ is the empty binding list \NIL, and inheritance in
the \dg\ is determined by subsumption.  The bilattice constructed as
functions from this \dg\ into \type{*active-bilattice*} is used
extensively by the prover, and is available to the user as
\type{bdg-to-truth-val}.

The binding \dg\ itself is known as \type{binding-dag}, and is constructed
using the following functions:
 \begin{description}
 \item [\itype{equal-binding(b1 b2)}] Checks to see if two binding
lists are effectively the same.
 \item [\itype{append-binding-lists(b1 b2)}] Computes the least
specific binding list that is more specific than both $b_1$ and $b_2$
or returns \NIL\ if no such binding list exists ($b_1$ and $b_2$ may
bind the same variable to different constants, for example).
 \item [\itype{binding-le(b1 b2)}] Returns \NIL\ unless $b_1$ is a more
specific binding that $b_2$, in which case a list of relative binding
lists is returned.
 \end{description}

The following global variables are also used:
 \begin{description}
 \item [\itype{binding-dag}] The \dg\ of binding lists.
 \item [\itype{bdg-to-truth-val}] The bilattice of functions from
the \dg\ of binding lists into \type{*active-bilattice*}.
 \end{description}

\section{Modal operators}
\label{s.modal}

An additional feature that makes the \mvl\ system unique is its
ability to handle {\em modal operators}.  The approach taken by \mvl\
to modal operators is that they are functions from the bilattice of
truth values to itself.  Thus, for example, the modal operator that
takes a sentence $p$ into the sentence, ``I know that $p$,''
corresponds to the bilattice function (on the first-order bilattice)
that takes $t$ and $\bot$ into $t$ (since we know $p$ if we know it
to be true or know it to be true and false both), and takes $f$ and
$u$ into $f$ (since we do {\em not\/} know $p$ if we either know it
to be false or do not know if it is true or false).  Further details
can be found in the paper, ``Bilattices and modal operators.''

A modal expression in \mvl\ is any sentence of the form 
 \[\type{(op p-1 \dots\ p-n)}\]
 where \type{op} is recognizable as a modal operator because it has a
non-\NIL\ value for its \type{modal-op} property in the currently
active invocation of \mvl.  In this case, the value of the
\type{modal-op} property should be an instance of the structure
\itype{modal-operator}.  This structure has the following fields:
 \begin{description}
 \item [\itype{name}] The name of the modal operator.  This is used
to identify modal operators in different bilattices that have the
same commonsense meaning (and can therefore be usefully combined when
the bilattices are combined as in the previous section).
 \item [\itype{function}] A function that accepts as arguments the
truth values of the sentences under the scope of the modal operator
and returns the truth value of the modal expression itself.
 \item [\itype{parameters}] If the modal operator accepts $n$
arguments, this should be a list of $n$ entries indicating whether
each argument is a truth value or a parameter that should be passed
directly to the modal operator function.  (An entry of \T\ indicates
that the associated argument is parametric; an entry of \NIL\
indicates that it isn't.)  Suppose, for example, that we are working
with a bilattice where the truth values are functions describing the
behaviors of various fluents over the course of time, and have a modal
operator \type{delay}, where the truth value assigned to
 \[\type{(delay t p)}\]
 is expected to be the same as the truth value assigned to $p$ but
delayed by an amount of time $t$.  Here, the first argument ($t$) is
a parameter and the second argument ($p$) is not.

If there are fewer than $n$ entries on the parameter list, then all of
the missing entries are assumed to be \NIL\ (so that the associated
arguments to the modal function are not parameters).  Note: All
parametric arguments should be bound when attempting to derive a modal
expression.  This is similar to the requirement that arguments to a
procedurally-attached predicate be bound when the procedure is called.
 \end{description}

\subsection{Modal operators supplied with the MVL bilattices}

All bilattices in \mvl\ are equipped with a modal operator \itype{idn}
that does not change the truth value of its argument at all; this
modal operator can be used to semantically mark points where it may be
desirable to interrupt the inference process.  In addition, if the
prover is invoked on \type{(idn p)} instead of on \type p, attempts
will be made to prove both $p$ and its negation $\neg p$.

The bilattices supplied as part of the \mvl\ system are equipped with
the following modal operators:

 \begin{description}
 \item [\type{*first-order-bilattice*}] This bilattice has two
modal operators: $L$, which maps $p$ into, ``necessarily $p$,'' or,
``I know that $p$,'' and $M$, which maps $p$ into ``possibly $p$.''
 \item [\type{*atms-bilattice*}]  No modal operators are supplied with
this or any of the other \atms-type bilattices.
 \item [\type{*default-bilattice*}] This bilattice and the hierarchical
default bilattice both inherit the modal operators $L$ and $M$ from the
first-order bilattice.
 \item [\type{*circumscription-bilattice*}] This bilattice does not
come equipped with any modal operators.
 \item [\type{*time-bilattice*}] This bilattice inherits the modal
operators $L$ and $M$ from the bilattice of basic values.  The
following three additional modal operators are also defined:
 \begin{enumerate}
 \item \itype{delay(t p)} This is a modal operator that pushes the
truth value of $p$ into the future by an amount of time $t$.  Thus if
$p$ is true at time 2, $\type{delay}(p,2)$ will be true at time 4.
 \item \itype{propagate} is a modal operator that corresponds to the frame
axiom.  If $p$ is true at time 2, unknown between times 3 and 20, and false
at time 21, then $\type{propagate}(p)$ will be true from times 2 until 20
and false subsequently.
 \item \itype{at(t p)} is a modal operator that returns the value taken
by the truth value $p$ at the time $t$.  (In actuality, a function has
to be returned; the function is the one that takes value $p(t)$ at all
times instead of just at $t$.)
 \end{enumerate}
 \end{description}

\subsection{Quantification and modal operators}

In addition to \type{idn}, there are two other modal operators that
are available for use with any bilattice, \itype{exists} and
\itype{forall}.  Although not really modal operators on the bilattice
being used, these {\em are\/} legitimate modal operators on the bilattice
\type{bdg-to-truth-val} that was described in Section \ref{s.binding-dag}.

The functions associated to these modal operators mimic the usual behaviors
given to the quantifiers.  Thus an expression such as
 \begin{equation}
\type{(exists ?x (foo ?x))}
\label{e.exists}
\end{equation}
 is evaluated by first computing the truth value that could be assigned
to \type{(foo ?x)} as an element of \type{bdg-to-truth-val}, so that
a truth value is obtained for every binding of \type{?x}.  This truth
value is then manipulated so that the values for any particular binding
of \type{?x} are disjoined (the usual interpretation of the existential
quantifier), and the result is returned as the truth value assigned to
\ref{e.exists}.  The modal operator \type{forall} is handled similarly
and corresponds to universal quantification.  The first (parametric)
argument to these operators can also be a list of variables, in which
case all of the variables are quantified over.

\subsection{Modal operators created automatically}

Finally, many of the functions that construct new bilattices are able
to construct modal operators on these bilattice in terms of the modal
operators on the bilattices they are using in the construction.  More
specifically:

\begin{enumerate}
 \item A product bilattice $A \times B$ will inherit a modal operator
from $A$ and $B$ if these component bilattices have modal operators
with the same name.
 \item A bilattice constructed using \type{append-bilattice} will have
modal operators corresponding to whatever modal operators were available
on the bilattice used in the construction.
 \item A bilattice constructed by splicing two bilattices $A$ and $B$
together will inherit the modal operators from the ``upper'' bilattice.
 \item A bilattice constructed using \type{hierarchy-bilattice} will
have modal operators corresponding to whatever modal operators were
available on the bilattice used in the construction.
 \item A bilattice constructed from a lattice will not normally have
any modal operators associated to it.
 \item The bilattice $B^D$ will inherit whatever modal operators the
bilattice $B$ has.
 \end{enumerate}

\section*{Acknowledgement}

Support for the development of \mvl\ has been provided by the Air
Force Office of Scientific Research, the Defense Advance Research
Projects Agency, General Dynamics, the National Science Foundation,
Rockwell International, and Rome Labs.

\clearpage

\appendix

\section{Summary of MVL functions and constants}

\begin{list}{}{\setlength{\labelwidth}{1.5in}\setlength{\labelsep}{.25in}
\setlength{\leftmargin}{1.75in}}

% *A*

\item [\itype{activate(th)}] Activate the theory \type{th} for reading
and writing.

\item [\itype{activate-best-first}] Activate best-first search by
setting \type{*node-evaluation-fn*} to \type{best-first-fn}.  Allowed
keywords: \type{initial-limit}, \type{limit-delta}, and
\type{final-limit}.

\item [\itype{activate-breadth-first}] Activate breadth-first search
by setting \type{*node-evaluation-fn*} to \type{depth-fn}.  Allowed
keywords: \type{initial-limit}, \type{limit-delta}, and
\type{final-limit}.

\item [\itype{activate-depth-first}] Activate depth-first search by
setting \type{*initial-limit*} to \NIL.  Allowed keyword:
\type{final-limit}.

\item [\itype{activate-theory-mechanism}] Activates the theory
mechanism.  Accepts as an optional argument either a one-character
bilattice descriptor, or a bilattice.

\item [\itype{*active-bilattice*}] The currently active bilattice.

\item [\itype{*active-theory*}] The theory that is currently active
for reading and writing.

\item [\itype{add-dag(d b f)}] Given the \dg-function $f$, add the
value $b$ at the point $d$.  This function works by side effect only.
There is a keyword \type{modifier} that, if supplied, should be a
function that takes two elements of the bilattice and returns a single
one.  Given such a function $m$, the \dg-function is changed so that
the new value at $d$ is $m(b,f(d))$, where $f(d)$ is the old value.
(This is similar to the same keyword in \type{dag-accumulate-fn}.)
The keyword \type{simplify} is also available.

\item[\itype{answer-binding(a)}] Returns the binding field of the
answer \type a.

\item[\itype{answer-value(a)}] Returns the value field of the answer
\type a.

\item [\itype{anytime-trace(pat)}] Include \type{pat} on \at, so that
any instance of \type{pat} will be traced.  The argument is optional;
if omitted, \type{(?*)} (of which everything is an instance) will be
used.

\item [\itype{anytime-untrace(pat)}] Remove any instance of
\type{pat} from \at.  The argument is optional; if omitted,
\type{(?*)} (of which everything is an instance) will be used.

\item [\itype{append-bilattice(b f)}] Returns the bilattice obtained
by appending the result of applying the function $f$ to the points of
the original bilattice $b$.

\item [\itype{append-binding-lists(b1 b2)}] Computes the least
specific binding list that is more specific than both $b_1$ and $b_2$
or returns \NIL\ if no such binding list exists ($b_1$ and $b_2$ may
bind the same variable to different constants, for example).

\item [\itype{*atms-bilattice*}] The bilattice corresponding to an \atms,
where the justification information is maintained in disjunctive normal
form as suggested by de Kleer.

% *B*

\item [\itype{(bagof}\type{ ?x ?p ?b)}] This relation holds for
$x$, $p$ and $b$ if and only if $b$ is a list of all $x$ satisfying
$p$.

\item [\itype{bc(q)}] Attempts to prove the query $q$ from information
in the database, looking only for a single answer.

\item [\itype{bc-trace(pat)}] Include \type{pat} on
\type{*bc-trace*}, so that any instance of \type{pat} will be
traced.  The argument is optional; if omitted, \type{(?*)} (of which
everything is an instance) will be used.

\item [\itype{bc-untrace(pat)}] Remove any instance of \type{pat}
from \type{*bc-trace*}.  The argument is optional; if omitted,
\type{(?*)} (of which everything is an instance) will be used.

\item [\itype{bc-unwatch(preds)}] Removes the given predicates from
the list of {\sc WATCH}ed predicates.  If no arguments are supplied,
all predicates are un{\sc watch}ed.

\item [\itype{bc-watch(preds)}] Accepts as arguments any number of
predicates that should be watched as the prover proceeds.  Each
predicate can also be a list of the form \type{(predicate depth)}
where \type{depth} is the number of successive ancestors to be
displayed when the prover encounters an instance of the predicate.
Returns a list of the predicates currently being {\sc WATCH}ed.

\item [\itype{bcs(q)}] Attempts to prove the query $q$ from information
in the database, looking for all possible answers.

\item [\itype{bdg-to-truth-val}] The bilattice of functions from
the \dg\ of binding lists into \type{*active-bilattice*}.

\item [\itype{bilattice(b)}] Add the bilattice $b$ to \type{*bilattices*}.

\item [\itype{*bilattices*}] A list of the bilattices known to the system.

\item [\itype{binding-dag}] The \dg\ of binding lists.

\item [\itype{binding-le(b1 b2)}] Returns \NIL\ unless $b_1$ is a more
specific binding that $b_2$, in which case a list of relative binding
lists is returned.

\item [\itype{bottom}] The truth value \type{bottom} in the active bilattice.

\item [\itype{*break-node*}] The number of a node at which the
proof process should be interrupted.

% *C*

\item [\itype{choose-task}]  Selects a task to work on next.  The 
current implementation of \mvl\ simply selects any task that might
have some impact on the final answer being computed.

\item [\itype{circum(p)}] Like \type{bc}, but looks for instances of
$p$ that follow from the circumscription axiom.  Allowed keywords:
\type{cutoffs} and \type{succeeds}.

\item [\itype{*circum-trace*}] If not \NIL, the circumscriptive proof
functions will output diagnostics describing their progress.

\item [\itype{circums(p)}] Like \type{bcs}, but looks for instances of
$p$ that follow from the circumscription axiom.  Allowed keywords:
\type{cutoffs} and \type{succeeds}.

\item [\itype{*circumscription-bilattice*}] The bilattice used to compute
the results of the circumscription axiom.

\item [\itype{cnf(p)}] Returns $p$ in conjunctive normal form, as a list
of conjuncts, where each conjunct is a list of disjuncts.  Thus
 \[\type{(CNF '(AND A B))}\]
returns
 \[\type{((A) (B))}\]

\item [\itype{*cnf-atms-bilattice*}] An alternative version of the \atms\
bilattice where the justification information is maintained in conjunctive
normal form.

\item [\itype{combine-1(f params f1 ... fn)}] This is an $n$-ary
version of \type{combine-fns}, where $f$ is an $n$-ary bilattice
function and there are $n$ \dg-functions supplied as well.  The
additional argument \type{params} indicates which of the $n$ arguments
to this function are not necessarily \dg-functions, but are instead
parameters expected by the overall function $f$.  This additional
argument is a list of \T\ (if the $n$th argument is a parameter) or
\NIL\ (if it isn't).  If the parameter list ends early, all of the
remaining entries are assumed to be \NIL.

\item [\itype{combine-fns(f1 f2 fn comp)}] This function takes the two
\dg-functions $f_1$ and $f_2$ and combines them using the function
\type{fn}, which should be a function that takes two elements of the
bilattice and returns a single one.  The comparison function
\type{comp} is used to determine if one of the two functions can be
returned directly; if, for example \type{(comp f1 f2)} holds, than we
simply return $f_1$, and similarly for $f_2$.  As an example, if the
combining function is the bilattice operation $\cdot$, the comparison
function is \type{k-le}.

\item [\itype{compile-mvl}] Compiles and loads the \mvl\ system.
On a Symbolics system, the keyword \type{reload} forces all of the
files to be reloaded and the keyword \type{recompile} forces them all
to be recompiled.

\item [\itype{*completed-database*}] If \T, find all exceptions
to default proofs.  If \NIL\ (the default), accept defaults proofs
with possible exceptions as complete.

\item [\itype{cont-analysis(analysis)}] This function continues the
investigation of a specific analysis, and can be used to resume a
proof effort that has been curtailed for some reason.  It returns two
values: an instance of the \type{answer} structure giving the results
of the proof attempts, and a modified version of the \type{analysis}
structure that can be used to continue the proof process if desired.

\item [\itype{cont-prover(proof)}] Resume the task associated to the
auxiliary information in the given instance of the \type{proof}
structure.  This function returns an instance of \type{bcanswer}
corresponding to the result obtained by the prover, and returns as
soon as a solution to the original query is found.

\item [\itype{contents}] Displays the contents of the database.  Keywords:
\type{cutoffs}, \type{props}.

\item [\itype{create-mvl-invocation}] This function sets up a new
copy of \mvl.  It accepts an optional argument that is an instance of
\mi; this can be used to reset the state of \mvl\ to a previous one.
If the optional argument is omitted, default values are used.  The
value returned is the \mvl\ invocation that was active before the
function was called, so that another call to
\type{create-mvl-invocation} can be used to restore the original
environment.

\item [\itype{cutoffs-not(cutoffs \&optional bilattice)}] Returns a
new cutoffs $c$ such that
 \[\type{(test x c)} = \type{(test (mvl-not x) cutoffs)}\]
 for all truth values $x$.  In other words, $x$ passes $c$ if and
only if $\neg x$ passed the cutoffs supplied originally.

% *D*

\item [\itype{dag-accumulate-fn(dag bilattice list)}] Given a \dg, a
bilattice, and a list of elements of the form
 \[\type{(dag-elt . bilattice-elt)}\]
a \dg-function is constructed that takes the given values at the
given elements of the \dg.

A keyword to \type{dag-accumulate-fn} is \type{modifier}.  This
function is used if the \type{d-list} contains the same \dg\ element
more than once, in which case the modifier function is called with
the new and old values (in that order) and the result is used as the
value of the \dg-function at that particular point.  The default
action is to simply overwrite the old value with the new one.  (In
other words, the default modifier function simply returns its first
argument while ignoring the second.)

\item [\itype{dag-bilattice(dag bilattice)}] Given a bilattice $B$
and a dag $D$, returns the bilattice $B^D$.

\item [\itype{dag-change-dag(f d-f)}] The argument \type{d-f} should
be a function from the \dg\ to itself; a modified version of the
original \dg-function $f$ is returned in which every point at which
the \dg-function takes a new value has had the function \type{d-f}
applied to it.  The bilattice argument is expected to be the bilattice
of functional truth values as opposed to the bilattice that is the
target of the functions.  There is an optional keyword
\type{simplify}; if supplied, the computed function is simplified
before being returned.  (See \type{simplify}.)

\item [\itype{dag-change-tv(f b-f)}] The argument \type{b-f} should be
a function from the bilattice to itself; a modified version of the
original \dg-function $f$ is returned in which every point in the
\dg-function has had the function \type{b-f} applied to it.  There is
an optional keyword \type{simplify}; if supplied, the computed
function is simplified before being returned.  (See \type{simplify}.)

\item [\itype{dag-list-of-answers(f)}] Given a \dg-function $f$, assumed
to be from the \dg\ of binding lists into the currently active bilattice, this
function constructs a list of answers corresponding to the points where
the \dg-function takes values.

\item [\itype{dag-prune(f pred)}] Given a \dg-function $f$ and a
predicate on the bilattice, this function prunes from $f$ all points
$d$ where $f(d)$ fails to satisfy the predicate.  There is an optional
keyword \type{modify-fn} that, if supplied, should be a function $m$
from the bilattice to itself.  In this case, the points where
$m(f(d))$ fails to satisfy the predicate are pruned.

\item [\itype{*default-bilattice*}] A bilattice that can be used for simple
default reasoning.

\item [\itype{delete-bdg(var bdg-list)}] Destructively modifies
\type{bdg-list}, removing any binding for the variable \type{var}.

\item [\itype{denotes(a)}] Returns the proposition associated with a
particular \lisp\ atom.

\item [\itype{dnf(p)}] Returns $p$ in disjunctive normal form, as a list
of disjuncts, where each disjunct is a list of conjuncts.  Thus
 \[\type{(DNF '(AND A B))}\]
returns
 \[\type{((A B))}\]

\item [\itype{(done exp-1 ... exp-n)}] This predicate, which is
useful for stashing only, causes \lisp\ to evaluate each of the subsequent
expressions.

\item [\itype{dot-with-not(x)}] Computes $x\cdot \neg x$.

\item [\itype{dot-with-star}] The function \type{dws} in the active bilattice.

% *E*

\item [\itype{empty}] Unstashes every sentence in the database.  Keyword:
\type{value}.

\item [\itype{empty-theory()}] Empties the currently active theory,
although this theory may be changed using the optional theory
argument.  The keyword \type{value} has a meaning similar to its
meaning as a keyword to \type{erase}.

\item[\itype{equal-answer(a1 a2)}] Returns \T\ if the two answers
\type{a1} and \type{a2} are equal, and \NIL\ otherwise.  The
functions used are \type{binding-equal} and \type{mvl-eq}, the
functions normally used to test binding lists and truth values for
equality.

\item [\itype{equal-binding(b1 b2)}] Checks to see if two binding lists are
effectively the same.

\item [\itype{*equivalence-translation*}] Controls the fashion in which
\type{normal-form} deals with the connective \type{<=>}.

\item [\itype{erase(p)}] Removes $p$ from the database and forward
chains.  Allowed keywords: \type{value}, \type{decrement-p} and
\type{cutoffs}.

\item [\itype{erase-from-theory(p)}] Removes the proposition $p$ from
the currently active theory, and forward chains on the result.
Allowed keywords: \type{theory} can be used to specify the theory from
which $p$ should be removed; \type{value}, \type{decrement-p} and
\type{cutoffs} have the same meaning as for \type{erase}.  (See
Section \ref{s.fc}.)

\item [\itype{existing-mvl-invocation}] This function returns a
single value that is an instance of the structure \mi.  It is useful
primarily for saving the current copy of \mvl\ so that it can be
restored later.

% *F*

\item [\itype{false}] The truth value \type{false} in the active bilattice.

\item [\itype{fc(p)}] Inserts the sentence $p$ into the database and
forward chains.  Allowed keywords: \type{value}, \type{increment-p},
\type{fc-increment-p}, \type{cutoffs} and \type{include-modal}.

\item [\itype{fc-trace(pat)}] Include \type{pat} on
\type{*fc-trace*}, so that any instance of \type{pat} will be
traced.  The argument is optional; if omitted, \type{(?*)} (of which
everything is an instance) will be used.

\item [\itype{fc-untrace(pat)}] Remove any instance of \type{pat}
from \type{*fc-trace*}.  The argument is optional; if omitted,
\type{(?*)} (of which everything is an instance) will be used.

\item [\itype{*final-limit*}] is the limit point beyond which the
search is abandoned, so that any node that evaluates to a value worse
than this will be pruned.  The default value of this parameter is
\NIL, indicating that no nodes are to be pruned.

\item [\itype{*first-order-atms-bilattice*}] This bilattice corresponds
to a fully first-order \atms, as opposed to the more conventional one
dealing only with propositional logic.

\item [\itype{*first-order-bilattice*}] The bilattice corresponding to
first-order logic.

\item [\itype{fn-and(f1 f2)}] This function conjoins the values of
$f_1$ and $f_2$ using the bilattice $\wedge$ operation.  The new
\dg-function is returned.

\item [\itype{fn-comp(f1 f2 fn)}] This functions compares the values
of $f_1$ and $f_2$ using the bilattice comparison function \type{fn}.
Either \T\ or \NIL\ is returned.

\item [\itype{fn-dot(f1 f2)}] This function combines the values of
$f_1$ and $f_2$ using the bilattice $\cdot$ operation.  The new
\dg-function is returned.

\item [\itype{fn-dws(f1 f2)}] This function combines the values of
$f_1$ and $f_2$ using the bilattice operation \type{dot-with-star}.
The new \dg-function is returned.

\item [\itype{fn-eq(f1 f2)}] This function compares the values of
$f_1$ and $f_2$ using the bilattice comparison operation.  Either \T\
or \NIL\ is returned.

\item [\itype{fn-not(f1)}] This function negates the values of $f_1$
using the bilattice $\neg$ operation.  The new \dg-function is
returned.

\item [\itype{fn-or(f1 f2)}] This function disjoins the values of
$f_1$ and $f_2$ using the bilattice $\vee$ operation.  The new
\dg-function is returned.

\item [\itype{fn-plus(f1 f2)}] This function adds the values of $f_1$
and $f_2$ using the bilattice $+$ operation.  The new \dg-function is
returned.

% *G*

\item [\itype{get-bdg(var bdg-list)}] Returns the value assigned to
the variable \type{var} by the binding list \type{bdg-list}.  If the
variable does not appear in the binding list, \NIL\ is returned.

\item [\itype{get-mvl (symbol indicator)}]  Like the Common Lisp
\type{get}, but if the value is a list, treats it as an a-list and
searches for an item whose car matches \type{*mvl-invocation*}.
The value taken can be changed with \type{setf}.

\item [\itype{get-val(d f)}] Returns the value of the function $f$ at
the point $d$ of the \dg.

\item [\itype{*global*}] The root of the \dg\ of theories.

\item [\itype{groundp(p)}] Returns \T\ if the proposition $p$ is
ground (i.e., contains no variables) and \NIL\ otherwise.

% *H*

\item [\itype{hierarchy-bilattice(b)}] Returns the bilattice constructed
as an infinite descending chain of copies of $b$.

\item [\itype{*hierarchy-bilattice*}] A bilattice similar to the
default bilattice but where the defaults are labelled with
integer-valued strengths.

% *I*

\item [\itype{*if-translation*}] Controls the fashion in which
\type{normal-form} deals with the connective \type{if}.

\item [\itype{includes(big little)}] This function is used to state that
the theory \type{little} is a subtheory of the theory \type{big}, and
should inherit from it.

\item [\itype{index(p)}] Updates the discrimination net as needed
to insert the proposition $p$ into the database.

\item [\itype{init-analysis(prop)}] The \mvl\ primitive upon which
the proving functions are based, this function accepts a proposition
and creates and returns an instance of the \type{analysis} structure.

\item [\itype{init-prover(p effort)}] This function initializes and
returns an instance of \type{proof} containing auxiliary information
about the proof process.

\item [\itype{*initial-limit*}] is used to indicate the type of search.
If \NIL\ (the default), then the values assigned to the nodes are not
used to order the search, and depth-first search is used.  If not
\NIL, \type{*initial-limit*} is typically 0.

\item [\itype{instp(p1 p2)}] If $p_2$ is an instance of $p_1$,
return a list of their most general unifiers; only variables in $p_1$
are bound.  Returns \NIL\ if $p_2$ is not an instance of $p_1$.

\item [\itype{invoke(fn p)}] This function invokes \type{fn} on the
sentence $p$.  In addition, any keyword arguments are passed to
\type{fn}, but with \type{\&allow-other-keys} so that the keywords
will be ignored if \type{fn} is not expecting them.

% *J*

% *K*

\item [\itype{k-ge(x y)}] Returns \T\ if $x \geq_k y$, and \NIL\ otherwise.

\item [\itype{k-gt(x y)}] Returns \T\ if $x >_k y$, and \NIL\ otherwise.

\item [\itype{k-le(x y)}] Returns \T\ if $x \leq_k y$, and \NIL\ otherwise.

\item [\itype{k-lt(x y)}] Returns \T\ if $x <_k y$, and \NIL\ otherwise.

\item [\itype{k-not-ge(x y)}] Returns \T\ if $x \not\geq_k y$, and
\NIL\ otherwise.

\item [\itype{k-not-gt(x y)}] Returns \T\ if $x \not>_k y$, and \NIL\
otherwise.

\item [\itype{k-not-le(x y)}] Returns \T\ if $x \not\leq_k y$, and
\NIL\ otherwise.

\item [\itype{k-not-lt(x y)}] Returns \T\ if $x \not<_k y$, and \NIL\
otherwise.

\item [\itype{(known p)}] Stashing this predicate simply stashes $p$;
unstashing it causes an error.  Looking it up looks up $p$.

% *L*

\item [\itype{lattice-to-bilattice(l)}] Returns the bilattice constructed
from the lattice $l$.

\item [\itype{lattice-to-dag(lattice)}] Construct a \dg\ from the
given lattice.

\item [\itype{lattice-to-dag-to-bilattice(lattice bilattice)}]  First
construct a \dg\ $D$ from the given lattice.  Then use the bilattice
$B$ to construct the bilattice $B^D$.

\item [\itype{*limit-delta*}] is used if \type{*initial-limit*} is
non-\NIL.  In this case, the nodes are ordered by increasing
 \[\frac{\type{value(node)}-\type{*initial-limit*}}{\type{*limit-delta*}}\]
 A typical value for *initial-limit* is 0, and for *limit-delta* is
1.  The value of this parameter is irrelevant if
\type{*initial-limit*} is \NIL.

\item [\itype{lisp-defined(fn)}] This macro is used to indicate that
the given function can be evaluated by a \lisp\ call.  It is also
used to cause \mvl\ to signal an error if an attempt is made to stash
or unstash an instance of the function.  \type{Lisp-defined} also
accepts an optional argument, which is an inverter function.  The
overall result of attempting to lookup an instance of the function
can be described as follows:
 \begin{enumerate}
 \item If none of the arguments to the function (i.e., all but the
last of the arguments of the relation) is a variable, then the function
is called and the result is unified with the final argument of the
relation.
 \item If the inverter function has been supplied and is equal to \T,
then the function is called and the result is unified with the final
argument of the relation even if some of the arguments are variables.
 \item If some of the arguments are variables but the result is not a
variable, then the inverter function is invoked if one has been
supplied.  The arguments passed to the inverter function are the
result (i.e., the last argument of the relation), and a list of all
of the arguments (i.e., all but the last argument of the relation).
If there is no inverter function, then \NIL\ is returned.  In
addition, if \wbli\ \index{*warn-on-bad-lisp-invocations*} is not
\NIL, the user is warned.
 \item If both some of the arguments and the result are variables, then
\NIL\ is returned.  In addition, if \wbli\ is not \NIL, the user is
warned.
 \end{enumerate}

\item [\itype{lisp-predicate(pred)}] This macro is used to indicate
that the given predicate can be evaluated by a \lisp\ call.  It is
also used to cause \mvl\ to signal an error if an attempt is made to
stash or unstash an instance of the function.

\item [\itype{load-logic}] Empty the database and make a new bilattice
active.  Accepts as an optional argument either a character designating
a bilattice, or the new bilattice itself.  If no argument is supplied,
the user is prompted for the new bilattice.

\item [\itype{load-mvl}] Loads the \mvl\ system.

\item [\itype{lookup(q)}] Searches the database for a sentence
matching the query $q$.  Allowed keywords: \type{cutoffs}, \type{inv},
\apa, \somo.  Returns two values: an answer and the source of the
match.  Returns \NIL\ if the query cannot be found in the database.

\item [\itype{lookup-call(fn p)}] This function is similar to
\type{invoke}, but modifies the values returned to make them agree
with those expected to be returned by \type{lookup}.  Unless the
function \type{fn} supplies them, the truth value is assumed to be
\type{true}.  The second value is \type{attach}.

\item [\itype{lookups(q)}] Similar to \type{lookup}, but looks
for {\em all\/} sentences matching the query.  Returns two lists of
values.

\item [\itype{lookups-call(fn p)}] The plural version of
\type{lookup-call}.

% *M*

\item [\itype{make-root(dag bilattice b)}] Returns a \dg-function that
takes the value $b$ at the root of the \dg.

\item [\itype{make-tv(dag bilattice d b)}] Returns a \dg-function that
takes the value $b$ at the point $d$.  If $d$ is not the root of the
\dg, the value at the root of the \dg\ is generally the point
\type{unknown} in the bilattice, but the user may supply the keyword
\type{root-fn} if he wishes, in which case the value at the root of
the \dg\ is the result of applying \type{root-fn} to the bilattice.

\item [\itype{matchp(pattern atom)}] Returns a list of elements of
the form \type{(b . d)} where $b$ is a most general binding list such
that \type{(plug pattern b)} is an instance of the proposition
labelled with \type{atom} and $d$ is a binding list for the variables
in the database proposition.  \NIL\ is returned if no such binding
list exists.  Thus if \type{P1} is \type{(TALL \#:?2)}, and \type{P2}
is \type{(TALL TOM)}, then
 \[\type{(matchp (TALL TOM) P1)}\]
returns a list consisting of the single entry 
 \[\type{(NIL . ((\#:?2 . TOM)))}\]
 The \type{car} of this entry is \NIL, indicating that no variables
in the query need to be bound to match \type{P1}; the \type{cdr} is
\type{((\#:?2 . TOM))} since the variable \type{\#:?2} in \type{P1}
needs to be bound to \type{TOM} to match the query.
 \[\type{(matchp (TALL ?X) P2)}\]
returns \type{((((?X . TOM)) . NIL))}.

\item [\itype{mvl-and}] The function \type{and} in the active bilattice.

\item [\itype{mvl-dot}] The function \type{dot} in the active bilattice.

\item [\itype{mvl-eq}] The function \type{eq} in the active bilattice.

\item [\itype{mvl-file(string)}] Returns the path for the \mvl\ file
associated with the given name.  The file name may include an
extension, one may be supplied as an optional argument, or
\type{.MVL} will be used by default.  There is an optional keyword
\type{no-error} indicating that even if the file cannot be found, the
path should be returned.

\item [\itype{*mvl-invocation*}] A global variable containing a unique
identifier for the currently active copy of \mvl.

\item [\itype{mvl-load(file)}] Loads the sentences in the given file into
the \mvl\ database.

\item [\itype{mvl-not}] The function \type{not} in the active bilattice.

\item [\itype{mvl-or}] The function \type{or} in the active bilattice.

\item [\itype{mvl-plug}] The function \type{plug} in the active bilattice.

\item [\itype{mvl-plus}] The function \type{plus} in the active bilattice.

\item [\itype{mvl-print}] This function accepts a single argument, a
truth value, and modifies it in a way that will cause it to be
pretty-printed.  The value returned by mvl-print should not be used
for any purpose other than printing.

\item [\itype{mvl-save(file)}] Writes the \mvl\ database out to the
given file.  This function also accepts the keywords \type{cutoffs}
and \type{props}, which are treated as for \type{contents} (see
Section \ref{s.contents}).  Finally, the keyword \type{if-exists} is
passed to the \type{with-open-file} macro.

\item [\itype{mvl-simplicity}] The function \type{simplicity} in the active
bilattice.

\item [\itype{mvl-test-file(string)}] Can be used to evaluate
various forms by the \mvl\ system and compare the results to
predefined values.  Both the forms and the expected answers are in
the \type{.TEST} file whose name is supplied to the function.
Returns the number of errors found.

\item [\itype{mvl-unload(file)}] Removes the sentences in the given file from
the \mvl\ database.

\item [\itype{mvl-vars}] The function \type{vars} in the active bilattice.

% *N*

\item [\itype{negate(p)}] Returns the negation of $p$.

\item [\type{new-$*$var()}]\index{new-star-var} Returns a new
temporary sequence variable.

\item [\type{new-?var()}]\index{new-query-var} Returns a new
temporary variable.

\item [\itype{new-val(a v)}] Sets the truth value of the sentence
associated with the atom $a$ to $v$.  If $v$ is the truth value
\type{unknown}, the sentence is removed from the database.

\item [\itype{normal-form(p)}] Returns an equivalent form of $p$
suitable for insertion into the database (i.e., with the connectives
\type{IFF}, \type{IF} and \type{<=>} reduced to \type{<=} and \type{=>}).

% *O*

% *P*

\item [\itype{plug(p bdg-list)}] Returns the result of instantiating
the proposition $p$ using the given binding list.  This function is
nondestructive.

\item [\itype{prfacts(atom)}] Displays the database facts concerning the
given atom.  Keyword: \type{cutoffs}.

\item [\itype{*probability-bilattice*}] An extended version of the \atms\
bilattice that maintains probabilistic information as well.

\item [\itype{product-bilattice(b1 b2)}] Returns the bilattice
that is the product of the bilattices $b_1$ and $b_2$.

\item [\itype{(property symbol value indicator)}] Stashing this
predicate sets the symbol's \type{indicator} property to \type{value}
for the current invocation of \mvl.  Unstashing it removes the
\type{indicator} property for the current invocation of \mvl\ if the
previous value of this property unifies with \type{value}.  (See the
definitions of \type{get-mvl} and \type{rem-mvl}.)

\item [\itype{props()}] Returns a list of all propositions about
which any information has been recorded.

% *Q*

% *R*

\item [\itype{release-mvl}] (Symbolics only)  Compiles and loads the
\mvl\ system, and declares the new version to be \type{released}.  The
keyword \type{no-recompile} can be used to skip the compilation step.

\item [\itype{rem-mvl (symbol indicator)}] Like the Common Lisp
\type{remprop}, but if the value is a list, treats it as an a-list
and removes the item whose car matches \type{*mvl-invocation*}.

% *S*

\item [\itype{samep(p1 p2)}] Returns \T\ if the propositions $p_1$ and
$p_2$ are identical (up to variable renaming), and \NIL\ otherwise.
If the optional \type{return-answer?} argument is supplied and the
propositions are the same, a binding list is returned instead of \T.

\item [\itype{*save-nodes*}] If \T\ (the default), nodes are
labelled with consecutive integers when they are displayed by the
first-order prover.

\item [\itype{set-search-parameters}] This function is the primitive
used to modify the parameters used to control the first-order theorem
prover's search.  It accepts as keywords \type{initial-limit},
\type{limit-delta}, and \type{final-limit}.

\item [\itype{simplify(f)}] In some cases, manipulating the
\dg-function $f$ may have cause it to contain redundancies, where the
value at some point is presented explicitly even though this value is
in fact the one that would normally be inherited from its parents.
This function returns a modification of $f$ where these irrelevant
entries have been removed.

\item [\itype{splice-bilattice(upper lower)}] Returns the bilattice
obtained by replacing the point \type{unknown} in the bilattice
\type{upper} with a copy of the bilattice \type{lower}.

\item [\itype{standardize-operators(p)}] Returns an equivalent form of $p$
with the connectives \type{<=>}, \type{IF} and \type{IFF} removed.
This function does not use the control variables
\type{*if-translation*} and \et, since it is used when proving
sentences, not when inserting them into the database.

\item [\itype{stash(p)}] Inserts $p$ into the database.  The allowed
keywords are \type{value}, \type{increment-p} and \apa.

\item [\itype{stash-value}] The function \type{stash-val} in the
active bilattice.

\item [\itype{state(p)}] Inserts $p$ into the database after
converting it to normal form.  Keywords: \type{value},
\type{increment-p}.

\item [\itype{state-in-theory(p)}] Adds the proposition $p$ to
the currently active theory.  Allowed keywords: \type{theory} can be
used to specify a different theory to which $p$ should be added;
\type{value} and \type{increment-p} have the same meaning as for
\type{state}.  (See Section \ref{s.database}; if these keywords are
not used, the function \type{state} can be used equally well, since
the default stash value with the theory mechanism activated also
inserts the proposition into the currently active theory.)

\item [\itype{std-cutoffs}] Cutoff information constructed from the
value of \type{cutoff-val} in the active bilattice.  See Section
\ref{s.cutoffs}.

% *T*

\item [\itype{t-ge(x y)}] Returns \T\ if $x \geq_t y$, and \NIL\ otherwise.

\item [\itype{t-gt(x y)}] Returns \T\ if $x >_t y$, and \NIL\ otherwise.

\item [\itype{t-le(x y)}] Returns \T\ if $x \leq_t y$, and \NIL\ otherwise.

\item [\itype{t-lt(x y)}] Returns \T\ if $x <_t y$, and \NIL\ otherwise.

\item [\itype{t-not-ge(x y)}] Returns \T\ if $x \not\geq_t y$, and
\NIL\ otherwise.

\item [\itype{t-not-gt(x y)}] Returns \T\ if $x \not>_t y$, and \NIL\
otherwise.

\item [\itype{t-not-le(x y)}] Returns \T\ if $x \not\leq_t y$, and
\NIL\ otherwise.

\item [\itype{t-not-lt(x y)}] Returns \T\ if $x \not<_t y$, and \NIL\
otherwise.

\item [\itype{test(value cutoffs \&optional bilattice)}] Returns \T\
if the truth value passes the test associated to \type{cutoffs}, and
\NIL\ otherwise.  The test is performed for the supplied bilattice,
or for the currently active one is the optional argument is omitted.

\item [\itype{test-mvl}] Tests the \mvl\ system by invoking the function
\type{mvl-test-file} on the file \type{MVL.TEST}.

\item [\itype{theory-active-p}] Returns \T\ or \NIL\ depending on
whether or not the theory mechanism has been activated.

\item [\itype{theory-contents()}] Displays the contents of the
currently active theory, which can be changed using the
keyword \type{theory}.  The keywords \type{cutoffs} and \type{props}
are allowed and are treated as for \type{contents}.  (See Section
\ref{s.contents}.)  There is an additional keyword,
\itype{this-theory-only}, which defaults to \T\ and, if supplied,
indicates that only sentences whose truth values {\em change\/} in the
given theory are to be displayed.

\item [\itype{theory-dag}] The \dg\ of theories.

\item [\itype{*time-bilattice*}] A compound bilattice that can be used
for reasoning about events occurring against a temporal background.

\item [\itype{true}] The truth value \type{true} in the active bilattice.

\item [\itype{true-in-theory}]  The truth value of a sentence that is
true in the currently-active theory.

\item [\itype{truth-value(a)}] Returns the truth value associated with
a particular \lisp\ atom.  Thus, if the result of asserting
\type{(not (short tom))} into the database was the atom \type{P38},
\type{(denotes 'p38)} would evaluate to \type{(short tom)}, and
\type{(truth-value 'p38)} would evaluate to \type{false}. 

% *U*

\item [\itype{unifyp(p1 p2)}] Returns a list of the most general
unifiers of $p_1$ and $p_2$, or \NIL\ if the propositions fail to
unify.

\item [\itype{unincludes(big little)}] This function is used to state
that the theory \type{little} is no longer a subtheory of the theory
\type{big}, and should not inherit from it.

\item [\itype{unindex(p)}] Removes from the discrimination net 
information relating to the proposition $p$.

\item [\itype{unknown}] The truth value \type{unknown} in the active
bilattice.

\item [\itype{(unknown p)}]  Stashing this predicate unstashes $p$;
unstashing it causes an error.  Looking it up fails if $p$
can be found in the database, and succeeds otherwise.

\item [\itype{unstash(p)}] Removes $p$ from the database.  Allowed
keywords: \type{value}, \type{decrement-p} and \apa.

\item [\itype{unstash-facts-with-atom(atom)}] Unstashes from the
database any facts concerning the given atom.  Keywords:
\type{value}, \type{decrement-p}, \apa.

\item [\itype{unstash-value}] The function \type{unstash-val} in the
active bilattice.

\item [\itype{unstate(p)}] Removes $p$ from the database after
converting it to normal form.  Keywords: \type{value},
\type{decrement-p}.

\item [\itype{unstate-from-theory(p)}] Removes the proposition $p$
from the currently active theory.  Allowed keywords:
\type{theory} can be used to specify a different theory from which
$p$ should be removed; \type{value} and \type{decrement-p} have the
same meaning as for \type{unstate}.  (See Section \ref{s.database}.)

% *V*

\item [\itype{(value symbol val)}] Stashing this predicate sets the
given symbol to the value \type{val}.  Unstashing it makes the symbol
unbound if its current value unifies with \type{val}.

\item [\itype{varp(x)}] Returns \type{T} if $x$ is a variable, and \NIL\
otherwise.

\item [\type{varp$*$(x)}]\index{varp-star} Returns \T\ if $x$ is a
sequence variable, and \NIL\ otherwise.

\item [\type{varp?(x)}]\index{varp-query} Returns \T\ if $x$ is a
standard (i.e., non-sequence) variable, and \NIL\ otherwise.

\item [\itype{vars-in(p)}] Returns a list of the variables appearing in $p$.

\item [\itype{vartype(x)}] Returns ? if $x$ is a standard variable,
$*$ if $x$ is a sequence variable, and \NIL\ is $x$ is not a
variable.

% *W*

\item [\itype{with-mvl-invocation}] This is a macro similar to the
Common Lisp macro \type{with-open-file}.  It accepts an optional
argument that is an instance of \mi\ and is followed by a body of
forms (presumably including a \type{load-logic}) that are evaluated
in a locally-bound copy of \mvl.

% *X*

% *Y*

% *Z*

\end{list}

\documentstyle{article}

\setlength{\topmargin}{0pt}
\setlength{\oddsidemargin}{0pt}
\setlength{\evensidemargin}{0pt}
\setlength{\textwidth}{6.5in}
\setlength{\textheight}{8.5in}
\setlength{\headheight}{0pt}

%here are some mathematical symbols
\def\emptyset{{\hbox{\rm \O}}}
\def\set#1{\{#1\}}
\newcommand{\type}[1]{\mbox{\tt #1}}
\newcommand{\str}{?$*$}
\def\R{\rlap{\rm I}\,{\rm R}}
\def\powset{{\cal P}}
\def\glb{{\mbox{\rm glb}}}\def\lub{{\mbox{\rm lub}}}

\newcommand{\nm}{non\-mon\-o\-ton\-ic}
\newcommand{\Nm}{Non\-mon\-o\-ton\-ic}

\newcommand{\prolog}{{\sc prolog}}\newcommand{\Prolog}{{\sc Prolog}}
\newcommand{\atms}{{\sc atms}}\newcommand{\Atms}{{\sc Atms}}

%this is where the title page begins

\def\draftline{}

\def\authors{Matthew L. Ginsberg}

\newcommand{\nodblspace}{\baselineskip=\normalbaselineskip}
\def\group{\nodblspace Department of Computer Science\\
Stanford University\\
Stanford, California 94305\\
\bigskip
Updated \today
}

\def\titlebottom{\vspace{1.55in}
{\large \authors}\\
\vspace{1.55in}
\group
\end{center}
\clearpage}

\def\xtitletop{\vspace*{1.35in}\begin{center}}
\def\xtitl#1#2{\xtitletop
{\Large\bf #2}\\
\titlebottom}

\newcommand{\xtitle}[2]
{\def\tphrase{#2}\begin{titlepage}\draftline\xtitl{#1}{#2}\end{titlepage}}

\newcommand{\mvl}{{\sc mvl}}
\newcommand{\Mvl}{{\sc Mvl}}
\newcommand{\lisp}{{\sc lisp}}
\newcommand{\Lisp}{{\sc Lisp}}
\newcommand{\dg}{{\sc dag}}
\newcommand{\Dg}{{\sc Dag}}
\newcommand{\T}{\type{T}}
\newcommand{\NIL}{\type{NIL}}
\newcommand{\apa}{{\tt allow-\-procedural-\-attachment}}
\newcommand{\at}{{\tt *anytime-\-trace*}}
\newcommand{\mi}{{\tt mvl-\-invocation}}
\newcommand{\wbli}{{\tt *warn-\-on-\-bad-\-lisp-\-invocations*}}
\newcommand{\somo}{{\tt succeed-\-on-\-modal-\-ops}}
\newcommand{\et}{{\tt *equivalence-\-translation*}}
\newcommand{\itype}[1]{\type{#1}\index{#1 **}}

\makeindex
%\renewcommand{\index}[1]{}

\begin{document}
\xtitle{}{User's Guide to the MVL System}

\tableofcontents
\clearpage

\section{Introduction}
\subsection{What you got}

The \mvl\ system is currently distributed via anonymous ftp from
t.stanford.edu.  A .DVI version of this file, the \mvl\ manual,
can be found in the \type{papers} directory on this machine.

The \type{mvl} directory contains the source code for the \mvl\ system
itself, which is written in standard Common Lisp.  It should be
possible to port the system to any Common Lisp without substantial
difficulty.  In addition, the \type{papers} directory contains .DVI
versions of the following papers that relate to the functioning of the
system:
 \begin{enumerate}
 \item ``Multivalued logics: A uniform approach to inference in
artificial intelligence.''  This paper gives a formal description of
multivalued logics and the \mvl\ system.
 \item ``Bilattices and modal operators.''  This paper describes the
formal connection between the \mvl\ system and modal operators.
 \end{enumerate}

The \mvl\ sources are public domain.  You are free to use, modify, and
distribute them as you see fit, provided that such use or distribution
is not for commercial purposes.  No warranties are made with regard to
the performance of the system, however, and no commitment of support
has been made.

\subsection{Overview}

This user's guide has been written to help you install and use the \mvl\
system.  It describes the basic structure of the system, and shows you
how to take advantage of the flexibility of multivalued inference.

Essentially, \mvl\ is just like any other logic programming system.
There is a database consisting of those sentences about which the
system knows something, and there are facilities for determining
whether or not some query follows from the information the database
contains.  Viewed in this fashion, the \mvl\ system is not much
different from a conventional logic programming language such as
\prolog.

What makes \mvl\ unique is that each sentence in the database is
labelled with a {\em truth value\/}\index{truth value}.  Thus,
instead of simply assuming a statement to be true and dumping it into
the database, we will {\em label\/} the statement as true when we
insert it.  Of course, we can use other labels as well: true by
default, true by virtue of some assumption, or even false or false by
default.  \Mvl\ considers these truth values when determining how to
respond to a \prolog-like query.

This manual is divided into two parts.  In the first part (Sections
\ref{s.fol1} -- \ref{s.support}), we describe the use of the system
where the truth values assigned to sentences are simply ``true'' or
``false.''  Used in this fashion, \mvl\ is not much different from
\prolog.

The second part of the manual shows you how to use logics involving
more complex truth values.  We examine the logics that have been
predefined as part of the \mvl\ system, and describe the various ways
in which you can define your own logics.  Several tools for this
purpose have been included as part of the \mvl\ system.  In addition,
we describe \mvl's facilities for handling modal operators, which
are operators that accept logical sentences as arguments.

\subsection{Why MVL is unique}

As mentioned in the previous section, what distinguishes \mvl\ from
other declarative programming languages is its support for multiple
truth values.  There are other differences, however, and these will
be discussed in detail in subsequent sections.  Here's a quick summary:
 \begin{enumerate}
 \item As mentioned, the principal feature offered by \mvl\ is its
support of multiple truth values.
 \item The basic theorem prover used by the system includes facilities
to detect loops and control recursion encountered in proof search.
The first-order theorem prover is discussed in Section \ref{s.fo},
although we will have nothing further to say about recursion control.
 \item The system can handle negative subgoals containing free variables
as discussed in Section \ref{s.effort}.
 \item The \mvl\ unifier includes an occurcheck and full support for
sequence variables.  This is discussed in Section \ref{s.rep}.
 \end{enumerate}

\section{Installation and testing}

The \mvl\ system is currently developed and maintained on a
Sparcstation running Allegro Common Lisp and is tested only
infrequently on other platforms.  Although the system has been written
in standard Common Lisp to maximize ease of portability and no
problems have been encountered when moving the system to a variety of
platforms, performance on other systems is not guaranteed.

\Mvl\ currently supports both ANSI Common Lisp and Common Lisp as
described in Steele's 1984 book.  In order to determine whether or not
the ANSI standard is conformed to, the system checks to see if
\type{in-package} is a macro or a function.  If the former, it is
assumed that the entire ANSI standard is supported; if the latter, it
is assumed that the language is as in Steele and various extensions
are defined.  It is expected that some future release of \mvl\ will
support the ANSI standard only.

\subsection{Sun-4}

In order to load the system on a Sparcstation running Allegro Common
Lisp, do the following:
 \begin{enumerate}

 \item (ANSI) Set up a logical pathname translation so that logical
pathnames beginning with \type{MVL} point to the direction where \mvl\
is kept.  (non-ANSI) Edit the file \type{mvldcl.lisp} so that the
\type{:default-pathname} field in the \mvl\ system definition points
to the directory where \mvl\ will be kept.  This directory should also
contain the following files:
 \begin{enumerate}\setlength{\itemsep}{0pt}\raggedright
 \item {\tt f1.mvl, recursion.mvl, def-test.mvl, circ-test-1.mvl,
circ-test-2.mvl, circ-test-3.mvl, circ-test-4.mvl, completed-db-1.mvl,
completed-db-2.mvl, modal-test.mvl, actions.mvl, shooting.mvl}
 \item {\tt mvl.test, bc.test, recursion.test, atms.test,
defaults.test, circumscription.test, completed-db.test, modal.test,
fc.test, temporal.test}
 \end{enumerate}
 \item Compile and load the file \type{mvldcl.lisp}.
 \item Execute the function
\type{(mvl:compile-mvl)}\index{compile-mvl} to compile and load the
\mvl\ system.  Using this function in the future will only compile 
those files that have been changed since last compiled.
 \item Invoke the function \type{(use-package 'mvl)} to access the
functions and symbols of the \mvl\ system (which is kept in its own
package; see Section \ref{s.package}).
 \end{enumerate}

Now invoke the function {\tt (test-mvl)}, which should return 0
(i.e., 0 errors encountered).  If it doesn't, the problem is probably
that \mvl\ is unable to locate the files mentioned in (1) above.

\subsection{Symbolics 3600 series}

The \mvl\ system will run on Symbolics Genera 7.2 or later.  It is
possible to run the system using Genera 7.0 or 7.1, although a bug in
the implementation of the Common Lisp \type{\&allow-other-keys}
keyword may cause the system to dump to the cold load stream if
certain predicates are procedurally attached.  (See Section
\ref{s.attach}.)

You should begin by editing the file \type{mvldcl.lisp} so that the
default pathname in the \mvl\ system declaration points to the
directory in which the files are actually kept.  You can then load
the \mvl\ system using either the Load System command or the Lisp
function \type{(mvl:load-mvl)}\index{load-mvl} that is defined in the
file \type{mvldcl.lisp}.  Now access the \mvl\ package using
\type{(use-package 'mvl)} and execute the function \itype{test-mvl}.
This function should return 0 (i.e., 0 errors encountered).  If it
doesn't, the problem is probably that \mvl\ is unable to find the
directory its test files are on (they're included in the distribution
tape).

\subsection{Other machines}

The \mvl\ system has also been run and tested on a wide variety of
other machines.  In order to load the system, follow the instructions
provided for Sun-4's.

\section{The MVL database}
\label{s.fol1}
\subsection{Representation}
\subsubsection{Structure of the database}
\label{s.rep}

As already remarked, the \mvl\ database consists of sentences labelled
with truth values.  In this section, we will describe the form of the
sentences themselves; in the next section, we comment briefly on the
truth values used to label them.

Sentences in the database are represented as \lisp\ s-expressions
using the prenex normal form \index{prenex normal form} of the
standard logical connectives \type{NOT}, \type{OR} and \type{AND}.
Thus the statement that Fred is tall and Harry is not would be
written as:
 \[\type{(AND (TALL FRED) (NOT (TALL HARRY)))}\]

Any \lisp\ atom can be used as a constant term, with the exception of
variables (see below) and the following reserved connectives:
 \[\type{NOT AND OR => <= <=> IF IFF}\]
The connective \type{<=} is used in backward-chaining rules;
 \[\type{(<= CONCLUSION PREMISE-1 ... PREMISE-N)}\]
is logically equivalent to
 \[\type{(IF (AND PREMISE-1 ... PREMISE-N) CONCLUSION)}\]
as is the forward-chaining version
 \[\type{(=> PREMISE-1 ... PREMISE-N CONCLUSION)}\]
 Although logically equivalent, these expressions behave differently
in practice.  The version
 \[\type{(IF (AND PREMISE-1 ... PREMISE-N) CONCLUSION)}\]
 in which the connective \type{IF} is used represents a pair of
rules, depending on the value of the parameter \itype{*if-translation*}.
If the value is \type{bc} (the default), then two backward-chaining
rules are produced, so that
 \[\type{(IF (A) (B))}\]
is equivalent to
 \[\type{(AND (<= (B) (A)) (<= (NOT (A)) (NOT (B))))}\]
 the conjunction of the original rule and its contraposition.
If the value of \type{*if-translation*} is \type{fc}, then
forward-chaining rules are produced and the translation is to
 \[\type{(AND (=> (A) (B)) (=> (NOT (B)) (NOT (A))))}\]
 Finally, if \type{*if-translation*} has the value \type{mix}, the
pair of noncontraposed rules
 \[\type{(AND (<= (B) (A)) (=> (A) (B)))}\]
 is produced.

Equivalences are represented using either \type{IFF} or \type{<=>}.
Thus
 \[\type{(IFF (FOO FRED) (GOO FRED))}\]
is equivalent to
 \[\type{(AND (IF (FOO FRED) (GOO FRED)) (IF (GOO FRED) (FOO FRED)))}\]
 Translations involving \type{<=>} depend on the value of the
parameter \et\index{*equivalence-translation*}.  As above, if the
value is \type{bc} (the default), then two backward-chaining rules
are produced, so that
 \[\type{(<=> (FOO FRED) (GOO FRED))}\]
translates to
 \[\type{(AND (<= (FOO FRED) (GOO FRED)) (<= (GOO FRED) (FOO FRED)))}\]
 If \et\ has the value \type{fc}, two forward-chaining rules are produced
so that the translation is to
 \[\type{(AND (=> (FOO FRED) (GOO FRED)) (=> (GOO FRED) (FOO FRED)))}\]
 Finally, if the value is \type{mix}, one forward-chaining rule and
one backward-chaining rule are produced:
 \[\type{(AND (=> (FOO FRED) (GOO FRED)) (<= (FOO FRED) (GOO FRED)))}\]

Variables \index{variables} are represented using \lisp\ atoms
beginning with the character \type{?} and are handled as in \prolog,
so that free variables in database statements are implicitly
universally quantified.  Again as in \prolog, existential
quantification is handled via Skolemization; thus the
backward-chaining sentence ``Every rich person is happy'' and the
sentence ``There is a rich person'' would be written as
 \begin{equation}
\type{(<= (HAPPY ?P) (PERSON ?P) (RICH ?P))}
\label{e.if}
\end{equation}
and
 \[\type{(AND (PERSON P0017) (RICH P0017))}\]
respectively.

The variables in database sentences are standardized apart from the
other variables used by the system.  This is done by using uninterned
variables in the actual representation of a database fact, so that
the internal representation of (\ref{e.if}) would be
 \[\type{(<= (HAPPY \#:?P) (PERSON \#:?P) (RICH \#:?P))}\]
 The original variables are generally reinstated when \mvl\ is asked
to display a database proposition.

\subsubsection{Sequence variables}

In addition to the simple variables discussed in the previous
section, \mvl\ supports the use of ``sequence
variables.''\index{sequence variables} These variables begin with the
characters \type{\str}, and will match not a single atom, but an
arbitrary (possibly empty) sequence of atoms.

Thus, for example,
 \[\type{(FOO A B)}\]
 will unify with
 \[\type{(FOO \str X)},\]
 returning the binding list
 \[\type{((\str X . (A B)))}\]
 in which \type{?*X} has been bound to the list \type{(A B)}.  This
binding list will presumably be displayed as
 \[\type{((\str X A B))}.\]
 Similarly, the most general unifier of $\type{(FOO)}$ and
$\type{(FOO \str X)}$ is given by $\type{((\str X))}$.

The presence of sequence variables may cause the most general unifier
of two patterns to be nonunique.  For example, the result of unifying
 \[\type{(FOO A B)}\]
 with
 \[\type{(FOO \str A \str B)}\]
 can be any one of the three answers
 \begin{eqnarray*}
& \type{((\str A A B) (\str B))} & \\
& \type{((\str A A) (\str B B))} & \\
& \type{((\str A) (\str B A B))} &
 \end{eqnarray*}

For this reason, the unification functions return not a single answer,
but a {\em list\/} of possible unifiers.

\subsubsection{Functions to manipulate database sentences}
\label{s.manipulate}

The following functions and global parameters are used to manipulate
sentences in the database:
 \begin{description}
 \item [\itype{varp(x)}] Returns \T\ if $x$ is a variable, and
\NIL\ otherwise.
 \item [\type{varp?(x)}]\index{varp-query} Returns \T\ if $x$ is a
standard (i.e., non-sequence) variable, and \NIL\ otherwise.
 \item [\type{varp$*$(x)}]\index{varp-star} Returns \T\ if $x$ is a
sequence variable, and \NIL\ otherwise.
 \item [\itype{vartype(x)}] Returns ? if $x$ is a standard variable,
$*$ if $x$ is a sequence variable, and \NIL\ is $x$ is not a
variable.
 \item [\type{new-?var()}]\index{new-query-var} Returns a new
temporary variable.
 \item [\type{new-$*$var()}]\index{new-star-var} Returns a new
temporary sequence variable.
 \item [\itype{normal-form(p)}] Returns an equivalent form of $p$
suitable for insertion into the database (i.e., with the connectives
\type{IFF}, \type{IF} and \type{<=>} reduced to \type{<=} and \type{=>}).
 \item [\itype{*if-translation*}] Controls the fashion in which
\type{normal-form} deals with the connective \type{if}.
 \item [\itype{*equivalence-translation*}] Controls the fashion in which
\type{normal-form} deals with the connective \type{<=>}.
 \item [\itype{negate(p)}] Returns the negation of $p$.
 \item [\itype{standardize-operators(p)}] Returns an equivalent form
of $p$ with the connectives \type{<=>}, \type{IF} and \type{IFF}
removed.  (Note: This function does not use the control variables
\type{*if-translation*} and \et, since it is used when proving
sentences, not when inserting them into the database.)
 \item [\itype{cnf(p)}] Returns $p$ in conjunctive normal form, as a
list of conjuncts, where each conjunct is a list of disjuncts.  Thus
 \[\type{(CNF '(AND A B))}\]
returns
 \[\type{((A) (B))}\]
\item [\itype{dnf(p)}] Returns $p$ in disjunctive normal form, as a list
of disjuncts, where each disjunct is a list of conjuncts.  Thus
 \[\type{(DNF '(AND A B))}\]
returns
 \[\type{((A B))}\]
\end{description}

\subsection{Truth values}

The truth values used to label the sentences in the \mvl\ database
are elements of a bilattice that can be selected by the user with
the function \type{load-logic}.  We will defer a discussion of these
truth values to the second portion of this manual, which discusses
the use of bilattices that label sentences as other than simply true
or false.  At this point, we merely remark that if it is our intention
to assert the sentence \type{(not (foo))}, this is actually accomplished
by asserting that the negated sentence \type{(foo)} has truth value
\type{false}.

\subsection{Database form and internals}
\label{s.database}

The database is maintained using a discrimination net, as suggested
by Charniak et.~al in their book on {\sc lisp} programming.  Any
particular sentence in the database, say \type{(tall tom)}, is
associated with a \lisp\ atom such as \type{P27} that is its
``proposition number.''  The propositions are numbered consecutively,
beginning with \type{P1}.  At any given point in time, the function
\type{(props)} can be invoked to return a list of all of the
propositions in the database.

Very few of the functions that are used to manipulate the internal form
of the database are likely to be needed by a user of the \mvl\ system.  The
following are the most useful:
 \begin{description}
 \item [\itype{props()}] This function returns a list of all
propositions about which any information has been recorded.
 \item [\itype{denotes(a)}] Returns the proposition associated with a
particular \lisp\ atom.
 \item [\itype{truth-value(a)}] Returns the truth value associated
with a particular \lisp\ atom.  Thus, if the result of asserting
\type{(not (short tom))} into the database was the atom \type{P38},
\type{(denotes 'p38)} would evaluate to \type{(short tom)}, and
\type{(truth-value 'p38)} would evaluate to \type{false}. 
 \item [\itype{index(p)}] Updates the discrimination net as
needed to insert the proposition $p$ into the database.
 \item [\itype{unindex(p)}] Removes from the discrimination net
information relating to the proposition $p$.
 \end{description}

\section{Quantification}

As we have remarked, quantification is treated in \mvl\ in a fashion
similar to that used by \prolog\ systems:  Free variables appearing
in the database are assumed to be universally quantified; free variables
appearing in queries are existentially quantified.  Skolemization is
used to handle existential quantification in database statements or
universal quantification in queries.

The use of this approach means that the answer to an existentially
quantified query will often be of a form that indicates how the
variables should be bound in the result.  In \mvl\, these binding
lists \index{binding lists} are represented as expressions of the
form
 \[\type{((var-1 . bdg-1) ... (var-n . bdg-n))}\]
 where \type{var-1} through \type{var-n} are variables and
\type{bdg-1} through \type{bdg-n} are their respective bindings.
If no variable is bound to a value, the empty binding list \NIL\
is used.

As an example, suppose that the database contains the single sentence
 \[\type{(LIKES COCO ?X)}\]
 indicating that Coco likes everything.  Now the query
 \[\type{(LIKES ?Y MATT)}\]
 will be interpreted as, ``Find something that likes Matt.''  The answer
returned will be
 \[\type{((?Y . COCO))}\]
since Coco likes Matt.  Since the queries
 \[\type{(LIKES COCO MATT)}\]
and
 \[\type{(LIKES COCO ?Z)}\]
 both succeed without binding a variable, the answer \NIL\ will be
returned for each.\footnote{In order to distinguish the empty binding
list \NIL\ from a failed database query, all queries return a binding
list as well as a truth value.  This is discussed in Section
\ref{s.answer}.}

\subsection{Functions used to manipulate binding lists}

A variety of functions within \mvl\ are used to construct and manipulate
binding lists.  Here are some of the more useful ones:
 \begin{description}
 \item [\itype{groundp(p)}] Returns \T\ if the proposition $p$ is
ground (i.e., contains no variables) and \NIL\ otherwise.
 \item [\itype{vars-in(p)}] Returns a list of the variables
appearing in $p$.
 \item [\itype{get-bdg(var bdg-list)}] Returns the value assigned to
the variable \type{var} by the binding list \type{bdg-list}.  If the
variable does not appear in the binding list, \NIL\ is returned.
 \item [\itype{delete-bdg(var bdg-list)}] Destructively modifies
\type{bdg-list}, removing any binding for the variable \type{var}.
 \item [\itype{plug(p bdg-list)}] Returns the result of instantiating
the proposition $p$ using the given binding list.  This function is
nondestructive.
 \item [\itype{samep(p1 p2)}] Returns \T\ if the propositions $p_1$
and $p_2$ are identical (up to variable renaming), and \NIL\
otherwise.  If the optional \type{return-answer?} argument is
supplied and the propositions are the same, a binding list is
returned instead of \T.
 \item [\itype{unifyp(p1 p2)}] Returns a list of the most general
unifiers of $p_1$ and $p_2$, or \NIL\ if the propositions fail to
unify.
 \item [\itype{instp(p1 p2)}] If $p_2$ is an instance of $p_1$,
return a list of their most general unifiers; only variables in
$p_1$ are bound.  Returns \NIL\ if $p_2$ is not an instance of $p_1$.
 \item [\itype{matchp(pattern atom)}] Returns a list of elements of
the form \type{(b . d)} where $b$ is a most general binding list such
that \type{(plug pattern b)} is an instance of the proposition
labelled with \type{atom} and $d$ is a binding list for the variables
in the database proposition.  \NIL\ is returned if no such binding
list exists.  Thus if \type{P1} is \type{(TALL \#:?2)}, and \type{P2}
is \type{(TALL TOM)}, then
 \[\type{(matchp (TALL TOM) P1)}\]
returns a list consisting of the single entry 
 \[\type{(NIL . ((\#:?2 . TOM)))}\]
 The \type{car} of this entry is \NIL, indicating that no variables
in the query need to be bound to match \type{P1}; the \type{cdr} is
\type{((\#:?2 . TOM))} since the variable \type{\#:?2} in \type{P1}
needs to be bound to \type{TOM} to match the query.
 \[\type{(matchp (TALL ?X) P2)}\]
returns \type{((((?X . TOM)) . NIL))}.
 \end{description}

Additional functions are described in Section \ref{s.binding-dag}.

\subsection{Negative subgoals containing free variables}
\label{s.effort}

Consider the following \prolog\ fragment:
\begin{verbatim}
     f(x) :- not(g(x)).
     g(A).
\end{verbatim}
 The interpretation we will assign to the implication
\type{f(x) :- not(g(x))} will be that it is universally quantified
over $x$, and therefore corresponds to the implication
 \[\forall x.[\neg g(x) \supset f(x)]\]
 Given this interpretation, what should the response be if we provide
the system with the query \type{f(x)}?

In light of the negation as failure rule of \prolog\ (which has analogs
in the \mvl\ system), we would expect the query to succeed for all values
of $x$ {\em except\/} $x=A$.

Most \prolog\ systems cannot respond to queries such as this one, since
a nonground negative subgoal is encountered during the proof process.
One approach to these queries, suggested by David Chan, is to do just
what we have suggested in the previous paragraph, computing those binding
lists for which the query fails.

Unfortunately, there are many queries that succeed in general, but
for which it is impractical to compute all of the bindings for which
the query fails.  This may be because there are infinitely many such
failures, or because the time needed to compute them explicitly is
prohibitive.  Situations such as this arise frequently in
applications of declarative programming techniques to planning
problems; here, we would like simply to return an indication that the
query ``succeeds in general.''  In the above example, we might
report that the query had succeeded without binding the variable $x$;
even though $x=A$ is an exception, the report of success is true in
general.

Both approaches are supported by \mvl, and the global variable
\type{*completed-database*} is used to distinguish between them.  If
this variable has the value \T, all exceptions are computed; if it
has the value \NIL, the other approach is taken.  The default value
is \NIL:
 \begin{description}
 \item [\itype{*completed-database*}] If \T, find all exceptions
to default proofs.  If \NIL\ (the default), accept defaults proofs
with possible exceptions as complete.
 \end{description}

\section{Modifying the MVL database}

\subsection{Keyboard input}

The simplest functions to modify the \mvl\ database simply accept
propositions from the user and either add or remove them from the
knowledge base.  The primitive function for adding to the database is
\itype{stash}; the function \itype{unstash} is used for removal.

The function \type{stash} accepts a proposition and up to three
keywords:
 \begin{description}
 \item [\itype{value}]  This keyword is used to supply a specific truth
value to be assigned to the proposition.  If no value is supplied, the
default value is used (this value is supplied when the logic is chosen;
see Section \ref{s.bilattice}).
 \item [\itype{increment-p}] If the proposition in question already
appears in the database, the user has a choice: The given truth value
can be used to replace the one currently appearing in the database,
or it can be combined with the existing value using the combination
function supplied with the logic being used.  The truth value is
replaced if \type{increment-p} is \NIL, and is updated if it is \T\
(the default).
 \item [\apa\index{allow-procedural-attachment}] This keyword,
which defaults to \T, controls whether or not it is possible to
procedurally attach the action of stashing the supplied proposition.
See Section \ref{s.attach}.
 \end{description}

Somewhat more useful than \type{stash} is the function \itype{state}.
This function is identical to \type{stash}, except it translates the
supplied sentence into a form suitable for inference by using the
function \itype{normal-form} described in Section \ref{s.manipulate}.
In addition, \type{state} cannot be procedurally attached, so the
keyword \apa\ is not allowed.

For removing sentences from the database, the functions \type{unstash}
and \itype{unstate} are used.

The function \type{unstash} accepts a proposition to be removed from
the database, and up to three keywords: \type{value},
\itype{decrement-p} and \apa.  The last of these is treated as for
\type{stash}.  In general, if no keywords are used, the effect is to
remove the proposition from the database.

If \type{decrement-p} is \NIL, then \type{unstash} behaves like
\type{stash}, except the default value used for unstashing is not
necessarily the same as the default value used for stashing.

If \type{decrement-p} is \T, then the new truth value assigned to
the proposition is given by
 \begin{equation}
o \cdot \neg v
\label{e.dot-star}
\end{equation}
 where $o$ is the old value and $v$ is the value supplied via the
\type{value} keyword.  The $\cdot$ and $\neg$ operations are as
described in Section \ref{s.bilattice}.

There is one additional function for removing sentences from the
database, \itype{empty}.  This function accepts no arguments, but
does take a single keyword, \type{value}.  The function calls
\type{unstash} for every sentence in the database.

All of the functions described in this section work by calling the
primitive function \itype{new-val}.  This function accepts as
arguments a proposition name and a truth value, and sets the truth
value of the proposition to the given value.  If the truth value is
\type{unknown}, indicating that nothing is known about the
proposition, it is removed from the database and the discrimination
net.

\subsection{Input from files}

It is also possible to enter a database from a disk file.  These files
should in general have the extension \type{.MVL}, and should be
located either on the currently active disk area or on the \mvl\
directory.  The files should consist simply of the sentences to be
added to the database.

In addition to consisting of such sentences, the files may contain the
keyword \type{:value}, followed by the value to be assigned to each
sentence, or the keyword \type{:function}, followed by a function to
be applied to each sentence in order to {\em compute\/} the value to
be assigned to it.  The keyword \type{:empty} empties the database
before loading the file.  Examples of these constructs can be found in
the various \type{.MVL} files included on the distribution tape.

The sentences are loaded with \itype{mvl-load}, which accepts a file
name and an optional keyword \type{empty}.  If the keyword is
supplied, the database is emptied before the file is read; in any
event, \type{state} is then applied to each sentence in the given
file.  The sentences in the file can be removed from the database
with the function \itype{mvl-unload}, which accepts a file name as its
argument.  These functions each call \type{mvl-file}, which accepts
a string and returns a path to the \type{.MVL} file associated to it.

Related to \type{mvl-load} is \itype{mvl-save}, which also accepts a
file name as an argument, but instead of loading a declarative
database from the file, writes the current database {\em to\/} the
file.

\subsection{The editor interface} \index{editor interface}
\subsubsection{Sun EMACS editor}

In addition to the file commands \type{mvl-load} and
\type{mvl-unload}, there is an interface to the {\sc emacs} editor on
Sparcstations.  For a file of mode \mvl, certain keys are bound as
follows (each of these commands also recognizes the
\type{:value} and \type{:function} keywords):

 \begin{description}
 \item [\type{meta-S}] Applies \type{state} to each s-expression in the
current region.
 \item [\type{meta-U}] Applies \type{unstate} to each s-expression in the
current region.
 \item [\itype{meta-s}] Applies \type{state} to each s-expression in the
current defun (i.e., s-expression beginning in column 1).
 \item [\itype{meta-u}] Applies \type{unstate} to each s-expression in the
current defun.
 \end{description}

In order to use these commands, certain modifications will need to be made
to your \type{.emacs} file.  The file \type{dot.emacs} is an example of
what to do.

\subsubsection{Symbolics ZMACS editor}

In addition to the file commands \type{mvl-load} and
\type{mvl-unload}, there is an interface to the {\sc zmacs} editor on
Symbolics 3600 series machines.  For a file of mode \mvl, certain
keys are bound as follows (each of these commands also recognizes the
\type{:value} and \type{:function} keywords):

 \begin{description}
 \item [\itype{super-s}] Applies \type{state} to each s-expression in the
current region.
 \item [\itype{super-u}] Applies \type{unstate} to each s-expression in the
current region.
 \item [\itype{super-f}] Applies \type{fc} to each s-expression in the
current region.  See Section \ref{s.fc}.
 \item [\itype{super-e}] Applies \type{erase} to each s-expression in the
current region.  See Section \ref{s.fc}.
 \end{description}

\subsection{Summary of database modification functions}

\begin{description}
 \item [\itype{stash(p)}] Inserts $p$ into the database.  Allowed
keywords: \type{value}, \type{increment-p} and \apa.
 \item [\itype{state(p)}] Inserts $p$ into the database after
converting it to normal form.  Keywords: \type{value},
\type{increment-p}.
 \item [\itype{unstash(p)}] Removes $p$ from the database.  Keywords:
\type{value}, \type{decrement-p}, \apa.
 \item [\itype{unstate(p)}] Removes $p$ from the database after
converting it to normal form.  Keywords: \type{value},
\type{decrement-p}.
 \item [\itype{empty}] Unstashes every sentence in the database.  Keyword:
\type{value}.
 \item [\itype{new-val(a v)}] Sets the truth value of the sentence
associated with the atom $a$ to $v$.  If $v$ is the truth value
\type{unknown}, the sentence is removed from the database.
 \item [\itype{mvl-file(string)}] Returns the path for the \mvl\ file
associated with the given name.  The file name may include an
extension, one may be supplied as an optional argument, or
\type{.MVL} will be used by default.  There is an optional keyword
\type{no-error} indicating that even if the file cannot be found, the
path should be returned.
 \item [\itype{mvl-load(file)}] Loads the sentences in the given file into
the \mvl\ database.
 \item [\itype{mvl-unload(file)}] Removes the sentences in the
given file from the \mvl\ database.
 \item [\itype{mvl-save(file)}] Writes the \mvl\ database out to the
given file.  This function also accepts the keywords \type{cutoffs}
and \type{props}, which are treated as for \type{contents} (see
Section \ref{s.contents}).  Finally, the keyword \type{if-exists} is
passed to the \type{with-open-file} macro.
 \end{description}

\section{Querying the MVL database}

\subsection{MVL responses to queries}
\label{s.answer}

Because of the fact that a particular fact may be stored in the
database with any of a variety of truth values, \mvl\ responds to a
query by returning not simply a binding list (as most \prolog\
interpreters do, for example), but by returning a binding list {\em
and a truth value}. 

In actuality, what is returned is an instance of the \type{answer}
structure\index{answer structure}.  This structure contains the
following fields:
 \begin{description}
 \item[\type{binding}] Bindings for some of the variables in the query.
 \item[\type{value}] A truth value.  The bound version of the query has
this truth value.
 \end{description}

The definition of this structure leads Common Lisp to define the
access functions \type{answer-binding} and \type{answer-value};
in addition, a function \type{equal-answer} is defined by \mvl\ to
test two answers for equality:
 \begin{description}
 \item[\itype{answer-binding(a)}] Returns the binding field of the
answer \type a.
 \item[\itype{answer-value(a)}] Returns the value field of the answer
\type a.
 \item[\itype{equal-answer(a1 a2)}] Returns \T\ if the two answers
\type{a1} and \type{a2} are equal, and \NIL\ otherwise.  The
functions used are \type{binding-equal} and \type{mvl-eq}, the
functions normally used to test binding lists and truth values for
equality.
 \end{description}

Thus, in response to the query
 \[\type{(flies ?x)}\]
 the response of an answer with binding \type{((?x . Tweety))} and
truth value \type{dt} would indicate that
 \[\type{(flies Tweety)}\] 
 has truth value \type{dt}.  In other words, Tweety's flying is a
conclusion that is true by default.  (See Section \ref{s.j} for a
description of default truth values in the \mvl\ system.)

As a side effect, this allows the user to distinguish a response of
\NIL\ (indicating that the query has failed) from a response in which
the binding is \NIL\ (indicating that the query has succeeded without
binding any variables).

\subsection{Searching for specific sentences}

There are two functions that are most commonly used to retrieve facts
that have been stored in the database, \itype{lookup} and
\itype{lookups}.  Each of these functions accepts as arguments a
proposition to search for and up to four keywords, and returns two
values.

The keywords and their meanings are as follows:
 \begin{description}
 \item [\itype{cutoffs}] This keyword is used to give conditions on
the sentences being considered as possible matches to the query.  Any
sentence that is considered as a match must have an associated truth
value \type{v} such that \type{(test v cutoffs)} evaluates to \T\ (see
Section \ref{s.cutoffs}).  The default value of this keyword is defined
by the bilattice in use (see Section \ref{s.bilattice}).
 \item [\itype{inv}] This keyword defaults to \NIL.  If it is \T,
however, any truth values found should be negated before they are
returned.
 \item [\apa] This is similar to the keyword for \type{stash}.
 \item [\itype{succeed-on-modal-ops}] If \T, then any attempt to look
up a proposition involving a modal operator will immediately succeed;
if \NIL\ (the default), then the result returned will be the result
of applying the modal operator to the truth values of the
propositions that are its arguments.  See Section \ref{s.modal}.
 \end{description}

The two functions \type{lookup} and \type{lookups} are identical
except that \type{lookup} looks for only a single match for the
supplied proposition, while \type{lookups} searches for all such
matches.  Thus while \type{lookup} returns two values, \type{lookups}
returns two {\em lists\/} of values.  Each entry on each of these
lists is analogous to the corresponding value returned by
\type{lookup}.  The values returned by \type{(lookup q)} are as
follows:
 \begin{enumerate}
 \item First, an \type{answer} is returned.  If this answer contains
the binding list \type b and the truth value \type v, then applying
the binding list to the query produces a result with the given truth
value.  In other words, \type{(plug q b)} is an instantiation of a
sentence appearing explicitly in the database with truth value \type
v.
 \item Second, one of the following three values is returned:
 \begin{enumerate}
 \item If the sentence was found by procedural attachment, the atom
\type{attach} is returned.
 \item If the sentence was found by applying a modal operator, the
atom \type{modal} is returned.
 \item Under other conditions, the instantiated version of the
proposition that resulted in the match is returned.
 \end{enumerate}
 \end{enumerate}

As already mentioned, \type{lookups} returns two lists of values.

Here are summaries of \type{lookup} and \type{lookups}:

 \begin{description}
 \item [\itype{lookup(q)}] Searches the database for a sentence
matching the query $q$.  Allowed keywords: \type{cutoffs}, \type{inv},
\apa, \somo.  Returns two values: an answer and the source of the
match.  Returns \NIL\ if the query cannot be found in the database.
 \item [\itype{lookups(q)}] Similar to \type{lookup}, but looks
for {\em all\/} sentences matching the query.  Returns two lists of
values.
 \end{description}

\subsection{Searching the entire database}
\label{s.contents}

In addition to searching for a particular sentence, it is also
possible to display large portions of the entire \mvl\ database.  The
basic functions for doing this are \type{contents}, which displays
the entire database, and \type{prfacts}, which accepts a \lisp\ atom
as an argument and displays all of those facts concerning that atom.

Each of these two functions also accepts a variety of keywords:
 \begin{description}
 \item [\itype{cutoffs}] As for \type{lookup}, it is possible to
supply a test that must be passed by the truth value of any database
sentence before that sentence is displayed.
 \end{description}
 In addition, \type{contents} accepts a keyword \type{props}, which
defaults to \type{(props)} (see Section \ref{s.database}) and is a
list of the sentences to be displayed.

Related to \type{prfacts} is \itype{unstash-facts-with-atom}, which
accepts an atom as an argument and removes from the database any
fact concerning that atom.  \type{Unstash-facts-with-atom} accepts
the same keywords as \type{unstash} does.

Here, then, are the functions described in this section:

\begin{description}
 \item [\itype{contents}] Displays the contents of the database.  Keywords:
\type{cutoffs}, \type{props}.
 \item [\itype{prfacts(atom)}] Displays the database facts concerning the
given atom.  Keyword: \type{cutoffs}.
 \item [\itype{unstash-facts-with-atom(atom)}] Unstashes from the
database any facts concerning the given atom.  Keywords:
\type{value}, \type{decrement-p}, \apa.
 \end{description}

\subsubsection{Command line interface}

\paragraph{Sun}
In addition to the above functions, on Sparcstations, the key
$\langle$meta$\rangle$-* may be typed to the command line processor
running from the editor.  The result is that any information in the \mvl\
database concerning the currently open predicate will be displayed to the
user.

\paragraph{Symbolics}
On Symbolics systems, the key $\langle$super-shift$\rangle$-A may be typed
to the command line processor.  The result is that any information in the
\mvl\ database concerning the currently open predicate will be displayed
to the user.  This feature is the \mvl\ analog to the
$\langle$ctrl-shift$\rangle$-A key, which displays information about the
arguments to the currently open function.

\section{The first-order prover}
\label{s.fo}

As discussed in the technical papers distributed with \mvl, inference
in \mvl\ proceeds using a conventional first-order theorem prover; in
this section, we will describe this prover and the interface to it.
Be warned that the information in this section is unlikely to be used
by other than the most sophisticated users of \mvl; most readers will
want to proceed directly to Section \ref{s.inf}.

In general, there will be a variety of invocations to the first-order
prover active at any particular time.  Each such invocation is
associated with a task that the multivalued prover needs to complete,
and each such task is a \lisp\ object that is an instance of the
\type{consider} structure.\footnote{In fact, each task is an instance
of a structure that {\em includes\/} the \type{consider} structure.
See the file \type{bc.lisp} for further details.}

The multiple invocations of the first-order prover are run in
parallel by \mvl.  To each such invocation is associated an instance
of the structure \type{proof} that is used to record temporary
information about the progress of the proof effort.  We will discuss
first the first-order prover itself, then the construction and
termination of the invocations of it, and finally the structure of
the multivalued proving tasks.

\subsection{Nature}
\label{s.simple}

The prover itself is a simple natural-deduction style complete
theorem prover.  Only inferences against the initial clauses of
database statements are allowed; the result of this is that after
applying \type{state} to a sentence of the form
 \[\type{(<= CONCLUSION PREMISE-1 ... PREMISE-N)}\]
 \type{PREMISE-1} will always be tried first when the prover attempts
to prove \type{CONCLUSION}.  In a similar way, stating the
forward-chaining rule
 \[\type{(=> PREMISE-1 ... PREMISE-N CONCLUSION)}\]
 will result in a situation where only \type{PREMISE-1} will trigger
forward chaining.

Control of the proof search is handled by a function that evaluates
the ``merit'' of each of the ongoing proof attempts, and then
attempts to select the best one.  The evaluation function is stored
in the value of the global constant \itype{*node-evaluation-fn*},
which defaults to \itype{best-first-fn}, a function that attempts to
determine the value of a proof by looking at the truth value being
accumulated.  An alternative to \type{best-first-fn} is the function
\itype{depth-fn}, which evaluates a node simply by its depth in the
proof tree.

Three global parameters are used to control the proof search:
 \begin{description}
 \item [\itype{*final-limit*}] is the limit point beyond which the
search is abandoned, so that any node that evaluates to a value worse
than this will be pruned.  The default value of this parameter is
\NIL, indicating that no nodes are to be pruned.
 \item [\itype{*initial-limit*}] is used to indicate the type of search.
If \NIL\ (the default), then the values assigned to the nodes are not
used to order the search, and depth-first search is used.
 \item [\itype{*limit-delta*}] is used if \type{*initial-limit*} is
non-\NIL.  In this case, the nodes are ordered by increasing
 \[\frac{\type{value(node)}-\type{*initial-limit*}}{\type{*limit-delta*}}\]
 A typical value for *initial-limit* is 0, and for *limit-delta* is
1.  The value of this parameter is irrelevant if
\type{*initial-limit*} is \NIL.
 \end{description}

These three parameters may be reset by the user:
 \begin{description}
 \item [\itype{set-search-parameters}] This function is the primitive
used to modify the parameters used to control the first-order theorem
prover's search.  It accepts as keywords \type{initial-limit},
\type{limit-delta}, and \type{final-limit}.
 \item [\itype{activate-breadth-first}] Activate breadth-first search
by setting \type{*node-evaluation-fn*} to \type{depth-fn}.  Allowed
keywords: \type{initial-limit}, \type{limit-delta}, and
\type{final-limit}.
 \item [\itype{activate-depth-first}] Activate depth-first search by
setting \type{*initial-limit*} to \NIL.  Allowed keyword:
\type{final-limit}.
 \item [\itype{activate-best-first}] Activate best-first search by
setting \type{*node-evaluation-fn*} to \type{best-first-fn}.  Allowed
keywords: \type{initial-limit}, \type{limit-delta}, and
\type{final-limit}.
 \end{description}

\subsection{Invoking the prover}

The prover is invoked using the function \type{init-prover}.  This
function accepts as arguments a proposition and a pointer to the task
responsible for this proof attempt.

When \type{init-prover} is invoked, an instance of \itype{proof}
is created and returned that contains auxiliary information about the
proof attempt.  Here are the fields in the structure:
 \begin{description}
 \item [\itype{active-nodes}] A list of the nodes that are still
under active consideration.
 \item [\itype{answers}] A list of answers that have been found.
 \item [\itype{indexp} and \itype{indexn}] These entries are used to
index the various subtasks generated by the prover.
 \item [\itype{initial-limit}, \itype{final-limit}, \itype{limit-delta}
and \itype{evaluation-fn}] These entries are the values of the search
control parameters used in this proof attempt.
 \item [\itype{task}] A pointer back to the task responsible
for this invocation of the first-order prover.
 \item [\itype{vars}] The variables that appear in the proposition
being proven.
 \end{description}

Each node is an instance of the structure \itype{goal}, which is
described in the file \type{first-order.lisp}.

To continue a proof process, the function \type{cont-prover} is used.
This function accepts as an argument an instance of \type{proof} that
contains the auxiliary information for the proof process being
resumed.  The value returned by \type{cont-prover} is a structure
describing the results of the proof process itself;
\type{cont-prover} returns as soon as the prover finds a single
solution to its query.  The value returned is an instance of the
\type{bcanswer} structure\index{\type{bcanswer} structure}; the
fields in this structure are as follows:
 \begin{description}
 \item[\type{answer}] This field contains a binding list and truth
value that the system has determined can be assigned to the given
instantiation of the original query.
 \item[\type{just}] This field contains a list of those database
sentences that were used in the proof.
 \end{description}

If the proof attempt fails, \NIL\ is returned.

Here is a summary of the functions described in this section:

\begin{description}
 \item [\itype{init-prover(p effort)}] This function
initializes and returns an instance of \type{proof} containing
auxiliary information about the proof process.
 \item [\itype{cont-prover(proof)}] Resume the task associated to the
auxiliary information in the given instance of the \type{proof}
structure.  This function returns an instance of \type{bcanswer}
corresponding to the result obtained by the prover, and returns as
soon as a solution to the original query is found.
 \end{description}

\subsection{The multivalued proof tasks}
\label{s.consider}

The individual proof attempts generated by a multivalued query are
instances of structures that include the structure
\type{consider}\index{consider structure}.  The exact details of this
and related structures are unlikely to be of much interest to a user
of the \mvl\ system, but we include them here in the interests of
completeness.  These are the fields in the
\type{consider} structure:
 \begin{description}
 \item [\itype{prop}] The proposition being investigated by this task.
 \item [\itype{parent}] The task for which the current task is a
child.
 \item [\itype{children}] Children of the current task.
 \item [\itype{truth-value}] The current best estimate of the truth value
that should be assigned to the proposition being considered.
 \item [\itype{relevant}] A list of answers that, if returned by this
task, are known to be relevant to the original query.
 \item [\itype{irrelevant}] A list of answers know {\em not\/} to be
relevant to the original query.
 \item [\itype{paths}] The \mvl\ system goes to a great deal of effort
to ensure that it does not consider proofs that will not affect the
eventual response to a user's query; intermediate information used in
this computation is stored in this field.  Further details can be found
in the file \type{bc.lisp} and the files to which it points.
 \end{description}

When the prover must select a task to work on next, it does so via the
following function:
 \begin{description}
 \item [\itype{choose-task}] Selects a task to work on next.  The
current implementation of \mvl\ simply selects any active task (see
Section \ref{s.bc}).
 \end{description}

\section{Inference}
\label{s.inf}

\subsection{Backward chaining}
\label{s.bc}

The information in Section \ref{s.fo} is more technical than most users
of \mvl\ are likely to need.  What is of much greater interest are those
database query tools that are like \type{lookup} and \type{lookups} but,
instead of simply looking for a fact in the database, attempt to derive
it as a consequence of other database information.  These functions use
those of the previous section as described in the accompanying technical
papers.

The functions accept three sorts of cutoff information, instead of the
simple information provided to \type{lookup} and \type{lookups}.  The
first, provided via the \type{cutoffs} keyword, is used to construct a
test to decide whether an answer computed by the prover is good enough
to be returned when the proof effort is complete.  The second,
provided via the \type{succeeds} keyword, indicates that an answer
should be returned as soon as it can be shown that the final answer
will pass the given test.  The final cutoff, supplied via the
\type{terminate} keyword, indicates that the answer should be returned
as soon as there is any indication that the final answer {\em might\/}
pass the given test.

As an example, suppose that we were trying to find a plan to achieve some
goal.  We might well be prepared to terminate our search as soon as we found
a plan that was {\em guaranteed\/} to achieve the goal; if we completed our
search {\em without\/} finding such a plan, we might then want to return
all plans that we {\em expected\/} would achieve the goal.  It is to
facilitate such possibilities that the backward-chaining functions expect
three sets of cutoff information.  The default value for \type{cutoffs}
is the standard one associated to the active bilattice (see Section
\ref{s.bilattice}); the default choice for \type{succeeds} corresponds
to complete information for the singular versions of the queries and
guarantees failure for the plural versions (since an answer might be
about to be discovered for a different instantiation of the query).
The default value for \type{terminate} always fails.

Here, then are the various backward-chaining functions.  They all
return two values: an answer (or list of answers for the plural
versions \type{bcs} and \type{consider}) and an instance of the
\type{analysis} structure describing tasks remaining to be considered
by the prover.  In addition, they all accept an optional
\itype{analysis} keyword that, if supplied, should be an instance of
the \type{analysis} structure that is to be continued instead of a
new one being created.  This is to allow the user to continue proof
efforts that have returned for some reason.
 \begin{description}
 \item [\itype{bc(q)}] Attempts to prove the query $q$ from information
in the database, looking only for a single answer.
 \item [\itype{bcs(q)}] Attempts to prove the query $q$ from information
in the database, looking for all possible answers.
 \end{description}

All of these functions work by maintaining an instance of the
\type{analysis} structure\index{\type{analysis} structure} that
includes a variety of information about the status of the current
proof attempt.  Here are the fields in this structure:
 \begin{description}
 \item [\type{prop}] The proposition being considered by the prover.
 \item [\type{cutoffs}] The cutoff information that must be satisfied
by any result before it is to be returned.
 \item [\type{succeeds}] The cutoff information that will allow the
prover to terminate early if an answer is guaranteed to pass the test.
 \item [\type{terminate}] The cutoff information that will allow the
prover to terminate early if an answer might pass the test.
 \item [\type{active-tasks}] A list of the proof tasks (see Section
\ref{s.consider}) that still need to be investigated.
 \item [\type{main-task}] The task that was created upon the 
original invocation of the prover.
 \item [\type{returned-answers}] A list of answers that have already
been returned to the user; this is to keep \mvl\ from compulsively
re-returning these answers if the \type{analysis} is reinvoked.
 \end{description}

The top-level functions using this structure are as follows:
 \begin{description}
 \item [\itype{init-analysis(prop)}] The \mvl\ primitive upon which
the proving functions are based, this function accepts a proposition
and creates and returns an instance of the \type{analysis} structure.
 \item [\itype{cont-analysis(analysis)}] This function continues the
investigation of a specific analysis, and can be used to resume a
proof effort that has been curtailed for some reason.  It returns two
values: an instance of the \type{answer} structure giving the results
of the proof attempts, and a modified version of the \type{analysis}
structure that can be used to continue the proof process if desired.
 \end{description}

The functions \type{init-analysis} and the user functions such as
\type{bc} all accept the following keywords:
 \begin{description}
 \item [\type{cutoffs}] Cutoff test as described above.
 \item [\type{succeeds}] Success test as described above.
 \item [\type{terminate}] Termination test as described above.
 \item [\type{initial-value}] The truth value assigned to the query in
the absence of database information.
 \item [\type{analysis}]  If supplied, use this analysis instead of
creating a new one.
 \end{description}

\subsection{Forward chaining}
\label{s.fc}

Forward chaining in \mvl\ is not as simple as it is in many other
logic programming languages because of the \nm\ nature of multivalued
inference.  It follows from this that inserting a new sentence into
the database may cause a previous conclusion to be retracted, and
this may in turn cause other sentences to become unsupported.

There are two functions that use this information to modify the
database.  The more common is called \itype{fc} and accepts as an
argument a proposition to be added to the database.  The proposition
is added and provided that it was not already in the database when the
forward chainer was invoked, its consequences are also added to the
database and forward chained upon.  The \type{fc} function accepts
the following keywords:
 \begin{description}
 \item [\itype{value}] This keyword is similar to that used by
\type{stash} and \type{state}.
 \item [\itype{increment-p}] This keyword is also similar to that used
by \type{stash} and \type{state}.  The default value is \NIL.
 \item [\itype{fc-increment-p}] When the information content of some
sentence in the database has {\em fallen}, should this drop be
propagated to other sentences that appear to depend on the sentence
in question?  The default value is \T.
 \item [\itype{cutoffs}] Only forward chain if the sentence being
inserted into the database has a truth value that passes the test
corresponding to \type{cutoffs}.  (See Section \ref{s.cutoffs}.)  The
default value used is the one associated with the active bilattice.
 \item [\itype{include-modal}] If \T\ (the default), then the forward
chainer also looks for rules that fire when the truth value of a
modal expression involving $p$ changes (where $p$ is the sentence
being forward chained upon).  See Section \ref{s.modal}.  The reason
this keyword is included is that searching for rules of this form is
computationally expensive.
 \end{description}

Slightly less common is \itype{erase}, the function used to remove a
proposition from the database.  This function accepts a database
sentence as an argument, and also the keywords \itype{value},
\itype{decrement-p} and \itype{cutoffs}, which are similar to the
corresponding keywords for \type{fc}.

Summarizing:
 \begin{description}
 \item [\itype{fc(p)}] Inserts the sentence $p$ into the database and
forward chains.  Allowed keywords: \type{value}, \type{increment-p},
\type{fc-increment-p}, \type{cutoffs} and \type{include-modal}.
 \item [\itype{erase(p)}] Removes $p$ from the database and forward
chains.  Allowed keywords: \type{value}, \type{decrement-p} and
\type{cutoffs}.
 \end{description}

\subsection{Debugging tools}

Debugging facilities are limited in \mvl.  The facilities that exist
describe the prover's activities during backward-chaining,
forward-chaining, or while the first-order theorem prover is looking
for a conventional proof of some proposition.

\subsubsection{Backward chaining}

Diagnostics are displayed during backward chaining using the two
global variables \itype{*bc-trace*} and \at\index{*anytime-trace*}.

Whenever a task is encountered that involves a proposition that is an
instance of some element of \type{*bc-trace*}, information concerning
the task will be displayed to the user.  The variable \at\ is
similar, but only reports when an {\em inference\/} task is
encountered.  An inference task is one that is trying to prove
something, as opposed to one analyzing a modal operator; this
distinction is described in greater detail in the file {\type
bc.lisp}.  In practice, the distinction between \type{*bc-trace*} and
\at\ is that the former should be used for debugging purposes, while
the latter should be used if you are interested in watching the
prover approximate and then refine its eventual answer to some
declarative query.

Although the user can manipulate \type{*bc-trace*} and \at\ directly,
there are four functions written to help do so:
 \begin{description}
 \item [\itype{bc-trace(pat)}] Include \type{pat} on
\type{*bc-trace*}, so that any instance of \type{pat} will be
traced.
 \item [\itype{bc-untrace(pat)}] Remove any instance of \type{pat}
from \type{*bc-trace*}.
 \item [\itype{anytime-trace(pat)}] Include \type{pat} on \at, so
that any instance of \type{pat} will be traced.
 \item [\itype{anytime-untrace(pat)}] Remove any instance of
\type{pat} from \at.
 \end{description}
 For all of these functions, the argument is optional.  If omitted,
\type{(?*)} (of which everything is an instance) is used.

\subsubsection{Forward chaining}

Tracing the activity of the forward chainer is essentially identical
to tracing the behavior of the backward chainer; the global variable
\type{*fc-trace*} is used and is manipulated using the functions:
 \begin{description}
 \item [\itype{fc-trace(pat)}] Include \type{pat} on
\type{*fc-trace*}, so that any instance of \type{pat} will be
traced.
 \item [\itype{fc-untrace(pat)}] Remove any instance of \type{pat}
from \type{*fc-trace*}.
 \end{description}
 Once again, the argument is optional.

\subsubsection{The first-order prover}

Tracing the activity of the first-order theorem prover itself is
somewhat different.  The facility for doing this is similar to the
Common Lisp \type{trace} facility, but instead of reporting when a
function is invoked, it reports when an attempt is made to prove an
instance of a particular predicate, and also when such an instance is
found successfully in the database.

The functions used to control this are \itype{bc-watch} and
\itype{bc-unwatch}.  They are similar to \type{trace} and
\type{untrace}:

 \begin{description}
 \item [\itype{bc-watch(preds)}] Accepts as arguments any number of
predicates that should be watched as the prover proceeds.  Each
predicate can also be a list of the form \type{(predicate depth)}
where \type{depth} is the number of successive ancestors to be
displayed when the prover encounters an instance of the predicate.
(This can be used to understand {\em why\/} a particular clause is
relevant to the proof being attempted.)  \type{bc-watch} returns a
list of the predicates currently being {\sc watch}ed.
 \item [\itype{bc-unwatch(preds)}] Removes the given predicates from
the list of {\sc watch}ed predicates.  If no arguments are supplied,
all predicates are un{\sc watch}ed.
 \end{description}

There are two additional features involved in tracing the activities
of the first-order prover.  The first is used whenever the global
parameter \itype{*save-nodes*} is \T\ (which it is by default).  When
the first-order prover displays information regarding a particular
proof node, the node is labelled with a number.  Subsequent nodes
will then use this label if they need to refer to the node in
question (e.g., as a parent).

If \type{*save-nodes*} is \T, then \type{*break-node*} can be set
to an integer; when a node is displayed that is labelled with this
integer, a \type{(break)} will be executed.  This is to allow the
user to interrupt the first-order prover if he needs to examine the
proof stack or some other \mvl\ internal.

 \begin{description}
 \item [\itype{*save-nodes*}] If \T\ (the default), nodes are
labelled with consecutive integers when they are displayed by the
first-order prover.
 \item [\itype{*break-node*}] The number of a node at which the
proof process should be interrupted.
 \end{description}

\section{Procedural attachment}
\label{s.attach}

In some cases, it may be desirable to store a fact in the database
{\em without\/} using the functions supplied in \mvl.  Suppose, for
example, that we have an $n+1$-ary relation $R$ that corresponds to
an $n$-ary function $f$, so that $R(x_1,\dots,x_n,y)$ holds just in
case $f(x_1,\dots,x_n)=y$.  Now, if we assert $R(c_1,\dots,c_n,d)$
into the database, where $d$ and the $c_i$ are constants, we should
probably also {\em remove\/} any sentences of the form
$R(c_1,\dots,c_n,d')$, where $d'\neq d$.  This sort of a situation
might arise if we were monitoring the temperature in a room, and
represented the observed temperature using a database sentence such
as
 \begin{equation}
\type{(TEMPERATURE ROOM 72)}
\label{e.temp}
\end{equation}

We handle this by informing \mvl\ that to stash a sentence of the
form \type{(TEMPERATURE ...)}, it should use not the usual function
\type{stash}, but some other user-defined function, say
\type{unique-stash}.  This new function would first \type{unstash}
any sentence that agreed with the given one except for the value of
the final term, and would then invoke \type{stash} on the original
sentence, but with \apa\ set to \NIL\ so that \type{unique-stash} was
not invoked a second time.

To modify the function used by \mvl\ to stash sentences such as this,
the user should simply place the desired function (in this case,
\type{unique-stash}) on the \type{stash} property of the relation in
question (in this case, \type{TEMPERATURE}).\footnote{As described in
Section \ref{s.invocation}, the \type{stash} property of a relation
can also be an association list, if multiple invocations of the
system are running in parallel.} The system will then automatically
invoke \type{unique-stash} instead of \type{stash} when a fact such
as (\ref{e.temp}) is asserted.

The functions within \mvl\ that can be modified in this fashion are
\type{stash}, \type{unstash}, \type{lookup} and \type{lookups}.  In
addition, the procedural attachments for \type{lookup} and
\type{lookups} are linked, so that if \type{lookup} has a procedural
attachment and \type{lookups} does not, a call to \type{lookups} will
invoke the procedural attachment for \type{lookup} and then pluralize
the result.  Similar remarks hold if \type{lookups} and not
\type{lookup} has a procedural attachment.

The procedural attachment is effected using the following functions:
 \begin{description}
 \item [\itype{invoke(fn p)}] This function invokes \type{fn} on the
sentence $p$.  In addition, any keyword arguments are passed to
\type{fn}, but with \type{\&allow-other-keys} so that the keywords
will be ignored if \type{fn} is not expecting them.
 \item [\itype{lookup-call(fn p)}] This function is similar to
\type{invoke}, but modifies the values returned to make them agree
with those expected to be returned by \type{lookup}.  Unless the
function \type{fn} supplies them, the truth value is assumed to be
\type{true}.  The second value is \type{attach}.
 \item [\itype{lookups-call(fn p)}] The plural version of
\type{lookup-call}.
 \end{description}

\subsection{Facilities to create procedural attachments}

Some facilities exist to create procedural attachments within the
\mvl\ system.  The following two macros are used for this purpose:
 \begin{description}
 \item [\itype{lisp-predicate(pred)}] This macro is used to indicate
that the given predicate can be evaluated by a \lisp\ call.  It will
also cause \mvl\ to signal an error if an attempt is made to stash
or unstash an instance of the function.
 \item [\itype{lisp-defined(fn)}] This macro is used to indicate that
the given function can be evaluated by a \lisp\ call.  It will also
cause \mvl\ to signal an error if an attempt is made to stash or
unstash an instance of the function.  \type{Lisp-defined} also
accepts an optional argument, which is an inverter function.  The
overall result of attempting to lookup an instance of the function
can be described as follows:
 \begin{enumerate}
 \item If none of the arguments to the function (i.e., all but the
last of the arguments of the relation) is a variable, then the function
is called and the result is unified with the final argument of the
relation.
 \item If the inverter function has been supplied and is equal to \T,
then the function is called and the result is unified with the final
argument of the relation even if some of the arguments are variables.
 \item If some of the arguments are variables but the result is not a
variable, then the inverter function is invoked if one has been
supplied.  The arguments passed to the inverter function are the
result (i.e., the last argument of the relation), and a list of all
of the arguments (i.e., all but the last argument of the relation).
If there is no inverter function, then \NIL\ is returned.  In
addition, if \wbli\ \index{*warn-on-bad-lisp-invocations*} is not
\NIL, the user is warned.  (The parameter \wbli\ defaults to \NIL.)
 \item If both some of the arguments and the result are variables, then
\NIL\ is returned.  In addition, if \wbli\ is not \NIL, the user is
warned.
 \end{enumerate}
 \end{description}

The current version of \mvl\ invokes \type{lisp-predicate} on the
following \lisp\ predicates: \type{numberp}, $<$, $>$, \type{equal},
\type{null}, \type{atom}, \type{listp}, \type{varp}, \type{varp?},
\type{varp$*$}, \type{groundp}, and $=$.  The procedural attachment
to $=$ attempts to unify its arguments.

The macro \type{lisp-defined} is invoked on the following:
\type{car}, \type{cdr}, \type{cadr}, \type{cons}, \type{list},
\type{list*}, \type{revappend}, \type{append}, \type{reverse},
\type{length}, \type{assoc}, \type{delete}, \type{+}, \type{*},
\type{-}, \type{/}, \type{floor}, \type{ceiling}, \type{truncate},
\type{round}, \type{float}, \type{mod}, \type{rem}, \type{max},
\type{min}, \type{sin}, \type{cos}, \type{atan}, \type{sqrt}, 
and \type{eval}.

\subsection{Other procedurally attached relations}

In addition to the \lisp\ functions and predicates described in the
previous section, the following relations have been defined and
are procedurally attached to suitable \lisp\ functions:
 \begin{description}
 \item [\itype{(bagof}\type{ ?x ?p ?b)}] This relation holds for $x$,
$p$ and $b$ if and only if $b$ is a list of all $x$ satisfying $p$.
 \item [\itype{(done exp-1 ... exp-n)}] This predicate, which is
useful for stashing only, causes \lisp\ to evaluate each of the subsequent
expressions.
 \item [\itype{(value symbol val)}] Stashing this predicate sets the
given symbol to the value \type{val}.  Unstashing it makes the symbol
unbound if its current value unifies with \type{val}.
 \item [\itype{(property symbol value indicator)}] Stashing this
predicate sets the symbol's \type{indicator} property to \type{value}
for the current invocation of \mvl.  Unstashing it removes the
\type{indicator} property for the current invocation of \mvl\ if the
previous value of this property unifies with \type{value}.  (See the
definitions of \type{get-mvl} and \type{rem-mvl} in Section
\ref{s.invocation}.)
 \item [\itype{(known p)}]  Stashing this predicate simply stashes $p$;
unstashing it causes an error.  Looking it up looks up $p$.
 \item [\itype{(unknown p)}]  Stashing this predicate unstashes $p$;
unstashing it causes an error.  Looking it up fails if $p$
can be found in the database, and succeeds otherwise.
 \item [\itype{(member item list)}] Tries to unify \type{item} with
the various elements in \type{list}.
 \end{description}

\section{Testing the MVL system}
\label{s.foln}

\subsection{As a diagnostic tool}

There are two functions available to help you to test both \mvl\ and
any applications you write using it.  The first, \type{test-mvl}, is
one you should already have run; if \mvl\ has been installed successfully,
it should return 0.

\type{Test-mvl} works by invoking the generic testing function
\type{mvl-test-file} on the file \type{mvl.test} in the \mvl\
directory.\footnote{\type{Test-mvl} works in its own invocation of
\mvl, so that the existing database is not modified.  This is not true
of \type{mvl-test-file}, however.} This file also contains
documentation regarding the details of the various tests being
executed.

In general, \type{mvl-test-file} accepts as input a file name, and
looks for a file with an extension \type{.TEST} on either the \mvl\
or the current directory.  It then works through the file by
evaluating the first s-expression it contains, and comparing the
value returned to the second s-expression.  If these two
s-expressions match, the next s-expression is evaluated, and so on.
If there is a mismatch, \type{mvl-test-file} reports the error;
eventually, it returns the total number of errors found in the entire
file.  When doing the comparison, \type{**} acts as a wildcard.
Thus, if the second s-expression is simply \type{**}, it will always
match; if it is \type{(** ** **)}, it will match whenever the first
expression returns a list of three values, and so on.

\begin{description}
 \item [\itype{test-mvl}] Tests the \mvl\ system by invoking the function
\type{mvl-test-file} on the file \type{MVL.TEST}.
 \item [\itype{mvl-test-file(string)}] Can be used to evaluate
various forms by the \mvl\ system and compare the results to
predefined values.  Both the forms and the expected answers are in
the \type{.TEST} file whose name is supplied to the function.
Returns the number of errors found.
 \end{description}

\subsection{As a tutorial}

The information used by the function \type{test-mvl} can also be used
as a tutorial in the \mvl\ system.  Take a look at the file
\type{mvl.test} that this function uses; you will see that the first
two lines are
 \begin{verbatim}
     (notime (mvl-test-file "bc" "backward chainer")
     0
\end{verbatim}
 This means that the function \type{mvl-test-file} will be invoked on
the file \type{bc.test}, and that this function should return 0
(i.e., 0 errors).  So let's have a look at \type{bc.test}:
 \begin{verbatim}
     (load-logic #\f)                    ;first-order logic
     **
\end{verbatim}
 Here, the system loads the bilattice corresponding to first-order
logic and is not interested in the value returned by the loading
operation.  Next, a declarative database corresponding to a full
adder is loaded (with the timer turned off).  The sentences in this
database can be found in the file \type{f1.mvl} and are also
described in Section 15 and Appendix B.2 of the paper ``Multivalued
logics: A uniform approach to inference in artificial intelligence.''

Next, we see in \type{mvl.test} the two lines
 \begin{verbatim}
     (get-bdg '?x (answer-binding (bc '(val (out 1 f1) 1 ?x))))
     OFF                   ;the first output (the sum line) should be off
\end{verbatim}
 We invoke the backward chainer on the sentence
 \[\type{(val (out 1 f1) 1 ?x))}\]
 to see if we can prove that the first output of the full adder takes
some value \type{?x}.  Then we get the binding of this variable to
see what the value actually is -- \type{OFF}, in this case.  The
remainder of the various test files can be analyzed similarly; many
of the examples in these files can also be found in the various
technical papers included in the \mvl\ distribution package.

\section{System support}
\label{s.support}

\subsection{Functions that compile and load the MVL system}

The following three functions can be used to manipulate the \mvl\
system as a whole:
 \begin{description}
 \item [\itype{load-mvl}] Loads the \mvl\ system.
 \item [\itype{compile-mvl}] Compiles and loads the \mvl\ system.  On
an Allegro system, the keyword \type{recompile} forces a recompile of
all of the files and the keyword \type{reload} is ignored (since the
files will all be reloaded anyway).  On a Symbolics system, the
keyword \type{reload} forces all of the files to be reloaded and the
keyword \type{recompile} forces them all to be recompiled.
 \item [\itype{release-mvl}] (Symbolics only)  Compiles and loads the
\mvl\ system, and declares the new version to be \type{released}.  The
keyword \type{no-recompile} can be used to skip the compilation step.
 \end{description}

\subsection{The MVL package}
\label{s.package}

In order to avoid name conflicts with existing code, the \mvl\ system
resides in its own package, \type{MVL}\index{MVL package}.  All
of the symbols and functions appearing in the appendix are exported
from this package and can therefore be accessed by invoking the
function \type{use-package(mvl)}.  When a new bilattice is loaded,
the modal operators it defines are also exported. 

To use the \mvl\ system, invoke the function \type{(use-package 'mvl)}
to make these symbols accessible.\footnote{As an alternative,
\type{(in-package 'mvl)} wil put you in the MVL package itself.}

\subsection{Running several copies of MVL in parallel}
\label{s.invocation}

Some users may wish to run several copies of the \mvl\ system
simultaneously, perhaps associating them with distinct {\sc lisp}
processes or with individual agents.  (The use of multiple copies of
\mvl\ to represent distinct components of a single agent's knowledge
goes against the grain of the multivalued approach and is not
recommended.)  In order to support this facility, the following
functions and macros have been defined:
 \begin{description}
 \item [\itype{existing-mvl-invocation}] This function returns a
single value that is an instance of the structure \mi.  It is useful
primarily for saving the current copy of \mvl\ so that it can be
restored later.
 \item [\itype{create-mvl-invocation}] This function sets up a new
copy of \mvl.  It accepts an optional argument that is an instance of
\mi; this can be used to reset the state of \mvl\ to a previous one.
If the optional argument is omitted, default values are used.  The
value returned is the \mvl\ invocation that was active before the
function was called, so that another call to
\type{create-mvl-invocation} can be used to restore the original
environment.
 \item [\itype{with-mvl-invocation}] This is a macro similar to the
Common Lisp macro \type{with-open-file}.  It accepts an optional
argument that is an instance of \mi\ and is followed by a body of
forms (presumably including a \type{load-logic}) that are evaluated
in a locally-bound copy of \mvl.
 \end{description}

These functions work by locally or globally rebinding the values of
the various parameters used by the \mvl\ system
(\type{*active-bilattice*}, \type{true} and so on).  The only
difficulty with this is that procedural attachments for a relation
$R$ use the property list of $R$, and this cannot be rebound.

In order to cater to this difficulty, there is an additional global
variable \itype{*mvl-invocation*} that labels the currently active
copy of \mvl.  To determine what procedure is attached to the
\type{stash} operation for $R$, the \type{stash} property of $R$ is
examined.  There are now the following possibilities:
 \begin{enumerate}
 \item If the value is \NIL, \type{stash} is not procedurally attached,
so the usual \mvl\ function is called.
 \item If the value is an atom, then it is assumed that this atom is
the name of a function that is the procedural attachment for \type{stash}
independent of the copy of \mvl\ that is currently active.
 \item If the value is a list, it is treated as an association list
and a cons whose car is \type{*mvl-invocation*} is searched
for.
 \begin{enumerate}
 \item If such a cons is found, its cdr is used as the procedural
attachment to \type{stash}. 
 \item If no such cons is found, but there is a cons whose car is \T,
then the cdr of that cons is used.
 \item If no cons has either \type{*mvl-invocation*} or \T\ as its
car, then \type{stash} is assumed to be unattached.
 \end{enumerate}
 \end{enumerate}
 The other functions described in Section \ref{s.attach} are handled
similarly.

The association lists described above are manipulated using the following
functions:
 \begin{description}
 \item [\itype{get-mvl (symbol indicator)}]  Like the Common Lisp
\type{get}, but if the value is a list, treats it as an a-list and
searches for an item whose car matches \type{*mvl-invocation*}.
The value taken can be changed with \type{setf}.
 \item [\itype{rem-mvl (symbol indicator)}] Like the Common Lisp
\type{remprop}, but if the value is a list, treats it as an a-list
and removes the item whose car matches \type{*mvl-invocation*}.
 \end{description}

\section{Selecting a bilattice}
\label{s.bilattice}

\subsection{Bilattice structure}

Of course, most of the power of the \mvl\ system comes not from the
features that were described in Sections \ref{s.fol1}--\ref{s.support}
of this manual, but from the fact that the user can change
representation language conveniently within the system.  He does this
by changing the {\em bilattice\/} that he is using.

A bilattice is the mathematical structure from which the truth values
used by the system are selected, and the formal ideas underlying this
approach are discussed at length in the accompanying technical
papers.  In the \mvl\ system, a bilattice is described as an instance
of the \type{bilattice} structure.  This structure has 27 fields; of
these, 17 are used to provide the formal information required by the
system, and the others are used to simplify the user interface.  Here
are 13 of the fields used to describe the bilattice itself:

\begin{description}
 \item [\itype{true}] A \lisp\ object that is the truth value assigned
to a sentence that is unconditionally true, such as the tautology
\type{(OR (P) (NOT (P)))}.
 \item [\itype{false}] A \lisp\ object that is the truth value assigned
to a sentence that is unconditionally false, such as the negation of a
tautology.
 \item [\itype{unknown}] A \lisp\ object that is the truth value assigned
to a sentence about which nothing is known.
 \item [\itype{bottom}] A \lisp\ object that is the truth value assigned
to a sentence that is both true and false.  If any sentence is assigned
this truth value (or can be shown to have it), it indicates the presence
of a contradiction in the database.
 \item [\itype{eq}] A \lisp\ function that takes two truth values as
arguments and returns \T\ if they are equal and \NIL\ if they are
not.  This function is needed because it may be possible to represent
a single element of the bilattice as a \lisp\ object in a variety of
fashions.
 \item [\itype{and}] A \lisp\ function that takes two truth values as
arguments and returns the truth value corresponding to their
conjunction.  If the two sentences $p$ and $q$ are labelled with
truth values $v_1$ and $v_2$ respectively, then the conjunction of
the truth values is the truth value that can be assigned to $p\wedge q$.
 \item [\itype{or}] A \lisp\ function that takes two truth values as
arguments and returns the truth value corresponding to their
disjunction.
 \item [\itype{not}] A \lisp\ function that takes a truth value as its
argument and returns the truth value corresponding to its negation.
 \item [\itype{plus}] A \lisp\ function that takes two truth values as
arguments and returns the truth value corresponding to their ``sum.''
If there is reason to believe that a sentence $p$ should be labelled
with truth value $v_1$, and independent reason to believe that it
should be labelled with truth value $v_2$, then the actual label is
given by the sum $v_1+v_2$.  This operation therefore corresponds to
combination of evidence.
 \item [\itype{dot}] A \lisp\ function that is the dual of the
\type{plus} operation.  It takes two truth values $v_1$ and $v_2$ as
arguments and returns the greatest truth value that could
subsequently be strengthened to a conclusion of either $v_1$ or
$v_2$.  This is a ``consensus'' operation.
 \item [\itype{kle}] A \lisp\ function that takes two truth values $x$
and $y$ and determines whether $x+y=y$.  Although this can also
be done using the bilattice $+$ and $=$ operations, this field can be
provided to give a more efficient implementation of this operation.
This field is optional.
 \item [\itype{tle}] A \lisp\ function that takes two truth values $x$
and $y$ and determines whether $x\vee y=y$.  This field is also
optional.
 \item [\itype{modal-ops}] A list of the modal operators associated
with the given bilattice.  See Section \ref{s.modal}.
 \end{description}

Three of the remaining formal slots are concerned with the incorporation
of binding information into truth values.  In some cases, when a
proposition is retrieved from the database, it is only an instance of
the database proposition that is used, and the truth value may have to
be modified to reflect that:
 \begin{description}
 \item [\itype{plug}]  This \lisp\ function accepts as arguments
a truth value and a binding list, and returns a potentially modified
truth value.  The default value is a function that leaves the truth
value unchanged.
 \end{description}

There is also a function that accepts a truth value and returns the
variables appearing within it:
 \begin{description}
 \item [\itype{vars}]  This \lisp\ function accepts a truth value
as its argument and returns a list of the variables appearing in it.
The default value is a function that returns \NIL\ independent of its
argument.
 \end{description}

The function \itype{absorb} is needed for subtle technical reasons
when working with bilattices with non-\NIL\ plugging functions;
further information on this can be found in the file \type{load.lisp}.
All three of these functions are \NIL\ if there is no interaction
between the bilattice truth values and the variables that they
contain.

The final formal field has to do with search control:
 \begin{description}
 \item [\itype{simplicity}] This \lisp\ function accepts a truth
value (which will frequently be the value obtained by conjoining the
premises in a particular proof attempt) and returns a figure
indicating the ``merit'' that should be assigned to the proof attempt
in question.  The value returned should be a nonempty list of
integers and is viewed lexicographically, with high values being
preferred to low ones.  The default value is a function that returns
the list \type{(0)} independent of the truth value supplied.
 \end{description}

The remaining ten fields in the \type{bilattice} structure are as
follows:
 \begin{description}
 \item [\itype{dws}] This \lisp\ function, ``dot with star,'' accepts
as arguments two truth values, $v_1$ and $v_2$.  The truth value returned
is $v_1$, modified so as to remove any contribution that could be
attributed to a sentence with truth value $v_2$.
 \item [\itype{stash-val}] A function that, given a sentence $p$,
returns the default truth value to be passed to \type{stash}.
 \item [\itype{unstash-val}] A function that, given a sentence $p$,
returns the default truth value to be passed to \type{unstash}.
 \item [\itype{cutoff-val}] The default cutoff point that must be
exceeded by any sentence if it is to be returned by \type{lookup}.
See Section \ref{s.cutoffs}.
 \item [\itype{print-fn}] This function is used to pretty-print
truth values in the associated bilattice.  Its three arguments
are a truth value, the bilattice in question, and a stream.
 \item [\itype{long-desc}] A long description of the bilattice, used
for input and output only.
 \item [\itype{short-desc}] A short description of the bilattice, also
used for input and output only.
 \item [\itype{char}] A one-character identifier for the bilattice.
This character can be used to change the bilattice using the function
\type{load-logic} described in Section \ref{s.change}.
 \item [\itype{type}] An overall description of the mathematical form
of the bilattice.  Perhaps it is a product of two other bilattices,
or formed from them in some other fashion.  This defaults to \type{basic}
if no other value is supplied; various methods for constructing new
bilattices from old ones are described in Section \ref{s.construct}. 
 \item [\itype{components}]  A list of the components used to construct
a bilattice whose \type{type} is not \type{basic}.
 \end{description}

\subsection{The active bilattice}

At any particular point in time, there will be a single bilattice
that has been selected by the user to be used in the application he
is developing.  This bilattice is stored as the value of
\type{*active-bilattice*}; in addition, many of the slots of the
active bilattice are directly available to the user.  Here they
are:

\begin{description}
 \item [\itype{*active-bilattice*}] The currently active bilattice.
 \item [\itype{true}] The truth value \type{true} in the active bilattice.
 \item [\itype{false}] The truth value \type{false} in the active bilattice.
 \item [\itype{unknown}] The truth value \type{unknown} in the active
bilattice.
 \item [\itype{bottom}] The truth value \type{bottom} in the active bilattice.
 \item [\itype{std-cutoffs}] Cutoff information constructed from the
value of \type{cutoff-val} in the active bilattice.  See Section
\ref{s.cutoffs}.
 \item [\itype{stash-value}] The function \type{stash-val} in the
active bilattice.
 \item [\itype{unstash-value}] The function \type{unstash-val} in the
active bilattice.
 \item [\itype{mvl-plus}] The function \type{plus} in the active bilattice.
 \item [\itype{mvl-dot}] The function \type{dot} in the active bilattice.
 \item [\itype{mvl-and}] The function \type{and} in the active bilattice.
 \item [\itype{mvl-or}] The function \type{or} in the active bilattice.
 \item [\itype{mvl-not}] The function \type{not} in the active bilattice.
 \item [\itype{mvl-eq}] The function \type{eq} in the active bilattice.
 \item [\itype{mvl-simplicity}] The function \type{simplicity} in the active
bilattice.
 \item [\itype{dot-with-star}] The function \type{dws} in the active 
bilattice.
 \item [\itype{mvl-vars}] The function \type{vars} in the active bilattice.
 \item [\itype{mvl-plug}] The function \type{plug} in the active bilattice.
 \item [\itype{mvl-print}] This function accepts a single argument, a truth
value, and modifies it in a way that will cause it to be pretty-printed.
The value returned by mvl-print should not be used for any purpose other
than printing.
 \end{description}

All of the functions above accept an optional argument that can be
used to evaluate their result using some bilattice other than the
active one.

\subsubsection{Functions derived from the active bilattice}

There are a variety of functions associated with the functions that
are defined when a new bilattice is made active.  In order to
simplify the description of these functions, we will denote
\type{mvl-and} by $\wedge$, \type{mvl-or} by $\vee$, \type{mvl-not}
by $\neg$, \type{mvl-plus} by $+$, and \type{mvl-dot} by $\cdot$.
The first auxiliary function defined is \type{dot-with-not}:
 \begin{description}
 \item [\itype{dot-with-not(x)}] Computes $x\cdot \neg x$.
 \end{description}

The other functions are constructed from the partial orders defined
using the bilattice operations.  We will write $x\leq_t y$ just in
case $x\wedge y = x$, for example, and $x\leq_k y$ just in case
$x\cdot y = x$.  (All of this notation follows that in the paper,
``Multivalued logics: A uniform approach to inference in artificial
intelligence''.)

Given these definitions, the following \lisp\ functions are defined:
 \begin{description}
 \item [\itype{k-le(x y)}] Returns \T\ if $x \leq_k y$, and \NIL\ otherwise.
 \item [\itype{k-ge(x y)}] Returns \T\ if $x \geq_k y$, and \NIL\ otherwise.
 \item [\itype{t-le(x y)}] Returns \T\ if $x \leq_t y$, and \NIL\ otherwise.
 \item [\itype{t-ge(x y)}] Returns \T\ if $x \geq_t y$, and \NIL\ otherwise.
 \item [\itype{k-lt(x y)}] Returns \T\ if $x <_k y$, and \NIL\ otherwise.
 \item [\itype{k-gt(x y)}] Returns \T\ if $x >_k y$, and \NIL\ otherwise.
 \item [\itype{t-lt(x y)}] Returns \T\ if $x <_t y$, and \NIL\ otherwise.
 \item [\itype{t-gt(x y)}] Returns \T\ if $x >_t y$, and \NIL\ otherwise.
 \item [\itype{k-not-le(x y)}] Returns \T\ if $x \not\leq_k y$, and
\NIL\ otherwise.
 \item [\itype{k-not-ge(x y)}] Returns \T\ if $x \not\geq_k y$, and
\NIL\ otherwise.
 \item [\itype{t-not-le(x y)}] Returns \T\ if $x \not\leq_t y$, and
\NIL\ otherwise.
 \item [\itype{t-not-ge(x y)}] Returns \T\ if $x \not\geq_t y$, and
\NIL\ otherwise.
 \item [\itype{k-not-lt(x y)}] Returns \T\ if $x \not<_k y$, and \NIL\
otherwise.
 \item [\itype{k-not-gt(x y)}] Returns \T\ if $x \not>_k y$, and \NIL\
otherwise.
 \item [\itype{t-not-lt(x y)}] Returns \T\ if $x \not<_t y$, and \NIL\
otherwise.
 \item [\itype{t-not-gt(x y)}] Returns \T\ if $x \not>_t y$, and \NIL\
otherwise.
 \end{description}

All of these functions accept an optional argument that is the
bilattice to be used in the comparison.  This argument defaults to
the currently active bilattice.

\subsubsection{Cutoff information}
\label{s.cutoffs}

We will refer to each of the 16 functions defined at the end of the
previous section (e.g., excluding \type{dot-with-not}) as a {\em tag}
\index{tag}.  Thus \type{k-le} is a tag, and so on.

We will now define a {\em simple test} to be a dotted pair
 \[\type{(v . tag)}\]
 where \type{v} is a truth value and \type{tag} is a tag.  We will
say that a truth value \type{x} {\em satisfies\/} the simple test
if and only if the function \type{tag}, when applied to the arguments
$x$ and $v$ (in that order), returns \T.

As an example, consider the simple test
 \[\type{(true . t-ge)}\]
 This test will be satisfied by any truth value that is $\geq_t
\type{true}$.  Since \type{true} itself is the only truth value
satisfying this test, the test will be satisfied for \type{true} and
no other truth values.  This is the test that is frequently used to
decide to terminate a proof effort.\footnote{Tests in the \mvl\ system
are in fact \lisp\ structures with \type{value} and \type{tag} fields;
we are using dotted lists in this manual only to make the presentation
clearer.}

Alternatively, the test used to determine whether or not to return
an answer after the theorem prover runs to completion is often given
by
 \[\type{(unknown . k-gt)}\]
 This test will succeed for any truth value about which more is known
that \type{unknown} -- in other words, for any truth value other than
\type{unknown} itself.

In general, the simple tests are combined into {\em cutoffs}.  A cutoff
is simply a list of lists of simple tests, and is interpreted as being
in conjunctive normal form.  Thus the cutoff corresponding to the simple
test $s$ is given simply by
 \[\type{((s))}\]

A cutoff corresponding to the conjunction of tests $s_1,\dots,s_n$ is given
by
 \[\type{((s-1) ... (s-n))}\]
 while that corresponding to their {\em disjunction\/} is given by
 \[\type{((s-1 ... s-n))}\]

In general, we can build up conjunctions of disjunctions of simple
tests in as complex a manner as we wish.  The cutoff \NIL\
corresponds to the empty conjunction and therefore always succeeds;
the cutoff \type{(NIL)} contains an empty disjunction and therefore
always fails.  For convenience, these two cutoffs are stored in the
global constants \itype{*always-succeeds*} and
\itype{*never-succeeds*} respectively.

Given a truth value $x$ and a cutoff $c$, the expression \type{(test x c)}
evaluates to \T\ if $x$ passes the test, and to \NIL\ if it fails it.
There is also a function that negates cutoffs:
 \begin{description}
 \item [\itype{test(value cutoffs \&optional bilattice)}] Returns \T\
if the truth value passes the test associated to \type{cutoffs}, and
\NIL\ otherwise.  The test is performed for the supplied bilattice,
or for the currently active one if the optional argument is omitted.
 \item [\itype{cutoffs-not(cutoffs \&optional bilattice)}] Returns a
new cutoffs $c$ such that
 \[\type{(test x c)} = \type{(test (mvl-not x) cutoffs)}\]
 for all truth values $x$.  In other words, $x$ passes $c$ if and
only if $\neg x$ passed the cutoffs supplied originally.
 \end{description}

\subsection{Selecting a new bilattice}
\label{s.change}

To change the bilattice being used, simply invoke the function
\type{load-logic}.  There are three ways to specify the new bilattice:
 \begin{enumerate}
 \item Do nothing.  You will then be asked to select the bilattice
from the list of bilattices known to the system.
 \item Supply the one-character identifier of the new bilattice as an
argument to \type{load-logic}.  The new bilattice must already be
known to the system.
 \item Supply the new bilattice itself as an argument to
\type{load-logic}.
 \end{enumerate}
 \type{Load-logic} empties the database before the bilattice is
changed.

The bilattices known to the system are kept in the global variable
\type{*bilattices*}.  To add a new bilattice to \type{*bilattices*},
invoke the function \type{bilattice} with the bilattice as an
argument.  This function will warn you if you are overwriting another
bilattice that has the same one-character designator.

Summarizing:
 \begin{description}
 \item [\itype{*bilattices*}] A list of the bilattices known to the system.
 \item [\itype{bilattice(b)}] Add the bilattice $b$ to \type{*bilattices*}.
 \item [\itype{load-logic}] Empty the database and make a new bilattice
active.  Accepts as an optional argument either a character designating
a bilattice, or the new bilattice itself.  If no argument is supplied,
the user is prompted for the new bilattice.
 \end{description}

\section{Bilattices available in the MVL system}

Here are the bilattices available in the \mvl\ system as distributed.
Each is discussed more fully in a subsequent section:
 \begin{description}
 \item [\itype{*first-order-bilattice*}] The bilattice corresponding to
first-order logic.
 \item [\itype{*atms-bilattice*}] The bilattice corresponding to an \atms,
where the justification information is maintained in disjunctive normal
form as suggested by de Kleer.
 \item [\itype{*cnf-atms-bilattice*}] An alternative version of the \atms\
bilattice where the justification information is maintained in conjunctive
normal form.
 \item [\itype{*default-bilattice*}] A bilattice that can be used for simple
default reasoning.
 \item [\itype{*hierarchy-bilattice*}] A bilattice similar to the
default bilattice but where the defaults are labelled with
integer-valued strengths.
 \item [\itype{*first-order-atms-bilattice*}] This bilattice corresponds
to a fully first-order \atms, as opposed to the more conventional one
dealing only with propositional logic.
 \item [\itype{*circumscription-bilattice*}] The bilattice used to compute
the results of the circumscription axiom.
 \item [\itype{*probability-bilattice*}] An extended version of the \atms\
bilattice that maintains probabilistic information as well.
 \item [\itype{*time-bilattice*}] A compound bilattice that can be used
for reasoning about events occurring against a temporal background.
 \end{description}

\subsection{First-order logic} \index{first-order logic}

The simplest bilattice is that corresponding to first-order logic, and
contains only the four truth values \type{true} (displayed as \type
T), \type{false} (displayed as \type F), \type{unknown} (\type U) and
\type{bottom} (\type{T/F}).  This bilattice is identified using the
character {\bf f}, and is active when \mvl\ is first loaded.

\subsection{Assumption-based truth maintenance}
\index{assumption-based truth maintenance}

Except for the \nm\ and probabilistic \atms\ bilattices, all of the
bilattices corresponding to \atms's have truth values that are dotted
pairs \type{(j . k)}, where $j$ is a justification for the sentence
being labelled, and $k$ is a justification for its negation.  The
bilattices differ in that the justifications are represented
differently.

\subsubsection{The ATMS bilattice}

For the simplest \atms\ bilattice, the justifications are simply
disjunctive normal form expressions.  Thus the justification
 \begin{equation}
\type{((P23 P37) (P1))}
\label{e.just}
\end{equation}
 is interpreted to mean that the sentence in question is a
consequence either of \type{P1}, or of \type{P23} and \type{P37} in
conjunction.

This bilattice is identified by the character {\bf a}.

\subsubsection{ATMS's in conjunctive normal form}

Alternatively, we can represent the justifications in conjunctive
normal form, as opposed to disjunctive normal form.  This representation
is useful, for example, in diagnosing multiple faults, since we want to
be able to explain an observed failure {\em either\/} in terms of the
collective failures of one set of components, or the collective failures
of another set, and so on.

This bilattice is identified by the character {\bf b}.

\subsubsection{First-order ATMS's}
\label{s.fo-atms}

It is also possible to extend these ideas slightly to obtain a fully
first-order \atms.  In order to do this, we need simply modify
justifications such as that appearing in (\ref{e.just}) so that the
primitive elements of the justification are not atoms referring to
propositions, but pairs
 \[\type{(atom . binding-list)}\]
 where the atom refers to a proposition and the binding list
indicates precisely what instantiation of that proposition is
actually needed to justify the sentence in question.

As an example, suppose that an assumption in our database is \type{P1},
given by
 \begin{equation}
\type{(FRIENDLY ?X)}
\label{e.assume}
\end{equation}
 and indicating that everything is friendly.  If we also have the rule
of inference
 \[\type{(<= (HAPPY ?X) (FRIENDLY ?X))}\]
 indicating that anything friendly is happy, and this rule is unconditionally
true (i.e., not an assumption), then the truth value assigned to the
query \type{(HAPPY HARRY)} will be
 \[\type{((((P1 . ((?X . HARRY))))))}\]
 indicating that Harry is happy by virtue of (\ref{e.assume}) applied
to $x=\mbox{Harry}$, and that there is no justification for the
conclusion that Harry is not happy.  (Recall that all of the \atms\
bilattices combine justification for and against each sentence into a
single truth value.)

This bilattice is identified by the character {\bf o}.

\subsubsection{Probabilistic ATMS's}
\label{s.prob}

Truth values in the probabilistic \atms\ bilattice consist of dotted
pairs, where the \type{car} of a truth value is simply an element of
the \atms\ bilattice and the \type{cdr} is the probability that the
given assumption is valid.

When entering an assumption, the user will be queried for the
probability in question; the probabilities of compound sentences are
then computed automatically by the system.  The algorithm used is
that described in Grnarov, Kleinrock and Gerla in the 1979
Proceedings of the Computer Networking Symposium.

The probabilistic \atms\ bilattice is identified by the character
{\bf p}.
  
\subsection{Default reasoning} \index{default reasoning}
\label{s.j}

The default bilattice consists of seven truth values.  In addition to
the usual elements \type{true}, \type{false}, \type{unknown} and
\type{bottom}, there are elements \type{dt}, \type{df} and \type{*}.

The truth value \type{dt} labels a conclusion that is true by default
(such as Tweety flying in the canonical example), while \type{df} labels
a conclusion that is {\em false\/} by default.  The final value,
\type{*}, labels conclusions that are both true {\em and\/} false
by default, as in the well-known Nixon-Republican-Quaker example.

This bilattice is identified by the character {\bf d}.

\subsubsection{Prioritized default reasoning}

Another version of the default bilattice allows for priorities among
the defaults.  In this bilattice, truth values are of the form
 \[\type{(n . x)}\]
 where $n$ is a nonnegative integer and $x$ is an element of the
first-order bilattice.  Truth values with $n=0$ are displayed as for
the first-order bilattice; with $n=1$ as for the default bilattice.
For $n>1$, we might see \type{D(3)T} or something like it.  This
particular truth value is the same as
 \[\type{(3 . true)}\]
 Truth values with small $n$ override those with larger $n$.
\type{Unknown} in this bilattice is given by the truth value \NIL.
The one-character identifier for the hierarchical default bilattice is
{\bf h}.

\subsection{Circumscription} \index{circumscription}
\label{s.circ}

The default bilattices are not able to deal with all of the
situations in which multiple extensions arise because single truth
values such as \type{*} label sentences that are true in some
extensions and false in others without providing any information
about precisely which extensions are which.  In order to handle this
difficulty, we need to work with a more complex bilattice.

One way to view the default bilattice is as a copy of the first-order
bilattice in which the point \type{unknown} is interpreted as, ``not
known with certainty,'' and is in fact replaced with another copy of
the first-order bilattice that is used to record uncertain
information (i.e., default truth, default falsity, or the presence of
a default contradiction).

The circumscriptive bilattice is formed similarly, except that
instead of simply recording default truth, falsity, and so on, we
record a {\em justification\/} for default conclusions.  This is done
by replacing the point \type{unknown} of the first-order bilattice
not with another copy of the first-order bilattice, but with a copy
of the first-order \atms\ bilattice.  (This idea is in fact a general
way to construct new bilattices from old ones, and is called a {\em
surgery}.  It is discussed in Section \ref{s.surgery}.)  The
designator for the circumscriptive bilattice is {\bf c}.
\type{Circum} and \type{circums} use this bilattice to accept a
sentence as an argument and attempt to prove it from the
circumscription axiom; they resemble \type{bc} and \type{bcs}.

The truth values in the circumscription bilattice corresponding to
certain information are of the form \type{(T . x)}, where $x$ is an
element of the first-order bilattice.  Those corresponding to uncertain
information are of the form \type{(NIL . y)}, where $y$ is an element
of the first-order \atms\ bilattice.

Finally, the circumscriptive theorem prover will output diagnostics
concerning its progress if the global variable \type{*circum-trace*}
has a non-\NIL\ value.

Here is a summary of the circumscriptive facilities:
 \begin{description} 
 \item [\itype{*circum-trace*}] If not \NIL, the circumscriptive proof
functions will output diagnostics describing their progress.
 \item [\itype{circum(p)}] Like \type{bc}, but looks for instances of
$p$ that follow from the circumscription axiom.  Allowed keywords:
\type{cutoffs} and \type{succeeds}.
 \item [\itype{circums(p)}] Like \type{bcs}, but looks for instances of
$p$ that follow from the circumscription axiom.  Allowed keywords:
\type{cutoffs} and \type{succeeds}.
 \end{description}

\subsection{Temporal reasoning}\index{temporal reasoning}

\type{*Time-bilattice*} is a bilattice of the truth values that might
be assigned to propositions that have a time-dependent nature, such as
whether or not a particular switch is on.  The truth values represent
the entire histories of the propositions, so that a single truth value
is a complete function from a set of time points (taken to be the
integers) into a set of ``base'' truth values (take to be elements of
the nonmonotonic {\sc atms} bilattice). This bilattice is discussed at
greater length in the paper, ``Computational considerations in
reasoning about action,'' and a sample use of this bilattice can be
found by examining the file \type{temporal.test} that is included with
\mvl.

\section{Facilities to create new bilattices}
\label{s.construct}

There are a variety of ways in which it is possible to construct
new bilattices from old ones, or from old ones in conjunction
with instances of some other structure.  In this section, we
describe the extent to which these facilities exist within the
\mvl\ system.

\subsection{Products}
\label{s.product}

Perhaps the simplest way in which a new bilattice can be constructed
is as the product of two existing bilattices.  If $A$ and $B$ are
bilattices, then a truth value in the product bilattice $A\times B$
will simply be a pair of truth values \type{(x . y)}, where $x$ is a
truth value in $A$ and $y$ is one in $B$.  All of the bilattice
operations on the product bilattice are computed pointwise using the
corresponding operations on the component bilattices $A$ and $B$.

The product construction is performed using the following function:
 \begin{description}
 \item [\itype{product-bilattice(b1 b2)}] Returns the bilattice
that is the product of the bilattices $b_1$ and $b_2$.
 \end{description}

\subsection{Functional attachments}

Another way in which to construct a bilattice is to start with a
bilattice $B$ and a function $f$ that maps $B$ into a fixed set $S$.
The elements of the new bilattice are simply pairs of the form
 \[\type{(x .~(f x))}\]
 where $x$ is an element of $B$.  The function $f$ here is simply
tagging along for the ride, and has no impact on the bilattice
operations.  It can be used to provide additional information about
the truth values, and is used to compute probabilistic information in
the bilattice {\bf p} described in Section \ref{s.prob}.

The relevant function is the following one:
 \begin{description}
 \item [\itype{append-bilattice(b f)}] Returns the bilattice obtained
by appending the result of applying the function $f$ to the points of
the original bilattice $b$.
 \end{description}

\subsection{Surgeries}
\label{s.surgery}

A somewhat more interesting construction is the one we described in
Section \ref{s.circ} as a {\em surgery}.  Here, we take one bilattice
$b_1$ and replace the point corresponding to \type{unknown} in this
bilattice with another bilattice $b_2$.  The truth values in the
resulting bilattice are either of the form \type{(T . x)} for a point
$x$ in the ``upper'' bilattice $b_1$, or \type{(NIL . y)} for a point
$y$ in the lower bilattice.  Truth values from $b_1$ take priority over
values from $b_2$.

To perform this construction generally, we have: 
 \begin{description}
 \item [\itype{splice-bilattice(upper lower)}] Returns the bilattice
obtained by replacing the point \type{unknown} in the bilattice
\type{upper} with a copy of the bilattice \type{lower}.
 \end{description}

As an example, the default bilattice {\bf d} of Section \ref{s.j} is
given by
 \[\type{(splice-bilattice *first-order-bilattice* *first-order-bilattice*)}\]
 The circumscriptive bilattice of Section \ref{s.circ} is given by
 \[\type{(splice-bilattice *first-order-bilattice*
*first-order-atms-bilattice*)}\]

\subsection{Hierarchies}

It is possible to iterate the surgery construction of the previous
section, producing a bilattice that is an infinite descending
``chain'' of copies of a basic bilattice $b$.  This is done by
the function \type{hierarchy-bilattice}:
 \begin{description}
 \item [\itype{hierarchy-bilattice(b)}] Returns the bilattice constructed
as an infinite descending chain of copies of $b$.
 \end{description}

Truth values in this bilattice have the form \type{(n . x)}, where
$x$ is an element of $b$, and $n$ is an integer giving the hierarchy
level.  The top hierarchy level of the bilattice has $n=0$; the
second level is $n=1$, and so on.  The point \type{unknown} is given
by \NIL.

This function is used to construct the hierarchical version of the
default bilattice.

\subsection{Lattices}

Bilattices can also be constructed from lattices.  A lattice is a set
equipped with a {\em single\/} partial order and greatest lower bound
and least upper bound operations, which we denote by $+$ and $\cdot$
respectively.

An element of the associated bilattice is a {\em pair\/} of elements
of the lattice, and the bilattice operations are given by:
 \begin{eqnarray*}
\neg(a,b) & = & (b,a) \\
(a,b) + (c,d) & = & (a+c,b+d) \\
(a,b) \cdot (c,d) & = & (a\cdot c,b\cdot d) \\
(a,b) \vee (c,d) & = & (a+c, b\cdot d) \\
(a,b) \wedge (c,d) & = & (a\cdot c,b+d)
 \end{eqnarray*}

Roughly speaking, the first component of the bilattice element labels
the ``truth'' of the sentence, and the second labels its falsity.  This
is how the \atms\ bilattices are constructed.

The function used is \type{lattice-to-bilattice}:

 \begin{description}
 \item [\itype{lattice-to-bilattice(l)}] Returns the bilattice constructed
from the lattice $l$.
 \end{description}

A lattice is a structure rather like a bilattice, but with fewer slots
(because it only has a single partial order).  The fields in an instance
of the \itype{lattice} structure are:

\begin{description}
 \item [\type{max}] A \lisp\ object that is the maximal element of
the lattice.
 \item [\type{min}] A \lisp\ object that is the minimal element of
the lattice.
 \item [\type{eq}] A \lisp\ function that takes two lattice elements
as arguments and returns \T\ if they are equal and \NIL\ if they are
not.
 \item [\type{and}] A \lisp\ function that takes two lattice elements
as arguments and returns their greatest lower bound.
 \item [\type{or}] A \lisp\ function that takes two lattice elements
as arguments and returns their least upper bound.
 \item [\type{le}] This optional field is a \lisp\ function that
accepts two arguments $x$ and $y$ and determines whether or not
$x\leq y$ under the partial order corresponding to the lattice.
 \item [\type{vars}] This \lisp\ function accepts a lattice element
as an argument and returns a list of the variables that it contains.
The default value returns \NIL\ independent of its argument.
 \item [\type{plug}] This \lisp\ function accepts as arguments a
lattice value and a binding list, and returns a potentially modified
lattice value.  The default value is a function that leaves the lattice
value unchanged.
 \item [\type{simplicity}] The lattice analog of the bilattice
\type{simplicity} slot.  When constructing a bilattice from a
lattice, the \type{simplicity} of a truth value $(a,b)$ is the result
of applying the simplicity function from the lattice to the lattice
element $a$.  The value $b$ is ignored in the simplicity computation.
 \item [\type{dws}] This \lisp\ function, ``dot with star,'' accepts
as arguments two lattice values, $v_1$ and $v_2$.  The lattice
element returned is $v_1$, modified so as to remove any contribution
that could be attributed to the lattice element $v_2$.
 \item [\type{stash-val}] The lattice element entering into
the bilattice function \type{stash-val}.
 \item [\type{long-desc}] A long description of the lattice, used
for input and output only.
 \item [\type{short-desc}] A short description of the lattice, also
used for input and output only.
 \item [\type{char}] A one-character identifier for the lattice.
 \end{description}

\subsection{Directed acyclic graphs}
\label{s.dag}

The final fashion in which \mvl\ supports the construction of new
bilattices is an extension of the product construction of Section
\ref{s.product}.  There, we formed the product of two bilattices
$A\times B$.  If the bilattices were identical, we would have formed
the product bilattice $B\times B$, or $B^2$.  Another way to view
this bilattice is as the collection of functions from the two point
set $\set{1,2}$ into the original bilattice $B$.  Since such a
function $f$ can be described as the pair $(f(1),f(2))$, there is no
difference between this description and the previous one.

Given this analogy, it is not hard to see that we can extend the
construction to the bilattice $B^S$ for an arbitrary set $S$, where
elements of this bilattice are functions from the set $S$ into $B$,
and all of the bilattice functions ($\wedge$, $\cdot$ and so on) are
once again computed pointwise.

The difficulty with this idea is that for large sets $S$, it is not
clear how we can represent a function from $S$ into $B$ in a
convenient way.  It turns out, however, that if $S$ has the structure
of a directed acyclic graph, or \dg, this problem can be overcome.

\subsubsection{Defining the DAG}

As a preliminary, we need some way to define the \dg\ itself.  In
\mvl, a \dg\ \index{dag} is an instance of the \type{dag} structure;
this structure has the following fields:
 \begin{description}
 \item [\itype{root}] The root of the \dg.
 \item [\itype{eq}] A function that accepts two elements of the \dg\
and returns \T\ if they are the same element and \NIL\ otherwise.
 \item [\itype{leq}] A function that accepts two elements of the \dg\
and returns \T\ if the first is farther from the root of the \dg\
than the second.  If not supplied by the user, this function can by
computed from the \type{dot} operation on the \dg.
 \item [\itype{dot}] A function that accepts two elements of the \dg\
and returns a list of their first common descendants, or \NIL\ if no
such descendant exists.
 \item [\itype{long}] A long description of the \dg.
 \item [\itype{short}] A short description of the \dg.
 \item [\itype{inherit}] A function describing the fashion in which
truth values are inherited within the \dg.  This function $f$ should take
four arguments: a bilattice $b$, two elements of the \dg\ $x$ and $y$,
and an element $z$ of the bilattice.  If $z$ is the element of the
bilattice associated to the \dg\ element $x$, then $f(b,x,y,z)$ will
be the value associated to the \dg\ element $y$ in the absence of
information to the contrary.  By default, this function ignores
its first three arguments and assigns $y$ the same value as $x$ has.
 \item [\type{vars}] This \lisp\ function accepts a \dg\ element
as an argument and returns a list of the variables that it contains.
The default value returns \NIL\ independent of its argument.
 \item [\type{plug}] This \lisp\ function accepts as arguments a
\dg\ value and a binding list, and returns a potentially modified
\dg\ value.  The default value is a function that leaves the \dg\
value unchanged.
 \item [\type{select}] In some applications (none of which is distributed
with \mvl), it appears that truth values should be inherited from only
{\em some\/} of the elements of the \dg\ which might potentially
impact the value of the function at some specific point.  The
selection function is given a point where the truth value is being
computed and a list of points that apparently contribute; it returns a
list of points that {\em actually\/} contribute.  By default, all
points contribute; most users of \mvl\ are very unlikely to need this
facility.
 \end{description}

\subsubsection{Constructing the bilattice}

Having selected a \dg\ $D$ and a bilattice $B$, we can now construct
the bilattice $B^D$.  Specific functions from $D$ to $B$ are described
as structures with the following slots:
 \begin{description}
 \item [\type{dag}] The domain of the function, $D$.
 \item [\type{bilattice}] The range of the function, $B$.
 \item [\type{list}] A list of elements of the form \type{(p . v)},
where \type p is a point where the value changes in a fashion {\em
other\/} than the one predicted by the inheritance function associated
with $D$ and \type v is the value taken there.
 \end{description}
 We will refer to these functions as {\em \dg-functions}.

There are a variety of functions available to manipulate \dg-functions,
and these are used to construct the bilattice $B^D$.  In some cases,
however, the user may want to access these functions directly.  Here
are some of them:
 \begin{description}
 \item [\itype{get-val(d f)}] Returns the value of the
function $f$ at the point $d$ of the \dg.
 \item [\itype{make-tv(dag bilattice d b)}] Returns a \dg-function that
takes the value $b$ at the point $d$.  If $d$ is not the root of the
\dg, the value at the root of the \dg\ is generally the point
\type{unknown} in the bilattice, but the user may supply the keyword
\type{root-fn} if he wishes, in which case the value at the root of
the \dg\ is the result of applying \type{root-fn} to the bilattice.
 \item [\itype{make-root(dag bilattice b)}] Returns a \dg-function that
takes the value $b$ at the root of the \dg.
 \item [\itype{dag-change-tv(f b-f)}] The argument
\type{b-f} should be a function from the bilattice to itself; a
modified version of the original \dg-function $f$ is returned in which
every point in the \dg-function has had the function \type{b-f}
applied to it.  There is an optional keyword \type{simplify}; if
supplied, the computed function is simplified before being returned.
(See \type{simplify}.)
 \item [\itype{dag-change-dag(f d-f)}] The argument
\type{d-f} should be a function from the \dg\ to itself; a modified
version of the original \dg-function $f$ is returned in which every
point at which the \dg-function takes a new value has had the function
\type{d-f} applied to it.  The bilattice argument is expected to be
the bilattice of functional truth values as opposed to the bilattice
that is the target of the functions.  There is an optional keyword
\type{simplify}; if supplied, the computed function is simplified
before being returned.  (See \type{simplify}.)
 \item [\itype{dag-list-of-answers(f)}] Given a \dg-function $f$, assumed
to be from the \dg\ of binding lists into the currently active bilattice, this
function constructs a list of answers corresponding to the points where
the \dg-function takes values.
 \item [\itype{dag-accumulate-fn(dag bilattice list)}] Given a \dg, a
bilattice, and a list of elements of the form
 \[\type{(dag-elt . bilattice-elt)}\]
a \dg-function is constructed that takes the given values at the
given elements of the \dg.

A keyword to \type{dag-accumulate-fn} is \type{modifier}.  This
function is used if the \type{d-list} contains the same \dg\ element
more than once, in which case the modifier function is called with
the new and old values (in that order) and the result is used as the
value of the \dg-function at that particular point.  The default
action is to simply overwrite the old value with the new one.  (In
other words, the default modifier function simply returns its first
argument while ignoring the second.)
 \item [\itype{dag-prune(f pred)}] Given a \dg-function
$f$ and a predicate on the bilattice, this function prunes from $f$
all points $d$ where $f(d)$ fails to satisfy the predicate.  There is
an optional keyword \type{modify-fn} that, if supplied, should be a
function $m$ from the bilattice to itself.  In this case, the points
where $m(f(d))$ fails to satisfy the predicate are pruned.
 \item [\itype{combine-fns(f1 f2 fn comp)}]
This function takes the two \dg-functions $f_1$ and $f_2$ and
combines them using the function \type{fn}, which should be a
function that takes two elements of the bilattice and returns a
single one.  The comparison function \type{comp} is used to determine
if one of the two functions can be returned directly; if, for example
\type{(comp f1 f2)} holds, than  we simply return $f_1$, and similarly
for $f_2$.  As an example, if the combining function is the bilattice
operation $\cdot$, the comparison function is \type{k-le}.
 \item [\itype{combine-1(f params f1 ... fn)}] This is an
$n$-ary version of \type{combine-fns}, where $f$ is an $n$-ary
bilattice function and there are $n$ \dg-functions supplied as well.
The additional argument \type{params} indicates which of the $n$ arguments
to this function are not necessarily \dg-functions, but are instead parameters
expected by the overall function $f$.  This additional argument is a list
of \T\ (if the $n$th argument is a parameter) or \NIL\ (if it isn't).
If the parameter list ends early, all of the remaining entries are
assumed to be \NIL.
 \item [\itype{add-dag(d b f)}] Given the \dg-function
$f$, add the value $b$ at the point $d$.  This function works by side
effect only.  There is a keyword \type{modifier} that, if supplied,
should be a function that takes two elements of the bilattice and
returns a single one.  Given such a function $m$, the \dg-function is
changed so that the new value at $d$ is $m(b,f(d))$, where $f(d)$ is
the old value.  (This is similar to the same keyword in
\type{dag-accumulate-fn}.)  The keyword \type{simplify} is also available.
 \item [\itype{simplify(f)}] In some cases,
manipulating the \dg-function $f$ may cause it to contain
redundancies, where the value at some point is presented explicitly
even though this value is in fact the one that would normally be
inherited from its parents.  This function returns a modification of
$f$ where these irrelevant entries have been removed.
 \item [\itype{fn-and(f1 f2)}] This function conjoins
the values of $f_1$ and $f_2$ using the bilattice $\wedge$ operation.
The new \dg-function is returned.
 \item [\itype{fn-or(f1 f2)}] This function disjoins
the values of $f_1$ and $f_2$ using the bilattice $\vee$ operation.
The new \dg-function is returned.
 \item [\itype{fn-plus(f1 f2)}] This function adds
the values of $f_1$ and $f_2$ using the bilattice $+$ operation.
The new \dg-function is returned.
 \item [\itype{fn-dot(f1 f2)}] This function combines
the values of $f_1$ and $f_2$ using the bilattice $\cdot$ operation.
The new \dg-function is returned.
 \item [\itype{fn-dws(f1 f2)}] This function combines
the values of $f_1$ and $f_2$ using the bilattice operation
\type{dot-with-star}.  The new \dg-function is returned.
 \item [\itype{fn-not(f1)}] This function negates the
values of $f_1$ using the bilattice $\neg$ operation.  The new
\dg-function is returned.
 \item [\itype{fn-eq(f1 f2)}] This function compares
the values of $f_1$ and $f_2$ using the bilattice comparison operation.
Either \T\ or \NIL\ is returned.
 \item [\itype{fn-comp(f1 f2 fn)}] This functions compares
the values of $f_1$ and $f_2$ using the bilattice comparison function
\type{fn}.  Either \T\ or \NIL\ is returned.
 \end{description}

Of course, the point of all of this is to enable us to define the
following function:
 \begin{description}
 \item [\itype{dag-bilattice(dag bilattice)}] Given a bilattice $B$
and a dag $D$, returns the bilattice $B^D$.
 \end{description}

The following functions are also available:
 \begin{description}
 \item [\itype{lattice-to-dag(lattice)}] Construct a \dg\ from the
given lattice.
 \item [\itype{lattice-to-dag-to-bilattice(lattice bilattice)}]  First
construct a \dg\ $D$ from the given lattice.  Then use the bilattice
$B$ to construct the bilattice $B^D$.
 \end{description}

\subsubsection{The DAG of theories}
\label{s.theory}

Two specific \dg s are provided for use in the \mvl\ system.  The first
uses a \dg\ corresponding to the existence of multiple theories within
a single database; the second uses a \dg\ to handle binding information.

The theory \dg\ consists of a hierarchy of theories.  In a single
database, for example, we might want to have rules that applied to
celestial bodies, rules applying to stars, rules applying only to red
giants or only to our sun, and so on.  The truth value of a rule
applicable to all celestial bodies should be inherited if we are
considering our sun (but it should be possible to overwrite it with
a different truth value if we wish), but a rule applicable to red
giants should not be inherited in this fashion.  So we see that our
\dg\ construction is well-suited for this sort of an application.

To activate the theory mechanism, simply invoke the \lisp\ function
\type{activate-theory-mechanism}.  This function replaces the active
bilattice, say $B$, with the bilattice constructed of functions from
the theory \dg\ into $B$.  The original bilattice can still be
accessed, since it is the second \type{component} of the newly formed
bilattice.  The function \type{theory-active-p} can be used to
determine whether or not the theory mechanism has been activated.

The \dg\ of theories is the value of the global variable
\type{theory-dag}, and has as its root the theory whose name is the
value of the constant \type{*global*}.  To indicate that a theory
\type{big} includes another theory \type{little} as a subtheory,
invoke the function 
 \[\type{(includes big little)}\]
 Thus, for example, one would introduce a new theory called
\type{theory} by writing
 \[\type{(includes *global* 'theory)}\]
 The effects of \type{includes} can be reversed by invoking the
function \type{unincludes}.

Note that since the theory \dg\ is used by the functions that
manipulate the bilattice $B^D$, it is important that this \dg\ be
maintained as new theories are introduced.  If \type{includes} is not
used to ``locate'' these new theories in the overall theory \dg\
itself, the theories will be inaccessible to many of the \mvl\
functions.

At any point in time, any new fact will be inserted into the database
as true in a particular theory; the theory into which it is inserted
is the value of \type{*active-theory*}.  The value of this variable
should not be changed directly, but should instead be modified by
invoking the function \type{activate} on the theory into which
subsequent facts should be written.  The variable
\type{true-in-theory} contains the truth value of a sentence that is
true in the currently active writing theory.

The theory \type{*active-theory*} is also the theory that is active
for reading at any point in time.  Of course, any parents of the
active theory are also active by virtue of the structure of the theory
\dg\ itself; if it is desired to read from multiple theories
simultaneously, a new theory should be created that is a child of all
of them.

Here is a summary of the functions described so far:
 \begin{description}
 \item [\itype{activate-theory-mechanism}] Activates the theory
mechanism.  Accepts as an optional argument either a one-character
bilattice descriptor, or a bilattice.
 \item [\itype{theory-active-p}] Returns \T\ or \NIL\ depending on
whether or not the theory mechanism has been activated.
 \item [\itype{theory-dag}] The \dg\ of theories.
 \item [\itype{*global*}] The root of the \dg\ of theories.
 \item [\itype{includes(big little)}] This function is used to state that
the theory \type{little} is a subtheory of the theory \type{big}, and
should inherit from it.
 \item [\itype{unincludes(big little)}] This function is used to state
that the theory \type{little} is no longer a subtheory of the theory
\type{big}, and should not inherit from it.
 \item [\itype{*active-theory*}] The theory that is currently active
for reading and writing.
 \item [\itype{activate(th)}] Activate the theory \type{th} for reading
and writing.
 \item [\itype{true-in-theory}]  The truth value of a sentence that is
true in the currently-active writing theory.
 \end{description}

Finally, five functions have been added to simplify the user
interface somewhat.  The function {\tt theory-\-contents} displays
the contents of a given theory, and \type{empty-theory} removes all
of the propositions in this theory.  To modify the contents of a
specific theory, the following functions can be used:
 \begin{description}
 \item [\itype{state-in-theory(p)}] Adds the proposition $p$ to the
currently active theory.  Allowed keywords: \type{theory} can
be used to specify a different theory to which $p$ should be added;
\type{value} and \type{increment-p} have the same meaning as for
\type{state}.  (See Section \ref{s.database}; if these keywords are
not used, the function \type{state} can be used equally well, since
the default stash value with the theory mechanism activated also
inserts the proposition into the currently active theory.)
 \item [\itype{unstate-from-theory(p)}] Removes the proposition $p$
from the currently active theory.  Allowed keywords:
\type{theory} can be used to specify a different theory from which
$p$ should be removed; \type{value} and \type{decrement-p} have the
same meaning as for \type{unstate}.  (See Section \ref{s.database}.)
 \item [\itype{erase-from-theory(p)}] Removes the proposition $p$
from the currently active theory, and forward chains on the result.
Allowed keywords: \type{theory} can be used to specify the theory from
which $p$ should be removed; \type{value}, \type{decrement-p} and
\type{cutoffs} have the same meaning as for \type{erase}.  (See
Section \ref{s.fc}.)
 \item [\itype{theory-contents()}] Displays the contents of the
currently active theory, which can be changed using the
keyword \type{theory}.  The keywords \type{cutoffs} and \type{props}
are allowed and are treated as for \type{contents}.  (See Section
\ref{s.contents}.)  There is an additional keyword,
\itype{this-theory-only}, which defaults to \T\ and, if supplied,
indicates that only sentences whose truth values {\em change\/} in the
given theory are to be displayed.  (This avoids displaying all
sentences in the global theory, for example.)
 \item [\itype{empty-theory()}] Empties the currently active
theory, although this theory may be changed using the keyword
\type{theory}.  The keyword \type{value} has a meaning similar to its
meaning as a keyword to \type{erase}.
 \end{description}

\subsubsection{The DAG of binding lists}
\label{s.binding-dag}

The other \dg\ supplied with the \mvl\ system is used for binding
lists since these, too, admit a natural inheritance structure.  The
root of the \dg\ is the empty binding list \NIL, and inheritance in
the \dg\ is determined by subsumption.  The bilattice constructed as
functions from this \dg\ into \type{*active-bilattice*} is used
extensively by the prover, and is available to the user as
\type{bdg-to-truth-val}.

The binding \dg\ itself is known as \type{binding-dag}, and is constructed
using the following functions:
 \begin{description}
 \item [\itype{equal-binding(b1 b2)}] Checks to see if two binding
lists are effectively the same.
 \item [\itype{append-binding-lists(b1 b2)}] Computes the least
specific binding list that is more specific than both $b_1$ and $b_2$
or returns \NIL\ if no such binding list exists ($b_1$ and $b_2$ may
bind the same variable to different constants, for example).
 \item [\itype{binding-le(b1 b2)}] Returns \NIL\ unless $b_1$ is a more
specific binding that $b_2$, in which case a list of relative binding
lists is returned.
 \end{description}

The following global variables are also used:
 \begin{description}
 \item [\itype{binding-dag}] The \dg\ of binding lists.
 \item [\itype{bdg-to-truth-val}] The bilattice of functions from
the \dg\ of binding lists into \type{*active-bilattice*}.
 \end{description}

\section{Modal operators}
\label{s.modal}

An additional feature that makes the \mvl\ system unique is its
ability to handle {\em modal operators}.  The approach taken by \mvl\
to modal operators is that they are functions from the bilattice of
truth values to itself.  Thus, for example, the modal operator that
takes a sentence $p$ into the sentence, ``I know that $p$,''
corresponds to the bilattice function (on the first-order bilattice)
that takes $t$ and $\bot$ into $t$ (since we know $p$ if we know it
to be true or know it to be true and false both), and takes $f$ and
$u$ into $f$ (since we do {\em not\/} know $p$ if we either know it
to be false or do not know if it is true or false).  Further details
can be found in the paper, ``Bilattices and modal operators.''

A modal expression in \mvl\ is any sentence of the form 
 \[\type{(op p-1 \dots\ p-n)}\]
 where \type{op} is recognizable as a modal operator because it has a
non-\NIL\ value for its \type{modal-op} property in the currently
active invocation of \mvl.  In this case, the value of the
\type{modal-op} property should be an instance of the structure
\itype{modal-operator}.  This structure has the following fields:
 \begin{description}
 \item [\itype{name}] The name of the modal operator.  This is used
to identify modal operators in different bilattices that have the
same commonsense meaning (and can therefore be usefully combined when
the bilattices are combined as in the previous section).
 \item [\itype{function}] A function that accepts as arguments the
truth values of the sentences under the scope of the modal operator
and returns the truth value of the modal expression itself.
 \item [\itype{parameters}] If the modal operator accepts $n$
arguments, this should be a list of $n$ entries indicating whether
each argument is a truth value or a parameter that should be passed
directly to the modal operator function.  (An entry of \T\ indicates
that the associated argument is parametric; an entry of \NIL\
indicates that it isn't.)  Suppose, for example, that we are working
with a bilattice where the truth values are functions describing the
behaviors of various fluents over the course of time, and have a modal
operator \type{delay}, where the truth value assigned to
 \[\type{(delay t p)}\]
 is expected to be the same as the truth value assigned to $p$ but
delayed by an amount of time $t$.  Here, the first argument ($t$) is
a parameter and the second argument ($p$) is not.

If there are fewer than $n$ entries on the parameter list, then all of
the missing entries are assumed to be \NIL\ (so that the associated
arguments to the modal function are not parameters).  Note: All
parametric arguments should be bound when attempting to derive a modal
expression.  This is similar to the requirement that arguments to a
procedurally-attached predicate be bound when the procedure is called.
 \end{description}

\subsection{Modal operators supplied with the MVL bilattices}

All bilattices in \mvl\ are equipped with a modal operator \itype{idn}
that does not change the truth value of its argument at all; this
modal operator can be used to semantically mark points where it may be
desirable to interrupt the inference process.  In addition, if the
prover is invoked on \type{(idn p)} instead of on \type p, attempts
will be made to prove both $p$ and its negation $\neg p$.

The bilattices supplied as part of the \mvl\ system are equipped with
the following modal operators:

 \begin{description}
 \item [\type{*first-order-bilattice*}] This bilattice has two
modal operators: $L$, which maps $p$ into, ``necessarily $p$,'' or,
``I know that $p$,'' and $M$, which maps $p$ into ``possibly $p$.''
 \item [\type{*atms-bilattice*}]  No modal operators are supplied with
this or any of the other \atms-type bilattices.
 \item [\type{*default-bilattice*}] This bilattice and the hierarchical
default bilattice both inherit the modal operators $L$ and $M$ from the
first-order bilattice.
 \item [\type{*circumscription-bilattice*}] This bilattice does not
come equipped with any modal operators.
 \item [\type{*time-bilattice*}] This bilattice inherits the modal
operators $L$ and $M$ from the bilattice of basic values.  The
following three additional modal operators are also defined:
 \begin{enumerate}
 \item \itype{delay(t p)} This is a modal operator that pushes the
truth value of $p$ into the future by an amount of time $t$.  Thus if
$p$ is true at time 2, $\type{delay}(p,2)$ will be true at time 4.
 \item \itype{propagate} is a modal operator that corresponds to the frame
axiom.  If $p$ is true at time 2, unknown between times 3 and 20, and false
at time 21, then $\type{propagate}(p)$ will be true from times 2 until 20
and false subsequently.
 \item \itype{at(t p)} is a modal operator that returns the value taken
by the truth value $p$ at the time $t$.  (In actuality, a function has
to be returned; the function is the one that takes value $p(t)$ at all
times instead of just at $t$.)
 \end{enumerate}
 \end{description}

\subsection{Quantification and modal operators}

In addition to \type{idn}, there are two other modal operators that
are available for use with any bilattice, \itype{exists} and
\itype{forall}.  Although not really modal operators on the bilattice
being used, these {\em are\/} legitimate modal operators on the bilattice
\type{bdg-to-truth-val} that was described in Section \ref{s.binding-dag}.

The functions associated to these modal operators mimic the usual behaviors
given to the quantifiers.  Thus an expression such as
 \begin{equation}
\type{(exists ?x (foo ?x))}
\label{e.exists}
\end{equation}
 is evaluated by first computing the truth value that could be assigned
to \type{(foo ?x)} as an element of \type{bdg-to-truth-val}, so that
a truth value is obtained for every binding of \type{?x}.  This truth
value is then manipulated so that the values for any particular binding
of \type{?x} are disjoined (the usual interpretation of the existential
quantifier), and the result is returned as the truth value assigned to
\ref{e.exists}.  The modal operator \type{forall} is handled similarly
and corresponds to universal quantification.  The first (parametric)
argument to these operators can also be a list of variables, in which
case all of the variables are quantified over.

\subsection{Modal operators created automatically}

Finally, many of the functions that construct new bilattices are able
to construct modal operators on these bilattice in terms of the modal
operators on the bilattices they are using in the construction.  More
specifically:

\begin{enumerate}
 \item A product bilattice $A \times B$ will inherit a modal operator
from $A$ and $B$ if these component bilattices have modal operators
with the same name.
 \item A bilattice constructed using \type{append-bilattice} will have
modal operators corresponding to whatever modal operators were available
on the bilattice used in the construction.
 \item A bilattice constructed by splicing two bilattices $A$ and $B$
together will inherit the modal operators from the ``upper'' bilattice.
 \item A bilattice constructed using \type{hierarchy-bilattice} will
have modal operators corresponding to whatever modal operators were
available on the bilattice used in the construction.
 \item A bilattice constructed from a lattice will not normally have
any modal operators associated to it.
 \item The bilattice $B^D$ will inherit whatever modal operators the
bilattice $B$ has.
 \end{enumerate}

\section*{Acknowledgement}

Support for the development of \mvl\ has been provided by the Air
Force Office of Scientific Research, the Defense Advance Research
Projects Agency, General Dynamics, the National Science Foundation,
Rockwell International, and Rome Labs.

\clearpage

\appendix

\section{Summary of MVL functions and constants}

\begin{list}{}{\setlength{\labelwidth}{1.5in}\setlength{\labelsep}{.25in}
\setlength{\leftmargin}{1.75in}}

% *A*

\item [\itype{activate(th)}] Activate the theory \type{th} for reading
and writing.

\item [\itype{activate-best-first}] Activate best-first search by
setting \type{*node-evaluation-fn*} to \type{best-first-fn}.  Allowed
keywords: \type{initial-limit}, \type{limit-delta}, and
\type{final-limit}.

\item [\itype{activate-breadth-first}] Activate breadth-first search
by setting \type{*node-evaluation-fn*} to \type{depth-fn}.  Allowed
keywords: \type{initial-limit}, \type{limit-delta}, and
\type{final-limit}.

\item [\itype{activate-depth-first}] Activate depth-first search by
setting \type{*initial-limit*} to \NIL.  Allowed keyword:
\type{final-limit}.

\item [\itype{activate-theory-mechanism}] Activates the theory
mechanism.  Accepts as an optional argument either a one-character
bilattice descriptor, or a bilattice.

\item [\itype{*active-bilattice*}] The currently active bilattice.

\item [\itype{*active-theory*}] The theory that is currently active
for reading and writing.

\item [\itype{add-dag(d b f)}] Given the \dg-function $f$, add the
value $b$ at the point $d$.  This function works by side effect only.
There is a keyword \type{modifier} that, if supplied, should be a
function that takes two elements of the bilattice and returns a single
one.  Given such a function $m$, the \dg-function is changed so that
the new value at $d$ is $m(b,f(d))$, where $f(d)$ is the old value.
(This is similar to the same keyword in \type{dag-accumulate-fn}.)
The keyword \type{simplify} is also available.

\item[\itype{answer-binding(a)}] Returns the binding field of the
answer \type a.

\item[\itype{answer-value(a)}] Returns the value field of the answer
\type a.

\item [\itype{anytime-trace(pat)}] Include \type{pat} on \at, so that
any instance of \type{pat} will be traced.  The argument is optional;
if omitted, \type{(?*)} (of which everything is an instance) will be
used.

\item [\itype{anytime-untrace(pat)}] Remove any instance of
\type{pat} from \at.  The argument is optional; if omitted,
\type{(?*)} (of which everything is an instance) will be used.

\item [\itype{append-bilattice(b f)}] Returns the bilattice obtained
by appending the result of applying the function $f$ to the points of
the original bilattice $b$.

\item [\itype{append-binding-lists(b1 b2)}] Computes the least
specific binding list that is more specific than both $b_1$ and $b_2$
or returns \NIL\ if no such binding list exists ($b_1$ and $b_2$ may
bind the same variable to different constants, for example).

\item [\itype{*atms-bilattice*}] The bilattice corresponding to an \atms,
where the justification information is maintained in disjunctive normal
form as suggested by de Kleer.

% *B*

\item [\itype{(bagof}\type{ ?x ?p ?b)}] This relation holds for
$x$, $p$ and $b$ if and only if $b$ is a list of all $x$ satisfying
$p$.

\item [\itype{bc(q)}] Attempts to prove the query $q$ from information
in the database, looking only for a single answer.

\item [\itype{bc-trace(pat)}] Include \type{pat} on
\type{*bc-trace*}, so that any instance of \type{pat} will be
traced.  The argument is optional; if omitted, \type{(?*)} (of which
everything is an instance) will be used.

\item [\itype{bc-untrace(pat)}] Remove any instance of \type{pat}
from \type{*bc-trace*}.  The argument is optional; if omitted,
\type{(?*)} (of which everything is an instance) will be used.

\item [\itype{bc-unwatch(preds)}] Removes the given predicates from
the list of {\sc WATCH}ed predicates.  If no arguments are supplied,
all predicates are un{\sc watch}ed.

\item [\itype{bc-watch(preds)}] Accepts as arguments any number of
predicates that should be watched as the prover proceeds.  Each
predicate can also be a list of the form \type{(predicate depth)}
where \type{depth} is the number of successive ancestors to be
displayed when the prover encounters an instance of the predicate.
Returns a list of the predicates currently being {\sc WATCH}ed.

\item [\itype{bcs(q)}] Attempts to prove the query $q$ from information
in the database, looking for all possible answers.

\item [\itype{bdg-to-truth-val}] The bilattice of functions from
the \dg\ of binding lists into \type{*active-bilattice*}.

\item [\itype{bilattice(b)}] Add the bilattice $b$ to \type{*bilattices*}.

\item [\itype{*bilattices*}] A list of the bilattices known to the system.

\item [\itype{binding-dag}] The \dg\ of binding lists.

\item [\itype{binding-le(b1 b2)}] Returns \NIL\ unless $b_1$ is a more
specific binding that $b_2$, in which case a list of relative binding
lists is returned.

\item [\itype{bottom}] The truth value \type{bottom} in the active bilattice.

\item [\itype{*break-node*}] The number of a node at which the
proof process should be interrupted.

% *C*

\item [\itype{choose-task}]  Selects a task to work on next.  The 
current implementation of \mvl\ simply selects any task that might
have some impact on the final answer being computed.

\item [\itype{circum(p)}] Like \type{bc}, but looks for instances of
$p$ that follow from the circumscription axiom.  Allowed keywords:
\type{cutoffs} and \type{succeeds}.

\item [\itype{*circum-trace*}] If not \NIL, the circumscriptive proof
functions will output diagnostics describing their progress.

\item [\itype{circums(p)}] Like \type{bcs}, but looks for instances of
$p$ that follow from the circumscription axiom.  Allowed keywords:
\type{cutoffs} and \type{succeeds}.

\item [\itype{*circumscription-bilattice*}] The bilattice used to compute
the results of the circumscription axiom.

\item [\itype{cnf(p)}] Returns $p$ in conjunctive normal form, as a list
of conjuncts, where each conjunct is a list of disjuncts.  Thus
 \[\type{(CNF '(AND A B))}\]
returns
 \[\type{((A) (B))}\]

\item [\itype{*cnf-atms-bilattice*}] An alternative version of the \atms\
bilattice where the justification information is maintained in conjunctive
normal form.

\item [\itype{combine-1(f params f1 ... fn)}] This is an $n$-ary
version of \type{combine-fns}, where $f$ is an $n$-ary bilattice
function and there are $n$ \dg-functions supplied as well.  The
additional argument \type{params} indicates which of the $n$ arguments
to this function are not necessarily \dg-functions, but are instead
parameters expected by the overall function $f$.  This additional
argument is a list of \T\ (if the $n$th argument is a parameter) or
\NIL\ (if it isn't).  If the parameter list ends early, all of the
remaining entries are assumed to be \NIL.

\item [\itype{combine-fns(f1 f2 fn comp)}] This function takes the two
\dg-functions $f_1$ and $f_2$ and combines them using the function
\type{fn}, which should be a function that takes two elements of the
bilattice and returns a single one.  The comparison function
\type{comp} is used to determine if one of the two functions can be
returned directly; if, for example \type{(comp f1 f2)} holds, than we
simply return $f_1$, and similarly for $f_2$.  As an example, if the
combining function is the bilattice operation $\cdot$, the comparison
function is \type{k-le}.

\item [\itype{compile-mvl}] Compiles and loads the \mvl\ system.
On a Symbolics system, the keyword \type{reload} forces all of the
files to be reloaded and the keyword \type{recompile} forces them all
to be recompiled.

\item [\itype{*completed-database*}] If \T, find all exceptions
to default proofs.  If \NIL\ (the default), accept defaults proofs
with possible exceptions as complete.

\item [\itype{cont-analysis(analysis)}] This function continues the
investigation of a specific analysis, and can be used to resume a
proof effort that has been curtailed for some reason.  It returns two
values: an instance of the \type{answer} structure giving the results
of the proof attempts, and a modified version of the \type{analysis}
structure that can be used to continue the proof process if desired.

\item [\itype{cont-prover(proof)}] Resume the task associated to the
auxiliary information in the given instance of the \type{proof}
structure.  This function returns an instance of \type{bcanswer}
corresponding to the result obtained by the prover, and returns as
soon as a solution to the original query is found.

\item [\itype{contents}] Displays the contents of the database.  Keywords:
\type{cutoffs}, \type{props}.

\item [\itype{create-mvl-invocation}] This function sets up a new
copy of \mvl.  It accepts an optional argument that is an instance of
\mi; this can be used to reset the state of \mvl\ to a previous one.
If the optional argument is omitted, default values are used.  The
value returned is the \mvl\ invocation that was active before the
function was called, so that another call to
\type{create-mvl-invocation} can be used to restore the original
environment.

\item [\itype{cutoffs-not(cutoffs \&optional bilattice)}] Returns a
new cutoffs $c$ such that
 \[\type{(test x c)} = \type{(test (mvl-not x) cutoffs)}\]
 for all truth values $x$.  In other words, $x$ passes $c$ if and
only if $\neg x$ passed the cutoffs supplied originally.

% *D*

\item [\itype{dag-accumulate-fn(dag bilattice list)}] Given a \dg, a
bilattice, and a list of elements of the form
 \[\type{(dag-elt . bilattice-elt)}\]
a \dg-function is constructed that takes the given values at the
given elements of the \dg.

A keyword to \type{dag-accumulate-fn} is \type{modifier}.  This
function is used if the \type{d-list} contains the same \dg\ element
more than once, in which case the modifier function is called with
the new and old values (in that order) and the result is used as the
value of the \dg-function at that particular point.  The default
action is to simply overwrite the old value with the new one.  (In
other words, the default modifier function simply returns its first
argument while ignoring the second.)

\item [\itype{dag-bilattice(dag bilattice)}] Given a bilattice $B$
and a dag $D$, returns the bilattice $B^D$.

\item [\itype{dag-change-dag(f d-f)}] The argument \type{d-f} should
be a function from the \dg\ to itself; a modified version of the
original \dg-function $f$ is returned in which every point at which
the \dg-function takes a new value has had the function \type{d-f}
applied to it.  The bilattice argument is expected to be the bilattice
of functional truth values as opposed to the bilattice that is the
target of the functions.  There is an optional keyword
\type{simplify}; if supplied, the computed function is simplified
before being returned.  (See \type{simplify}.)

\item [\itype{dag-change-tv(f b-f)}] The argument \type{b-f} should be
a function from the bilattice to itself; a modified version of the
original \dg-function $f$ is returned in which every point in the
\dg-function has had the function \type{b-f} applied to it.  There is
an optional keyword \type{simplify}; if supplied, the computed
function is simplified before being returned.  (See \type{simplify}.)

\item [\itype{dag-list-of-answers(f)}] Given a \dg-function $f$, assumed
to be from the \dg\ of binding lists into the currently active bilattice, this
function constructs a list of answers corresponding to the points where
the \dg-function takes values.

\item [\itype{dag-prune(f pred)}] Given a \dg-function $f$ and a
predicate on the bilattice, this function prunes from $f$ all points
$d$ where $f(d)$ fails to satisfy the predicate.  There is an optional
keyword \type{modify-fn} that, if supplied, should be a function $m$
from the bilattice to itself.  In this case, the points where
$m(f(d))$ fails to satisfy the predicate are pruned.

\item [\itype{*default-bilattice*}] A bilattice that can be used for simple
default reasoning.

\item [\itype{delete-bdg(var bdg-list)}] Destructively modifies
\type{bdg-list}, removing any binding for the variable \type{var}.

\item [\itype{denotes(a)}] Returns the proposition associated with a
particular \lisp\ atom.

\item [\itype{dnf(p)}] Returns $p$ in disjunctive normal form, as a list
of disjuncts, where each disjunct is a list of conjuncts.  Thus
 \[\type{(DNF '(AND A B))}\]
returns
 \[\type{((A B))}\]

\item [\itype{(done exp-1 ... exp-n)}] This predicate, which is
useful for stashing only, causes \lisp\ to evaluate each of the subsequent
expressions.

\item [\itype{dot-with-not(x)}] Computes $x\cdot \neg x$.

\item [\itype{dot-with-star}] The function \type{dws} in the active bilattice.

% *E*

\item [\itype{empty}] Unstashes every sentence in the database.  Keyword:
\type{value}.

\item [\itype{empty-theory()}] Empties the currently active theory,
although this theory may be changed using the optional theory
argument.  The keyword \type{value} has a meaning similar to its
meaning as a keyword to \type{erase}.

\item[\itype{equal-answer(a1 a2)}] Returns \T\ if the two answers
\type{a1} and \type{a2} are equal, and \NIL\ otherwise.  The
functions used are \type{binding-equal} and \type{mvl-eq}, the
functions normally used to test binding lists and truth values for
equality.

\item [\itype{equal-binding(b1 b2)}] Checks to see if two binding lists are
effectively the same.

\item [\itype{*equivalence-translation*}] Controls the fashion in which
\type{normal-form} deals with the connective \type{<=>}.

\item [\itype{erase(p)}] Removes $p$ from the database and forward
chains.  Allowed keywords: \type{value}, \type{decrement-p} and
\type{cutoffs}.

\item [\itype{erase-from-theory(p)}] Removes the proposition $p$ from
the currently active theory, and forward chains on the result.
Allowed keywords: \type{theory} can be used to specify the theory from
which $p$ should be removed; \type{value}, \type{decrement-p} and
\type{cutoffs} have the same meaning as for \type{erase}.  (See
Section \ref{s.fc}.)

\item [\itype{existing-mvl-invocation}] This function returns a
single value that is an instance of the structure \mi.  It is useful
primarily for saving the current copy of \mvl\ so that it can be
restored later.

% *F*

\item [\itype{false}] The truth value \type{false} in the active bilattice.

\item [\itype{fc(p)}] Inserts the sentence $p$ into the database and
forward chains.  Allowed keywords: \type{value}, \type{increment-p},
\type{fc-increment-p}, \type{cutoffs} and \type{include-modal}.

\item [\itype{fc-trace(pat)}] Include \type{pat} on
\type{*fc-trace*}, so that any instance of \type{pat} will be
traced.  The argument is optional; if omitted, \type{(?*)} (of which
everything is an instance) will be used.

\item [\itype{fc-untrace(pat)}] Remove any instance of \type{pat}
from \type{*fc-trace*}.  The argument is optional; if omitted,
\type{(?*)} (of which everything is an instance) will be used.

\item [\itype{*final-limit*}] is the limit point beyond which the
search is abandoned, so that any node that evaluates to a value worse
than this will be pruned.  The default value of this parameter is
\NIL, indicating that no nodes are to be pruned.

\item [\itype{*first-order-atms-bilattice*}] This bilattice corresponds
to a fully first-order \atms, as opposed to the more conventional one
dealing only with propositional logic.

\item [\itype{*first-order-bilattice*}] The bilattice corresponding to
first-order logic.

\item [\itype{fn-and(f1 f2)}] This function conjoins the values of
$f_1$ and $f_2$ using the bilattice $\wedge$ operation.  The new
\dg-function is returned.

\item [\itype{fn-comp(f1 f2 fn)}] This functions compares the values
of $f_1$ and $f_2$ using the bilattice comparison function \type{fn}.
Either \T\ or \NIL\ is returned.

\item [\itype{fn-dot(f1 f2)}] This function combines the values of
$f_1$ and $f_2$ using the bilattice $\cdot$ operation.  The new
\dg-function is returned.

\item [\itype{fn-dws(f1 f2)}] This function combines the values of
$f_1$ and $f_2$ using the bilattice operation \type{dot-with-star}.
The new \dg-function is returned.

\item [\itype{fn-eq(f1 f2)}] This function compares the values of
$f_1$ and $f_2$ using the bilattice comparison operation.  Either \T\
or \NIL\ is returned.

\item [\itype{fn-not(f1)}] This function negates the values of $f_1$
using the bilattice $\neg$ operation.  The new \dg-function is
returned.

\item [\itype{fn-or(f1 f2)}] This function disjoins the values of
$f_1$ and $f_2$ using the bilattice $\vee$ operation.  The new
\dg-function is returned.

\item [\itype{fn-plus(f1 f2)}] This function adds the values of $f_1$
and $f_2$ using the bilattice $+$ operation.  The new \dg-function is
returned.

% *G*

\item [\itype{get-bdg(var bdg-list)}] Returns the value assigned to
the variable \type{var} by the binding list \type{bdg-list}.  If the
variable does not appear in the binding list, \NIL\ is returned.

\item [\itype{get-mvl (symbol indicator)}]  Like the Common Lisp
\type{get}, but if the value is a list, treats it as an a-list and
searches for an item whose car matches \type{*mvl-invocation*}.
The value taken can be changed with \type{setf}.

\item [\itype{get-val(d f)}] Returns the value of the function $f$ at
the point $d$ of the \dg.

\item [\itype{*global*}] The root of the \dg\ of theories.

\item [\itype{groundp(p)}] Returns \T\ if the proposition $p$ is
ground (i.e., contains no variables) and \NIL\ otherwise.

% *H*

\item [\itype{hierarchy-bilattice(b)}] Returns the bilattice constructed
as an infinite descending chain of copies of $b$.

\item [\itype{*hierarchy-bilattice*}] A bilattice similar to the
default bilattice but where the defaults are labelled with
integer-valued strengths.

% *I*

\item [\itype{*if-translation*}] Controls the fashion in which
\type{normal-form} deals with the connective \type{if}.

\item [\itype{includes(big little)}] This function is used to state that
the theory \type{little} is a subtheory of the theory \type{big}, and
should inherit from it.

\item [\itype{index(p)}] Updates the discrimination net as needed
to insert the proposition $p$ into the database.

\item [\itype{init-analysis(prop)}] The \mvl\ primitive upon which
the proving functions are based, this function accepts a proposition
and creates and returns an instance of the \type{analysis} structure.

\item [\itype{init-prover(p effort)}] This function initializes and
returns an instance of \type{proof} containing auxiliary information
about the proof process.

\item [\itype{*initial-limit*}] is used to indicate the type of search.
If \NIL\ (the default), then the values assigned to the nodes are not
used to order the search, and depth-first search is used.  If not
\NIL, \type{*initial-limit*} is typically 0.

\item [\itype{instp(p1 p2)}] If $p_2$ is an instance of $p_1$,
return a list of their most general unifiers; only variables in $p_1$
are bound.  Returns \NIL\ if $p_2$ is not an instance of $p_1$.

\item [\itype{invoke(fn p)}] This function invokes \type{fn} on the
sentence $p$.  In addition, any keyword arguments are passed to
\type{fn}, but with \type{\&allow-other-keys} so that the keywords
will be ignored if \type{fn} is not expecting them.

% *J*

% *K*

\item [\itype{k-ge(x y)}] Returns \T\ if $x \geq_k y$, and \NIL\ otherwise.

\item [\itype{k-gt(x y)}] Returns \T\ if $x >_k y$, and \NIL\ otherwise.

\item [\itype{k-le(x y)}] Returns \T\ if $x \leq_k y$, and \NIL\ otherwise.

\item [\itype{k-lt(x y)}] Returns \T\ if $x <_k y$, and \NIL\ otherwise.

\item [\itype{k-not-ge(x y)}] Returns \T\ if $x \not\geq_k y$, and
\NIL\ otherwise.

\item [\itype{k-not-gt(x y)}] Returns \T\ if $x \not>_k y$, and \NIL\
otherwise.

\item [\itype{k-not-le(x y)}] Returns \T\ if $x \not\leq_k y$, and
\NIL\ otherwise.

\item [\itype{k-not-lt(x y)}] Returns \T\ if $x \not<_k y$, and \NIL\
otherwise.

\item [\itype{(known p)}] Stashing this predicate simply stashes $p$;
unstashing it causes an error.  Looking it up looks up $p$.

% *L*

\item [\itype{lattice-to-bilattice(l)}] Returns the bilattice constructed
from the lattice $l$.

\item [\itype{lattice-to-dag(lattice)}] Construct a \dg\ from the
given lattice.

\item [\itype{lattice-to-dag-to-bilattice(lattice bilattice)}]  First
construct a \dg\ $D$ from the given lattice.  Then use the bilattice
$B$ to construct the bilattice $B^D$.

\item [\itype{*limit-delta*}] is used if \type{*initial-limit*} is
non-\NIL.  In this case, the nodes are ordered by increasing
 \[\frac{\type{value(node)}-\type{*initial-limit*}}{\type{*limit-delta*}}\]
 A typical value for *initial-limit* is 0, and for *limit-delta* is
1.  The value of this parameter is irrelevant if
\type{*initial-limit*} is \NIL.

\item [\itype{lisp-defined(fn)}] This macro is used to indicate that
the given function can be evaluated by a \lisp\ call.  It is also
used to cause \mvl\ to signal an error if an attempt is made to stash
or unstash an instance of the function.  \type{Lisp-defined} also
accepts an optional argument, which is an inverter function.  The
overall result of attempting to lookup an instance of the function
can be described as follows:
 \begin{enumerate}
 \item If none of the arguments to the function (i.e., all but the
last of the arguments of the relation) is a variable, then the function
is called and the result is unified with the final argument of the
relation.
 \item If the inverter function has been supplied and is equal to \T,
then the function is called and the result is unified with the final
argument of the relation even if some of the arguments are variables.
 \item If some of the arguments are variables but the result is not a
variable, then the inverter function is invoked if one has been
supplied.  The arguments passed to the inverter function are the
result (i.e., the last argument of the relation), and a list of all
of the arguments (i.e., all but the last argument of the relation).
If there is no inverter function, then \NIL\ is returned.  In
addition, if \wbli\ \index{*warn-on-bad-lisp-invocations*} is not
\NIL, the user is warned.
 \item If both some of the arguments and the result are variables, then
\NIL\ is returned.  In addition, if \wbli\ is not \NIL, the user is
warned.
 \end{enumerate}

\item [\itype{lisp-predicate(pred)}] This macro is used to indicate
that the given predicate can be evaluated by a \lisp\ call.  It is
also used to cause \mvl\ to signal an error if an attempt is made to
stash or unstash an instance of the function.

\item [\itype{load-logic}] Empty the database and make a new bilattice
active.  Accepts as an optional argument either a character designating
a bilattice, or the new bilattice itself.  If no argument is supplied,
the user is prompted for the new bilattice.

\item [\itype{load-mvl}] Loads the \mvl\ system.

\item [\itype{lookup(q)}] Searches the database for a sentence
matching the query $q$.  Allowed keywords: \type{cutoffs}, \type{inv},
\apa, \somo.  Returns two values: an answer and the source of the
match.  Returns \NIL\ if the query cannot be found in the database.

\item [\itype{lookup-call(fn p)}] This function is similar to
\type{invoke}, but modifies the values returned to make them agree
with those expected to be returned by \type{lookup}.  Unless the
function \type{fn} supplies them, the truth value is assumed to be
\type{true}.  The second value is \type{attach}.

\item [\itype{lookups(q)}] Similar to \type{lookup}, but looks
for {\em all\/} sentences matching the query.  Returns two lists of
values.

\item [\itype{lookups-call(fn p)}] The plural version of
\type{lookup-call}.

% *M*

\item [\itype{make-root(dag bilattice b)}] Returns a \dg-function that
takes the value $b$ at the root of the \dg.

\item [\itype{make-tv(dag bilattice d b)}] Returns a \dg-function that
takes the value $b$ at the point $d$.  If $d$ is not the root of the
\dg, the value at the root of the \dg\ is generally the point
\type{unknown} in the bilattice, but the user may supply the keyword
\type{root-fn} if he wishes, in which case the value at the root of
the \dg\ is the result of applying \type{root-fn} to the bilattice.

\item [\itype{matchp(pattern atom)}] Returns a list of elements of
the form \type{(b . d)} where $b$ is a most general binding list such
that \type{(plug pattern b)} is an instance of the proposition
labelled with \type{atom} and $d$ is a binding list for the variables
in the database proposition.  \NIL\ is returned if no such binding
list exists.  Thus if \type{P1} is \type{(TALL \#:?2)}, and \type{P2}
is \type{(TALL TOM)}, then
 \[\type{(matchp (TALL TOM) P1)}\]
returns a list consisting of the single entry 
 \[\type{(NIL . ((\#:?2 . TOM)))}\]
 The \type{car} of this entry is \NIL, indicating that no variables
in the query need to be bound to match \type{P1}; the \type{cdr} is
\type{((\#:?2 . TOM))} since the variable \type{\#:?2} in \type{P1}
needs to be bound to \type{TOM} to match the query.
 \[\type{(matchp (TALL ?X) P2)}\]
returns \type{((((?X . TOM)) . NIL))}.

\item [\itype{mvl-and}] The function \type{and} in the active bilattice.

\item [\itype{mvl-dot}] The function \type{dot} in the active bilattice.

\item [\itype{mvl-eq}] The function \type{eq} in the active bilattice.

\item [\itype{mvl-file(string)}] Returns the path for the \mvl\ file
associated with the given name.  The file name may include an
extension, one may be supplied as an optional argument, or
\type{.MVL} will be used by default.  There is an optional keyword
\type{no-error} indicating that even if the file cannot be found, the
path should be returned.

\item [\itype{*mvl-invocation*}] A global variable containing a unique
identifier for the currently active copy of \mvl.

\item [\itype{mvl-load(file)}] Loads the sentences in the given file into
the \mvl\ database.

\item [\itype{mvl-not}] The function \type{not} in the active bilattice.

\item [\itype{mvl-or}] The function \type{or} in the active bilattice.

\item [\itype{mvl-plug}] The function \type{plug} in the active bilattice.

\item [\itype{mvl-plus}] The function \type{plus} in the active bilattice.

\item [\itype{mvl-print}] This function accepts a single argument, a
truth value, and modifies it in a way that will cause it to be
pretty-printed.  The value returned by mvl-print should not be used
for any purpose other than printing.

\item [\itype{mvl-save(file)}] Writes the \mvl\ database out to the
given file.  This function also accepts the keywords \type{cutoffs}
and \type{props}, which are treated as for \type{contents} (see
Section \ref{s.contents}).  Finally, the keyword \type{if-exists} is
passed to the \type{with-open-file} macro.

\item [\itype{mvl-simplicity}] The function \type{simplicity} in the active
bilattice.

\item [\itype{mvl-test-file(string)}] Can be used to evaluate
various forms by the \mvl\ system and compare the results to
predefined values.  Both the forms and the expected answers are in
the \type{.TEST} file whose name is supplied to the function.
Returns the number of errors found.

\item [\itype{mvl-unload(file)}] Removes the sentences in the given file from
the \mvl\ database.

\item [\itype{mvl-vars}] The function \type{vars} in the active bilattice.

% *N*

\item [\itype{negate(p)}] Returns the negation of $p$.

\item [\type{new-$*$var()}]\index{new-star-var} Returns a new
temporary sequence variable.

\item [\type{new-?var()}]\index{new-query-var} Returns a new
temporary variable.

\item [\itype{new-val(a v)}] Sets the truth value of the sentence
associated with the atom $a$ to $v$.  If $v$ is the truth value
\type{unknown}, the sentence is removed from the database.

\item [\itype{normal-form(p)}] Returns an equivalent form of $p$
suitable for insertion into the database (i.e., with the connectives
\type{IFF}, \type{IF} and \type{<=>} reduced to \type{<=} and \type{=>}).

% *O*

% *P*

\item [\itype{plug(p bdg-list)}] Returns the result of instantiating
the proposition $p$ using the given binding list.  This function is
nondestructive.

\item [\itype{prfacts(atom)}] Displays the database facts concerning the
given atom.  Keyword: \type{cutoffs}.

\item [\itype{*probability-bilattice*}] An extended version of the \atms\
bilattice that maintains probabilistic information as well.

\item [\itype{product-bilattice(b1 b2)}] Returns the bilattice
that is the product of the bilattices $b_1$ and $b_2$.

\item [\itype{(property symbol value indicator)}] Stashing this
predicate sets the symbol's \type{indicator} property to \type{value}
for the current invocation of \mvl.  Unstashing it removes the
\type{indicator} property for the current invocation of \mvl\ if the
previous value of this property unifies with \type{value}.  (See the
definitions of \type{get-mvl} and \type{rem-mvl}.)

\item [\itype{props()}] Returns a list of all propositions about
which any information has been recorded.

% *Q*

% *R*

\item [\itype{release-mvl}] (Symbolics only)  Compiles and loads the
\mvl\ system, and declares the new version to be \type{released}.  The
keyword \type{no-recompile} can be used to skip the compilation step.

\item [\itype{rem-mvl (symbol indicator)}] Like the Common Lisp
\type{remprop}, but if the value is a list, treats it as an a-list
and removes the item whose car matches \type{*mvl-invocation*}.

% *S*

\item [\itype{samep(p1 p2)}] Returns \T\ if the propositions $p_1$ and
$p_2$ are identical (up to variable renaming), and \NIL\ otherwise.
If the optional \type{return-answer?} argument is supplied and the
propositions are the same, a binding list is returned instead of \T.

\item [\itype{*save-nodes*}] If \T\ (the default), nodes are
labelled with consecutive integers when they are displayed by the
first-order prover.

\item [\itype{set-search-parameters}] This function is the primitive
used to modify the parameters used to control the first-order theorem
prover's search.  It accepts as keywords \type{initial-limit},
\type{limit-delta}, and \type{final-limit}.

\item [\itype{simplify(f)}] In some cases, manipulating the
\dg-function $f$ may have cause it to contain redundancies, where the
value at some point is presented explicitly even though this value is
in fact the one that would normally be inherited from its parents.
This function returns a modification of $f$ where these irrelevant
entries have been removed.

\item [\itype{splice-bilattice(upper lower)}] Returns the bilattice
obtained by replacing the point \type{unknown} in the bilattice
\type{upper} with a copy of the bilattice \type{lower}.

\item [\itype{standardize-operators(p)}] Returns an equivalent form of $p$
with the connectives \type{<=>}, \type{IF} and \type{IFF} removed.
This function does not use the control variables
\type{*if-translation*} and \et, since it is used when proving
sentences, not when inserting them into the database.

\item [\itype{stash(p)}] Inserts $p$ into the database.  The allowed
keywords are \type{value}, \type{increment-p} and \apa.

\item [\itype{stash-value}] The function \type{stash-val} in the
active bilattice.

\item [\itype{state(p)}] Inserts $p$ into the database after
converting it to normal form.  Keywords: \type{value},
\type{increment-p}.

\item [\itype{state-in-theory(p)}] Adds the proposition $p$ to
the currently active theory.  Allowed keywords: \type{theory} can be
used to specify a different theory to which $p$ should be added;
\type{value} and \type{increment-p} have the same meaning as for
\type{state}.  (See Section \ref{s.database}; if these keywords are
not used, the function \type{state} can be used equally well, since
the default stash value with the theory mechanism activated also
inserts the proposition into the currently active theory.)

\item [\itype{std-cutoffs}] Cutoff information constructed from the
value of \type{cutoff-val} in the active bilattice.  See Section
\ref{s.cutoffs}.

% *T*

\item [\itype{t-ge(x y)}] Returns \T\ if $x \geq_t y$, and \NIL\ otherwise.

\item [\itype{t-gt(x y)}] Returns \T\ if $x >_t y$, and \NIL\ otherwise.

\item [\itype{t-le(x y)}] Returns \T\ if $x \leq_t y$, and \NIL\ otherwise.

\item [\itype{t-lt(x y)}] Returns \T\ if $x <_t y$, and \NIL\ otherwise.

\item [\itype{t-not-ge(x y)}] Returns \T\ if $x \not\geq_t y$, and
\NIL\ otherwise.

\item [\itype{t-not-gt(x y)}] Returns \T\ if $x \not>_t y$, and \NIL\
otherwise.

\item [\itype{t-not-le(x y)}] Returns \T\ if $x \not\leq_t y$, and
\NIL\ otherwise.

\item [\itype{t-not-lt(x y)}] Returns \T\ if $x \not<_t y$, and \NIL\
otherwise.

\item [\itype{test(value cutoffs \&optional bilattice)}] Returns \T\
if the truth value passes the test associated to \type{cutoffs}, and
\NIL\ otherwise.  The test is performed for the supplied bilattice,
or for the currently active one is the optional argument is omitted.

\item [\itype{test-mvl}] Tests the \mvl\ system by invoking the function
\type{mvl-test-file} on the file \type{MVL.TEST}.

\item [\itype{theory-active-p}] Returns \T\ or \NIL\ depending on
whether or not the theory mechanism has been activated.

\item [\itype{theory-contents()}] Displays the contents of the
currently active theory, which can be changed using the
keyword \type{theory}.  The keywords \type{cutoffs} and \type{props}
are allowed and are treated as for \type{contents}.  (See Section
\ref{s.contents}.)  There is an additional keyword,
\itype{this-theory-only}, which defaults to \T\ and, if supplied,
indicates that only sentences whose truth values {\em change\/} in the
given theory are to be displayed.

\item [\itype{theory-dag}] The \dg\ of theories.

\item [\itype{*time-bilattice*}] A compound bilattice that can be used
for reasoning about events occurring against a temporal background.

\item [\itype{true}] The truth value \type{true} in the active bilattice.

\item [\itype{true-in-theory}]  The truth value of a sentence that is
true in the currently-active theory.

\item [\itype{truth-value(a)}] Returns the truth value associated with
a particular \lisp\ atom.  Thus, if the result of asserting
\type{(not (short tom))} into the database was the atom \type{P38},
\type{(denotes 'p38)} would evaluate to \type{(short tom)}, and
\type{(truth-value 'p38)} would evaluate to \type{false}. 

% *U*

\item [\itype{unifyp(p1 p2)}] Returns a list of the most general
unifiers of $p_1$ and $p_2$, or \NIL\ if the propositions fail to
unify.

\item [\itype{unincludes(big little)}] This function is used to state
that the theory \type{little} is no longer a subtheory of the theory
\type{big}, and should not inherit from it.

\item [\itype{unindex(p)}] Removes from the discrimination net 
information relating to the proposition $p$.

\item [\itype{unknown}] The truth value \type{unknown} in the active
bilattice.

\item [\itype{(unknown p)}]  Stashing this predicate unstashes $p$;
unstashing it causes an error.  Looking it up fails if $p$
can be found in the database, and succeeds otherwise.

\item [\itype{unstash(p)}] Removes $p$ from the database.  Allowed
keywords: \type{value}, \type{decrement-p} and \apa.

\item [\itype{unstash-facts-with-atom(atom)}] Unstashes from the
database any facts concerning the given atom.  Keywords:
\type{value}, \type{decrement-p}, \apa.

\item [\itype{unstash-value}] The function \type{unstash-val} in the
active bilattice.

\item [\itype{unstate(p)}] Removes $p$ from the database after
converting it to normal form.  Keywords: \type{value},
\type{decrement-p}.

\item [\itype{unstate-from-theory(p)}] Removes the proposition $p$
from the currently active theory.  Allowed keywords:
\type{theory} can be used to specify a different theory from which
$p$ should be removed; \type{value} and \type{decrement-p} have the
same meaning as for \type{unstate}.  (See Section \ref{s.database}.)

% *V*

\item [\itype{(value symbol val)}] Stashing this predicate sets the
given symbol to the value \type{val}.  Unstashing it makes the symbol
unbound if its current value unifies with \type{val}.

\item [\itype{varp(x)}] Returns \type{T} if $x$ is a variable, and \NIL\
otherwise.

\item [\type{varp$*$(x)}]\index{varp-star} Returns \T\ if $x$ is a
sequence variable, and \NIL\ otherwise.

\item [\type{varp?(x)}]\index{varp-query} Returns \T\ if $x$ is a
standard (i.e., non-sequence) variable, and \NIL\ otherwise.

\item [\itype{vars-in(p)}] Returns a list of the variables appearing in $p$.

\item [\itype{vartype(x)}] Returns ? if $x$ is a standard variable,
$*$ if $x$ is a sequence variable, and \NIL\ is $x$ is not a
variable.

% *W*

\item [\itype{with-mvl-invocation}] This is a macro similar to the
Common Lisp macro \type{with-open-file}.  It accepts an optional
argument that is an instance of \mi\ and is followed by a body of
forms (presumably including a \type{load-logic}) that are evaluated
in a locally-bound copy of \mvl.

% *X*

% *Y*

% *Z*

\end{list}

\documentstyle{article}

\setlength{\topmargin}{0pt}
\setlength{\oddsidemargin}{0pt}
\setlength{\evensidemargin}{0pt}
\setlength{\textwidth}{6.5in}
\setlength{\textheight}{8.5in}
\setlength{\headheight}{0pt}

%here are some mathematical symbols
\def\emptyset{{\hbox{\rm \O}}}
\def\set#1{\{#1\}}
\newcommand{\type}[1]{\mbox{\tt #1}}
\newcommand{\str}{?$*$}
\def\R{\rlap{\rm I}\,{\rm R}}
\def\powset{{\cal P}}
\def\glb{{\mbox{\rm glb}}}\def\lub{{\mbox{\rm lub}}}

\newcommand{\nm}{non\-mon\-o\-ton\-ic}
\newcommand{\Nm}{Non\-mon\-o\-ton\-ic}

\newcommand{\prolog}{{\sc prolog}}\newcommand{\Prolog}{{\sc Prolog}}
\newcommand{\atms}{{\sc atms}}\newcommand{\Atms}{{\sc Atms}}

%this is where the title page begins

\def\draftline{}

\def\authors{Matthew L. Ginsberg}

\newcommand{\nodblspace}{\baselineskip=\normalbaselineskip}
\def\group{\nodblspace Department of Computer Science\\
Stanford University\\
Stanford, California 94305\\
\bigskip
Updated \today
}

\def\titlebottom{\vspace{1.55in}
{\large \authors}\\
\vspace{1.55in}
\group
\end{center}
\clearpage}

\def\xtitletop{\vspace*{1.35in}\begin{center}}
\def\xtitl#1#2{\xtitletop
{\Large\bf #2}\\
\titlebottom}

\newcommand{\xtitle}[2]
{\def\tphrase{#2}\begin{titlepage}\draftline\xtitl{#1}{#2}\end{titlepage}}

\newcommand{\mvl}{{\sc mvl}}
\newcommand{\Mvl}{{\sc Mvl}}
\newcommand{\lisp}{{\sc lisp}}
\newcommand{\Lisp}{{\sc Lisp}}
\newcommand{\dg}{{\sc dag}}
\newcommand{\Dg}{{\sc Dag}}
\newcommand{\T}{\type{T}}
\newcommand{\NIL}{\type{NIL}}
\newcommand{\apa}{{\tt allow-\-procedural-\-attachment}}
\newcommand{\at}{{\tt *anytime-\-trace*}}
\newcommand{\mi}{{\tt mvl-\-invocation}}
\newcommand{\wbli}{{\tt *warn-\-on-\-bad-\-lisp-\-invocations*}}
\newcommand{\somo}{{\tt succeed-\-on-\-modal-\-ops}}
\newcommand{\et}{{\tt *equivalence-\-translation*}}
\newcommand{\itype}[1]{\type{#1}\index{#1 **}}

\makeindex
%\renewcommand{\index}[1]{}

\begin{document}
\xtitle{}{User's Guide to the MVL System}

\tableofcontents
\clearpage

\section{Introduction}
\subsection{What you got}

The \mvl\ system is currently distributed via anonymous ftp from
t.stanford.edu.  A .DVI version of this file, the \mvl\ manual,
can be found in the \type{papers} directory on this machine.

The \type{mvl} directory contains the source code for the \mvl\ system
itself, which is written in standard Common Lisp.  It should be
possible to port the system to any Common Lisp without substantial
difficulty.  In addition, the \type{papers} directory contains .DVI
versions of the following papers that relate to the functioning of the
system:
 \begin{enumerate}
 \item ``Multivalued logics: A uniform approach to inference in
artificial intelligence.''  This paper gives a formal description of
multivalued logics and the \mvl\ system.
 \item ``Bilattices and modal operators.''  This paper describes the
formal connection between the \mvl\ system and modal operators.
 \end{enumerate}

The \mvl\ sources are public domain.  You are free to use, modify, and
distribute them as you see fit, provided that such use or distribution
is not for commercial purposes.  No warranties are made with regard to
the performance of the system, however, and no commitment of support
has been made.

\subsection{Overview}

This user's guide has been written to help you install and use the \mvl\
system.  It describes the basic structure of the system, and shows you
how to take advantage of the flexibility of multivalued inference.

Essentially, \mvl\ is just like any other logic programming system.
There is a database consisting of those sentences about which the
system knows something, and there are facilities for determining
whether or not some query follows from the information the database
contains.  Viewed in this fashion, the \mvl\ system is not much
different from a conventional logic programming language such as
\prolog.

What makes \mvl\ unique is that each sentence in the database is
labelled with a {\em truth value\/}\index{truth value}.  Thus,
instead of simply assuming a statement to be true and dumping it into
the database, we will {\em label\/} the statement as true when we
insert it.  Of course, we can use other labels as well: true by
default, true by virtue of some assumption, or even false or false by
default.  \Mvl\ considers these truth values when determining how to
respond to a \prolog-like query.

This manual is divided into two parts.  In the first part (Sections
\ref{s.fol1} -- \ref{s.support}), we describe the use of the system
where the truth values assigned to sentences are simply ``true'' or
``false.''  Used in this fashion, \mvl\ is not much different from
\prolog.

The second part of the manual shows you how to use logics involving
more complex truth values.  We examine the logics that have been
predefined as part of the \mvl\ system, and describe the various ways
in which you can define your own logics.  Several tools for this
purpose have been included as part of the \mvl\ system.  In addition,
we describe \mvl's facilities for handling modal operators, which
are operators that accept logical sentences as arguments.

\subsection{Why MVL is unique}

As mentioned in the previous section, what distinguishes \mvl\ from
other declarative programming languages is its support for multiple
truth values.  There are other differences, however, and these will
be discussed in detail in subsequent sections.  Here's a quick summary:
 \begin{enumerate}
 \item As mentioned, the principal feature offered by \mvl\ is its
support of multiple truth values.
 \item The basic theorem prover used by the system includes facilities
to detect loops and control recursion encountered in proof search.
The first-order theorem prover is discussed in Section \ref{s.fo},
although we will have nothing further to say about recursion control.
 \item The system can handle negative subgoals containing free variables
as discussed in Section \ref{s.effort}.
 \item The \mvl\ unifier includes an occurcheck and full support for
sequence variables.  This is discussed in Section \ref{s.rep}.
 \end{enumerate}

\section{Installation and testing}

The \mvl\ system is currently developed and maintained on a
Sparcstation running Allegro Common Lisp and is tested only
infrequently on other platforms.  Although the system has been written
in standard Common Lisp to maximize ease of portability and no
problems have been encountered when moving the system to a variety of
platforms, performance on other systems is not guaranteed.

\Mvl\ currently supports both ANSI Common Lisp and Common Lisp as
described in Steele's 1984 book.  In order to determine whether or not
the ANSI standard is conformed to, the system checks to see if
\type{in-package} is a macro or a function.  If the former, it is
assumed that the entire ANSI standard is supported; if the latter, it
is assumed that the language is as in Steele and various extensions
are defined.  It is expected that some future release of \mvl\ will
support the ANSI standard only.

\subsection{Sun-4}

In order to load the system on a Sparcstation running Allegro Common
Lisp, do the following:
 \begin{enumerate}

 \item (ANSI) Set up a logical pathname translation so that logical
pathnames beginning with \type{MVL} point to the direction where \mvl\
is kept.  (non-ANSI) Edit the file \type{mvldcl.lisp} so that the
\type{:default-pathname} field in the \mvl\ system definition points
to the directory where \mvl\ will be kept.  This directory should also
contain the following files:
 \begin{enumerate}\setlength{\itemsep}{0pt}\raggedright
 \item {\tt f1.mvl, recursion.mvl, def-test.mvl, circ-test-1.mvl,
circ-test-2.mvl, circ-test-3.mvl, circ-test-4.mvl, completed-db-1.mvl,
completed-db-2.mvl, modal-test.mvl, actions.mvl, shooting.mvl}
 \item {\tt mvl.test, bc.test, recursion.test, atms.test,
defaults.test, circumscription.test, completed-db.test, modal.test,
fc.test, temporal.test}
 \end{enumerate}
 \item Compile and load the file \type{mvldcl.lisp}.
 \item Execute the function
\type{(mvl:compile-mvl)}\index{compile-mvl} to compile and load the
\mvl\ system.  Using this function in the future will only compile 
those files that have been changed since last compiled.
 \item Invoke the function \type{(use-package 'mvl)} to access the
functions and symbols of the \mvl\ system (which is kept in its own
package; see Section \ref{s.package}).
 \end{enumerate}

Now invoke the function {\tt (test-mvl)}, which should return 0
(i.e., 0 errors encountered).  If it doesn't, the problem is probably
that \mvl\ is unable to locate the files mentioned in (1) above.

\subsection{Symbolics 3600 series}

The \mvl\ system will run on Symbolics Genera 7.2 or later.  It is
possible to run the system using Genera 7.0 or 7.1, although a bug in
the implementation of the Common Lisp \type{\&allow-other-keys}
keyword may cause the system to dump to the cold load stream if
certain predicates are procedurally attached.  (See Section
\ref{s.attach}.)

You should begin by editing the file \type{mvldcl.lisp} so that the
default pathname in the \mvl\ system declaration points to the
directory in which the files are actually kept.  You can then load
the \mvl\ system using either the Load System command or the Lisp
function \type{(mvl:load-mvl)}\index{load-mvl} that is defined in the
file \type{mvldcl.lisp}.  Now access the \mvl\ package using
\type{(use-package 'mvl)} and execute the function \itype{test-mvl}.
This function should return 0 (i.e., 0 errors encountered).  If it
doesn't, the problem is probably that \mvl\ is unable to find the
directory its test files are on (they're included in the distribution
tape).

\subsection{Other machines}

The \mvl\ system has also been run and tested on a wide variety of
other machines.  In order to load the system, follow the instructions
provided for Sun-4's.

\section{The MVL database}
\label{s.fol1}
\subsection{Representation}
\subsubsection{Structure of the database}
\label{s.rep}

As already remarked, the \mvl\ database consists of sentences labelled
with truth values.  In this section, we will describe the form of the
sentences themselves; in the next section, we comment briefly on the
truth values used to label them.

Sentences in the database are represented as \lisp\ s-expressions
using the prenex normal form \index{prenex normal form} of the
standard logical connectives \type{NOT}, \type{OR} and \type{AND}.
Thus the statement that Fred is tall and Harry is not would be
written as:
 \[\type{(AND (TALL FRED) (NOT (TALL HARRY)))}\]

Any \lisp\ atom can be used as a constant term, with the exception of
variables (see below) and the following reserved connectives:
 \[\type{NOT AND OR => <= <=> IF IFF}\]
The connective \type{<=} is used in backward-chaining rules;
 \[\type{(<= CONCLUSION PREMISE-1 ... PREMISE-N)}\]
is logically equivalent to
 \[\type{(IF (AND PREMISE-1 ... PREMISE-N) CONCLUSION)}\]
as is the forward-chaining version
 \[\type{(=> PREMISE-1 ... PREMISE-N CONCLUSION)}\]
 Although logically equivalent, these expressions behave differently
in practice.  The version
 \[\type{(IF (AND PREMISE-1 ... PREMISE-N) CONCLUSION)}\]
 in which the connective \type{IF} is used represents a pair of
rules, depending on the value of the parameter \itype{*if-translation*}.
If the value is \type{bc} (the default), then two backward-chaining
rules are produced, so that
 \[\type{(IF (A) (B))}\]
is equivalent to
 \[\type{(AND (<= (B) (A)) (<= (NOT (A)) (NOT (B))))}\]
 the conjunction of the original rule and its contraposition.
If the value of \type{*if-translation*} is \type{fc}, then
forward-chaining rules are produced and the translation is to
 \[\type{(AND (=> (A) (B)) (=> (NOT (B)) (NOT (A))))}\]
 Finally, if \type{*if-translation*} has the value \type{mix}, the
pair of noncontraposed rules
 \[\type{(AND (<= (B) (A)) (=> (A) (B)))}\]
 is produced.

Equivalences are represented using either \type{IFF} or \type{<=>}.
Thus
 \[\type{(IFF (FOO FRED) (GOO FRED))}\]
is equivalent to
 \[\type{(AND (IF (FOO FRED) (GOO FRED)) (IF (GOO FRED) (FOO FRED)))}\]
 Translations involving \type{<=>} depend on the value of the
parameter \et\index{*equivalence-translation*}.  As above, if the
value is \type{bc} (the default), then two backward-chaining rules
are produced, so that
 \[\type{(<=> (FOO FRED) (GOO FRED))}\]
translates to
 \[\type{(AND (<= (FOO FRED) (GOO FRED)) (<= (GOO FRED) (FOO FRED)))}\]
 If \et\ has the value \type{fc}, two forward-chaining rules are produced
so that the translation is to
 \[\type{(AND (=> (FOO FRED) (GOO FRED)) (=> (GOO FRED) (FOO FRED)))}\]
 Finally, if the value is \type{mix}, one forward-chaining rule and
one backward-chaining rule are produced:
 \[\type{(AND (=> (FOO FRED) (GOO FRED)) (<= (FOO FRED) (GOO FRED)))}\]

Variables \index{variables} are represented using \lisp\ atoms
beginning with the character \type{?} and are handled as in \prolog,
so that free variables in database statements are implicitly
universally quantified.  Again as in \prolog, existential
quantification is handled via Skolemization; thus the
backward-chaining sentence ``Every rich person is happy'' and the
sentence ``There is a rich person'' would be written as
 \begin{equation}
\type{(<= (HAPPY ?P) (PERSON ?P) (RICH ?P))}
\label{e.if}
\end{equation}
and
 \[\type{(AND (PERSON P0017) (RICH P0017))}\]
respectively.

The variables in database sentences are standardized apart from the
other variables used by the system.  This is done by using uninterned
variables in the actual representation of a database fact, so that
the internal representation of (\ref{e.if}) would be
 \[\type{(<= (HAPPY \#:?P) (PERSON \#:?P) (RICH \#:?P))}\]
 The original variables are generally reinstated when \mvl\ is asked
to display a database proposition.

\subsubsection{Sequence variables}

In addition to the simple variables discussed in the previous
section, \mvl\ supports the use of ``sequence
variables.''\index{sequence variables} These variables begin with the
characters \type{\str}, and will match not a single atom, but an
arbitrary (possibly empty) sequence of atoms.

Thus, for example,
 \[\type{(FOO A B)}\]
 will unify with
 \[\type{(FOO \str X)},\]
 returning the binding list
 \[\type{((\str X . (A B)))}\]
 in which \type{?*X} has been bound to the list \type{(A B)}.  This
binding list will presumably be displayed as
 \[\type{((\str X A B))}.\]
 Similarly, the most general unifier of $\type{(FOO)}$ and
$\type{(FOO \str X)}$ is given by $\type{((\str X))}$.

The presence of sequence variables may cause the most general unifier
of two patterns to be nonunique.  For example, the result of unifying
 \[\type{(FOO A B)}\]
 with
 \[\type{(FOO \str A \str B)}\]
 can be any one of the three answers
 \begin{eqnarray*}
& \type{((\str A A B) (\str B))} & \\
& \type{((\str A A) (\str B B))} & \\
& \type{((\str A) (\str B A B))} &
 \end{eqnarray*}

For this reason, the unification functions return not a single answer,
but a {\em list\/} of possible unifiers.

\subsubsection{Functions to manipulate database sentences}
\label{s.manipulate}

The following functions and global parameters are used to manipulate
sentences in the database:
 \begin{description}
 \item [\itype{varp(x)}] Returns \T\ if $x$ is a variable, and
\NIL\ otherwise.
 \item [\type{varp?(x)}]\index{varp-query} Returns \T\ if $x$ is a
standard (i.e., non-sequence) variable, and \NIL\ otherwise.
 \item [\type{varp$*$(x)}]\index{varp-star} Returns \T\ if $x$ is a
sequence variable, and \NIL\ otherwise.
 \item [\itype{vartype(x)}] Returns ? if $x$ is a standard variable,
$*$ if $x$ is a sequence variable, and \NIL\ is $x$ is not a
variable.
 \item [\type{new-?var()}]\index{new-query-var} Returns a new
temporary variable.
 \item [\type{new-$*$var()}]\index{new-star-var} Returns a new
temporary sequence variable.
 \item [\itype{normal-form(p)}] Returns an equivalent form of $p$
suitable for insertion into the database (i.e., with the connectives
\type{IFF}, \type{IF} and \type{<=>} reduced to \type{<=} and \type{=>}).
 \item [\itype{*if-translation*}] Controls the fashion in which
\type{normal-form} deals with the connective \type{if}.
 \item [\itype{*equivalence-translation*}] Controls the fashion in which
\type{normal-form} deals with the connective \type{<=>}.
 \item [\itype{negate(p)}] Returns the negation of $p$.
 \item [\itype{standardize-operators(p)}] Returns an equivalent form
of $p$ with the connectives \type{<=>}, \type{IF} and \type{IFF}
removed.  (Note: This function does not use the control variables
\type{*if-translation*} and \et, since it is used when proving
sentences, not when inserting them into the database.)
 \item [\itype{cnf(p)}] Returns $p$ in conjunctive normal form, as a
list of conjuncts, where each conjunct is a list of disjuncts.  Thus
 \[\type{(CNF '(AND A B))}\]
returns
 \[\type{((A) (B))}\]
\item [\itype{dnf(p)}] Returns $p$ in disjunctive normal form, as a list
of disjuncts, where each disjunct is a list of conjuncts.  Thus
 \[\type{(DNF '(AND A B))}\]
returns
 \[\type{((A B))}\]
\end{description}

\subsection{Truth values}

The truth values used to label the sentences in the \mvl\ database
are elements of a bilattice that can be selected by the user with
the function \type{load-logic}.  We will defer a discussion of these
truth values to the second portion of this manual, which discusses
the use of bilattices that label sentences as other than simply true
or false.  At this point, we merely remark that if it is our intention
to assert the sentence \type{(not (foo))}, this is actually accomplished
by asserting that the negated sentence \type{(foo)} has truth value
\type{false}.

\subsection{Database form and internals}
\label{s.database}

The database is maintained using a discrimination net, as suggested
by Charniak et.~al in their book on {\sc lisp} programming.  Any
particular sentence in the database, say \type{(tall tom)}, is
associated with a \lisp\ atom such as \type{P27} that is its
``proposition number.''  The propositions are numbered consecutively,
beginning with \type{P1}.  At any given point in time, the function
\type{(props)} can be invoked to return a list of all of the
propositions in the database.

Very few of the functions that are used to manipulate the internal form
of the database are likely to be needed by a user of the \mvl\ system.  The
following are the most useful:
 \begin{description}
 \item [\itype{props()}] This function returns a list of all
propositions about which any information has been recorded.
 \item [\itype{denotes(a)}] Returns the proposition associated with a
particular \lisp\ atom.
 \item [\itype{truth-value(a)}] Returns the truth value associated
with a particular \lisp\ atom.  Thus, if the result of asserting
\type{(not (short tom))} into the database was the atom \type{P38},
\type{(denotes 'p38)} would evaluate to \type{(short tom)}, and
\type{(truth-value 'p38)} would evaluate to \type{false}. 
 \item [\itype{index(p)}] Updates the discrimination net as
needed to insert the proposition $p$ into the database.
 \item [\itype{unindex(p)}] Removes from the discrimination net
information relating to the proposition $p$.
 \end{description}

\section{Quantification}

As we have remarked, quantification is treated in \mvl\ in a fashion
similar to that used by \prolog\ systems:  Free variables appearing
in the database are assumed to be universally quantified; free variables
appearing in queries are existentially quantified.  Skolemization is
used to handle existential quantification in database statements or
universal quantification in queries.

The use of this approach means that the answer to an existentially
quantified query will often be of a form that indicates how the
variables should be bound in the result.  In \mvl\, these binding
lists \index{binding lists} are represented as expressions of the
form
 \[\type{((var-1 . bdg-1) ... (var-n . bdg-n))}\]
 where \type{var-1} through \type{var-n} are variables and
\type{bdg-1} through \type{bdg-n} are their respective bindings.
If no variable is bound to a value, the empty binding list \NIL\
is used.

As an example, suppose that the database contains the single sentence
 \[\type{(LIKES COCO ?X)}\]
 indicating that Coco likes everything.  Now the query
 \[\type{(LIKES ?Y MATT)}\]
 will be interpreted as, ``Find something that likes Matt.''  The answer
returned will be
 \[\type{((?Y . COCO))}\]
since Coco likes Matt.  Since the queries
 \[\type{(LIKES COCO MATT)}\]
and
 \[\type{(LIKES COCO ?Z)}\]
 both succeed without binding a variable, the answer \NIL\ will be
returned for each.\footnote{In order to distinguish the empty binding
list \NIL\ from a failed database query, all queries return a binding
list as well as a truth value.  This is discussed in Section
\ref{s.answer}.}

\subsection{Functions used to manipulate binding lists}

A variety of functions within \mvl\ are used to construct and manipulate
binding lists.  Here are some of the more useful ones:
 \begin{description}
 \item [\itype{groundp(p)}] Returns \T\ if the proposition $p$ is
ground (i.e., contains no variables) and \NIL\ otherwise.
 \item [\itype{vars-in(p)}] Returns a list of the variables
appearing in $p$.
 \item [\itype{get-bdg(var bdg-list)}] Returns the value assigned to
the variable \type{var} by the binding list \type{bdg-list}.  If the
variable does not appear in the binding list, \NIL\ is returned.
 \item [\itype{delete-bdg(var bdg-list)}] Destructively modifies
\type{bdg-list}, removing any binding for the variable \type{var}.
 \item [\itype{plug(p bdg-list)}] Returns the result of instantiating
the proposition $p$ using the given binding list.  This function is
nondestructive.
 \item [\itype{samep(p1 p2)}] Returns \T\ if the propositions $p_1$
and $p_2$ are identical (up to variable renaming), and \NIL\
otherwise.  If the optional \type{return-answer?} argument is
supplied and the propositions are the same, a binding list is
returned instead of \T.
 \item [\itype{unifyp(p1 p2)}] Returns a list of the most general
unifiers of $p_1$ and $p_2$, or \NIL\ if the propositions fail to
unify.
 \item [\itype{instp(p1 p2)}] If $p_2$ is an instance of $p_1$,
return a list of their most general unifiers; only variables in
$p_1$ are bound.  Returns \NIL\ if $p_2$ is not an instance of $p_1$.
 \item [\itype{matchp(pattern atom)}] Returns a list of elements of
the form \type{(b . d)} where $b$ is a most general binding list such
that \type{(plug pattern b)} is an instance of the proposition
labelled with \type{atom} and $d$ is a binding list for the variables
in the database proposition.  \NIL\ is returned if no such binding
list exists.  Thus if \type{P1} is \type{(TALL \#:?2)}, and \type{P2}
is \type{(TALL TOM)}, then
 \[\type{(matchp (TALL TOM) P1)}\]
returns a list consisting of the single entry 
 \[\type{(NIL . ((\#:?2 . TOM)))}\]
 The \type{car} of this entry is \NIL, indicating that no variables
in the query need to be bound to match \type{P1}; the \type{cdr} is
\type{((\#:?2 . TOM))} since the variable \type{\#:?2} in \type{P1}
needs to be bound to \type{TOM} to match the query.
 \[\type{(matchp (TALL ?X) P2)}\]
returns \type{((((?X . TOM)) . NIL))}.
 \end{description}

Additional functions are described in Section \ref{s.binding-dag}.

\subsection{Negative subgoals containing free variables}
\label{s.effort}

Consider the following \prolog\ fragment:
\begin{verbatim}
     f(x) :- not(g(x)).
     g(A).
\end{verbatim}
 The interpretation we will assign to the implication
\type{f(x) :- not(g(x))} will be that it is universally quantified
over $x$, and therefore corresponds to the implication
 \[\forall x.[\neg g(x) \supset f(x)]\]
 Given this interpretation, what should the response be if we provide
the system with the query \type{f(x)}?

In light of the negation as failure rule of \prolog\ (which has analogs
in the \mvl\ system), we would expect the query to succeed for all values
of $x$ {\em except\/} $x=A$.

Most \prolog\ systems cannot respond to queries such as this one, since
a nonground negative subgoal is encountered during the proof process.
One approach to these queries, suggested by David Chan, is to do just
what we have suggested in the previous paragraph, computing those binding
lists for which the query fails.

Unfortunately, there are many queries that succeed in general, but
for which it is impractical to compute all of the bindings for which
the query fails.  This may be because there are infinitely many such
failures, or because the time needed to compute them explicitly is
prohibitive.  Situations such as this arise frequently in
applications of declarative programming techniques to planning
problems; here, we would like simply to return an indication that the
query ``succeeds in general.''  In the above example, we might
report that the query had succeeded without binding the variable $x$;
even though $x=A$ is an exception, the report of success is true in
general.

Both approaches are supported by \mvl, and the global variable
\type{*completed-database*} is used to distinguish between them.  If
this variable has the value \T, all exceptions are computed; if it
has the value \NIL, the other approach is taken.  The default value
is \NIL:
 \begin{description}
 \item [\itype{*completed-database*}] If \T, find all exceptions
to default proofs.  If \NIL\ (the default), accept defaults proofs
with possible exceptions as complete.
 \end{description}

\section{Modifying the MVL database}

\subsection{Keyboard input}

The simplest functions to modify the \mvl\ database simply accept
propositions from the user and either add or remove them from the
knowledge base.  The primitive function for adding to the database is
\itype{stash}; the function \itype{unstash} is used for removal.

The function \type{stash} accepts a proposition and up to three
keywords:
 \begin{description}
 \item [\itype{value}]  This keyword is used to supply a specific truth
value to be assigned to the proposition.  If no value is supplied, the
default value is used (this value is supplied when the logic is chosen;
see Section \ref{s.bilattice}).
 \item [\itype{increment-p}] If the proposition in question already
appears in the database, the user has a choice: The given truth value
can be used to replace the one currently appearing in the database,
or it can be combined with the existing value using the combination
function supplied with the logic being used.  The truth value is
replaced if \type{increment-p} is \NIL, and is updated if it is \T\
(the default).
 \item [\apa\index{allow-procedural-attachment}] This keyword,
which defaults to \T, controls whether or not it is possible to
procedurally attach the action of stashing the supplied proposition.
See Section \ref{s.attach}.
 \end{description}

Somewhat more useful than \type{stash} is the function \itype{state}.
This function is identical to \type{stash}, except it translates the
supplied sentence into a form suitable for inference by using the
function \itype{normal-form} described in Section \ref{s.manipulate}.
In addition, \type{state} cannot be procedurally attached, so the
keyword \apa\ is not allowed.

For removing sentences from the database, the functions \type{unstash}
and \itype{unstate} are used.

The function \type{unstash} accepts a proposition to be removed from
the database, and up to three keywords: \type{value},
\itype{decrement-p} and \apa.  The last of these is treated as for
\type{stash}.  In general, if no keywords are used, the effect is to
remove the proposition from the database.

If \type{decrement-p} is \NIL, then \type{unstash} behaves like
\type{stash}, except the default value used for unstashing is not
necessarily the same as the default value used for stashing.

If \type{decrement-p} is \T, then the new truth value assigned to
the proposition is given by
 \begin{equation}
o \cdot \neg v
\label{e.dot-star}
\end{equation}
 where $o$ is the old value and $v$ is the value supplied via the
\type{value} keyword.  The $\cdot$ and $\neg$ operations are as
described in Section \ref{s.bilattice}.

There is one additional function for removing sentences from the
database, \itype{empty}.  This function accepts no arguments, but
does take a single keyword, \type{value}.  The function calls
\type{unstash} for every sentence in the database.

All of the functions described in this section work by calling the
primitive function \itype{new-val}.  This function accepts as
arguments a proposition name and a truth value, and sets the truth
value of the proposition to the given value.  If the truth value is
\type{unknown}, indicating that nothing is known about the
proposition, it is removed from the database and the discrimination
net.

\subsection{Input from files}

It is also possible to enter a database from a disk file.  These files
should in general have the extension \type{.MVL}, and should be
located either on the currently active disk area or on the \mvl\
directory.  The files should consist simply of the sentences to be
added to the database.

In addition to consisting of such sentences, the files may contain the
keyword \type{:value}, followed by the value to be assigned to each
sentence, or the keyword \type{:function}, followed by a function to
be applied to each sentence in order to {\em compute\/} the value to
be assigned to it.  The keyword \type{:empty} empties the database
before loading the file.  Examples of these constructs can be found in
the various \type{.MVL} files included on the distribution tape.

The sentences are loaded with \itype{mvl-load}, which accepts a file
name and an optional keyword \type{empty}.  If the keyword is
supplied, the database is emptied before the file is read; in any
event, \type{state} is then applied to each sentence in the given
file.  The sentences in the file can be removed from the database
with the function \itype{mvl-unload}, which accepts a file name as its
argument.  These functions each call \type{mvl-file}, which accepts
a string and returns a path to the \type{.MVL} file associated to it.

Related to \type{mvl-load} is \itype{mvl-save}, which also accepts a
file name as an argument, but instead of loading a declarative
database from the file, writes the current database {\em to\/} the
file.

\subsection{The editor interface} \index{editor interface}
\subsubsection{Sun EMACS editor}

In addition to the file commands \type{mvl-load} and
\type{mvl-unload}, there is an interface to the {\sc emacs} editor on
Sparcstations.  For a file of mode \mvl, certain keys are bound as
follows (each of these commands also recognizes the
\type{:value} and \type{:function} keywords):

 \begin{description}
 \item [\type{meta-S}] Applies \type{state} to each s-expression in the
current region.
 \item [\type{meta-U}] Applies \type{unstate} to each s-expression in the
current region.
 \item [\itype{meta-s}] Applies \type{state} to each s-expression in the
current defun (i.e., s-expression beginning in column 1).
 \item [\itype{meta-u}] Applies \type{unstate} to each s-expression in the
current defun.
 \end{description}

In order to use these commands, certain modifications will need to be made
to your \type{.emacs} file.  The file \type{dot.emacs} is an example of
what to do.

\subsubsection{Symbolics ZMACS editor}

In addition to the file commands \type{mvl-load} and
\type{mvl-unload}, there is an interface to the {\sc zmacs} editor on
Symbolics 3600 series machines.  For a file of mode \mvl, certain
keys are bound as follows (each of these commands also recognizes the
\type{:value} and \type{:function} keywords):

 \begin{description}
 \item [\itype{super-s}] Applies \type{state} to each s-expression in the
current region.
 \item [\itype{super-u}] Applies \type{unstate} to each s-expression in the
current region.
 \item [\itype{super-f}] Applies \type{fc} to each s-expression in the
current region.  See Section \ref{s.fc}.
 \item [\itype{super-e}] Applies \type{erase} to each s-expression in the
current region.  See Section \ref{s.fc}.
 \end{description}

\subsection{Summary of database modification functions}

\begin{description}
 \item [\itype{stash(p)}] Inserts $p$ into the database.  Allowed
keywords: \type{value}, \type{increment-p} and \apa.
 \item [\itype{state(p)}] Inserts $p$ into the database after
converting it to normal form.  Keywords: \type{value},
\type{increment-p}.
 \item [\itype{unstash(p)}] Removes $p$ from the database.  Keywords:
\type{value}, \type{decrement-p}, \apa.
 \item [\itype{unstate(p)}] Removes $p$ from the database after
converting it to normal form.  Keywords: \type{value},
\type{decrement-p}.
 \item [\itype{empty}] Unstashes every sentence in the database.  Keyword:
\type{value}.
 \item [\itype{new-val(a v)}] Sets the truth value of the sentence
associated with the atom $a$ to $v$.  If $v$ is the truth value
\type{unknown}, the sentence is removed from the database.
 \item [\itype{mvl-file(string)}] Returns the path for the \mvl\ file
associated with the given name.  The file name may include an
extension, one may be supplied as an optional argument, or
\type{.MVL} will be used by default.  There is an optional keyword
\type{no-error} indicating that even if the file cannot be found, the
path should be returned.
 \item [\itype{mvl-load(file)}] Loads the sentences in the given file into
the \mvl\ database.
 \item [\itype{mvl-unload(file)}] Removes the sentences in the
given file from the \mvl\ database.
 \item [\itype{mvl-save(file)}] Writes the \mvl\ database out to the
given file.  This function also accepts the keywords \type{cutoffs}
and \type{props}, which are treated as for \type{contents} (see
Section \ref{s.contents}).  Finally, the keyword \type{if-exists} is
passed to the \type{with-open-file} macro.
 \end{description}

\section{Querying the MVL database}

\subsection{MVL responses to queries}
\label{s.answer}

Because of the fact that a particular fact may be stored in the
database with any of a variety of truth values, \mvl\ responds to a
query by returning not simply a binding list (as most \prolog\
interpreters do, for example), but by returning a binding list {\em
and a truth value}. 

In actuality, what is returned is an instance of the \type{answer}
structure\index{answer structure}.  This structure contains the
following fields:
 \begin{description}
 \item[\type{binding}] Bindings for some of the variables in the query.
 \item[\type{value}] A truth value.  The bound version of the query has
this truth value.
 \end{description}

The definition of this structure leads Common Lisp to define the
access functions \type{answer-binding} and \type{answer-value};
in addition, a function \type{equal-answer} is defined by \mvl\ to
test two answers for equality:
 \begin{description}
 \item[\itype{answer-binding(a)}] Returns the binding field of the
answer \type a.
 \item[\itype{answer-value(a)}] Returns the value field of the answer
\type a.
 \item[\itype{equal-answer(a1 a2)}] Returns \T\ if the two answers
\type{a1} and \type{a2} are equal, and \NIL\ otherwise.  The
functions used are \type{binding-equal} and \type{mvl-eq}, the
functions normally used to test binding lists and truth values for
equality.
 \end{description}

Thus, in response to the query
 \[\type{(flies ?x)}\]
 the response of an answer with binding \type{((?x . Tweety))} and
truth value \type{dt} would indicate that
 \[\type{(flies Tweety)}\] 
 has truth value \type{dt}.  In other words, Tweety's flying is a
conclusion that is true by default.  (See Section \ref{s.j} for a
description of default truth values in the \mvl\ system.)

As a side effect, this allows the user to distinguish a response of
\NIL\ (indicating that the query has failed) from a response in which
the binding is \NIL\ (indicating that the query has succeeded without
binding any variables).

\subsection{Searching for specific sentences}

There are two functions that are most commonly used to retrieve facts
that have been stored in the database, \itype{lookup} and
\itype{lookups}.  Each of these functions accepts as arguments a
proposition to search for and up to four keywords, and returns two
values.

The keywords and their meanings are as follows:
 \begin{description}
 \item [\itype{cutoffs}] This keyword is used to give conditions on
the sentences being considered as possible matches to the query.  Any
sentence that is considered as a match must have an associated truth
value \type{v} such that \type{(test v cutoffs)} evaluates to \T\ (see
Section \ref{s.cutoffs}).  The default value of this keyword is defined
by the bilattice in use (see Section \ref{s.bilattice}).
 \item [\itype{inv}] This keyword defaults to \NIL.  If it is \T,
however, any truth values found should be negated before they are
returned.
 \item [\apa] This is similar to the keyword for \type{stash}.
 \item [\itype{succeed-on-modal-ops}] If \T, then any attempt to look
up a proposition involving a modal operator will immediately succeed;
if \NIL\ (the default), then the result returned will be the result
of applying the modal operator to the truth values of the
propositions that are its arguments.  See Section \ref{s.modal}.
 \end{description}

The two functions \type{lookup} and \type{lookups} are identical
except that \type{lookup} looks for only a single match for the
supplied proposition, while \type{lookups} searches for all such
matches.  Thus while \type{lookup} returns two values, \type{lookups}
returns two {\em lists\/} of values.  Each entry on each of these
lists is analogous to the corresponding value returned by
\type{lookup}.  The values returned by \type{(lookup q)} are as
follows:
 \begin{enumerate}
 \item First, an \type{answer} is returned.  If this answer contains
the binding list \type b and the truth value \type v, then applying
the binding list to the query produces a result with the given truth
value.  In other words, \type{(plug q b)} is an instantiation of a
sentence appearing explicitly in the database with truth value \type
v.
 \item Second, one of the following three values is returned:
 \begin{enumerate}
 \item If the sentence was found by procedural attachment, the atom
\type{attach} is returned.
 \item If the sentence was found by applying a modal operator, the
atom \type{modal} is returned.
 \item Under other conditions, the instantiated version of the
proposition that resulted in the match is returned.
 \end{enumerate}
 \end{enumerate}

As already mentioned, \type{lookups} returns two lists of values.

Here are summaries of \type{lookup} and \type{lookups}:

 \begin{description}
 \item [\itype{lookup(q)}] Searches the database for a sentence
matching the query $q$.  Allowed keywords: \type{cutoffs}, \type{inv},
\apa, \somo.  Returns two values: an answer and the source of the
match.  Returns \NIL\ if the query cannot be found in the database.
 \item [\itype{lookups(q)}] Similar to \type{lookup}, but looks
for {\em all\/} sentences matching the query.  Returns two lists of
values.
 \end{description}

\subsection{Searching the entire database}
\label{s.contents}

In addition to searching for a particular sentence, it is also
possible to display large portions of the entire \mvl\ database.  The
basic functions for doing this are \type{contents}, which displays
the entire database, and \type{prfacts}, which accepts a \lisp\ atom
as an argument and displays all of those facts concerning that atom.

Each of these two functions also accepts a variety of keywords:
 \begin{description}
 \item [\itype{cutoffs}] As for \type{lookup}, it is possible to
supply a test that must be passed by the truth value of any database
sentence before that sentence is displayed.
 \end{description}
 In addition, \type{contents} accepts a keyword \type{props}, which
defaults to \type{(props)} (see Section \ref{s.database}) and is a
list of the sentences to be displayed.

Related to \type{prfacts} is \itype{unstash-facts-with-atom}, which
accepts an atom as an argument and removes from the database any
fact concerning that atom.  \type{Unstash-facts-with-atom} accepts
the same keywords as \type{unstash} does.

Here, then, are the functions described in this section:

\begin{description}
 \item [\itype{contents}] Displays the contents of the database.  Keywords:
\type{cutoffs}, \type{props}.
 \item [\itype{prfacts(atom)}] Displays the database facts concerning the
given atom.  Keyword: \type{cutoffs}.
 \item [\itype{unstash-facts-with-atom(atom)}] Unstashes from the
database any facts concerning the given atom.  Keywords:
\type{value}, \type{decrement-p}, \apa.
 \end{description}

\subsubsection{Command line interface}

\paragraph{Sun}
In addition to the above functions, on Sparcstations, the key
$\langle$meta$\rangle$-* may be typed to the command line processor
running from the editor.  The result is that any information in the \mvl\
database concerning the currently open predicate will be displayed to the
user.

\paragraph{Symbolics}
On Symbolics systems, the key $\langle$super-shift$\rangle$-A may be typed
to the command line processor.  The result is that any information in the
\mvl\ database concerning the currently open predicate will be displayed
to the user.  This feature is the \mvl\ analog to the
$\langle$ctrl-shift$\rangle$-A key, which displays information about the
arguments to the currently open function.

\section{The first-order prover}
\label{s.fo}

As discussed in the technical papers distributed with \mvl, inference
in \mvl\ proceeds using a conventional first-order theorem prover; in
this section, we will describe this prover and the interface to it.
Be warned that the information in this section is unlikely to be used
by other than the most sophisticated users of \mvl; most readers will
want to proceed directly to Section \ref{s.inf}.

In general, there will be a variety of invocations to the first-order
prover active at any particular time.  Each such invocation is
associated with a task that the multivalued prover needs to complete,
and each such task is a \lisp\ object that is an instance of the
\type{consider} structure.\footnote{In fact, each task is an instance
of a structure that {\em includes\/} the \type{consider} structure.
See the file \type{bc.lisp} for further details.}

The multiple invocations of the first-order prover are run in
parallel by \mvl.  To each such invocation is associated an instance
of the structure \type{proof} that is used to record temporary
information about the progress of the proof effort.  We will discuss
first the first-order prover itself, then the construction and
termination of the invocations of it, and finally the structure of
the multivalued proving tasks.

\subsection{Nature}
\label{s.simple}

The prover itself is a simple natural-deduction style complete
theorem prover.  Only inferences against the initial clauses of
database statements are allowed; the result of this is that after
applying \type{state} to a sentence of the form
 \[\type{(<= CONCLUSION PREMISE-1 ... PREMISE-N)}\]
 \type{PREMISE-1} will always be tried first when the prover attempts
to prove \type{CONCLUSION}.  In a similar way, stating the
forward-chaining rule
 \[\type{(=> PREMISE-1 ... PREMISE-N CONCLUSION)}\]
 will result in a situation where only \type{PREMISE-1} will trigger
forward chaining.

Control of the proof search is handled by a function that evaluates
the ``merit'' of each of the ongoing proof attempts, and then
attempts to select the best one.  The evaluation function is stored
in the value of the global constant \itype{*node-evaluation-fn*},
which defaults to \itype{best-first-fn}, a function that attempts to
determine the value of a proof by looking at the truth value being
accumulated.  An alternative to \type{best-first-fn} is the function
\itype{depth-fn}, which evaluates a node simply by its depth in the
proof tree.

Three global parameters are used to control the proof search:
 \begin{description}
 \item [\itype{*final-limit*}] is the limit point beyond which the
search is abandoned, so that any node that evaluates to a value worse
than this will be pruned.  The default value of this parameter is
\NIL, indicating that no nodes are to be pruned.
 \item [\itype{*initial-limit*}] is used to indicate the type of search.
If \NIL\ (the default), then the values assigned to the nodes are not
used to order the search, and depth-first search is used.
 \item [\itype{*limit-delta*}] is used if \type{*initial-limit*} is
non-\NIL.  In this case, the nodes are ordered by increasing
 \[\frac{\type{value(node)}-\type{*initial-limit*}}{\type{*limit-delta*}}\]
 A typical value for *initial-limit* is 0, and for *limit-delta* is
1.  The value of this parameter is irrelevant if
\type{*initial-limit*} is \NIL.
 \end{description}

These three parameters may be reset by the user:
 \begin{description}
 \item [\itype{set-search-parameters}] This function is the primitive
used to modify the parameters used to control the first-order theorem
prover's search.  It accepts as keywords \type{initial-limit},
\type{limit-delta}, and \type{final-limit}.
 \item [\itype{activate-breadth-first}] Activate breadth-first search
by setting \type{*node-evaluation-fn*} to \type{depth-fn}.  Allowed
keywords: \type{initial-limit}, \type{limit-delta}, and
\type{final-limit}.
 \item [\itype{activate-depth-first}] Activate depth-first search by
setting \type{*initial-limit*} to \NIL.  Allowed keyword:
\type{final-limit}.
 \item [\itype{activate-best-first}] Activate best-first search by
setting \type{*node-evaluation-fn*} to \type{best-first-fn}.  Allowed
keywords: \type{initial-limit}, \type{limit-delta}, and
\type{final-limit}.
 \end{description}

\subsection{Invoking the prover}

The prover is invoked using the function \type{init-prover}.  This
function accepts as arguments a proposition and a pointer to the task
responsible for this proof attempt.

When \type{init-prover} is invoked, an instance of \itype{proof}
is created and returned that contains auxiliary information about the
proof attempt.  Here are the fields in the structure:
 \begin{description}
 \item [\itype{active-nodes}] A list of the nodes that are still
under active consideration.
 \item [\itype{answers}] A list of answers that have been found.
 \item [\itype{indexp} and \itype{indexn}] These entries are used to
index the various subtasks generated by the prover.
 \item [\itype{initial-limit}, \itype{final-limit}, \itype{limit-delta}
and \itype{evaluation-fn}] These entries are the values of the search
control parameters used in this proof attempt.
 \item [\itype{task}] A pointer back to the task responsible
for this invocation of the first-order prover.
 \item [\itype{vars}] The variables that appear in the proposition
being proven.
 \end{description}

Each node is an instance of the structure \itype{goal}, which is
described in the file \type{first-order.lisp}.

To continue a proof process, the function \type{cont-prover} is used.
This function accepts as an argument an instance of \type{proof} that
contains the auxiliary information for the proof process being
resumed.  The value returned by \type{cont-prover} is a structure
describing the results of the proof process itself;
\type{cont-prover} returns as soon as the prover finds a single
solution to its query.  The value returned is an instance of the
\type{bcanswer} structure\index{\type{bcanswer} structure}; the
fields in this structure are as follows:
 \begin{description}
 \item[\type{answer}] This field contains a binding list and truth
value that the system has determined can be assigned to the given
instantiation of the original query.
 \item[\type{just}] This field contains a list of those database
sentences that were used in the proof.
 \end{description}

If the proof attempt fails, \NIL\ is returned.

Here is a summary of the functions described in this section:

\begin{description}
 \item [\itype{init-prover(p effort)}] This function
initializes and returns an instance of \type{proof} containing
auxiliary information about the proof process.
 \item [\itype{cont-prover(proof)}] Resume the task associated to the
auxiliary information in the given instance of the \type{proof}
structure.  This function returns an instance of \type{bcanswer}
corresponding to the result obtained by the prover, and returns as
soon as a solution to the original query is found.
 \end{description}

\subsection{The multivalued proof tasks}
\label{s.consider}

The individual proof attempts generated by a multivalued query are
instances of structures that include the structure
\type{consider}\index{consider structure}.  The exact details of this
and related structures are unlikely to be of much interest to a user
of the \mvl\ system, but we include them here in the interests of
completeness.  These are the fields in the
\type{consider} structure:
 \begin{description}
 \item [\itype{prop}] The proposition being investigated by this task.
 \item [\itype{parent}] The task for which the current task is a
child.
 \item [\itype{children}] Children of the current task.
 \item [\itype{truth-value}] The current best estimate of the truth value
that should be assigned to the proposition being considered.
 \item [\itype{relevant}] A list of answers that, if returned by this
task, are known to be relevant to the original query.
 \item [\itype{irrelevant}] A list of answers know {\em not\/} to be
relevant to the original query.
 \item [\itype{paths}] The \mvl\ system goes to a great deal of effort
to ensure that it does not consider proofs that will not affect the
eventual response to a user's query; intermediate information used in
this computation is stored in this field.  Further details can be found
in the file \type{bc.lisp} and the files to which it points.
 \end{description}

When the prover must select a task to work on next, it does so via the
following function:
 \begin{description}
 \item [\itype{choose-task}] Selects a task to work on next.  The
current implementation of \mvl\ simply selects any active task (see
Section \ref{s.bc}).
 \end{description}

\section{Inference}
\label{s.inf}

\subsection{Backward chaining}
\label{s.bc}

The information in Section \ref{s.fo} is more technical than most users
of \mvl\ are likely to need.  What is of much greater interest are those
database query tools that are like \type{lookup} and \type{lookups} but,
instead of simply looking for a fact in the database, attempt to derive
it as a consequence of other database information.  These functions use
those of the previous section as described in the accompanying technical
papers.

The functions accept three sorts of cutoff information, instead of the
simple information provided to \type{lookup} and \type{lookups}.  The
first, provided via the \type{cutoffs} keyword, is used to construct a
test to decide whether an answer computed by the prover is good enough
to be returned when the proof effort is complete.  The second,
provided via the \type{succeeds} keyword, indicates that an answer
should be returned as soon as it can be shown that the final answer
will pass the given test.  The final cutoff, supplied via the
\type{terminate} keyword, indicates that the answer should be returned
as soon as there is any indication that the final answer {\em might\/}
pass the given test.

As an example, suppose that we were trying to find a plan to achieve some
goal.  We might well be prepared to terminate our search as soon as we found
a plan that was {\em guaranteed\/} to achieve the goal; if we completed our
search {\em without\/} finding such a plan, we might then want to return
all plans that we {\em expected\/} would achieve the goal.  It is to
facilitate such possibilities that the backward-chaining functions expect
three sets of cutoff information.  The default value for \type{cutoffs}
is the standard one associated to the active bilattice (see Section
\ref{s.bilattice}); the default choice for \type{succeeds} corresponds
to complete information for the singular versions of the queries and
guarantees failure for the plural versions (since an answer might be
about to be discovered for a different instantiation of the query).
The default value for \type{terminate} always fails.

Here, then are the various backward-chaining functions.  They all
return two values: an answer (or list of answers for the plural
versions \type{bcs} and \type{consider}) and an instance of the
\type{analysis} structure describing tasks remaining to be considered
by the prover.  In addition, they all accept an optional
\itype{analysis} keyword that, if supplied, should be an instance of
the \type{analysis} structure that is to be continued instead of a
new one being created.  This is to allow the user to continue proof
efforts that have returned for some reason.
 \begin{description}
 \item [\itype{bc(q)}] Attempts to prove the query $q$ from information
in the database, looking only for a single answer.
 \item [\itype{bcs(q)}] Attempts to prove the query $q$ from information
in the database, looking for all possible answers.
 \end{description}

All of these functions work by maintaining an instance of the
\type{analysis} structure\index{\type{analysis} structure} that
includes a variety of information about the status of the current
proof attempt.  Here are the fields in this structure:
 \begin{description}
 \item [\type{prop}] The proposition being considered by the prover.
 \item [\type{cutoffs}] The cutoff information that must be satisfied
by any result before it is to be returned.
 \item [\type{succeeds}] The cutoff information that will allow the
prover to terminate early if an answer is guaranteed to pass the test.
 \item [\type{terminate}] The cutoff information that will allow the
prover to terminate early if an answer might pass the test.
 \item [\type{active-tasks}] A list of the proof tasks (see Section
\ref{s.consider}) that still need to be investigated.
 \item [\type{main-task}] The task that was created upon the 
original invocation of the prover.
 \item [\type{returned-answers}] A list of answers that have already
been returned to the user; this is to keep \mvl\ from compulsively
re-returning these answers if the \type{analysis} is reinvoked.
 \end{description}

The top-level functions using this structure are as follows:
 \begin{description}
 \item [\itype{init-analysis(prop)}] The \mvl\ primitive upon which
the proving functions are based, this function accepts a proposition
and creates and returns an instance of the \type{analysis} structure.
 \item [\itype{cont-analysis(analysis)}] This function continues the
investigation of a specific analysis, and can be used to resume a
proof effort that has been curtailed for some reason.  It returns two
values: an instance of the \type{answer} structure giving the results
of the proof attempts, and a modified version of the \type{analysis}
structure that can be used to continue the proof process if desired.
 \end{description}

The functions \type{init-analysis} and the user functions such as
\type{bc} all accept the following keywords:
 \begin{description}
 \item [\type{cutoffs}] Cutoff test as described above.
 \item [\type{succeeds}] Success test as described above.
 \item [\type{terminate}] Termination test as described above.
 \item [\type{initial-value}] The truth value assigned to the query in
the absence of database information.
 \item [\type{analysis}]  If supplied, use this analysis instead of
creating a new one.
 \end{description}

\subsection{Forward chaining}
\label{s.fc}

Forward chaining in \mvl\ is not as simple as it is in many other
logic programming languages because of the \nm\ nature of multivalued
inference.  It follows from this that inserting a new sentence into
the database may cause a previous conclusion to be retracted, and
this may in turn cause other sentences to become unsupported.

There are two functions that use this information to modify the
database.  The more common is called \itype{fc} and accepts as an
argument a proposition to be added to the database.  The proposition
is added and provided that it was not already in the database when the
forward chainer was invoked, its consequences are also added to the
database and forward chained upon.  The \type{fc} function accepts
the following keywords:
 \begin{description}
 \item [\itype{value}] This keyword is similar to that used by
\type{stash} and \type{state}.
 \item [\itype{increment-p}] This keyword is also similar to that used
by \type{stash} and \type{state}.  The default value is \NIL.
 \item [\itype{fc-increment-p}] When the information content of some
sentence in the database has {\em fallen}, should this drop be
propagated to other sentences that appear to depend on the sentence
in question?  The default value is \T.
 \item [\itype{cutoffs}] Only forward chain if the sentence being
inserted into the database has a truth value that passes the test
corresponding to \type{cutoffs}.  (See Section \ref{s.cutoffs}.)  The
default value used is the one associated with the active bilattice.
 \item [\itype{include-modal}] If \T\ (the default), then the forward
chainer also looks for rules that fire when the truth value of a
modal expression involving $p$ changes (where $p$ is the sentence
being forward chained upon).  See Section \ref{s.modal}.  The reason
this keyword is included is that searching for rules of this form is
computationally expensive.
 \end{description}

Slightly less common is \itype{erase}, the function used to remove a
proposition from the database.  This function accepts a database
sentence as an argument, and also the keywords \itype{value},
\itype{decrement-p} and \itype{cutoffs}, which are similar to the
corresponding keywords for \type{fc}.

Summarizing:
 \begin{description}
 \item [\itype{fc(p)}] Inserts the sentence $p$ into the database and
forward chains.  Allowed keywords: \type{value}, \type{increment-p},
\type{fc-increment-p}, \type{cutoffs} and \type{include-modal}.
 \item [\itype{erase(p)}] Removes $p$ from the database and forward
chains.  Allowed keywords: \type{value}, \type{decrement-p} and
\type{cutoffs}.
 \end{description}

\subsection{Debugging tools}

Debugging facilities are limited in \mvl.  The facilities that exist
describe the prover's activities during backward-chaining,
forward-chaining, or while the first-order theorem prover is looking
for a conventional proof of some proposition.

\subsubsection{Backward chaining}

Diagnostics are displayed during backward chaining using the two
global variables \itype{*bc-trace*} and \at\index{*anytime-trace*}.

Whenever a task is encountered that involves a proposition that is an
instance of some element of \type{*bc-trace*}, information concerning
the task will be displayed to the user.  The variable \at\ is
similar, but only reports when an {\em inference\/} task is
encountered.  An inference task is one that is trying to prove
something, as opposed to one analyzing a modal operator; this
distinction is described in greater detail in the file {\type
bc.lisp}.  In practice, the distinction between \type{*bc-trace*} and
\at\ is that the former should be used for debugging purposes, while
the latter should be used if you are interested in watching the
prover approximate and then refine its eventual answer to some
declarative query.

Although the user can manipulate \type{*bc-trace*} and \at\ directly,
there are four functions written to help do so:
 \begin{description}
 \item [\itype{bc-trace(pat)}] Include \type{pat} on
\type{*bc-trace*}, so that any instance of \type{pat} will be
traced.
 \item [\itype{bc-untrace(pat)}] Remove any instance of \type{pat}
from \type{*bc-trace*}.
 \item [\itype{anytime-trace(pat)}] Include \type{pat} on \at, so
that any instance of \type{pat} will be traced.
 \item [\itype{anytime-untrace(pat)}] Remove any instance of
\type{pat} from \at.
 \end{description}
 For all of these functions, the argument is optional.  If omitted,
\type{(?*)} (of which everything is an instance) is used.

\subsubsection{Forward chaining}

Tracing the activity of the forward chainer is essentially identical
to tracing the behavior of the backward chainer; the global variable
\type{*fc-trace*} is used and is manipulated using the functions:
 \begin{description}
 \item [\itype{fc-trace(pat)}] Include \type{pat} on
\type{*fc-trace*}, so that any instance of \type{pat} will be
traced.
 \item [\itype{fc-untrace(pat)}] Remove any instance of \type{pat}
from \type{*fc-trace*}.
 \end{description}
 Once again, the argument is optional.

\subsubsection{The first-order prover}

Tracing the activity of the first-order theorem prover itself is
somewhat different.  The facility for doing this is similar to the
Common Lisp \type{trace} facility, but instead of reporting when a
function is invoked, it reports when an attempt is made to prove an
instance of a particular predicate, and also when such an instance is
found successfully in the database.

The functions used to control this are \itype{bc-watch} and
\itype{bc-unwatch}.  They are similar to \type{trace} and
\type{untrace}:

 \begin{description}
 \item [\itype{bc-watch(preds)}] Accepts as arguments any number of
predicates that should be watched as the prover proceeds.  Each
predicate can also be a list of the form \type{(predicate depth)}
where \type{depth} is the number of successive ancestors to be
displayed when the prover encounters an instance of the predicate.
(This can be used to understand {\em why\/} a particular clause is
relevant to the proof being attempted.)  \type{bc-watch} returns a
list of the predicates currently being {\sc watch}ed.
 \item [\itype{bc-unwatch(preds)}] Removes the given predicates from
the list of {\sc watch}ed predicates.  If no arguments are supplied,
all predicates are un{\sc watch}ed.
 \end{description}

There are two additional features involved in tracing the activities
of the first-order prover.  The first is used whenever the global
parameter \itype{*save-nodes*} is \T\ (which it is by default).  When
the first-order prover displays information regarding a particular
proof node, the node is labelled with a number.  Subsequent nodes
will then use this label if they need to refer to the node in
question (e.g., as a parent).

If \type{*save-nodes*} is \T, then \type{*break-node*} can be set
to an integer; when a node is displayed that is labelled with this
integer, a \type{(break)} will be executed.  This is to allow the
user to interrupt the first-order prover if he needs to examine the
proof stack or some other \mvl\ internal.

 \begin{description}
 \item [\itype{*save-nodes*}] If \T\ (the default), nodes are
labelled with consecutive integers when they are displayed by the
first-order prover.
 \item [\itype{*break-node*}] The number of a node at which the
proof process should be interrupted.
 \end{description}

\section{Procedural attachment}
\label{s.attach}

In some cases, it may be desirable to store a fact in the database
{\em without\/} using the functions supplied in \mvl.  Suppose, for
example, that we have an $n+1$-ary relation $R$ that corresponds to
an $n$-ary function $f$, so that $R(x_1,\dots,x_n,y)$ holds just in
case $f(x_1,\dots,x_n)=y$.  Now, if we assert $R(c_1,\dots,c_n,d)$
into the database, where $d$ and the $c_i$ are constants, we should
probably also {\em remove\/} any sentences of the form
$R(c_1,\dots,c_n,d')$, where $d'\neq d$.  This sort of a situation
might arise if we were monitoring the temperature in a room, and
represented the observed temperature using a database sentence such
as
 \begin{equation}
\type{(TEMPERATURE ROOM 72)}
\label{e.temp}
\end{equation}

We handle this by informing \mvl\ that to stash a sentence of the
form \type{(TEMPERATURE ...)}, it should use not the usual function
\type{stash}, but some other user-defined function, say
\type{unique-stash}.  This new function would first \type{unstash}
any sentence that agreed with the given one except for the value of
the final term, and would then invoke \type{stash} on the original
sentence, but with \apa\ set to \NIL\ so that \type{unique-stash} was
not invoked a second time.

To modify the function used by \mvl\ to stash sentences such as this,
the user should simply place the desired function (in this case,
\type{unique-stash}) on the \type{stash} property of the relation in
question (in this case, \type{TEMPERATURE}).\footnote{As described in
Section \ref{s.invocation}, the \type{stash} property of a relation
can also be an association list, if multiple invocations of the
system are running in parallel.} The system will then automatically
invoke \type{unique-stash} instead of \type{stash} when a fact such
as (\ref{e.temp}) is asserted.

The functions within \mvl\ that can be modified in this fashion are
\type{stash}, \type{unstash}, \type{lookup} and \type{lookups}.  In
addition, the procedural attachments for \type{lookup} and
\type{lookups} are linked, so that if \type{lookup} has a procedural
attachment and \type{lookups} does not, a call to \type{lookups} will
invoke the procedural attachment for \type{lookup} and then pluralize
the result.  Similar remarks hold if \type{lookups} and not
\type{lookup} has a procedural attachment.

The procedural attachment is effected using the following functions:
 \begin{description}
 \item [\itype{invoke(fn p)}] This function invokes \type{fn} on the
sentence $p$.  In addition, any keyword arguments are passed to
\type{fn}, but with \type{\&allow-other-keys} so that the keywords
will be ignored if \type{fn} is not expecting them.
 \item [\itype{lookup-call(fn p)}] This function is similar to
\type{invoke}, but modifies the values returned to make them agree
with those expected to be returned by \type{lookup}.  Unless the
function \type{fn} supplies them, the truth value is assumed to be
\type{true}.  The second value is \type{attach}.
 \item [\itype{lookups-call(fn p)}] The plural version of
\type{lookup-call}.
 \end{description}

\subsection{Facilities to create procedural attachments}

Some facilities exist to create procedural attachments within the
\mvl\ system.  The following two macros are used for this purpose:
 \begin{description}
 \item [\itype{lisp-predicate(pred)}] This macro is used to indicate
that the given predicate can be evaluated by a \lisp\ call.  It will
also cause \mvl\ to signal an error if an attempt is made to stash
or unstash an instance of the function.
 \item [\itype{lisp-defined(fn)}] This macro is used to indicate that
the given function can be evaluated by a \lisp\ call.  It will also
cause \mvl\ to signal an error if an attempt is made to stash or
unstash an instance of the function.  \type{Lisp-defined} also
accepts an optional argument, which is an inverter function.  The
overall result of attempting to lookup an instance of the function
can be described as follows:
 \begin{enumerate}
 \item If none of the arguments to the function (i.e., all but the
last of the arguments of the relation) is a variable, then the function
is called and the result is unified with the final argument of the
relation.
 \item If the inverter function has been supplied and is equal to \T,
then the function is called and the result is unified with the final
argument of the relation even if some of the arguments are variables.
 \item If some of the arguments are variables but the result is not a
variable, then the inverter function is invoked if one has been
supplied.  The arguments passed to the inverter function are the
result (i.e., the last argument of the relation), and a list of all
of the arguments (i.e., all but the last argument of the relation).
If there is no inverter function, then \NIL\ is returned.  In
addition, if \wbli\ \index{*warn-on-bad-lisp-invocations*} is not
\NIL, the user is warned.  (The parameter \wbli\ defaults to \NIL.)
 \item If both some of the arguments and the result are variables, then
\NIL\ is returned.  In addition, if \wbli\ is not \NIL, the user is
warned.
 \end{enumerate}
 \end{description}

The current version of \mvl\ invokes \type{lisp-predicate} on the
following \lisp\ predicates: \type{numberp}, $<$, $>$, \type{equal},
\type{null}, \type{atom}, \type{listp}, \type{varp}, \type{varp?},
\type{varp$*$}, \type{groundp}, and $=$.  The procedural attachment
to $=$ attempts to unify its arguments.

The macro \type{lisp-defined} is invoked on the following:
\type{car}, \type{cdr}, \type{cadr}, \type{cons}, \type{list},
\type{list*}, \type{revappend}, \type{append}, \type{reverse},
\type{length}, \type{assoc}, \type{delete}, \type{+}, \type{*},
\type{-}, \type{/}, \type{floor}, \type{ceiling}, \type{truncate},
\type{round}, \type{float}, \type{mod}, \type{rem}, \type{max},
\type{min}, \type{sin}, \type{cos}, \type{atan}, \type{sqrt}, 
and \type{eval}.

\subsection{Other procedurally attached relations}

In addition to the \lisp\ functions and predicates described in the
previous section, the following relations have been defined and
are procedurally attached to suitable \lisp\ functions:
 \begin{description}
 \item [\itype{(bagof}\type{ ?x ?p ?b)}] This relation holds for $x$,
$p$ and $b$ if and only if $b$ is a list of all $x$ satisfying $p$.
 \item [\itype{(done exp-1 ... exp-n)}] This predicate, which is
useful for stashing only, causes \lisp\ to evaluate each of the subsequent
expressions.
 \item [\itype{(value symbol val)}] Stashing this predicate sets the
given symbol to the value \type{val}.  Unstashing it makes the symbol
unbound if its current value unifies with \type{val}.
 \item [\itype{(property symbol value indicator)}] Stashing this
predicate sets the symbol's \type{indicator} property to \type{value}
for the current invocation of \mvl.  Unstashing it removes the
\type{indicator} property for the current invocation of \mvl\ if the
previous value of this property unifies with \type{value}.  (See the
definitions of \type{get-mvl} and \type{rem-mvl} in Section
\ref{s.invocation}.)
 \item [\itype{(known p)}]  Stashing this predicate simply stashes $p$;
unstashing it causes an error.  Looking it up looks up $p$.
 \item [\itype{(unknown p)}]  Stashing this predicate unstashes $p$;
unstashing it causes an error.  Looking it up fails if $p$
can be found in the database, and succeeds otherwise.
 \item [\itype{(member item list)}] Tries to unify \type{item} with
the various elements in \type{list}.
 \end{description}

\section{Testing the MVL system}
\label{s.foln}

\subsection{As a diagnostic tool}

There are two functions available to help you to test both \mvl\ and
any applications you write using it.  The first, \type{test-mvl}, is
one you should already have run; if \mvl\ has been installed successfully,
it should return 0.

\type{Test-mvl} works by invoking the generic testing function
\type{mvl-test-file} on the file \type{mvl.test} in the \mvl\
directory.\footnote{\type{Test-mvl} works in its own invocation of
\mvl, so that the existing database is not modified.  This is not true
of \type{mvl-test-file}, however.} This file also contains
documentation regarding the details of the various tests being
executed.

In general, \type{mvl-test-file} accepts as input a file name, and
looks for a file with an extension \type{.TEST} on either the \mvl\
or the current directory.  It then works through the file by
evaluating the first s-expression it contains, and comparing the
value returned to the second s-expression.  If these two
s-expressions match, the next s-expression is evaluated, and so on.
If there is a mismatch, \type{mvl-test-file} reports the error;
eventually, it returns the total number of errors found in the entire
file.  When doing the comparison, \type{**} acts as a wildcard.
Thus, if the second s-expression is simply \type{**}, it will always
match; if it is \type{(** ** **)}, it will match whenever the first
expression returns a list of three values, and so on.

\begin{description}
 \item [\itype{test-mvl}] Tests the \mvl\ system by invoking the function
\type{mvl-test-file} on the file \type{MVL.TEST}.
 \item [\itype{mvl-test-file(string)}] Can be used to evaluate
various forms by the \mvl\ system and compare the results to
predefined values.  Both the forms and the expected answers are in
the \type{.TEST} file whose name is supplied to the function.
Returns the number of errors found.
 \end{description}

\subsection{As a tutorial}

The information used by the function \type{test-mvl} can also be used
as a tutorial in the \mvl\ system.  Take a look at the file
\type{mvl.test} that this function uses; you will see that the first
two lines are
 \begin{verbatim}
     (notime (mvl-test-file "bc" "backward chainer")
     0
\end{verbatim}
 This means that the function \type{mvl-test-file} will be invoked on
the file \type{bc.test}, and that this function should return 0
(i.e., 0 errors).  So let's have a look at \type{bc.test}:
 \begin{verbatim}
     (load-logic #\f)                    ;first-order logic
     **
\end{verbatim}
 Here, the system loads the bilattice corresponding to first-order
logic and is not interested in the value returned by the loading
operation.  Next, a declarative database corresponding to a full
adder is loaded (with the timer turned off).  The sentences in this
database can be found in the file \type{f1.mvl} and are also
described in Section 15 and Appendix B.2 of the paper ``Multivalued
logics: A uniform approach to inference in artificial intelligence.''

Next, we see in \type{mvl.test} the two lines
 \begin{verbatim}
     (get-bdg '?x (answer-binding (bc '(val (out 1 f1) 1 ?x))))
     OFF                   ;the first output (the sum line) should be off
\end{verbatim}
 We invoke the backward chainer on the sentence
 \[\type{(val (out 1 f1) 1 ?x))}\]
 to see if we can prove that the first output of the full adder takes
some value \type{?x}.  Then we get the binding of this variable to
see what the value actually is -- \type{OFF}, in this case.  The
remainder of the various test files can be analyzed similarly; many
of the examples in these files can also be found in the various
technical papers included in the \mvl\ distribution package.

\section{System support}
\label{s.support}

\subsection{Functions that compile and load the MVL system}

The following three functions can be used to manipulate the \mvl\
system as a whole:
 \begin{description}
 \item [\itype{load-mvl}] Loads the \mvl\ system.
 \item [\itype{compile-mvl}] Compiles and loads the \mvl\ system.  On
an Allegro system, the keyword \type{recompile} forces a recompile of
all of the files and the keyword \type{reload} is ignored (since the
files will all be reloaded anyway).  On a Symbolics system, the
keyword \type{reload} forces all of the files to be reloaded and the
keyword \type{recompile} forces them all to be recompiled.
 \item [\itype{release-mvl}] (Symbolics only)  Compiles and loads the
\mvl\ system, and declares the new version to be \type{released}.  The
keyword \type{no-recompile} can be used to skip the compilation step.
 \end{description}

\subsection{The MVL package}
\label{s.package}

In order to avoid name conflicts with existing code, the \mvl\ system
resides in its own package, \type{MVL}\index{MVL package}.  All
of the symbols and functions appearing in the appendix are exported
from this package and can therefore be accessed by invoking the
function \type{use-package(mvl)}.  When a new bilattice is loaded,
the modal operators it defines are also exported. 

To use the \mvl\ system, invoke the function \type{(use-package 'mvl)}
to make these symbols accessible.\footnote{As an alternative,
\type{(in-package 'mvl)} wil put you in the MVL package itself.}

\subsection{Running several copies of MVL in parallel}
\label{s.invocation}

Some users may wish to run several copies of the \mvl\ system
simultaneously, perhaps associating them with distinct {\sc lisp}
processes or with individual agents.  (The use of multiple copies of
\mvl\ to represent distinct components of a single agent's knowledge
goes against the grain of the multivalued approach and is not
recommended.)  In order to support this facility, the following
functions and macros have been defined:
 \begin{description}
 \item [\itype{existing-mvl-invocation}] This function returns a
single value that is an instance of the structure \mi.  It is useful
primarily for saving the current copy of \mvl\ so that it can be
restored later.
 \item [\itype{create-mvl-invocation}] This function sets up a new
copy of \mvl.  It accepts an optional argument that is an instance of
\mi; this can be used to reset the state of \mvl\ to a previous one.
If the optional argument is omitted, default values are used.  The
value returned is the \mvl\ invocation that was active before the
function was called, so that another call to
\type{create-mvl-invocation} can be used to restore the original
environment.
 \item [\itype{with-mvl-invocation}] This is a macro similar to the
Common Lisp macro \type{with-open-file}.  It accepts an optional
argument that is an instance of \mi\ and is followed by a body of
forms (presumably including a \type{load-logic}) that are evaluated
in a locally-bound copy of \mvl.
 \end{description}

These functions work by locally or globally rebinding the values of
the various parameters used by the \mvl\ system
(\type{*active-bilattice*}, \type{true} and so on).  The only
difficulty with this is that procedural attachments for a relation
$R$ use the property list of $R$, and this cannot be rebound.

In order to cater to this difficulty, there is an additional global
variable \itype{*mvl-invocation*} that labels the currently active
copy of \mvl.  To determine what procedure is attached to the
\type{stash} operation for $R$, the \type{stash} property of $R$ is
examined.  There are now the following possibilities:
 \begin{enumerate}
 \item If the value is \NIL, \type{stash} is not procedurally attached,
so the usual \mvl\ function is called.
 \item If the value is an atom, then it is assumed that this atom is
the name of a function that is the procedural attachment for \type{stash}
independent of the copy of \mvl\ that is currently active.
 \item If the value is a list, it is treated as an association list
and a cons whose car is \type{*mvl-invocation*} is searched
for.
 \begin{enumerate}
 \item If such a cons is found, its cdr is used as the procedural
attachment to \type{stash}. 
 \item If no such cons is found, but there is a cons whose car is \T,
then the cdr of that cons is used.
 \item If no cons has either \type{*mvl-invocation*} or \T\ as its
car, then \type{stash} is assumed to be unattached.
 \end{enumerate}
 \end{enumerate}
 The other functions described in Section \ref{s.attach} are handled
similarly.

The association lists described above are manipulated using the following
functions:
 \begin{description}
 \item [\itype{get-mvl (symbol indicator)}]  Like the Common Lisp
\type{get}, but if the value is a list, treats it as an a-list and
searches for an item whose car matches \type{*mvl-invocation*}.
The value taken can be changed with \type{setf}.
 \item [\itype{rem-mvl (symbol indicator)}] Like the Common Lisp
\type{remprop}, but if the value is a list, treats it as an a-list
and removes the item whose car matches \type{*mvl-invocation*}.
 \end{description}

\section{Selecting a bilattice}
\label{s.bilattice}

\subsection{Bilattice structure}

Of course, most of the power of the \mvl\ system comes not from the
features that were described in Sections \ref{s.fol1}--\ref{s.support}
of this manual, but from the fact that the user can change
representation language conveniently within the system.  He does this
by changing the {\em bilattice\/} that he is using.

A bilattice is the mathematical structure from which the truth values
used by the system are selected, and the formal ideas underlying this
approach are discussed at length in the accompanying technical
papers.  In the \mvl\ system, a bilattice is described as an instance
of the \type{bilattice} structure.  This structure has 27 fields; of
these, 17 are used to provide the formal information required by the
system, and the others are used to simplify the user interface.  Here
are 13 of the fields used to describe the bilattice itself:

\begin{description}
 \item [\itype{true}] A \lisp\ object that is the truth value assigned
to a sentence that is unconditionally true, such as the tautology
\type{(OR (P) (NOT (P)))}.
 \item [\itype{false}] A \lisp\ object that is the truth value assigned
to a sentence that is unconditionally false, such as the negation of a
tautology.
 \item [\itype{unknown}] A \lisp\ object that is the truth value assigned
to a sentence about which nothing is known.
 \item [\itype{bottom}] A \lisp\ object that is the truth value assigned
to a sentence that is both true and false.  If any sentence is assigned
this truth value (or can be shown to have it), it indicates the presence
of a contradiction in the database.
 \item [\itype{eq}] A \lisp\ function that takes two truth values as
arguments and returns \T\ if they are equal and \NIL\ if they are
not.  This function is needed because it may be possible to represent
a single element of the bilattice as a \lisp\ object in a variety of
fashions.
 \item [\itype{and}] A \lisp\ function that takes two truth values as
arguments and returns the truth value corresponding to their
conjunction.  If the two sentences $p$ and $q$ are labelled with
truth values $v_1$ and $v_2$ respectively, then the conjunction of
the truth values is the truth value that can be assigned to $p\wedge q$.
 \item [\itype{or}] A \lisp\ function that takes two truth values as
arguments and returns the truth value corresponding to their
disjunction.
 \item [\itype{not}] A \lisp\ function that takes a truth value as its
argument and returns the truth value corresponding to its negation.
 \item [\itype{plus}] A \lisp\ function that takes two truth values as
arguments and returns the truth value corresponding to their ``sum.''
If there is reason to believe that a sentence $p$ should be labelled
with truth value $v_1$, and independent reason to believe that it
should be labelled with truth value $v_2$, then the actual label is
given by the sum $v_1+v_2$.  This operation therefore corresponds to
combination of evidence.
 \item [\itype{dot}] A \lisp\ function that is the dual of the
\type{plus} operation.  It takes two truth values $v_1$ and $v_2$ as
arguments and returns the greatest truth value that could
subsequently be strengthened to a conclusion of either $v_1$ or
$v_2$.  This is a ``consensus'' operation.
 \item [\itype{kle}] A \lisp\ function that takes two truth values $x$
and $y$ and determines whether $x+y=y$.  Although this can also
be done using the bilattice $+$ and $=$ operations, this field can be
provided to give a more efficient implementation of this operation.
This field is optional.
 \item [\itype{tle}] A \lisp\ function that takes two truth values $x$
and $y$ and determines whether $x\vee y=y$.  This field is also
optional.
 \item [\itype{modal-ops}] A list of the modal operators associated
with the given bilattice.  See Section \ref{s.modal}.
 \end{description}

Three of the remaining formal slots are concerned with the incorporation
of binding information into truth values.  In some cases, when a
proposition is retrieved from the database, it is only an instance of
the database proposition that is used, and the truth value may have to
be modified to reflect that:
 \begin{description}
 \item [\itype{plug}]  This \lisp\ function accepts as arguments
a truth value and a binding list, and returns a potentially modified
truth value.  The default value is a function that leaves the truth
value unchanged.
 \end{description}

There is also a function that accepts a truth value and returns the
variables appearing within it:
 \begin{description}
 \item [\itype{vars}]  This \lisp\ function accepts a truth value
as its argument and returns a list of the variables appearing in it.
The default value is a function that returns \NIL\ independent of its
argument.
 \end{description}

The function \itype{absorb} is needed for subtle technical reasons
when working with bilattices with non-\NIL\ plugging functions;
further information on this can be found in the file \type{load.lisp}.
All three of these functions are \NIL\ if there is no interaction
between the bilattice truth values and the variables that they
contain.

The final formal field has to do with search control:
 \begin{description}
 \item [\itype{simplicity}] This \lisp\ function accepts a truth
value (which will frequently be the value obtained by conjoining the
premises in a particular proof attempt) and returns a figure
indicating the ``merit'' that should be assigned to the proof attempt
in question.  The value returned should be a nonempty list of
integers and is viewed lexicographically, with high values being
preferred to low ones.  The default value is a function that returns
the list \type{(0)} independent of the truth value supplied.
 \end{description}

The remaining ten fields in the \type{bilattice} structure are as
follows:
 \begin{description}
 \item [\itype{dws}] This \lisp\ function, ``dot with star,'' accepts
as arguments two truth values, $v_1$ and $v_2$.  The truth value returned
is $v_1$, modified so as to remove any contribution that could be
attributed to a sentence with truth value $v_2$.
 \item [\itype{stash-val}] A function that, given a sentence $p$,
returns the default truth value to be passed to \type{stash}.
 \item [\itype{unstash-val}] A function that, given a sentence $p$,
returns the default truth value to be passed to \type{unstash}.
 \item [\itype{cutoff-val}] The default cutoff point that must be
exceeded by any sentence if it is to be returned by \type{lookup}.
See Section \ref{s.cutoffs}.
 \item [\itype{print-fn}] This function is used to pretty-print
truth values in the associated bilattice.  Its three arguments
are a truth value, the bilattice in question, and a stream.
 \item [\itype{long-desc}] A long description of the bilattice, used
for input and output only.
 \item [\itype{short-desc}] A short description of the bilattice, also
used for input and output only.
 \item [\itype{char}] A one-character identifier for the bilattice.
This character can be used to change the bilattice using the function
\type{load-logic} described in Section \ref{s.change}.
 \item [\itype{type}] An overall description of the mathematical form
of the bilattice.  Perhaps it is a product of two other bilattices,
or formed from them in some other fashion.  This defaults to \type{basic}
if no other value is supplied; various methods for constructing new
bilattices from old ones are described in Section \ref{s.construct}. 
 \item [\itype{components}]  A list of the components used to construct
a bilattice whose \type{type} is not \type{basic}.
 \end{description}

\subsection{The active bilattice}

At any particular point in time, there will be a single bilattice
that has been selected by the user to be used in the application he
is developing.  This bilattice is stored as the value of
\type{*active-bilattice*}; in addition, many of the slots of the
active bilattice are directly available to the user.  Here they
are:

\begin{description}
 \item [\itype{*active-bilattice*}] The currently active bilattice.
 \item [\itype{true}] The truth value \type{true} in the active bilattice.
 \item [\itype{false}] The truth value \type{false} in the active bilattice.
 \item [\itype{unknown}] The truth value \type{unknown} in the active
bilattice.
 \item [\itype{bottom}] The truth value \type{bottom} in the active bilattice.
 \item [\itype{std-cutoffs}] Cutoff information constructed from the
value of \type{cutoff-val} in the active bilattice.  See Section
\ref{s.cutoffs}.
 \item [\itype{stash-value}] The function \type{stash-val} in the
active bilattice.
 \item [\itype{unstash-value}] The function \type{unstash-val} in the
active bilattice.
 \item [\itype{mvl-plus}] The function \type{plus} in the active bilattice.
 \item [\itype{mvl-dot}] The function \type{dot} in the active bilattice.
 \item [\itype{mvl-and}] The function \type{and} in the active bilattice.
 \item [\itype{mvl-or}] The function \type{or} in the active bilattice.
 \item [\itype{mvl-not}] The function \type{not} in the active bilattice.
 \item [\itype{mvl-eq}] The function \type{eq} in the active bilattice.
 \item [\itype{mvl-simplicity}] The function \type{simplicity} in the active
bilattice.
 \item [\itype{dot-with-star}] The function \type{dws} in the active 
bilattice.
 \item [\itype{mvl-vars}] The function \type{vars} in the active bilattice.
 \item [\itype{mvl-plug}] The function \type{plug} in the active bilattice.
 \item [\itype{mvl-print}] This function accepts a single argument, a truth
value, and modifies it in a way that will cause it to be pretty-printed.
The value returned by mvl-print should not be used for any purpose other
than printing.
 \end{description}

All of the functions above accept an optional argument that can be
used to evaluate their result using some bilattice other than the
active one.

\subsubsection{Functions derived from the active bilattice}

There are a variety of functions associated with the functions that
are defined when a new bilattice is made active.  In order to
simplify the description of these functions, we will denote
\type{mvl-and} by $\wedge$, \type{mvl-or} by $\vee$, \type{mvl-not}
by $\neg$, \type{mvl-plus} by $+$, and \type{mvl-dot} by $\cdot$.
The first auxiliary function defined is \type{dot-with-not}:
 \begin{description}
 \item [\itype{dot-with-not(x)}] Computes $x\cdot \neg x$.
 \end{description}

The other functions are constructed from the partial orders defined
using the bilattice operations.  We will write $x\leq_t y$ just in
case $x\wedge y = x$, for example, and $x\leq_k y$ just in case
$x\cdot y = x$.  (All of this notation follows that in the paper,
``Multivalued logics: A uniform approach to inference in artificial
intelligence''.)

Given these definitions, the following \lisp\ functions are defined:
 \begin{description}
 \item [\itype{k-le(x y)}] Returns \T\ if $x \leq_k y$, and \NIL\ otherwise.
 \item [\itype{k-ge(x y)}] Returns \T\ if $x \geq_k y$, and \NIL\ otherwise.
 \item [\itype{t-le(x y)}] Returns \T\ if $x \leq_t y$, and \NIL\ otherwise.
 \item [\itype{t-ge(x y)}] Returns \T\ if $x \geq_t y$, and \NIL\ otherwise.
 \item [\itype{k-lt(x y)}] Returns \T\ if $x <_k y$, and \NIL\ otherwise.
 \item [\itype{k-gt(x y)}] Returns \T\ if $x >_k y$, and \NIL\ otherwise.
 \item [\itype{t-lt(x y)}] Returns \T\ if $x <_t y$, and \NIL\ otherwise.
 \item [\itype{t-gt(x y)}] Returns \T\ if $x >_t y$, and \NIL\ otherwise.
 \item [\itype{k-not-le(x y)}] Returns \T\ if $x \not\leq_k y$, and
\NIL\ otherwise.
 \item [\itype{k-not-ge(x y)}] Returns \T\ if $x \not\geq_k y$, and
\NIL\ otherwise.
 \item [\itype{t-not-le(x y)}] Returns \T\ if $x \not\leq_t y$, and
\NIL\ otherwise.
 \item [\itype{t-not-ge(x y)}] Returns \T\ if $x \not\geq_t y$, and
\NIL\ otherwise.
 \item [\itype{k-not-lt(x y)}] Returns \T\ if $x \not<_k y$, and \NIL\
otherwise.
 \item [\itype{k-not-gt(x y)}] Returns \T\ if $x \not>_k y$, and \NIL\
otherwise.
 \item [\itype{t-not-lt(x y)}] Returns \T\ if $x \not<_t y$, and \NIL\
otherwise.
 \item [\itype{t-not-gt(x y)}] Returns \T\ if $x \not>_t y$, and \NIL\
otherwise.
 \end{description}

All of these functions accept an optional argument that is the
bilattice to be used in the comparison.  This argument defaults to
the currently active bilattice.

\subsubsection{Cutoff information}
\label{s.cutoffs}

We will refer to each of the 16 functions defined at the end of the
previous section (e.g., excluding \type{dot-with-not}) as a {\em tag}
\index{tag}.  Thus \type{k-le} is a tag, and so on.

We will now define a {\em simple test} to be a dotted pair
 \[\type{(v . tag)}\]
 where \type{v} is a truth value and \type{tag} is a tag.  We will
say that a truth value \type{x} {\em satisfies\/} the simple test
if and only if the function \type{tag}, when applied to the arguments
$x$ and $v$ (in that order), returns \T.

As an example, consider the simple test
 \[\type{(true . t-ge)}\]
 This test will be satisfied by any truth value that is $\geq_t
\type{true}$.  Since \type{true} itself is the only truth value
satisfying this test, the test will be satisfied for \type{true} and
no other truth values.  This is the test that is frequently used to
decide to terminate a proof effort.\footnote{Tests in the \mvl\ system
are in fact \lisp\ structures with \type{value} and \type{tag} fields;
we are using dotted lists in this manual only to make the presentation
clearer.}

Alternatively, the test used to determine whether or not to return
an answer after the theorem prover runs to completion is often given
by
 \[\type{(unknown . k-gt)}\]
 This test will succeed for any truth value about which more is known
that \type{unknown} -- in other words, for any truth value other than
\type{unknown} itself.

In general, the simple tests are combined into {\em cutoffs}.  A cutoff
is simply a list of lists of simple tests, and is interpreted as being
in conjunctive normal form.  Thus the cutoff corresponding to the simple
test $s$ is given simply by
 \[\type{((s))}\]

A cutoff corresponding to the conjunction of tests $s_1,\dots,s_n$ is given
by
 \[\type{((s-1) ... (s-n))}\]
 while that corresponding to their {\em disjunction\/} is given by
 \[\type{((s-1 ... s-n))}\]

In general, we can build up conjunctions of disjunctions of simple
tests in as complex a manner as we wish.  The cutoff \NIL\
corresponds to the empty conjunction and therefore always succeeds;
the cutoff \type{(NIL)} contains an empty disjunction and therefore
always fails.  For convenience, these two cutoffs are stored in the
global constants \itype{*always-succeeds*} and
\itype{*never-succeeds*} respectively.

Given a truth value $x$ and a cutoff $c$, the expression \type{(test x c)}
evaluates to \T\ if $x$ passes the test, and to \NIL\ if it fails it.
There is also a function that negates cutoffs:
 \begin{description}
 \item [\itype{test(value cutoffs \&optional bilattice)}] Returns \T\
if the truth value passes the test associated to \type{cutoffs}, and
\NIL\ otherwise.  The test is performed for the supplied bilattice,
or for the currently active one if the optional argument is omitted.
 \item [\itype{cutoffs-not(cutoffs \&optional bilattice)}] Returns a
new cutoffs $c$ such that
 \[\type{(test x c)} = \type{(test (mvl-not x) cutoffs)}\]
 for all truth values $x$.  In other words, $x$ passes $c$ if and
only if $\neg x$ passed the cutoffs supplied originally.
 \end{description}

\subsection{Selecting a new bilattice}
\label{s.change}

To change the bilattice being used, simply invoke the function
\type{load-logic}.  There are three ways to specify the new bilattice:
 \begin{enumerate}
 \item Do nothing.  You will then be asked to select the bilattice
from the list of bilattices known to the system.
 \item Supply the one-character identifier of the new bilattice as an
argument to \type{load-logic}.  The new bilattice must already be
known to the system.
 \item Supply the new bilattice itself as an argument to
\type{load-logic}.
 \end{enumerate}
 \type{Load-logic} empties the database before the bilattice is
changed.

The bilattices known to the system are kept in the global variable
\type{*bilattices*}.  To add a new bilattice to \type{*bilattices*},
invoke the function \type{bilattice} with the bilattice as an
argument.  This function will warn you if you are overwriting another
bilattice that has the same one-character designator.

Summarizing:
 \begin{description}
 \item [\itype{*bilattices*}] A list of the bilattices known to the system.
 \item [\itype{bilattice(b)}] Add the bilattice $b$ to \type{*bilattices*}.
 \item [\itype{load-logic}] Empty the database and make a new bilattice
active.  Accepts as an optional argument either a character designating
a bilattice, or the new bilattice itself.  If no argument is supplied,
the user is prompted for the new bilattice.
 \end{description}

\section{Bilattices available in the MVL system}

Here are the bilattices available in the \mvl\ system as distributed.
Each is discussed more fully in a subsequent section:
 \begin{description}
 \item [\itype{*first-order-bilattice*}] The bilattice corresponding to
first-order logic.
 \item [\itype{*atms-bilattice*}] The bilattice corresponding to an \atms,
where the justification information is maintained in disjunctive normal
form as suggested by de Kleer.
 \item [\itype{*cnf-atms-bilattice*}] An alternative version of the \atms\
bilattice where the justification information is maintained in conjunctive
normal form.
 \item [\itype{*default-bilattice*}] A bilattice that can be used for simple
default reasoning.
 \item [\itype{*hierarchy-bilattice*}] A bilattice similar to the
default bilattice but where the defaults are labelled with
integer-valued strengths.
 \item [\itype{*first-order-atms-bilattice*}] This bilattice corresponds
to a fully first-order \atms, as opposed to the more conventional one
dealing only with propositional logic.
 \item [\itype{*circumscription-bilattice*}] The bilattice used to compute
the results of the circumscription axiom.
 \item [\itype{*probability-bilattice*}] An extended version of the \atms\
bilattice that maintains probabilistic information as well.
 \item [\itype{*time-bilattice*}] A compound bilattice that can be used
for reasoning about events occurring against a temporal background.
 \end{description}

\subsection{First-order logic} \index{first-order logic}

The simplest bilattice is that corresponding to first-order logic, and
contains only the four truth values \type{true} (displayed as \type
T), \type{false} (displayed as \type F), \type{unknown} (\type U) and
\type{bottom} (\type{T/F}).  This bilattice is identified using the
character {\bf f}, and is active when \mvl\ is first loaded.

\subsection{Assumption-based truth maintenance}
\index{assumption-based truth maintenance}

Except for the \nm\ and probabilistic \atms\ bilattices, all of the
bilattices corresponding to \atms's have truth values that are dotted
pairs \type{(j . k)}, where $j$ is a justification for the sentence
being labelled, and $k$ is a justification for its negation.  The
bilattices differ in that the justifications are represented
differently.

\subsubsection{The ATMS bilattice}

For the simplest \atms\ bilattice, the justifications are simply
disjunctive normal form expressions.  Thus the justification
 \begin{equation}
\type{((P23 P37) (P1))}
\label{e.just}
\end{equation}
 is interpreted to mean that the sentence in question is a
consequence either of \type{P1}, or of \type{P23} and \type{P37} in
conjunction.

This bilattice is identified by the character {\bf a}.

\subsubsection{ATMS's in conjunctive normal form}

Alternatively, we can represent the justifications in conjunctive
normal form, as opposed to disjunctive normal form.  This representation
is useful, for example, in diagnosing multiple faults, since we want to
be able to explain an observed failure {\em either\/} in terms of the
collective failures of one set of components, or the collective failures
of another set, and so on.

This bilattice is identified by the character {\bf b}.

\subsubsection{First-order ATMS's}
\label{s.fo-atms}

It is also possible to extend these ideas slightly to obtain a fully
first-order \atms.  In order to do this, we need simply modify
justifications such as that appearing in (\ref{e.just}) so that the
primitive elements of the justification are not atoms referring to
propositions, but pairs
 \[\type{(atom . binding-list)}\]
 where the atom refers to a proposition and the binding list
indicates precisely what instantiation of that proposition is
actually needed to justify the sentence in question.

As an example, suppose that an assumption in our database is \type{P1},
given by
 \begin{equation}
\type{(FRIENDLY ?X)}
\label{e.assume}
\end{equation}
 and indicating that everything is friendly.  If we also have the rule
of inference
 \[\type{(<= (HAPPY ?X) (FRIENDLY ?X))}\]
 indicating that anything friendly is happy, and this rule is unconditionally
true (i.e., not an assumption), then the truth value assigned to the
query \type{(HAPPY HARRY)} will be
 \[\type{((((P1 . ((?X . HARRY))))))}\]
 indicating that Harry is happy by virtue of (\ref{e.assume}) applied
to $x=\mbox{Harry}$, and that there is no justification for the
conclusion that Harry is not happy.  (Recall that all of the \atms\
bilattices combine justification for and against each sentence into a
single truth value.)

This bilattice is identified by the character {\bf o}.

\subsubsection{Probabilistic ATMS's}
\label{s.prob}

Truth values in the probabilistic \atms\ bilattice consist of dotted
pairs, where the \type{car} of a truth value is simply an element of
the \atms\ bilattice and the \type{cdr} is the probability that the
given assumption is valid.

When entering an assumption, the user will be queried for the
probability in question; the probabilities of compound sentences are
then computed automatically by the system.  The algorithm used is
that described in Grnarov, Kleinrock and Gerla in the 1979
Proceedings of the Computer Networking Symposium.

The probabilistic \atms\ bilattice is identified by the character
{\bf p}.
  
\subsection{Default reasoning} \index{default reasoning}
\label{s.j}

The default bilattice consists of seven truth values.  In addition to
the usual elements \type{true}, \type{false}, \type{unknown} and
\type{bottom}, there are elements \type{dt}, \type{df} and \type{*}.

The truth value \type{dt} labels a conclusion that is true by default
(such as Tweety flying in the canonical example), while \type{df} labels
a conclusion that is {\em false\/} by default.  The final value,
\type{*}, labels conclusions that are both true {\em and\/} false
by default, as in the well-known Nixon-Republican-Quaker example.

This bilattice is identified by the character {\bf d}.

\subsubsection{Prioritized default reasoning}

Another version of the default bilattice allows for priorities among
the defaults.  In this bilattice, truth values are of the form
 \[\type{(n . x)}\]
 where $n$ is a nonnegative integer and $x$ is an element of the
first-order bilattice.  Truth values with $n=0$ are displayed as for
the first-order bilattice; with $n=1$ as for the default bilattice.
For $n>1$, we might see \type{D(3)T} or something like it.  This
particular truth value is the same as
 \[\type{(3 . true)}\]
 Truth values with small $n$ override those with larger $n$.
\type{Unknown} in this bilattice is given by the truth value \NIL.
The one-character identifier for the hierarchical default bilattice is
{\bf h}.

\subsection{Circumscription} \index{circumscription}
\label{s.circ}

The default bilattices are not able to deal with all of the
situations in which multiple extensions arise because single truth
values such as \type{*} label sentences that are true in some
extensions and false in others without providing any information
about precisely which extensions are which.  In order to handle this
difficulty, we need to work with a more complex bilattice.

One way to view the default bilattice is as a copy of the first-order
bilattice in which the point \type{unknown} is interpreted as, ``not
known with certainty,'' and is in fact replaced with another copy of
the first-order bilattice that is used to record uncertain
information (i.e., default truth, default falsity, or the presence of
a default contradiction).

The circumscriptive bilattice is formed similarly, except that
instead of simply recording default truth, falsity, and so on, we
record a {\em justification\/} for default conclusions.  This is done
by replacing the point \type{unknown} of the first-order bilattice
not with another copy of the first-order bilattice, but with a copy
of the first-order \atms\ bilattice.  (This idea is in fact a general
way to construct new bilattices from old ones, and is called a {\em
surgery}.  It is discussed in Section \ref{s.surgery}.)  The
designator for the circumscriptive bilattice is {\bf c}.
\type{Circum} and \type{circums} use this bilattice to accept a
sentence as an argument and attempt to prove it from the
circumscription axiom; they resemble \type{bc} and \type{bcs}.

The truth values in the circumscription bilattice corresponding to
certain information are of the form \type{(T . x)}, where $x$ is an
element of the first-order bilattice.  Those corresponding to uncertain
information are of the form \type{(NIL . y)}, where $y$ is an element
of the first-order \atms\ bilattice.

Finally, the circumscriptive theorem prover will output diagnostics
concerning its progress if the global variable \type{*circum-trace*}
has a non-\NIL\ value.

Here is a summary of the circumscriptive facilities:
 \begin{description} 
 \item [\itype{*circum-trace*}] If not \NIL, the circumscriptive proof
functions will output diagnostics describing their progress.
 \item [\itype{circum(p)}] Like \type{bc}, but looks for instances of
$p$ that follow from the circumscription axiom.  Allowed keywords:
\type{cutoffs} and \type{succeeds}.
 \item [\itype{circums(p)}] Like \type{bcs}, but looks for instances of
$p$ that follow from the circumscription axiom.  Allowed keywords:
\type{cutoffs} and \type{succeeds}.
 \end{description}

\subsection{Temporal reasoning}\index{temporal reasoning}

\type{*Time-bilattice*} is a bilattice of the truth values that might
be assigned to propositions that have a time-dependent nature, such as
whether or not a particular switch is on.  The truth values represent
the entire histories of the propositions, so that a single truth value
is a complete function from a set of time points (taken to be the
integers) into a set of ``base'' truth values (take to be elements of
the nonmonotonic {\sc atms} bilattice). This bilattice is discussed at
greater length in the paper, ``Computational considerations in
reasoning about action,'' and a sample use of this bilattice can be
found by examining the file \type{temporal.test} that is included with
\mvl.

\section{Facilities to create new bilattices}
\label{s.construct}

There are a variety of ways in which it is possible to construct
new bilattices from old ones, or from old ones in conjunction
with instances of some other structure.  In this section, we
describe the extent to which these facilities exist within the
\mvl\ system.

\subsection{Products}
\label{s.product}

Perhaps the simplest way in which a new bilattice can be constructed
is as the product of two existing bilattices.  If $A$ and $B$ are
bilattices, then a truth value in the product bilattice $A\times B$
will simply be a pair of truth values \type{(x . y)}, where $x$ is a
truth value in $A$ and $y$ is one in $B$.  All of the bilattice
operations on the product bilattice are computed pointwise using the
corresponding operations on the component bilattices $A$ and $B$.

The product construction is performed using the following function:
 \begin{description}
 \item [\itype{product-bilattice(b1 b2)}] Returns the bilattice
that is the product of the bilattices $b_1$ and $b_2$.
 \end{description}

\subsection{Functional attachments}

Another way in which to construct a bilattice is to start with a
bilattice $B$ and a function $f$ that maps $B$ into a fixed set $S$.
The elements of the new bilattice are simply pairs of the form
 \[\type{(x .~(f x))}\]
 where $x$ is an element of $B$.  The function $f$ here is simply
tagging along for the ride, and has no impact on the bilattice
operations.  It can be used to provide additional information about
the truth values, and is used to compute probabilistic information in
the bilattice {\bf p} described in Section \ref{s.prob}.

The relevant function is the following one:
 \begin{description}
 \item [\itype{append-bilattice(b f)}] Returns the bilattice obtained
by appending the result of applying the function $f$ to the points of
the original bilattice $b$.
 \end{description}

\subsection{Surgeries}
\label{s.surgery}

A somewhat more interesting construction is the one we described in
Section \ref{s.circ} as a {\em surgery}.  Here, we take one bilattice
$b_1$ and replace the point corresponding to \type{unknown} in this
bilattice with another bilattice $b_2$.  The truth values in the
resulting bilattice are either of the form \type{(T . x)} for a point
$x$ in the ``upper'' bilattice $b_1$, or \type{(NIL . y)} for a point
$y$ in the lower bilattice.  Truth values from $b_1$ take priority over
values from $b_2$.

To perform this construction generally, we have: 
 \begin{description}
 \item [\itype{splice-bilattice(upper lower)}] Returns the bilattice
obtained by replacing the point \type{unknown} in the bilattice
\type{upper} with a copy of the bilattice \type{lower}.
 \end{description}

As an example, the default bilattice {\bf d} of Section \ref{s.j} is
given by
 \[\type{(splice-bilattice *first-order-bilattice* *first-order-bilattice*)}\]
 The circumscriptive bilattice of Section \ref{s.circ} is given by
 \[\type{(splice-bilattice *first-order-bilattice*
*first-order-atms-bilattice*)}\]

\subsection{Hierarchies}

It is possible to iterate the surgery construction of the previous
section, producing a bilattice that is an infinite descending
``chain'' of copies of a basic bilattice $b$.  This is done by
the function \type{hierarchy-bilattice}:
 \begin{description}
 \item [\itype{hierarchy-bilattice(b)}] Returns the bilattice constructed
as an infinite descending chain of copies of $b$.
 \end{description}

Truth values in this bilattice have the form \type{(n . x)}, where
$x$ is an element of $b$, and $n$ is an integer giving the hierarchy
level.  The top hierarchy level of the bilattice has $n=0$; the
second level is $n=1$, and so on.  The point \type{unknown} is given
by \NIL.

This function is used to construct the hierarchical version of the
default bilattice.

\subsection{Lattices}

Bilattices can also be constructed from lattices.  A lattice is a set
equipped with a {\em single\/} partial order and greatest lower bound
and least upper bound operations, which we denote by $+$ and $\cdot$
respectively.

An element of the associated bilattice is a {\em pair\/} of elements
of the lattice, and the bilattice operations are given by:
 \begin{eqnarray*}
\neg(a,b) & = & (b,a) \\
(a,b) + (c,d) & = & (a+c,b+d) \\
(a,b) \cdot (c,d) & = & (a\cdot c,b\cdot d) \\
(a,b) \vee (c,d) & = & (a+c, b\cdot d) \\
(a,b) \wedge (c,d) & = & (a\cdot c,b+d)
 \end{eqnarray*}

Roughly speaking, the first component of the bilattice element labels
the ``truth'' of the sentence, and the second labels its falsity.  This
is how the \atms\ bilattices are constructed.

The function used is \type{lattice-to-bilattice}:

 \begin{description}
 \item [\itype{lattice-to-bilattice(l)}] Returns the bilattice constructed
from the lattice $l$.
 \end{description}

A lattice is a structure rather like a bilattice, but with fewer slots
(because it only has a single partial order).  The fields in an instance
of the \itype{lattice} structure are:

\begin{description}
 \item [\type{max}] A \lisp\ object that is the maximal element of
the lattice.
 \item [\type{min}] A \lisp\ object that is the minimal element of
the lattice.
 \item [\type{eq}] A \lisp\ function that takes two lattice elements
as arguments and returns \T\ if they are equal and \NIL\ if they are
not.
 \item [\type{and}] A \lisp\ function that takes two lattice elements
as arguments and returns their greatest lower bound.
 \item [\type{or}] A \lisp\ function that takes two lattice elements
as arguments and returns their least upper bound.
 \item [\type{le}] This optional field is a \lisp\ function that
accepts two arguments $x$ and $y$ and determines whether or not
$x\leq y$ under the partial order corresponding to the lattice.
 \item [\type{vars}] This \lisp\ function accepts a lattice element
as an argument and returns a list of the variables that it contains.
The default value returns \NIL\ independent of its argument.
 \item [\type{plug}] This \lisp\ function accepts as arguments a
lattice value and a binding list, and returns a potentially modified
lattice value.  The default value is a function that leaves the lattice
value unchanged.
 \item [\type{simplicity}] The lattice analog of the bilattice
\type{simplicity} slot.  When constructing a bilattice from a
lattice, the \type{simplicity} of a truth value $(a,b)$ is the result
of applying the simplicity function from the lattice to the lattice
element $a$.  The value $b$ is ignored in the simplicity computation.
 \item [\type{dws}] This \lisp\ function, ``dot with star,'' accepts
as arguments two lattice values, $v_1$ and $v_2$.  The lattice
element returned is $v_1$, modified so as to remove any contribution
that could be attributed to the lattice element $v_2$.
 \item [\type{stash-val}] The lattice element entering into
the bilattice function \type{stash-val}.
 \item [\type{long-desc}] A long description of the lattice, used
for input and output only.
 \item [\type{short-desc}] A short description of the lattice, also
used for input and output only.
 \item [\type{char}] A one-character identifier for the lattice.
 \end{description}

\subsection{Directed acyclic graphs}
\label{s.dag}

The final fashion in which \mvl\ supports the construction of new
bilattices is an extension of the product construction of Section
\ref{s.product}.  There, we formed the product of two bilattices
$A\times B$.  If the bilattices were identical, we would have formed
the product bilattice $B\times B$, or $B^2$.  Another way to view
this bilattice is as the collection of functions from the two point
set $\set{1,2}$ into the original bilattice $B$.  Since such a
function $f$ can be described as the pair $(f(1),f(2))$, there is no
difference between this description and the previous one.

Given this analogy, it is not hard to see that we can extend the
construction to the bilattice $B^S$ for an arbitrary set $S$, where
elements of this bilattice are functions from the set $S$ into $B$,
and all of the bilattice functions ($\wedge$, $\cdot$ and so on) are
once again computed pointwise.

The difficulty with this idea is that for large sets $S$, it is not
clear how we can represent a function from $S$ into $B$ in a
convenient way.  It turns out, however, that if $S$ has the structure
of a directed acyclic graph, or \dg, this problem can be overcome.

\subsubsection{Defining the DAG}

As a preliminary, we need some way to define the \dg\ itself.  In
\mvl, a \dg\ \index{dag} is an instance of the \type{dag} structure;
this structure has the following fields:
 \begin{description}
 \item [\itype{root}] The root of the \dg.
 \item [\itype{eq}] A function that accepts two elements of the \dg\
and returns \T\ if they are the same element and \NIL\ otherwise.
 \item [\itype{leq}] A function that accepts two elements of the \dg\
and returns \T\ if the first is farther from the root of the \dg\
than the second.  If not supplied by the user, this function can by
computed from the \type{dot} operation on the \dg.
 \item [\itype{dot}] A function that accepts two elements of the \dg\
and returns a list of their first common descendants, or \NIL\ if no
such descendant exists.
 \item [\itype{long}] A long description of the \dg.
 \item [\itype{short}] A short description of the \dg.
 \item [\itype{inherit}] A function describing the fashion in which
truth values are inherited within the \dg.  This function $f$ should take
four arguments: a bilattice $b$, two elements of the \dg\ $x$ and $y$,
and an element $z$ of the bilattice.  If $z$ is the element of the
bilattice associated to the \dg\ element $x$, then $f(b,x,y,z)$ will
be the value associated to the \dg\ element $y$ in the absence of
information to the contrary.  By default, this function ignores
its first three arguments and assigns $y$ the same value as $x$ has.
 \item [\type{vars}] This \lisp\ function accepts a \dg\ element
as an argument and returns a list of the variables that it contains.
The default value returns \NIL\ independent of its argument.
 \item [\type{plug}] This \lisp\ function accepts as arguments a
\dg\ value and a binding list, and returns a potentially modified
\dg\ value.  The default value is a function that leaves the \dg\
value unchanged.
 \item [\type{select}] In some applications (none of which is distributed
with \mvl), it appears that truth values should be inherited from only
{\em some\/} of the elements of the \dg\ which might potentially
impact the value of the function at some specific point.  The
selection function is given a point where the truth value is being
computed and a list of points that apparently contribute; it returns a
list of points that {\em actually\/} contribute.  By default, all
points contribute; most users of \mvl\ are very unlikely to need this
facility.
 \end{description}

\subsubsection{Constructing the bilattice}

Having selected a \dg\ $D$ and a bilattice $B$, we can now construct
the bilattice $B^D$.  Specific functions from $D$ to $B$ are described
as structures with the following slots:
 \begin{description}
 \item [\type{dag}] The domain of the function, $D$.
 \item [\type{bilattice}] The range of the function, $B$.
 \item [\type{list}] A list of elements of the form \type{(p . v)},
where \type p is a point where the value changes in a fashion {\em
other\/} than the one predicted by the inheritance function associated
with $D$ and \type v is the value taken there.
 \end{description}
 We will refer to these functions as {\em \dg-functions}.

There are a variety of functions available to manipulate \dg-functions,
and these are used to construct the bilattice $B^D$.  In some cases,
however, the user may want to access these functions directly.  Here
are some of them:
 \begin{description}
 \item [\itype{get-val(d f)}] Returns the value of the
function $f$ at the point $d$ of the \dg.
 \item [\itype{make-tv(dag bilattice d b)}] Returns a \dg-function that
takes the value $b$ at the point $d$.  If $d$ is not the root of the
\dg, the value at the root of the \dg\ is generally the point
\type{unknown} in the bilattice, but the user may supply the keyword
\type{root-fn} if he wishes, in which case the value at the root of
the \dg\ is the result of applying \type{root-fn} to the bilattice.
 \item [\itype{make-root(dag bilattice b)}] Returns a \dg-function that
takes the value $b$ at the root of the \dg.
 \item [\itype{dag-change-tv(f b-f)}] The argument
\type{b-f} should be a function from the bilattice to itself; a
modified version of the original \dg-function $f$ is returned in which
every point in the \dg-function has had the function \type{b-f}
applied to it.  There is an optional keyword \type{simplify}; if
supplied, the computed function is simplified before being returned.
(See \type{simplify}.)
 \item [\itype{dag-change-dag(f d-f)}] The argument
\type{d-f} should be a function from the \dg\ to itself; a modified
version of the original \dg-function $f$ is returned in which every
point at which the \dg-function takes a new value has had the function
\type{d-f} applied to it.  The bilattice argument is expected to be
the bilattice of functional truth values as opposed to the bilattice
that is the target of the functions.  There is an optional keyword
\type{simplify}; if supplied, the computed function is simplified
before being returned.  (See \type{simplify}.)
 \item [\itype{dag-list-of-answers(f)}] Given a \dg-function $f$, assumed
to be from the \dg\ of binding lists into the currently active bilattice, this
function constructs a list of answers corresponding to the points where
the \dg-function takes values.
 \item [\itype{dag-accumulate-fn(dag bilattice list)}] Given a \dg, a
bilattice, and a list of elements of the form
 \[\type{(dag-elt . bilattice-elt)}\]
a \dg-function is constructed that takes the given values at the
given elements of the \dg.

A keyword to \type{dag-accumulate-fn} is \type{modifier}.  This
function is used if the \type{d-list} contains the same \dg\ element
more than once, in which case the modifier function is called with
the new and old values (in that order) and the result is used as the
value of the \dg-function at that particular point.  The default
action is to simply overwrite the old value with the new one.  (In
other words, the default modifier function simply returns its first
argument while ignoring the second.)
 \item [\itype{dag-prune(f pred)}] Given a \dg-function
$f$ and a predicate on the bilattice, this function prunes from $f$
all points $d$ where $f(d)$ fails to satisfy the predicate.  There is
an optional keyword \type{modify-fn} that, if supplied, should be a
function $m$ from the bilattice to itself.  In this case, the points
where $m(f(d))$ fails to satisfy the predicate are pruned.
 \item [\itype{combine-fns(f1 f2 fn comp)}]
This function takes the two \dg-functions $f_1$ and $f_2$ and
combines them using the function \type{fn}, which should be a
function that takes two elements of the bilattice and returns a
single one.  The comparison function \type{comp} is used to determine
if one of the two functions can be returned directly; if, for example
\type{(comp f1 f2)} holds, than  we simply return $f_1$, and similarly
for $f_2$.  As an example, if the combining function is the bilattice
operation $\cdot$, the comparison function is \type{k-le}.
 \item [\itype{combine-1(f params f1 ... fn)}] This is an
$n$-ary version of \type{combine-fns}, where $f$ is an $n$-ary
bilattice function and there are $n$ \dg-functions supplied as well.
The additional argument \type{params} indicates which of the $n$ arguments
to this function are not necessarily \dg-functions, but are instead parameters
expected by the overall function $f$.  This additional argument is a list
of \T\ (if the $n$th argument is a parameter) or \NIL\ (if it isn't).
If the parameter list ends early, all of the remaining entries are
assumed to be \NIL.
 \item [\itype{add-dag(d b f)}] Given the \dg-function
$f$, add the value $b$ at the point $d$.  This function works by side
effect only.  There is a keyword \type{modifier} that, if supplied,
should be a function that takes two elements of the bilattice and
returns a single one.  Given such a function $m$, the \dg-function is
changed so that the new value at $d$ is $m(b,f(d))$, where $f(d)$ is
the old value.  (This is similar to the same keyword in
\type{dag-accumulate-fn}.)  The keyword \type{simplify} is also available.
 \item [\itype{simplify(f)}] In some cases,
manipulating the \dg-function $f$ may cause it to contain
redundancies, where the value at some point is presented explicitly
even though this value is in fact the one that would normally be
inherited from its parents.  This function returns a modification of
$f$ where these irrelevant entries have been removed.
 \item [\itype{fn-and(f1 f2)}] This function conjoins
the values of $f_1$ and $f_2$ using the bilattice $\wedge$ operation.
The new \dg-function is returned.
 \item [\itype{fn-or(f1 f2)}] This function disjoins
the values of $f_1$ and $f_2$ using the bilattice $\vee$ operation.
The new \dg-function is returned.
 \item [\itype{fn-plus(f1 f2)}] This function adds
the values of $f_1$ and $f_2$ using the bilattice $+$ operation.
The new \dg-function is returned.
 \item [\itype{fn-dot(f1 f2)}] This function combines
the values of $f_1$ and $f_2$ using the bilattice $\cdot$ operation.
The new \dg-function is returned.
 \item [\itype{fn-dws(f1 f2)}] This function combines
the values of $f_1$ and $f_2$ using the bilattice operation
\type{dot-with-star}.  The new \dg-function is returned.
 \item [\itype{fn-not(f1)}] This function negates the
values of $f_1$ using the bilattice $\neg$ operation.  The new
\dg-function is returned.
 \item [\itype{fn-eq(f1 f2)}] This function compares
the values of $f_1$ and $f_2$ using the bilattice comparison operation.
Either \T\ or \NIL\ is returned.
 \item [\itype{fn-comp(f1 f2 fn)}] This functions compares
the values of $f_1$ and $f_2$ using the bilattice comparison function
\type{fn}.  Either \T\ or \NIL\ is returned.
 \end{description}

Of course, the point of all of this is to enable us to define the
following function:
 \begin{description}
 \item [\itype{dag-bilattice(dag bilattice)}] Given a bilattice $B$
and a dag $D$, returns the bilattice $B^D$.
 \end{description}

The following functions are also available:
 \begin{description}
 \item [\itype{lattice-to-dag(lattice)}] Construct a \dg\ from the
given lattice.
 \item [\itype{lattice-to-dag-to-bilattice(lattice bilattice)}]  First
construct a \dg\ $D$ from the given lattice.  Then use the bilattice
$B$ to construct the bilattice $B^D$.
 \end{description}

\subsubsection{The DAG of theories}
\label{s.theory}

Two specific \dg s are provided for use in the \mvl\ system.  The first
uses a \dg\ corresponding to the existence of multiple theories within
a single database; the second uses a \dg\ to handle binding information.

The theory \dg\ consists of a hierarchy of theories.  In a single
database, for example, we might want to have rules that applied to
celestial bodies, rules applying to stars, rules applying only to red
giants or only to our sun, and so on.  The truth value of a rule
applicable to all celestial bodies should be inherited if we are
considering our sun (but it should be possible to overwrite it with
a different truth value if we wish), but a rule applicable to red
giants should not be inherited in this fashion.  So we see that our
\dg\ construction is well-suited for this sort of an application.

To activate the theory mechanism, simply invoke the \lisp\ function
\type{activate-theory-mechanism}.  This function replaces the active
bilattice, say $B$, with the bilattice constructed of functions from
the theory \dg\ into $B$.  The original bilattice can still be
accessed, since it is the second \type{component} of the newly formed
bilattice.  The function \type{theory-active-p} can be used to
determine whether or not the theory mechanism has been activated.

The \dg\ of theories is the value of the global variable
\type{theory-dag}, and has as its root the theory whose name is the
value of the constant \type{*global*}.  To indicate that a theory
\type{big} includes another theory \type{little} as a subtheory,
invoke the function 
 \[\type{(includes big little)}\]
 Thus, for example, one would introduce a new theory called
\type{theory} by writing
 \[\type{(includes *global* 'theory)}\]
 The effects of \type{includes} can be reversed by invoking the
function \type{unincludes}.

Note that since the theory \dg\ is used by the functions that
manipulate the bilattice $B^D$, it is important that this \dg\ be
maintained as new theories are introduced.  If \type{includes} is not
used to ``locate'' these new theories in the overall theory \dg\
itself, the theories will be inaccessible to many of the \mvl\
functions.

At any point in time, any new fact will be inserted into the database
as true in a particular theory; the theory into which it is inserted
is the value of \type{*active-theory*}.  The value of this variable
should not be changed directly, but should instead be modified by
invoking the function \type{activate} on the theory into which
subsequent facts should be written.  The variable
\type{true-in-theory} contains the truth value of a sentence that is
true in the currently active writing theory.

The theory \type{*active-theory*} is also the theory that is active
for reading at any point in time.  Of course, any parents of the
active theory are also active by virtue of the structure of the theory
\dg\ itself; if it is desired to read from multiple theories
simultaneously, a new theory should be created that is a child of all
of them.

Here is a summary of the functions described so far:
 \begin{description}
 \item [\itype{activate-theory-mechanism}] Activates the theory
mechanism.  Accepts as an optional argument either a one-character
bilattice descriptor, or a bilattice.
 \item [\itype{theory-active-p}] Returns \T\ or \NIL\ depending on
whether or not the theory mechanism has been activated.
 \item [\itype{theory-dag}] The \dg\ of theories.
 \item [\itype{*global*}] The root of the \dg\ of theories.
 \item [\itype{includes(big little)}] This function is used to state that
the theory \type{little} is a subtheory of the theory \type{big}, and
should inherit from it.
 \item [\itype{unincludes(big little)}] This function is used to state
that the theory \type{little} is no longer a subtheory of the theory
\type{big}, and should not inherit from it.
 \item [\itype{*active-theory*}] The theory that is currently active
for reading and writing.
 \item [\itype{activate(th)}] Activate the theory \type{th} for reading
and writing.
 \item [\itype{true-in-theory}]  The truth value of a sentence that is
true in the currently-active writing theory.
 \end{description}

Finally, five functions have been added to simplify the user
interface somewhat.  The function {\tt theory-\-contents} displays
the contents of a given theory, and \type{empty-theory} removes all
of the propositions in this theory.  To modify the contents of a
specific theory, the following functions can be used:
 \begin{description}
 \item [\itype{state-in-theory(p)}] Adds the proposition $p$ to the
currently active theory.  Allowed keywords: \type{theory} can
be used to specify a different theory to which $p$ should be added;
\type{value} and \type{increment-p} have the same meaning as for
\type{state}.  (See Section \ref{s.database}; if these keywords are
not used, the function \type{state} can be used equally well, since
the default stash value with the theory mechanism activated also
inserts the proposition into the currently active theory.)
 \item [\itype{unstate-from-theory(p)}] Removes the proposition $p$
from the currently active theory.  Allowed keywords:
\type{theory} can be used to specify a different theory from which
$p$ should be removed; \type{value} and \type{decrement-p} have the
same meaning as for \type{unstate}.  (See Section \ref{s.database}.)
 \item [\itype{erase-from-theory(p)}] Removes the proposition $p$
from the currently active theory, and forward chains on the result.
Allowed keywords: \type{theory} can be used to specify the theory from
which $p$ should be removed; \type{value}, \type{decrement-p} and
\type{cutoffs} have the same meaning as for \type{erase}.  (See
Section \ref{s.fc}.)
 \item [\itype{theory-contents()}] Displays the contents of the
currently active theory, which can be changed using the
keyword \type{theory}.  The keywords \type{cutoffs} and \type{props}
are allowed and are treated as for \type{contents}.  (See Section
\ref{s.contents}.)  There is an additional keyword,
\itype{this-theory-only}, which defaults to \T\ and, if supplied,
indicates that only sentences whose truth values {\em change\/} in the
given theory are to be displayed.  (This avoids displaying all
sentences in the global theory, for example.)
 \item [\itype{empty-theory()}] Empties the currently active
theory, although this theory may be changed using the keyword
\type{theory}.  The keyword \type{value} has a meaning similar to its
meaning as a keyword to \type{erase}.
 \end{description}

\subsubsection{The DAG of binding lists}
\label{s.binding-dag}

The other \dg\ supplied with the \mvl\ system is used for binding
lists since these, too, admit a natural inheritance structure.  The
root of the \dg\ is the empty binding list \NIL, and inheritance in
the \dg\ is determined by subsumption.  The bilattice constructed as
functions from this \dg\ into \type{*active-bilattice*} is used
extensively by the prover, and is available to the user as
\type{bdg-to-truth-val}.

The binding \dg\ itself is known as \type{binding-dag}, and is constructed
using the following functions:
 \begin{description}
 \item [\itype{equal-binding(b1 b2)}] Checks to see if two binding
lists are effectively the same.
 \item [\itype{append-binding-lists(b1 b2)}] Computes the least
specific binding list that is more specific than both $b_1$ and $b_2$
or returns \NIL\ if no such binding list exists ($b_1$ and $b_2$ may
bind the same variable to different constants, for example).
 \item [\itype{binding-le(b1 b2)}] Returns \NIL\ unless $b_1$ is a more
specific binding that $b_2$, in which case a list of relative binding
lists is returned.
 \end{description}

The following global variables are also used:
 \begin{description}
 \item [\itype{binding-dag}] The \dg\ of binding lists.
 \item [\itype{bdg-to-truth-val}] The bilattice of functions from
the \dg\ of binding lists into \type{*active-bilattice*}.
 \end{description}

\section{Modal operators}
\label{s.modal}

An additional feature that makes the \mvl\ system unique is its
ability to handle {\em modal operators}.  The approach taken by \mvl\
to modal operators is that they are functions from the bilattice of
truth values to itself.  Thus, for example, the modal operator that
takes a sentence $p$ into the sentence, ``I know that $p$,''
corresponds to the bilattice function (on the first-order bilattice)
that takes $t$ and $\bot$ into $t$ (since we know $p$ if we know it
to be true or know it to be true and false both), and takes $f$ and
$u$ into $f$ (since we do {\em not\/} know $p$ if we either know it
to be false or do not know if it is true or false).  Further details
can be found in the paper, ``Bilattices and modal operators.''

A modal expression in \mvl\ is any sentence of the form 
 \[\type{(op p-1 \dots\ p-n)}\]
 where \type{op} is recognizable as a modal operator because it has a
non-\NIL\ value for its \type{modal-op} property in the currently
active invocation of \mvl.  In this case, the value of the
\type{modal-op} property should be an instance of the structure
\itype{modal-operator}.  This structure has the following fields:
 \begin{description}
 \item [\itype{name}] The name of the modal operator.  This is used
to identify modal operators in different bilattices that have the
same commonsense meaning (and can therefore be usefully combined when
the bilattices are combined as in the previous section).
 \item [\itype{function}] A function that accepts as arguments the
truth values of the sentences under the scope of the modal operator
and returns the truth value of the modal expression itself.
 \item [\itype{parameters}] If the modal operator accepts $n$
arguments, this should be a list of $n$ entries indicating whether
each argument is a truth value or a parameter that should be passed
directly to the modal operator function.  (An entry of \T\ indicates
that the associated argument is parametric; an entry of \NIL\
indicates that it isn't.)  Suppose, for example, that we are working
with a bilattice where the truth values are functions describing the
behaviors of various fluents over the course of time, and have a modal
operator \type{delay}, where the truth value assigned to
 \[\type{(delay t p)}\]
 is expected to be the same as the truth value assigned to $p$ but
delayed by an amount of time $t$.  Here, the first argument ($t$) is
a parameter and the second argument ($p$) is not.

If there are fewer than $n$ entries on the parameter list, then all of
the missing entries are assumed to be \NIL\ (so that the associated
arguments to the modal function are not parameters).  Note: All
parametric arguments should be bound when attempting to derive a modal
expression.  This is similar to the requirement that arguments to a
procedurally-attached predicate be bound when the procedure is called.
 \end{description}

\subsection{Modal operators supplied with the MVL bilattices}

All bilattices in \mvl\ are equipped with a modal operator \itype{idn}
that does not change the truth value of its argument at all; this
modal operator can be used to semantically mark points where it may be
desirable to interrupt the inference process.  In addition, if the
prover is invoked on \type{(idn p)} instead of on \type p, attempts
will be made to prove both $p$ and its negation $\neg p$.

The bilattices supplied as part of the \mvl\ system are equipped with
the following modal operators:

 \begin{description}
 \item [\type{*first-order-bilattice*}] This bilattice has two
modal operators: $L$, which maps $p$ into, ``necessarily $p$,'' or,
``I know that $p$,'' and $M$, which maps $p$ into ``possibly $p$.''
 \item [\type{*atms-bilattice*}]  No modal operators are supplied with
this or any of the other \atms-type bilattices.
 \item [\type{*default-bilattice*}] This bilattice and the hierarchical
default bilattice both inherit the modal operators $L$ and $M$ from the
first-order bilattice.
 \item [\type{*circumscription-bilattice*}] This bilattice does not
come equipped with any modal operators.
 \item [\type{*time-bilattice*}] This bilattice inherits the modal
operators $L$ and $M$ from the bilattice of basic values.  The
following three additional modal operators are also defined:
 \begin{enumerate}
 \item \itype{delay(t p)} This is a modal operator that pushes the
truth value of $p$ into the future by an amount of time $t$.  Thus if
$p$ is true at time 2, $\type{delay}(p,2)$ will be true at time 4.
 \item \itype{propagate} is a modal operator that corresponds to the frame
axiom.  If $p$ is true at time 2, unknown between times 3 and 20, and false
at time 21, then $\type{propagate}(p)$ will be true from times 2 until 20
and false subsequently.
 \item \itype{at(t p)} is a modal operator that returns the value taken
by the truth value $p$ at the time $t$.  (In actuality, a function has
to be returned; the function is the one that takes value $p(t)$ at all
times instead of just at $t$.)
 \end{enumerate}
 \end{description}

\subsection{Quantification and modal operators}

In addition to \type{idn}, there are two other modal operators that
are available for use with any bilattice, \itype{exists} and
\itype{forall}.  Although not really modal operators on the bilattice
being used, these {\em are\/} legitimate modal operators on the bilattice
\type{bdg-to-truth-val} that was described in Section \ref{s.binding-dag}.

The functions associated to these modal operators mimic the usual behaviors
given to the quantifiers.  Thus an expression such as
 \begin{equation}
\type{(exists ?x (foo ?x))}
\label{e.exists}
\end{equation}
 is evaluated by first computing the truth value that could be assigned
to \type{(foo ?x)} as an element of \type{bdg-to-truth-val}, so that
a truth value is obtained for every binding of \type{?x}.  This truth
value is then manipulated so that the values for any particular binding
of \type{?x} are disjoined (the usual interpretation of the existential
quantifier), and the result is returned as the truth value assigned to
\ref{e.exists}.  The modal operator \type{forall} is handled similarly
and corresponds to universal quantification.  The first (parametric)
argument to these operators can also be a list of variables, in which
case all of the variables are quantified over.

\subsection{Modal operators created automatically}

Finally, many of the functions that construct new bilattices are able
to construct modal operators on these bilattice in terms of the modal
operators on the bilattices they are using in the construction.  More
specifically:

\begin{enumerate}
 \item A product bilattice $A \times B$ will inherit a modal operator
from $A$ and $B$ if these component bilattices have modal operators
with the same name.
 \item A bilattice constructed using \type{append-bilattice} will have
modal operators corresponding to whatever modal operators were available
on the bilattice used in the construction.
 \item A bilattice constructed by splicing two bilattices $A$ and $B$
together will inherit the modal operators from the ``upper'' bilattice.
 \item A bilattice constructed using \type{hierarchy-bilattice} will
have modal operators corresponding to whatever modal operators were
available on the bilattice used in the construction.
 \item A bilattice constructed from a lattice will not normally have
any modal operators associated to it.
 \item The bilattice $B^D$ will inherit whatever modal operators the
bilattice $B$ has.
 \end{enumerate}

\section*{Acknowledgement}

Support for the development of \mvl\ has been provided by the Air
Force Office of Scientific Research, the Defense Advance Research
Projects Agency, General Dynamics, the National Science Foundation,
Rockwell International, and Rome Labs.

\clearpage

\appendix

\section{Summary of MVL functions and constants}

\begin{list}{}{\setlength{\labelwidth}{1.5in}\setlength{\labelsep}{.25in}
\setlength{\leftmargin}{1.75in}}

% *A*

\item [\itype{activate(th)}] Activate the theory \type{th} for reading
and writing.

\item [\itype{activate-best-first}] Activate best-first search by
setting \type{*node-evaluation-fn*} to \type{best-first-fn}.  Allowed
keywords: \type{initial-limit}, \type{limit-delta}, and
\type{final-limit}.

\item [\itype{activate-breadth-first}] Activate breadth-first search
by setting \type{*node-evaluation-fn*} to \type{depth-fn}.  Allowed
keywords: \type{initial-limit}, \type{limit-delta}, and
\type{final-limit}.

\item [\itype{activate-depth-first}] Activate depth-first search by
setting \type{*initial-limit*} to \NIL.  Allowed keyword:
\type{final-limit}.

\item [\itype{activate-theory-mechanism}] Activates the theory
mechanism.  Accepts as an optional argument either a one-character
bilattice descriptor, or a bilattice.

\item [\itype{*active-bilattice*}] The currently active bilattice.

\item [\itype{*active-theory*}] The theory that is currently active
for reading and writing.

\item [\itype{add-dag(d b f)}] Given the \dg-function $f$, add the
value $b$ at the point $d$.  This function works by side effect only.
There is a keyword \type{modifier} that, if supplied, should be a
function that takes two elements of the bilattice and returns a single
one.  Given such a function $m$, the \dg-function is changed so that
the new value at $d$ is $m(b,f(d))$, where $f(d)$ is the old value.
(This is similar to the same keyword in \type{dag-accumulate-fn}.)
The keyword \type{simplify} is also available.

\item[\itype{answer-binding(a)}] Returns the binding field of the
answer \type a.

\item[\itype{answer-value(a)}] Returns the value field of the answer
\type a.

\item [\itype{anytime-trace(pat)}] Include \type{pat} on \at, so that
any instance of \type{pat} will be traced.  The argument is optional;
if omitted, \type{(?*)} (of which everything is an instance) will be
used.

\item [\itype{anytime-untrace(pat)}] Remove any instance of
\type{pat} from \at.  The argument is optional; if omitted,
\type{(?*)} (of which everything is an instance) will be used.

\item [\itype{append-bilattice(b f)}] Returns the bilattice obtained
by appending the result of applying the function $f$ to the points of
the original bilattice $b$.

\item [\itype{append-binding-lists(b1 b2)}] Computes the least
specific binding list that is more specific than both $b_1$ and $b_2$
or returns \NIL\ if no such binding list exists ($b_1$ and $b_2$ may
bind the same variable to different constants, for example).

\item [\itype{*atms-bilattice*}] The bilattice corresponding to an \atms,
where the justification information is maintained in disjunctive normal
form as suggested by de Kleer.

% *B*

\item [\itype{(bagof}\type{ ?x ?p ?b)}] This relation holds for
$x$, $p$ and $b$ if and only if $b$ is a list of all $x$ satisfying
$p$.

\item [\itype{bc(q)}] Attempts to prove the query $q$ from information
in the database, looking only for a single answer.

\item [\itype{bc-trace(pat)}] Include \type{pat} on
\type{*bc-trace*}, so that any instance of \type{pat} will be
traced.  The argument is optional; if omitted, \type{(?*)} (of which
everything is an instance) will be used.

\item [\itype{bc-untrace(pat)}] Remove any instance of \type{pat}
from \type{*bc-trace*}.  The argument is optional; if omitted,
\type{(?*)} (of which everything is an instance) will be used.

\item [\itype{bc-unwatch(preds)}] Removes the given predicates from
the list of {\sc WATCH}ed predicates.  If no arguments are supplied,
all predicates are un{\sc watch}ed.

\item [\itype{bc-watch(preds)}] Accepts as arguments any number of
predicates that should be watched as the prover proceeds.  Each
predicate can also be a list of the form \type{(predicate depth)}
where \type{depth} is the number of successive ancestors to be
displayed when the prover encounters an instance of the predicate.
Returns a list of the predicates currently being {\sc WATCH}ed.

\item [\itype{bcs(q)}] Attempts to prove the query $q$ from information
in the database, looking for all possible answers.

\item [\itype{bdg-to-truth-val}] The bilattice of functions from
the \dg\ of binding lists into \type{*active-bilattice*}.

\item [\itype{bilattice(b)}] Add the bilattice $b$ to \type{*bilattices*}.

\item [\itype{*bilattices*}] A list of the bilattices known to the system.

\item [\itype{binding-dag}] The \dg\ of binding lists.

\item [\itype{binding-le(b1 b2)}] Returns \NIL\ unless $b_1$ is a more
specific binding that $b_2$, in which case a list of relative binding
lists is returned.

\item [\itype{bottom}] The truth value \type{bottom} in the active bilattice.

\item [\itype{*break-node*}] The number of a node at which the
proof process should be interrupted.

% *C*

\item [\itype{choose-task}]  Selects a task to work on next.  The 
current implementation of \mvl\ simply selects any task that might
have some impact on the final answer being computed.

\item [\itype{circum(p)}] Like \type{bc}, but looks for instances of
$p$ that follow from the circumscription axiom.  Allowed keywords:
\type{cutoffs} and \type{succeeds}.

\item [\itype{*circum-trace*}] If not \NIL, the circumscriptive proof
functions will output diagnostics describing their progress.

\item [\itype{circums(p)}] Like \type{bcs}, but looks for instances of
$p$ that follow from the circumscription axiom.  Allowed keywords:
\type{cutoffs} and \type{succeeds}.

\item [\itype{*circumscription-bilattice*}] The bilattice used to compute
the results of the circumscription axiom.

\item [\itype{cnf(p)}] Returns $p$ in conjunctive normal form, as a list
of conjuncts, where each conjunct is a list of disjuncts.  Thus
 \[\type{(CNF '(AND A B))}\]
returns
 \[\type{((A) (B))}\]

\item [\itype{*cnf-atms-bilattice*}] An alternative version of the \atms\
bilattice where the justification information is maintained in conjunctive
normal form.

\item [\itype{combine-1(f params f1 ... fn)}] This is an $n$-ary
version of \type{combine-fns}, where $f$ is an $n$-ary bilattice
function and there are $n$ \dg-functions supplied as well.  The
additional argument \type{params} indicates which of the $n$ arguments
to this function are not necessarily \dg-functions, but are instead
parameters expected by the overall function $f$.  This additional
argument is a list of \T\ (if the $n$th argument is a parameter) or
\NIL\ (if it isn't).  If the parameter list ends early, all of the
remaining entries are assumed to be \NIL.

\item [\itype{combine-fns(f1 f2 fn comp)}] This function takes the two
\dg-functions $f_1$ and $f_2$ and combines them using the function
\type{fn}, which should be a function that takes two elements of the
bilattice and returns a single one.  The comparison function
\type{comp} is used to determine if one of the two functions can be
returned directly; if, for example \type{(comp f1 f2)} holds, than we
simply return $f_1$, and similarly for $f_2$.  As an example, if the
combining function is the bilattice operation $\cdot$, the comparison
function is \type{k-le}.

\item [\itype{compile-mvl}] Compiles and loads the \mvl\ system.
On a Symbolics system, the keyword \type{reload} forces all of the
files to be reloaded and the keyword \type{recompile} forces them all
to be recompiled.

\item [\itype{*completed-database*}] If \T, find all exceptions
to default proofs.  If \NIL\ (the default), accept defaults proofs
with possible exceptions as complete.

\item [\itype{cont-analysis(analysis)}] This function continues the
investigation of a specific analysis, and can be used to resume a
proof effort that has been curtailed for some reason.  It returns two
values: an instance of the \type{answer} structure giving the results
of the proof attempts, and a modified version of the \type{analysis}
structure that can be used to continue the proof process if desired.

\item [\itype{cont-prover(proof)}] Resume the task associated to the
auxiliary information in the given instance of the \type{proof}
structure.  This function returns an instance of \type{bcanswer}
corresponding to the result obtained by the prover, and returns as
soon as a solution to the original query is found.

\item [\itype{contents}] Displays the contents of the database.  Keywords:
\type{cutoffs}, \type{props}.

\item [\itype{create-mvl-invocation}] This function sets up a new
copy of \mvl.  It accepts an optional argument that is an instance of
\mi; this can be used to reset the state of \mvl\ to a previous one.
If the optional argument is omitted, default values are used.  The
value returned is the \mvl\ invocation that was active before the
function was called, so that another call to
\type{create-mvl-invocation} can be used to restore the original
environment.

\item [\itype{cutoffs-not(cutoffs \&optional bilattice)}] Returns a
new cutoffs $c$ such that
 \[\type{(test x c)} = \type{(test (mvl-not x) cutoffs)}\]
 for all truth values $x$.  In other words, $x$ passes $c$ if and
only if $\neg x$ passed the cutoffs supplied originally.

% *D*

\item [\itype{dag-accumulate-fn(dag bilattice list)}] Given a \dg, a
bilattice, and a list of elements of the form
 \[\type{(dag-elt . bilattice-elt)}\]
a \dg-function is constructed that takes the given values at the
given elements of the \dg.

A keyword to \type{dag-accumulate-fn} is \type{modifier}.  This
function is used if the \type{d-list} contains the same \dg\ element
more than once, in which case the modifier function is called with
the new and old values (in that order) and the result is used as the
value of the \dg-function at that particular point.  The default
action is to simply overwrite the old value with the new one.  (In
other words, the default modifier function simply returns its first
argument while ignoring the second.)

\item [\itype{dag-bilattice(dag bilattice)}] Given a bilattice $B$
and a dag $D$, returns the bilattice $B^D$.

\item [\itype{dag-change-dag(f d-f)}] The argument \type{d-f} should
be a function from the \dg\ to itself; a modified version of the
original \dg-function $f$ is returned in which every point at which
the \dg-function takes a new value has had the function \type{d-f}
applied to it.  The bilattice argument is expected to be the bilattice
of functional truth values as opposed to the bilattice that is the
target of the functions.  There is an optional keyword
\type{simplify}; if supplied, the computed function is simplified
before being returned.  (See \type{simplify}.)

\item [\itype{dag-change-tv(f b-f)}] The argument \type{b-f} should be
a function from the bilattice to itself; a modified version of the
original \dg-function $f$ is returned in which every point in the
\dg-function has had the function \type{b-f} applied to it.  There is
an optional keyword \type{simplify}; if supplied, the computed
function is simplified before being returned.  (See \type{simplify}.)

\item [\itype{dag-list-of-answers(f)}] Given a \dg-function $f$, assumed
to be from the \dg\ of binding lists into the currently active bilattice, this
function constructs a list of answers corresponding to the points where
the \dg-function takes values.

\item [\itype{dag-prune(f pred)}] Given a \dg-function $f$ and a
predicate on the bilattice, this function prunes from $f$ all points
$d$ where $f(d)$ fails to satisfy the predicate.  There is an optional
keyword \type{modify-fn} that, if supplied, should be a function $m$
from the bilattice to itself.  In this case, the points where
$m(f(d))$ fails to satisfy the predicate are pruned.

\item [\itype{*default-bilattice*}] A bilattice that can be used for simple
default reasoning.

\item [\itype{delete-bdg(var bdg-list)}] Destructively modifies
\type{bdg-list}, removing any binding for the variable \type{var}.

\item [\itype{denotes(a)}] Returns the proposition associated with a
particular \lisp\ atom.

\item [\itype{dnf(p)}] Returns $p$ in disjunctive normal form, as a list
of disjuncts, where each disjunct is a list of conjuncts.  Thus
 \[\type{(DNF '(AND A B))}\]
returns
 \[\type{((A B))}\]

\item [\itype{(done exp-1 ... exp-n)}] This predicate, which is
useful for stashing only, causes \lisp\ to evaluate each of the subsequent
expressions.

\item [\itype{dot-with-not(x)}] Computes $x\cdot \neg x$.

\item [\itype{dot-with-star}] The function \type{dws} in the active bilattice.

% *E*

\item [\itype{empty}] Unstashes every sentence in the database.  Keyword:
\type{value}.

\item [\itype{empty-theory()}] Empties the currently active theory,
although this theory may be changed using the optional theory
argument.  The keyword \type{value} has a meaning similar to its
meaning as a keyword to \type{erase}.

\item[\itype{equal-answer(a1 a2)}] Returns \T\ if the two answers
\type{a1} and \type{a2} are equal, and \NIL\ otherwise.  The
functions used are \type{binding-equal} and \type{mvl-eq}, the
functions normally used to test binding lists and truth values for
equality.

\item [\itype{equal-binding(b1 b2)}] Checks to see if two binding lists are
effectively the same.

\item [\itype{*equivalence-translation*}] Controls the fashion in which
\type{normal-form} deals with the connective \type{<=>}.

\item [\itype{erase(p)}] Removes $p$ from the database and forward
chains.  Allowed keywords: \type{value}, \type{decrement-p} and
\type{cutoffs}.

\item [\itype{erase-from-theory(p)}] Removes the proposition $p$ from
the currently active theory, and forward chains on the result.
Allowed keywords: \type{theory} can be used to specify the theory from
which $p$ should be removed; \type{value}, \type{decrement-p} and
\type{cutoffs} have the same meaning as for \type{erase}.  (See
Section \ref{s.fc}.)

\item [\itype{existing-mvl-invocation}] This function returns a
single value that is an instance of the structure \mi.  It is useful
primarily for saving the current copy of \mvl\ so that it can be
restored later.

% *F*

\item [\itype{false}] The truth value \type{false} in the active bilattice.

\item [\itype{fc(p)}] Inserts the sentence $p$ into the database and
forward chains.  Allowed keywords: \type{value}, \type{increment-p},
\type{fc-increment-p}, \type{cutoffs} and \type{include-modal}.

\item [\itype{fc-trace(pat)}] Include \type{pat} on
\type{*fc-trace*}, so that any instance of \type{pat} will be
traced.  The argument is optional; if omitted, \type{(?*)} (of which
everything is an instance) will be used.

\item [\itype{fc-untrace(pat)}] Remove any instance of \type{pat}
from \type{*fc-trace*}.  The argument is optional; if omitted,
\type{(?*)} (of which everything is an instance) will be used.

\item [\itype{*final-limit*}] is the limit point beyond which the
search is abandoned, so that any node that evaluates to a value worse
than this will be pruned.  The default value of this parameter is
\NIL, indicating that no nodes are to be pruned.

\item [\itype{*first-order-atms-bilattice*}] This bilattice corresponds
to a fully first-order \atms, as opposed to the more conventional one
dealing only with propositional logic.

\item [\itype{*first-order-bilattice*}] The bilattice corresponding to
first-order logic.

\item [\itype{fn-and(f1 f2)}] This function conjoins the values of
$f_1$ and $f_2$ using the bilattice $\wedge$ operation.  The new
\dg-function is returned.

\item [\itype{fn-comp(f1 f2 fn)}] This functions compares the values
of $f_1$ and $f_2$ using the bilattice comparison function \type{fn}.
Either \T\ or \NIL\ is returned.

\item [\itype{fn-dot(f1 f2)}] This function combines the values of
$f_1$ and $f_2$ using the bilattice $\cdot$ operation.  The new
\dg-function is returned.

\item [\itype{fn-dws(f1 f2)}] This function combines the values of
$f_1$ and $f_2$ using the bilattice operation \type{dot-with-star}.
The new \dg-function is returned.

\item [\itype{fn-eq(f1 f2)}] This function compares the values of
$f_1$ and $f_2$ using the bilattice comparison operation.  Either \T\
or \NIL\ is returned.

\item [\itype{fn-not(f1)}] This function negates the values of $f_1$
using the bilattice $\neg$ operation.  The new \dg-function is
returned.

\item [\itype{fn-or(f1 f2)}] This function disjoins the values of
$f_1$ and $f_2$ using the bilattice $\vee$ operation.  The new
\dg-function is returned.

\item [\itype{fn-plus(f1 f2)}] This function adds the values of $f_1$
and $f_2$ using the bilattice $+$ operation.  The new \dg-function is
returned.

% *G*

\item [\itype{get-bdg(var bdg-list)}] Returns the value assigned to
the variable \type{var} by the binding list \type{bdg-list}.  If the
variable does not appear in the binding list, \NIL\ is returned.

\item [\itype{get-mvl (symbol indicator)}]  Like the Common Lisp
\type{get}, but if the value is a list, treats it as an a-list and
searches for an item whose car matches \type{*mvl-invocation*}.
The value taken can be changed with \type{setf}.

\item [\itype{get-val(d f)}] Returns the value of the function $f$ at
the point $d$ of the \dg.

\item [\itype{*global*}] The root of the \dg\ of theories.

\item [\itype{groundp(p)}] Returns \T\ if the proposition $p$ is
ground (i.e., contains no variables) and \NIL\ otherwise.

% *H*

\item [\itype{hierarchy-bilattice(b)}] Returns the bilattice constructed
as an infinite descending chain of copies of $b$.

\item [\itype{*hierarchy-bilattice*}] A bilattice similar to the
default bilattice but where the defaults are labelled with
integer-valued strengths.

% *I*

\item [\itype{*if-translation*}] Controls the fashion in which
\type{normal-form} deals with the connective \type{if}.

\item [\itype{includes(big little)}] This function is used to state that
the theory \type{little} is a subtheory of the theory \type{big}, and
should inherit from it.

\item [\itype{index(p)}] Updates the discrimination net as needed
to insert the proposition $p$ into the database.

\item [\itype{init-analysis(prop)}] The \mvl\ primitive upon which
the proving functions are based, this function accepts a proposition
and creates and returns an instance of the \type{analysis} structure.

\item [\itype{init-prover(p effort)}] This function initializes and
returns an instance of \type{proof} containing auxiliary information
about the proof process.

\item [\itype{*initial-limit*}] is used to indicate the type of search.
If \NIL\ (the default), then the values assigned to the nodes are not
used to order the search, and depth-first search is used.  If not
\NIL, \type{*initial-limit*} is typically 0.

\item [\itype{instp(p1 p2)}] If $p_2$ is an instance of $p_1$,
return a list of their most general unifiers; only variables in $p_1$
are bound.  Returns \NIL\ if $p_2$ is not an instance of $p_1$.

\item [\itype{invoke(fn p)}] This function invokes \type{fn} on the
sentence $p$.  In addition, any keyword arguments are passed to
\type{fn}, but with \type{\&allow-other-keys} so that the keywords
will be ignored if \type{fn} is not expecting them.

% *J*

% *K*

\item [\itype{k-ge(x y)}] Returns \T\ if $x \geq_k y$, and \NIL\ otherwise.

\item [\itype{k-gt(x y)}] Returns \T\ if $x >_k y$, and \NIL\ otherwise.

\item [\itype{k-le(x y)}] Returns \T\ if $x \leq_k y$, and \NIL\ otherwise.

\item [\itype{k-lt(x y)}] Returns \T\ if $x <_k y$, and \NIL\ otherwise.

\item [\itype{k-not-ge(x y)}] Returns \T\ if $x \not\geq_k y$, and
\NIL\ otherwise.

\item [\itype{k-not-gt(x y)}] Returns \T\ if $x \not>_k y$, and \NIL\
otherwise.

\item [\itype{k-not-le(x y)}] Returns \T\ if $x \not\leq_k y$, and
\NIL\ otherwise.

\item [\itype{k-not-lt(x y)}] Returns \T\ if $x \not<_k y$, and \NIL\
otherwise.

\item [\itype{(known p)}] Stashing this predicate simply stashes $p$;
unstashing it causes an error.  Looking it up looks up $p$.

% *L*

\item [\itype{lattice-to-bilattice(l)}] Returns the bilattice constructed
from the lattice $l$.

\item [\itype{lattice-to-dag(lattice)}] Construct a \dg\ from the
given lattice.

\item [\itype{lattice-to-dag-to-bilattice(lattice bilattice)}]  First
construct a \dg\ $D$ from the given lattice.  Then use the bilattice
$B$ to construct the bilattice $B^D$.

\item [\itype{*limit-delta*}] is used if \type{*initial-limit*} is
non-\NIL.  In this case, the nodes are ordered by increasing
 \[\frac{\type{value(node)}-\type{*initial-limit*}}{\type{*limit-delta*}}\]
 A typical value for *initial-limit* is 0, and for *limit-delta* is
1.  The value of this parameter is irrelevant if
\type{*initial-limit*} is \NIL.

\item [\itype{lisp-defined(fn)}] This macro is used to indicate that
the given function can be evaluated by a \lisp\ call.  It is also
used to cause \mvl\ to signal an error if an attempt is made to stash
or unstash an instance of the function.  \type{Lisp-defined} also
accepts an optional argument, which is an inverter function.  The
overall result of attempting to lookup an instance of the function
can be described as follows:
 \begin{enumerate}
 \item If none of the arguments to the function (i.e., all but the
last of the arguments of the relation) is a variable, then the function
is called and the result is unified with the final argument of the
relation.
 \item If the inverter function has been supplied and is equal to \T,
then the function is called and the result is unified with the final
argument of the relation even if some of the arguments are variables.
 \item If some of the arguments are variables but the result is not a
variable, then the inverter function is invoked if one has been
supplied.  The arguments passed to the inverter function are the
result (i.e., the last argument of the relation), and a list of all
of the arguments (i.e., all but the last argument of the relation).
If there is no inverter function, then \NIL\ is returned.  In
addition, if \wbli\ \index{*warn-on-bad-lisp-invocations*} is not
\NIL, the user is warned.
 \item If both some of the arguments and the result are variables, then
\NIL\ is returned.  In addition, if \wbli\ is not \NIL, the user is
warned.
 \end{enumerate}

\item [\itype{lisp-predicate(pred)}] This macro is used to indicate
that the given predicate can be evaluated by a \lisp\ call.  It is
also used to cause \mvl\ to signal an error if an attempt is made to
stash or unstash an instance of the function.

\item [\itype{load-logic}] Empty the database and make a new bilattice
active.  Accepts as an optional argument either a character designating
a bilattice, or the new bilattice itself.  If no argument is supplied,
the user is prompted for the new bilattice.

\item [\itype{load-mvl}] Loads the \mvl\ system.

\item [\itype{lookup(q)}] Searches the database for a sentence
matching the query $q$.  Allowed keywords: \type{cutoffs}, \type{inv},
\apa, \somo.  Returns two values: an answer and the source of the
match.  Returns \NIL\ if the query cannot be found in the database.

\item [\itype{lookup-call(fn p)}] This function is similar to
\type{invoke}, but modifies the values returned to make them agree
with those expected to be returned by \type{lookup}.  Unless the
function \type{fn} supplies them, the truth value is assumed to be
\type{true}.  The second value is \type{attach}.

\item [\itype{lookups(q)}] Similar to \type{lookup}, but looks
for {\em all\/} sentences matching the query.  Returns two lists of
values.

\item [\itype{lookups-call(fn p)}] The plural version of
\type{lookup-call}.

% *M*

\item [\itype{make-root(dag bilattice b)}] Returns a \dg-function that
takes the value $b$ at the root of the \dg.

\item [\itype{make-tv(dag bilattice d b)}] Returns a \dg-function that
takes the value $b$ at the point $d$.  If $d$ is not the root of the
\dg, the value at the root of the \dg\ is generally the point
\type{unknown} in the bilattice, but the user may supply the keyword
\type{root-fn} if he wishes, in which case the value at the root of
the \dg\ is the result of applying \type{root-fn} to the bilattice.

\item [\itype{matchp(pattern atom)}] Returns a list of elements of
the form \type{(b . d)} where $b$ is a most general binding list such
that \type{(plug pattern b)} is an instance of the proposition
labelled with \type{atom} and $d$ is a binding list for the variables
in the database proposition.  \NIL\ is returned if no such binding
list exists.  Thus if \type{P1} is \type{(TALL \#:?2)}, and \type{P2}
is \type{(TALL TOM)}, then
 \[\type{(matchp (TALL TOM) P1)}\]
returns a list consisting of the single entry 
 \[\type{(NIL . ((\#:?2 . TOM)))}\]
 The \type{car} of this entry is \NIL, indicating that no variables
in the query need to be bound to match \type{P1}; the \type{cdr} is
\type{((\#:?2 . TOM))} since the variable \type{\#:?2} in \type{P1}
needs to be bound to \type{TOM} to match the query.
 \[\type{(matchp (TALL ?X) P2)}\]
returns \type{((((?X . TOM)) . NIL))}.

\item [\itype{mvl-and}] The function \type{and} in the active bilattice.

\item [\itype{mvl-dot}] The function \type{dot} in the active bilattice.

\item [\itype{mvl-eq}] The function \type{eq} in the active bilattice.

\item [\itype{mvl-file(string)}] Returns the path for the \mvl\ file
associated with the given name.  The file name may include an
extension, one may be supplied as an optional argument, or
\type{.MVL} will be used by default.  There is an optional keyword
\type{no-error} indicating that even if the file cannot be found, the
path should be returned.

\item [\itype{*mvl-invocation*}] A global variable containing a unique
identifier for the currently active copy of \mvl.

\item [\itype{mvl-load(file)}] Loads the sentences in the given file into
the \mvl\ database.

\item [\itype{mvl-not}] The function \type{not} in the active bilattice.

\item [\itype{mvl-or}] The function \type{or} in the active bilattice.

\item [\itype{mvl-plug}] The function \type{plug} in the active bilattice.

\item [\itype{mvl-plus}] The function \type{plus} in the active bilattice.

\item [\itype{mvl-print}] This function accepts a single argument, a
truth value, and modifies it in a way that will cause it to be
pretty-printed.  The value returned by mvl-print should not be used
for any purpose other than printing.

\item [\itype{mvl-save(file)}] Writes the \mvl\ database out to the
given file.  This function also accepts the keywords \type{cutoffs}
and \type{props}, which are treated as for \type{contents} (see
Section \ref{s.contents}).  Finally, the keyword \type{if-exists} is
passed to the \type{with-open-file} macro.

\item [\itype{mvl-simplicity}] The function \type{simplicity} in the active
bilattice.

\item [\itype{mvl-test-file(string)}] Can be used to evaluate
various forms by the \mvl\ system and compare the results to
predefined values.  Both the forms and the expected answers are in
the \type{.TEST} file whose name is supplied to the function.
Returns the number of errors found.

\item [\itype{mvl-unload(file)}] Removes the sentences in the given file from
the \mvl\ database.

\item [\itype{mvl-vars}] The function \type{vars} in the active bilattice.

% *N*

\item [\itype{negate(p)}] Returns the negation of $p$.

\item [\type{new-$*$var()}]\index{new-star-var} Returns a new
temporary sequence variable.

\item [\type{new-?var()}]\index{new-query-var} Returns a new
temporary variable.

\item [\itype{new-val(a v)}] Sets the truth value of the sentence
associated with the atom $a$ to $v$.  If $v$ is the truth value
\type{unknown}, the sentence is removed from the database.

\item [\itype{normal-form(p)}] Returns an equivalent form of $p$
suitable for insertion into the database (i.e., with the connectives
\type{IFF}, \type{IF} and \type{<=>} reduced to \type{<=} and \type{=>}).

% *O*

% *P*

\item [\itype{plug(p bdg-list)}] Returns the result of instantiating
the proposition $p$ using the given binding list.  This function is
nondestructive.

\item [\itype{prfacts(atom)}] Displays the database facts concerning the
given atom.  Keyword: \type{cutoffs}.

\item [\itype{*probability-bilattice*}] An extended version of the \atms\
bilattice that maintains probabilistic information as well.

\item [\itype{product-bilattice(b1 b2)}] Returns the bilattice
that is the product of the bilattices $b_1$ and $b_2$.

\item [\itype{(property symbol value indicator)}] Stashing this
predicate sets the symbol's \type{indicator} property to \type{value}
for the current invocation of \mvl.  Unstashing it removes the
\type{indicator} property for the current invocation of \mvl\ if the
previous value of this property unifies with \type{value}.  (See the
definitions of \type{get-mvl} and \type{rem-mvl}.)

\item [\itype{props()}] Returns a list of all propositions about
which any information has been recorded.

% *Q*

% *R*

\item [\itype{release-mvl}] (Symbolics only)  Compiles and loads the
\mvl\ system, and declares the new version to be \type{released}.  The
keyword \type{no-recompile} can be used to skip the compilation step.

\item [\itype{rem-mvl (symbol indicator)}] Like the Common Lisp
\type{remprop}, but if the value is a list, treats it as an a-list
and removes the item whose car matches \type{*mvl-invocation*}.

% *S*

\item [\itype{samep(p1 p2)}] Returns \T\ if the propositions $p_1$ and
$p_2$ are identical (up to variable renaming), and \NIL\ otherwise.
If the optional \type{return-answer?} argument is supplied and the
propositions are the same, a binding list is returned instead of \T.

\item [\itype{*save-nodes*}] If \T\ (the default), nodes are
labelled with consecutive integers when they are displayed by the
first-order prover.

\item [\itype{set-search-parameters}] This function is the primitive
used to modify the parameters used to control the first-order theorem
prover's search.  It accepts as keywords \type{initial-limit},
\type{limit-delta}, and \type{final-limit}.

\item [\itype{simplify(f)}] In some cases, manipulating the
\dg-function $f$ may have cause it to contain redundancies, where the
value at some point is presented explicitly even though this value is
in fact the one that would normally be inherited from its parents.
This function returns a modification of $f$ where these irrelevant
entries have been removed.

\item [\itype{splice-bilattice(upper lower)}] Returns the bilattice
obtained by replacing the point \type{unknown} in the bilattice
\type{upper} with a copy of the bilattice \type{lower}.

\item [\itype{standardize-operators(p)}] Returns an equivalent form of $p$
with the connectives \type{<=>}, \type{IF} and \type{IFF} removed.
This function does not use the control variables
\type{*if-translation*} and \et, since it is used when proving
sentences, not when inserting them into the database.

\item [\itype{stash(p)}] Inserts $p$ into the database.  The allowed
keywords are \type{value}, \type{increment-p} and \apa.

\item [\itype{stash-value}] The function \type{stash-val} in the
active bilattice.

\item [\itype{state(p)}] Inserts $p$ into the database after
converting it to normal form.  Keywords: \type{value},
\type{increment-p}.

\item [\itype{state-in-theory(p)}] Adds the proposition $p$ to
the currently active theory.  Allowed keywords: \type{theory} can be
used to specify a different theory to which $p$ should be added;
\type{value} and \type{increment-p} have the same meaning as for
\type{state}.  (See Section \ref{s.database}; if these keywords are
not used, the function \type{state} can be used equally well, since
the default stash value with the theory mechanism activated also
inserts the proposition into the currently active theory.)

\item [\itype{std-cutoffs}] Cutoff information constructed from the
value of \type{cutoff-val} in the active bilattice.  See Section
\ref{s.cutoffs}.

% *T*

\item [\itype{t-ge(x y)}] Returns \T\ if $x \geq_t y$, and \NIL\ otherwise.

\item [\itype{t-gt(x y)}] Returns \T\ if $x >_t y$, and \NIL\ otherwise.

\item [\itype{t-le(x y)}] Returns \T\ if $x \leq_t y$, and \NIL\ otherwise.

\item [\itype{t-lt(x y)}] Returns \T\ if $x <_t y$, and \NIL\ otherwise.

\item [\itype{t-not-ge(x y)}] Returns \T\ if $x \not\geq_t y$, and
\NIL\ otherwise.

\item [\itype{t-not-gt(x y)}] Returns \T\ if $x \not>_t y$, and \NIL\
otherwise.

\item [\itype{t-not-le(x y)}] Returns \T\ if $x \not\leq_t y$, and
\NIL\ otherwise.

\item [\itype{t-not-lt(x y)}] Returns \T\ if $x \not<_t y$, and \NIL\
otherwise.

\item [\itype{test(value cutoffs \&optional bilattice)}] Returns \T\
if the truth value passes the test associated to \type{cutoffs}, and
\NIL\ otherwise.  The test is performed for the supplied bilattice,
or for the currently active one is the optional argument is omitted.

\item [\itype{test-mvl}] Tests the \mvl\ system by invoking the function
\type{mvl-test-file} on the file \type{MVL.TEST}.

\item [\itype{theory-active-p}] Returns \T\ or \NIL\ depending on
whether or not the theory mechanism has been activated.

\item [\itype{theory-contents()}] Displays the contents of the
currently active theory, which can be changed using the
keyword \type{theory}.  The keywords \type{cutoffs} and \type{props}
are allowed and are treated as for \type{contents}.  (See Section
\ref{s.contents}.)  There is an additional keyword,
\itype{this-theory-only}, which defaults to \T\ and, if supplied,
indicates that only sentences whose truth values {\em change\/} in the
given theory are to be displayed.

\item [\itype{theory-dag}] The \dg\ of theories.

\item [\itype{*time-bilattice*}] A compound bilattice that can be used
for reasoning about events occurring against a temporal background.

\item [\itype{true}] The truth value \type{true} in the active bilattice.

\item [\itype{true-in-theory}]  The truth value of a sentence that is
true in the currently-active theory.

\item [\itype{truth-value(a)}] Returns the truth value associated with
a particular \lisp\ atom.  Thus, if the result of asserting
\type{(not (short tom))} into the database was the atom \type{P38},
\type{(denotes 'p38)} would evaluate to \type{(short tom)}, and
\type{(truth-value 'p38)} would evaluate to \type{false}. 

% *U*

\item [\itype{unifyp(p1 p2)}] Returns a list of the most general
unifiers of $p_1$ and $p_2$, or \NIL\ if the propositions fail to
unify.

\item [\itype{unincludes(big little)}] This function is used to state
that the theory \type{little} is no longer a subtheory of the theory
\type{big}, and should not inherit from it.

\item [\itype{unindex(p)}] Removes from the discrimination net 
information relating to the proposition $p$.

\item [\itype{unknown}] The truth value \type{unknown} in the active
bilattice.

\item [\itype{(unknown p)}]  Stashing this predicate unstashes $p$;
unstashing it causes an error.  Looking it up fails if $p$
can be found in the database, and succeeds otherwise.

\item [\itype{unstash(p)}] Removes $p$ from the database.  Allowed
keywords: \type{value}, \type{decrement-p} and \apa.

\item [\itype{unstash-facts-with-atom(atom)}] Unstashes from the
database any facts concerning the given atom.  Keywords:
\type{value}, \type{decrement-p}, \apa.

\item [\itype{unstash-value}] The function \type{unstash-val} in the
active bilattice.

\item [\itype{unstate(p)}] Removes $p$ from the database after
converting it to normal form.  Keywords: \type{value},
\type{decrement-p}.

\item [\itype{unstate-from-theory(p)}] Removes the proposition $p$
from the currently active theory.  Allowed keywords:
\type{theory} can be used to specify a different theory from which
$p$ should be removed; \type{value} and \type{decrement-p} have the
same meaning as for \type{unstate}.  (See Section \ref{s.database}.)

% *V*

\item [\itype{(value symbol val)}] Stashing this predicate sets the
given symbol to the value \type{val}.  Unstashing it makes the symbol
unbound if its current value unifies with \type{val}.

\item [\itype{varp(x)}] Returns \type{T} if $x$ is a variable, and \NIL\
otherwise.

\item [\type{varp$*$(x)}]\index{varp-star} Returns \T\ if $x$ is a
sequence variable, and \NIL\ otherwise.

\item [\type{varp?(x)}]\index{varp-query} Returns \T\ if $x$ is a
standard (i.e., non-sequence) variable, and \NIL\ otherwise.

\item [\itype{vars-in(p)}] Returns a list of the variables appearing in $p$.

\item [\itype{vartype(x)}] Returns ? if $x$ is a standard variable,
$*$ if $x$ is a sequence variable, and \NIL\ is $x$ is not a
variable.

% *W*

\item [\itype{with-mvl-invocation}] This is a macro similar to the
Common Lisp macro \type{with-open-file}.  It accepts an optional
argument that is an instance of \mi\ and is followed by a body of
forms (presumably including a \type{load-logic}) that are evaluated
in a locally-bound copy of \mvl.

% *X*

% *Y*

% *Z*

\end{list}

\documentstyle{article}

\setlength{\topmargin}{0pt}
\setlength{\oddsidemargin}{0pt}
\setlength{\evensidemargin}{0pt}
\setlength{\textwidth}{6.5in}
\setlength{\textheight}{8.5in}
\setlength{\headheight}{0pt}

%here are some mathematical symbols
\def\emptyset{{\hbox{\rm \O}}}
\def\set#1{\{#1\}}
\newcommand{\type}[1]{\mbox{\tt #1}}
\newcommand{\str}{?$*$}
\def\R{\rlap{\rm I}\,{\rm R}}
\def\powset{{\cal P}}
\def\glb{{\mbox{\rm glb}}}\def\lub{{\mbox{\rm lub}}}

\newcommand{\nm}{non\-mon\-o\-ton\-ic}
\newcommand{\Nm}{Non\-mon\-o\-ton\-ic}

\newcommand{\prolog}{{\sc prolog}}\newcommand{\Prolog}{{\sc Prolog}}
\newcommand{\atms}{{\sc atms}}\newcommand{\Atms}{{\sc Atms}}

%this is where the title page begins

\def\draftline{}

\def\authors{Matthew L. Ginsberg}

\newcommand{\nodblspace}{\baselineskip=\normalbaselineskip}
\def\group{\nodblspace Department of Computer Science\\
Stanford University\\
Stanford, California 94305\\
\bigskip
Updated \today
}

\def\titlebottom{\vspace{1.55in}
{\large \authors}\\
\vspace{1.55in}
\group
\end{center}
\clearpage}

\def\xtitletop{\vspace*{1.35in}\begin{center}}
\def\xtitl#1#2{\xtitletop
{\Large\bf #2}\\
\titlebottom}

\newcommand{\xtitle}[2]
{\def\tphrase{#2}\begin{titlepage}\draftline\xtitl{#1}{#2}\end{titlepage}}

\newcommand{\mvl}{{\sc mvl}}
\newcommand{\Mvl}{{\sc Mvl}}
\newcommand{\lisp}{{\sc lisp}}
\newcommand{\Lisp}{{\sc Lisp}}
\newcommand{\dg}{{\sc dag}}
\newcommand{\Dg}{{\sc Dag}}
\newcommand{\T}{\type{T}}
\newcommand{\NIL}{\type{NIL}}
\newcommand{\apa}{{\tt allow-\-procedural-\-attachment}}
\newcommand{\at}{{\tt *anytime-\-trace*}}
\newcommand{\mi}{{\tt mvl-\-invocation}}
\newcommand{\wbli}{{\tt *warn-\-on-\-bad-\-lisp-\-invocations*}}
\newcommand{\somo}{{\tt succeed-\-on-\-modal-\-ops}}
\newcommand{\et}{{\tt *equivalence-\-translation*}}
\newcommand{\itype}[1]{\type{#1}\index{#1 **}}

\makeindex
%\renewcommand{\index}[1]{}

\begin{document}
\xtitle{}{User's Guide to the MVL System}

\tableofcontents
\clearpage

\section{Introduction}
\subsection{What you got}

The \mvl\ system is currently distributed via anonymous ftp from
t.stanford.edu.  A .DVI version of this file, the \mvl\ manual,
can be found in the \type{papers} directory on this machine.

The \type{mvl} directory contains the source code for the \mvl\ system
itself, which is written in standard Common Lisp.  It should be
possible to port the system to any Common Lisp without substantial
difficulty.  In addition, the \type{papers} directory contains .DVI
versions of the following papers that relate to the functioning of the
system:
 \begin{enumerate}
 \item ``Multivalued logics: A uniform approach to inference in
artificial intelligence.''  This paper gives a formal description of
multivalued logics and the \mvl\ system.
 \item ``Bilattices and modal operators.''  This paper describes the
formal connection between the \mvl\ system and modal operators.
 \end{enumerate}

The \mvl\ sources are public domain.  You are free to use, modify, and
distribute them as you see fit, provided that such use or distribution
is not for commercial purposes.  No warranties are made with regard to
the performance of the system, however, and no commitment of support
has been made.

\subsection{Overview}

This user's guide has been written to help you install and use the \mvl\
system.  It describes the basic structure of the system, and shows you
how to take advantage of the flexibility of multivalued inference.

Essentially, \mvl\ is just like any other logic programming system.
There is a database consisting of those sentences about which the
system knows something, and there are facilities for determining
whether or not some query follows from the information the database
contains.  Viewed in this fashion, the \mvl\ system is not much
different from a conventional logic programming language such as
\prolog.

What makes \mvl\ unique is that each sentence in the database is
labelled with a {\em truth value\/}\index{truth value}.  Thus,
instead of simply assuming a statement to be true and dumping it into
the database, we will {\em label\/} the statement as true when we
insert it.  Of course, we can use other labels as well: true by
default, true by virtue of some assumption, or even false or false by
default.  \Mvl\ considers these truth values when determining how to
respond to a \prolog-like query.

This manual is divided into two parts.  In the first part (Sections
\ref{s.fol1} -- \ref{s.support}), we describe the use of the system
where the truth values assigned to sentences are simply ``true'' or
``false.''  Used in this fashion, \mvl\ is not much different from
\prolog.

The second part of the manual shows you how to use logics involving
more complex truth values.  We examine the logics that have been
predefined as part of the \mvl\ system, and describe the various ways
in which you can define your own logics.  Several tools for this
purpose have been included as part of the \mvl\ system.  In addition,
we describe \mvl's facilities for handling modal operators, which
are operators that accept logical sentences as arguments.

\subsection{Why MVL is unique}

As mentioned in the previous section, what distinguishes \mvl\ from
other declarative programming languages is its support for multiple
truth values.  There are other differences, however, and these will
be discussed in detail in subsequent sections.  Here's a quick summary:
 \begin{enumerate}
 \item As mentioned, the principal feature offered by \mvl\ is its
support of multiple truth values.
 \item The basic theorem prover used by the system includes facilities
to detect loops and control recursion encountered in proof search.
The first-order theorem prover is discussed in Section \ref{s.fo},
although we will have nothing further to say about recursion control.
 \item The system can handle negative subgoals containing free variables
as discussed in Section \ref{s.effort}.
 \item The \mvl\ unifier includes an occurcheck and full support for
sequence variables.  This is discussed in Section \ref{s.rep}.
 \end{enumerate}

\section{Installation and testing}

The \mvl\ system is currently developed and maintained on a
Sparcstation running Allegro Common Lisp and is tested only
infrequently on other platforms.  Although the system has been written
in standard Common Lisp to maximize ease of portability and no
problems have been encountered when moving the system to a variety of
platforms, performance on other systems is not guaranteed.

\Mvl\ currently supports both ANSI Common Lisp and Common Lisp as
described in Steele's 1984 book.  In order to determine whether or not
the ANSI standard is conformed to, the system checks to see if
\type{in-package} is a macro or a function.  If the former, it is
assumed that the entire ANSI standard is supported; if the latter, it
is assumed that the language is as in Steele and various extensions
are defined.  It is expected that some future release of \mvl\ will
support the ANSI standard only.

\subsection{Sun-4}

In order to load the system on a Sparcstation running Allegro Common
Lisp, do the following:
 \begin{enumerate}

 \item (ANSI) Set up a logical pathname translation so that logical
pathnames beginning with \type{MVL} point to the direction where \mvl\
is kept.  (non-ANSI) Edit the file \type{mvldcl.lisp} so that the
\type{:default-pathname} field in the \mvl\ system definition points
to the directory where \mvl\ will be kept.  This directory should also
contain the following files:
 \begin{enumerate}\setlength{\itemsep}{0pt}\raggedright
 \item {\tt f1.mvl, recursion.mvl, def-test.mvl, circ-test-1.mvl,
circ-test-2.mvl, circ-test-3.mvl, circ-test-4.mvl, completed-db-1.mvl,
completed-db-2.mvl, modal-test.mvl, actions.mvl, shooting.mvl}
 \item {\tt mvl.test, bc.test, recursion.test, atms.test,
defaults.test, circumscription.test, completed-db.test, modal.test,
fc.test, temporal.test}
 \end{enumerate}
 \item Compile and load the file \type{mvldcl.lisp}.
 \item Execute the function
\type{(mvl:compile-mvl)}\index{compile-mvl} to compile and load the
\mvl\ system.  Using this function in the future will only compile 
those files that have been changed since last compiled.
 \item Invoke the function \type{(use-package 'mvl)} to access the
functions and symbols of the \mvl\ system (which is kept in its own
package; see Section \ref{s.package}).
 \end{enumerate}

Now invoke the function {\tt (test-mvl)}, which should return 0
(i.e., 0 errors encountered).  If it doesn't, the problem is probably
that \mvl\ is unable to locate the files mentioned in (1) above.

\subsection{Symbolics 3600 series}

The \mvl\ system will run on Symbolics Genera 7.2 or later.  It is
possible to run the system using Genera 7.0 or 7.1, although a bug in
the implementation of the Common Lisp \type{\&allow-other-keys}
keyword may cause the system to dump to the cold load stream if
certain predicates are procedurally attached.  (See Section
\ref{s.attach}.)

You should begin by editing the file \type{mvldcl.lisp} so that the
default pathname in the \mvl\ system declaration points to the
directory in which the files are actually kept.  You can then load
the \mvl\ system using either the Load System command or the Lisp
function \type{(mvl:load-mvl)}\index{load-mvl} that is defined in the
file \type{mvldcl.lisp}.  Now access the \mvl\ package using
\type{(use-package 'mvl)} and execute the function \itype{test-mvl}.
This function should return 0 (i.e., 0 errors encountered).  If it
doesn't, the problem is probably that \mvl\ is unable to find the
directory its test files are on (they're included in the distribution
tape).

\subsection{Other machines}

The \mvl\ system has also been run and tested on a wide variety of
other machines.  In order to load the system, follow the instructions
provided for Sun-4's.

\section{The MVL database}
\label{s.fol1}
\subsection{Representation}
\subsubsection{Structure of the database}
\label{s.rep}

As already remarked, the \mvl\ database consists of sentences labelled
with truth values.  In this section, we will describe the form of the
sentences themselves; in the next section, we comment briefly on the
truth values used to label them.

Sentences in the database are represented as \lisp\ s-expressions
using the prenex normal form \index{prenex normal form} of the
standard logical connectives \type{NOT}, \type{OR} and \type{AND}.
Thus the statement that Fred is tall and Harry is not would be
written as:
 \[\type{(AND (TALL FRED) (NOT (TALL HARRY)))}\]

Any \lisp\ atom can be used as a constant term, with the exception of
variables (see below) and the following reserved connectives:
 \[\type{NOT AND OR => <= <=> IF IFF}\]
The connective \type{<=} is used in backward-chaining rules;
 \[\type{(<= CONCLUSION PREMISE-1 ... PREMISE-N)}\]
is logically equivalent to
 \[\type{(IF (AND PREMISE-1 ... PREMISE-N) CONCLUSION)}\]
as is the forward-chaining version
 \[\type{(=> PREMISE-1 ... PREMISE-N CONCLUSION)}\]
 Although logically equivalent, these expressions behave differently
in practice.  The version
 \[\type{(IF (AND PREMISE-1 ... PREMISE-N) CONCLUSION)}\]
 in which the connective \type{IF} is used represents a pair of
rules, depending on the value of the parameter \itype{*if-translation*}.
If the value is \type{bc} (the default), then two backward-chaining
rules are produced, so that
 \[\type{(IF (A) (B))}\]
is equivalent to
 \[\type{(AND (<= (B) (A)) (<= (NOT (A)) (NOT (B))))}\]
 the conjunction of the original rule and its contraposition.
If the value of \type{*if-translation*} is \type{fc}, then
forward-chaining rules are produced and the translation is to
 \[\type{(AND (=> (A) (B)) (=> (NOT (B)) (NOT (A))))}\]
 Finally, if \type{*if-translation*} has the value \type{mix}, the
pair of noncontraposed rules
 \[\type{(AND (<= (B) (A)) (=> (A) (B)))}\]
 is produced.

Equivalences are represented using either \type{IFF} or \type{<=>}.
Thus
 \[\type{(IFF (FOO FRED) (GOO FRED))}\]
is equivalent to
 \[\type{(AND (IF (FOO FRED) (GOO FRED)) (IF (GOO FRED) (FOO FRED)))}\]
 Translations involving \type{<=>} depend on the value of the
parameter \et\index{*equivalence-translation*}.  As above, if the
value is \type{bc} (the default), then two backward-chaining rules
are produced, so that
 \[\type{(<=> (FOO FRED) (GOO FRED))}\]
translates to
 \[\type{(AND (<= (FOO FRED) (GOO FRED)) (<= (GOO FRED) (FOO FRED)))}\]
 If \et\ has the value \type{fc}, two forward-chaining rules are produced
so that the translation is to
 \[\type{(AND (=> (FOO FRED) (GOO FRED)) (=> (GOO FRED) (FOO FRED)))}\]
 Finally, if the value is \type{mix}, one forward-chaining rule and
one backward-chaining rule are produced:
 \[\type{(AND (=> (FOO FRED) (GOO FRED)) (<= (FOO FRED) (GOO FRED)))}\]

Variables \index{variables} are represented using \lisp\ atoms
beginning with the character \type{?} and are handled as in \prolog,
so that free variables in database statements are implicitly
universally quantified.  Again as in \prolog, existential
quantification is handled via Skolemization; thus the
backward-chaining sentence ``Every rich person is happy'' and the
sentence ``There is a rich person'' would be written as
 \begin{equation}
\type{(<= (HAPPY ?P) (PERSON ?P) (RICH ?P))}
\label{e.if}
\end{equation}
and
 \[\type{(AND (PERSON P0017) (RICH P0017))}\]
respectively.

The variables in database sentences are standardized apart from the
other variables used by the system.  This is done by using uninterned
variables in the actual representation of a database fact, so that
the internal representation of (\ref{e.if}) would be
 \[\type{(<= (HAPPY \#:?P) (PERSON \#:?P) (RICH \#:?P))}\]
 The original variables are generally reinstated when \mvl\ is asked
to display a database proposition.

\subsubsection{Sequence variables}

In addition to the simple variables discussed in the previous
section, \mvl\ supports the use of ``sequence
variables.''\index{sequence variables} These variables begin with the
characters \type{\str}, and will match not a single atom, but an
arbitrary (possibly empty) sequence of atoms.

Thus, for example,
 \[\type{(FOO A B)}\]
 will unify with
 \[\type{(FOO \str X)},\]
 returning the binding list
 \[\type{((\str X . (A B)))}\]
 in which \type{?*X} has been bound to the list \type{(A B)}.  This
binding list will presumably be displayed as
 \[\type{((\str X A B))}.\]
 Similarly, the most general unifier of $\type{(FOO)}$ and
$\type{(FOO \str X)}$ is given by $\type{((\str X))}$.

The presence of sequence variables may cause the most general unifier
of two patterns to be nonunique.  For example, the result of unifying
 \[\type{(FOO A B)}\]
 with
 \[\type{(FOO \str A \str B)}\]
 can be any one of the three answers
 \begin{eqnarray*}
& \type{((\str A A B) (\str B))} & \\
& \type{((\str A A) (\str B B))} & \\
& \type{((\str A) (\str B A B))} &
 \end{eqnarray*}

For this reason, the unification functions return not a single answer,
but a {\em list\/} of possible unifiers.

\subsubsection{Functions to manipulate database sentences}
\label{s.manipulate}

The following functions and global parameters are used to manipulate
sentences in the database:
 \begin{description}
 \item [\itype{varp(x)}] Returns \T\ if $x$ is a variable, and
\NIL\ otherwise.
 \item [\type{varp?(x)}]\index{varp-query} Returns \T\ if $x$ is a
standard (i.e., non-sequence) variable, and \NIL\ otherwise.
 \item [\type{varp$*$(x)}]\index{varp-star} Returns \T\ if $x$ is a
sequence variable, and \NIL\ otherwise.
 \item [\itype{vartype(x)}] Returns ? if $x$ is a standard variable,
$*$ if $x$ is a sequence variable, and \NIL\ is $x$ is not a
variable.
 \item [\type{new-?var()}]\index{new-query-var} Returns a new
temporary variable.
 \item [\type{new-$*$var()}]\index{new-star-var} Returns a new
temporary sequence variable.
 \item [\itype{normal-form(p)}] Returns an equivalent form of $p$
suitable for insertion into the database (i.e., with the connectives
\type{IFF}, \type{IF} and \type{<=>} reduced to \type{<=} and \type{=>}).
 \item [\itype{*if-translation*}] Controls the fashion in which
\type{normal-form} deals with the connective \type{if}.
 \item [\itype{*equivalence-translation*}] Controls the fashion in which
\type{normal-form} deals with the connective \type{<=>}.
 \item [\itype{negate(p)}] Returns the negation of $p$.
 \item [\itype{standardize-operators(p)}] Returns an equivalent form
of $p$ with the connectives \type{<=>}, \type{IF} and \type{IFF}
removed.  (Note: This function does not use the control variables
\type{*if-translation*} and \et, since it is used when proving
sentences, not when inserting them into the database.)
 \item [\itype{cnf(p)}] Returns $p$ in conjunctive normal form, as a
list of conjuncts, where each conjunct is a list of disjuncts.  Thus
 \[\type{(CNF '(AND A B))}\]
returns
 \[\type{((A) (B))}\]
\item [\itype{dnf(p)}] Returns $p$ in disjunctive normal form, as a list
of disjuncts, where each disjunct is a list of conjuncts.  Thus
 \[\type{(DNF '(AND A B))}\]
returns
 \[\type{((A B))}\]
\end{description}

\subsection{Truth values}

The truth values used to label the sentences in the \mvl\ database
are elements of a bilattice that can be selected by the user with
the function \type{load-logic}.  We will defer a discussion of these
truth values to the second portion of this manual, which discusses
the use of bilattices that label sentences as other than simply true
or false.  At this point, we merely remark that if it is our intention
to assert the sentence \type{(not (foo))}, this is actually accomplished
by asserting that the negated sentence \type{(foo)} has truth value
\type{false}.

\subsection{Database form and internals}
\label{s.database}

The database is maintained using a discrimination net, as suggested
by Charniak et.~al in their book on {\sc lisp} programming.  Any
particular sentence in the database, say \type{(tall tom)}, is
associated with a \lisp\ atom such as \type{P27} that is its
``proposition number.''  The propositions are numbered consecutively,
beginning with \type{P1}.  At any given point in time, the function
\type{(props)} can be invoked to return a list of all of the
propositions in the database.

Very few of the functions that are used to manipulate the internal form
of the database are likely to be needed by a user of the \mvl\ system.  The
following are the most useful:
 \begin{description}
 \item [\itype{props()}] This function returns a list of all
propositions about which any information has been recorded.
 \item [\itype{denotes(a)}] Returns the proposition associated with a
particular \lisp\ atom.
 \item [\itype{truth-value(a)}] Returns the truth value associated
with a particular \lisp\ atom.  Thus, if the result of asserting
\type{(not (short tom))} into the database was the atom \type{P38},
\type{(denotes 'p38)} would evaluate to \type{(short tom)}, and
\type{(truth-value 'p38)} would evaluate to \type{false}. 
 \item [\itype{index(p)}] Updates the discrimination net as
needed to insert the proposition $p$ into the database.
 \item [\itype{unindex(p)}] Removes from the discrimination net
information relating to the proposition $p$.
 \end{description}

\section{Quantification}

As we have remarked, quantification is treated in \mvl\ in a fashion
similar to that used by \prolog\ systems:  Free variables appearing
in the database are assumed to be universally quantified; free variables
appearing in queries are existentially quantified.  Skolemization is
used to handle existential quantification in database statements or
universal quantification in queries.

The use of this approach means that the answer to an existentially
quantified query will often be of a form that indicates how the
variables should be bound in the result.  In \mvl\, these binding
lists \index{binding lists} are represented as expressions of the
form
 \[\type{((var-1 . bdg-1) ... (var-n . bdg-n))}\]
 where \type{var-1} through \type{var-n} are variables and
\type{bdg-1} through \type{bdg-n} are their respective bindings.
If no variable is bound to a value, the empty binding list \NIL\
is used.

As an example, suppose that the database contains the single sentence
 \[\type{(LIKES COCO ?X)}\]
 indicating that Coco likes everything.  Now the query
 \[\type{(LIKES ?Y MATT)}\]
 will be interpreted as, ``Find something that likes Matt.''  The answer
returned will be
 \[\type{((?Y . COCO))}\]
since Coco likes Matt.  Since the queries
 \[\type{(LIKES COCO MATT)}\]
and
 \[\type{(LIKES COCO ?Z)}\]
 both succeed without binding a variable, the answer \NIL\ will be
returned for each.\footnote{In order to distinguish the empty binding
list \NIL\ from a failed database query, all queries return a binding
list as well as a truth value.  This is discussed in Section
\ref{s.answer}.}

\subsection{Functions used to manipulate binding lists}

A variety of functions within \mvl\ are used to construct and manipulate
binding lists.  Here are some of the more useful ones:
 \begin{description}
 \item [\itype{groundp(p)}] Returns \T\ if the proposition $p$ is
ground (i.e., contains no variables) and \NIL\ otherwise.
 \item [\itype{vars-in(p)}] Returns a list of the variables
appearing in $p$.
 \item [\itype{get-bdg(var bdg-list)}] Returns the value assigned to
the variable \type{var} by the binding list \type{bdg-list}.  If the
variable does not appear in the binding list, \NIL\ is returned.
 \item [\itype{delete-bdg(var bdg-list)}] Destructively modifies
\type{bdg-list}, removing any binding for the variable \type{var}.
 \item [\itype{plug(p bdg-list)}] Returns the result of instantiating
the proposition $p$ using the given binding list.  This function is
nondestructive.
 \item [\itype{samep(p1 p2)}] Returns \T\ if the propositions $p_1$
and $p_2$ are identical (up to variable renaming), and \NIL\
otherwise.  If the optional \type{return-answer?} argument is
supplied and the propositions are the same, a binding list is
returned instead of \T.
 \item [\itype{unifyp(p1 p2)}] Returns a list of the most general
unifiers of $p_1$ and $p_2$, or \NIL\ if the propositions fail to
unify.
 \item [\itype{instp(p1 p2)}] If $p_2$ is an instance of $p_1$,
return a list of their most general unifiers; only variables in
$p_1$ are bound.  Returns \NIL\ if $p_2$ is not an instance of $p_1$.
 \item [\itype{matchp(pattern atom)}] Returns a list of elements of
the form \type{(b . d)} where $b$ is a most general binding list such
that \type{(plug pattern b)} is an instance of the proposition
labelled with \type{atom} and $d$ is a binding list for the variables
in the database proposition.  \NIL\ is returned if no such binding
list exists.  Thus if \type{P1} is \type{(TALL \#:?2)}, and \type{P2}
is \type{(TALL TOM)}, then
 \[\type{(matchp (TALL TOM) P1)}\]
returns a list consisting of the single entry 
 \[\type{(NIL . ((\#:?2 . TOM)))}\]
 The \type{car} of this entry is \NIL, indicating that no variables
in the query need to be bound to match \type{P1}; the \type{cdr} is
\type{((\#:?2 . TOM))} since the variable \type{\#:?2} in \type{P1}
needs to be bound to \type{TOM} to match the query.
 \[\type{(matchp (TALL ?X) P2)}\]
returns \type{((((?X . TOM)) . NIL))}.
 \end{description}

Additional functions are described in Section \ref{s.binding-dag}.

\subsection{Negative subgoals containing free variables}
\label{s.effort}

Consider the following \prolog\ fragment:
\begin{verbatim}
     f(x) :- not(g(x)).
     g(A).
\end{verbatim}
 The interpretation we will assign to the implication
\type{f(x) :- not(g(x))} will be that it is universally quantified
over $x$, and therefore corresponds to the implication
 \[\forall x.[\neg g(x) \supset f(x)]\]
 Given this interpretation, what should the response be if we provide
the system with the query \type{f(x)}?

In light of the negation as failure rule of \prolog\ (which has analogs
in the \mvl\ system), we would expect the query to succeed for all values
of $x$ {\em except\/} $x=A$.

Most \prolog\ systems cannot respond to queries such as this one, since
a nonground negative subgoal is encountered during the proof process.
One approach to these queries, suggested by David Chan, is to do just
what we have suggested in the previous paragraph, computing those binding
lists for which the query fails.

Unfortunately, there are many queries that succeed in general, but
for which it is impractical to compute all of the bindings for which
the query fails.  This may be because there are infinitely many such
failures, or because the time needed to compute them explicitly is
prohibitive.  Situations such as this arise frequently in
applications of declarative programming techniques to planning
problems; here, we would like simply to return an indication that the
query ``succeeds in general.''  In the above example, we might
report that the query had succeeded without binding the variable $x$;
even though $x=A$ is an exception, the report of success is true in
general.

Both approaches are supported by \mvl, and the global variable
\type{*completed-database*} is used to distinguish between them.  If
this variable has the value \T, all exceptions are computed; if it
has the value \NIL, the other approach is taken.  The default value
is \NIL:
 \begin{description}
 \item [\itype{*completed-database*}] If \T, find all exceptions
to default proofs.  If \NIL\ (the default), accept defaults proofs
with possible exceptions as complete.
 \end{description}

\section{Modifying the MVL database}

\subsection{Keyboard input}

The simplest functions to modify the \mvl\ database simply accept
propositions from the user and either add or remove them from the
knowledge base.  The primitive function for adding to the database is
\itype{stash}; the function \itype{unstash} is used for removal.

The function \type{stash} accepts a proposition and up to three
keywords:
 \begin{description}
 \item [\itype{value}]  This keyword is used to supply a specific truth
value to be assigned to the proposition.  If no value is supplied, the
default value is used (this value is supplied when the logic is chosen;
see Section \ref{s.bilattice}).
 \item [\itype{increment-p}] If the proposition in question already
appears in the database, the user has a choice: The given truth value
can be used to replace the one currently appearing in the database,
or it can be combined with the existing value using the combination
function supplied with the logic being used.  The truth value is
replaced if \type{increment-p} is \NIL, and is updated if it is \T\
(the default).
 \item [\apa\index{allow-procedural-attachment}] This keyword,
which defaults to \T, controls whether or not it is possible to
procedurally attach the action of stashing the supplied proposition.
See Section \ref{s.attach}.
 \end{description}

Somewhat more useful than \type{stash} is the function \itype{state}.
This function is identical to \type{stash}, except it translates the
supplied sentence into a form suitable for inference by using the
function \itype{normal-form} described in Section \ref{s.manipulate}.
In addition, \type{state} cannot be procedurally attached, so the
keyword \apa\ is not allowed.

For removing sentences from the database, the functions \type{unstash}
and \itype{unstate} are used.

The function \type{unstash} accepts a proposition to be removed from
the database, and up to three keywords: \type{value},
\itype{decrement-p} and \apa.  The last of these is treated as for
\type{stash}.  In general, if no keywords are used, the effect is to
remove the proposition from the database.

If \type{decrement-p} is \NIL, then \type{unstash} behaves like
\type{stash}, except the default value used for unstashing is not
necessarily the same as the default value used for stashing.

If \type{decrement-p} is \T, then the new truth value assigned to
the proposition is given by
 \begin{equation}
o \cdot \neg v
\label{e.dot-star}
\end{equation}
 where $o$ is the old value and $v$ is the value supplied via the
\type{value} keyword.  The $\cdot$ and $\neg$ operations are as
described in Section \ref{s.bilattice}.

There is one additional function for removing sentences from the
database, \itype{empty}.  This function accepts no arguments, but
does take a single keyword, \type{value}.  The function calls
\type{unstash} for every sentence in the database.

All of the functions described in this section work by calling the
primitive function \itype{new-val}.  This function accepts as
arguments a proposition name and a truth value, and sets the truth
value of the proposition to the given value.  If the truth value is
\type{unknown}, indicating that nothing is known about the
proposition, it is removed from the database and the discrimination
net.

\subsection{Input from files}

It is also possible to enter a database from a disk file.  These files
should in general have the extension \type{.MVL}, and should be
located either on the currently active disk area or on the \mvl\
directory.  The files should consist simply of the sentences to be
added to the database.

In addition to consisting of such sentences, the files may contain the
keyword \type{:value}, followed by the value to be assigned to each
sentence, or the keyword \type{:function}, followed by a function to
be applied to each sentence in order to {\em compute\/} the value to
be assigned to it.  The keyword \type{:empty} empties the database
before loading the file.  Examples of these constructs can be found in
the various \type{.MVL} files included on the distribution tape.

The sentences are loaded with \itype{mvl-load}, which accepts a file
name and an optional keyword \type{empty}.  If the keyword is
supplied, the database is emptied before the file is read; in any
event, \type{state} is then applied to each sentence in the given
file.  The sentences in the file can be removed from the database
with the function \itype{mvl-unload}, which accepts a file name as its
argument.  These functions each call \type{mvl-file}, which accepts
a string and returns a path to the \type{.MVL} file associated to it.

Related to \type{mvl-load} is \itype{mvl-save}, which also accepts a
file name as an argument, but instead of loading a declarative
database from the file, writes the current database {\em to\/} the
file.

\subsection{The editor interface} \index{editor interface}
\subsubsection{Sun EMACS editor}

In addition to the file commands \type{mvl-load} and
\type{mvl-unload}, there is an interface to the {\sc emacs} editor on
Sparcstations.  For a file of mode \mvl, certain keys are bound as
follows (each of these commands also recognizes the
\type{:value} and \type{:function} keywords):

 \begin{description}
 \item [\type{meta-S}] Applies \type{state} to each s-expression in the
current region.
 \item [\type{meta-U}] Applies \type{unstate} to each s-expression in the
current region.
 \item [\itype{meta-s}] Applies \type{state} to each s-expression in the
current defun (i.e., s-expression beginning in column 1).
 \item [\itype{meta-u}] Applies \type{unstate} to each s-expression in the
current defun.
 \end{description}

In order to use these commands, certain modifications will need to be made
to your \type{.emacs} file.  The file \type{dot.emacs} is an example of
what to do.

\subsubsection{Symbolics ZMACS editor}

In addition to the file commands \type{mvl-load} and
\type{mvl-unload}, there is an interface to the {\sc zmacs} editor on
Symbolics 3600 series machines.  For a file of mode \mvl, certain
keys are bound as follows (each of these commands also recognizes the
\type{:value} and \type{:function} keywords):

 \begin{description}
 \item [\itype{super-s}] Applies \type{state} to each s-expression in the
current region.
 \item [\itype{super-u}] Applies \type{unstate} to each s-expression in the
current region.
 \item [\itype{super-f}] Applies \type{fc} to each s-expression in the
current region.  See Section \ref{s.fc}.
 \item [\itype{super-e}] Applies \type{erase} to each s-expression in the
current region.  See Section \ref{s.fc}.
 \end{description}

\subsection{Summary of database modification functions}

\begin{description}
 \item [\itype{stash(p)}] Inserts $p$ into the database.  Allowed
keywords: \type{value}, \type{increment-p} and \apa.
 \item [\itype{state(p)}] Inserts $p$ into the database after
converting it to normal form.  Keywords: \type{value},
\type{increment-p}.
 \item [\itype{unstash(p)}] Removes $p$ from the database.  Keywords:
\type{value}, \type{decrement-p}, \apa.
 \item [\itype{unstate(p)}] Removes $p$ from the database after
converting it to normal form.  Keywords: \type{value},
\type{decrement-p}.
 \item [\itype{empty}] Unstashes every sentence in the database.  Keyword:
\type{value}.
 \item [\itype{new-val(a v)}] Sets the truth value of the sentence
associated with the atom $a$ to $v$.  If $v$ is the truth value
\type{unknown}, the sentence is removed from the database.
 \item [\itype{mvl-file(string)}] Returns the path for the \mvl\ file
associated with the given name.  The file name may include an
extension, one may be supplied as an optional argument, or
\type{.MVL} will be used by default.  There is an optional keyword
\type{no-error} indicating that even if the file cannot be found, the
path should be returned.
 \item [\itype{mvl-load(file)}] Loads the sentences in the given file into
the \mvl\ database.
 \item [\itype{mvl-unload(file)}] Removes the sentences in the
given file from the \mvl\ database.
 \item [\itype{mvl-save(file)}] Writes the \mvl\ database out to the
given file.  This function also accepts the keywords \type{cutoffs}
and \type{props}, which are treated as for \type{contents} (see
Section \ref{s.contents}).  Finally, the keyword \type{if-exists} is
passed to the \type{with-open-file} macro.
 \end{description}

\section{Querying the MVL database}

\subsection{MVL responses to queries}
\label{s.answer}

Because of the fact that a particular fact may be stored in the
database with any of a variety of truth values, \mvl\ responds to a
query by returning not simply a binding list (as most \prolog\
interpreters do, for example), but by returning a binding list {\em
and a truth value}. 

In actuality, what is returned is an instance of the \type{answer}
structure\index{answer structure}.  This structure contains the
following fields:
 \begin{description}
 \item[\type{binding}] Bindings for some of the variables in the query.
 \item[\type{value}] A truth value.  The bound version of the query has
this truth value.
 \end{description}

The definition of this structure leads Common Lisp to define the
access functions \type{answer-binding} and \type{answer-value};
in addition, a function \type{equal-answer} is defined by \mvl\ to
test two answers for equality:
 \begin{description}
 \item[\itype{answer-binding(a)}] Returns the binding field of the
answer \type a.
 \item[\itype{answer-value(a)}] Returns the value field of the answer
\type a.
 \item[\itype{equal-answer(a1 a2)}] Returns \T\ if the two answers
\type{a1} and \type{a2} are equal, and \NIL\ otherwise.  The
functions used are \type{binding-equal} and \type{mvl-eq}, the
functions normally used to test binding lists and truth values for
equality.
 \end{description}

Thus, in response to the query
 \[\type{(flies ?x)}\]
 the response of an answer with binding \type{((?x . Tweety))} and
truth value \type{dt} would indicate that
 \[\type{(flies Tweety)}\] 
 has truth value \type{dt}.  In other words, Tweety's flying is a
conclusion that is true by default.  (See Section \ref{s.j} for a
description of default truth values in the \mvl\ system.)

As a side effect, this allows the user to distinguish a response of
\NIL\ (indicating that the query has failed) from a response in which
the binding is \NIL\ (indicating that the query has succeeded without
binding any variables).

\subsection{Searching for specific sentences}

There are two functions that are most commonly used to retrieve facts
that have been stored in the database, \itype{lookup} and
\itype{lookups}.  Each of these functions accepts as arguments a
proposition to search for and up to four keywords, and returns two
values.

The keywords and their meanings are as follows:
 \begin{description}
 \item [\itype{cutoffs}] This keyword is used to give conditions on
the sentences being considered as possible matches to the query.  Any
sentence that is considered as a match must have an associated truth
value \type{v} such that \type{(test v cutoffs)} evaluates to \T\ (see
Section \ref{s.cutoffs}).  The default value of this keyword is defined
by the bilattice in use (see Section \ref{s.bilattice}).
 \item [\itype{inv}] This keyword defaults to \NIL.  If it is \T,
however, any truth values found should be negated before they are
returned.
 \item [\apa] This is similar to the keyword for \type{stash}.
 \item [\itype{succeed-on-modal-ops}] If \T, then any attempt to look
up a proposition involving a modal operator will immediately succeed;
if \NIL\ (the default), then the result returned will be the result
of applying the modal operator to the truth values of the
propositions that are its arguments.  See Section \ref{s.modal}.
 \end{description}

The two functions \type{lookup} and \type{lookups} are identical
except that \type{lookup} looks for only a single match for the
supplied proposition, while \type{lookups} searches for all such
matches.  Thus while \type{lookup} returns two values, \type{lookups}
returns two {\em lists\/} of values.  Each entry on each of these
lists is analogous to the corresponding value returned by
\type{lookup}.  The values returned by \type{(lookup q)} are as
follows:
 \begin{enumerate}
 \item First, an \type{answer} is returned.  If this answer contains
the binding list \type b and the truth value \type v, then applying
the binding list to the query produces a result with the given truth
value.  In other words, \type{(plug q b)} is an instantiation of a
sentence appearing explicitly in the database with truth value \type
v.
 \item Second, one of the following three values is returned:
 \begin{enumerate}
 \item If the sentence was found by procedural attachment, the atom
\type{attach} is returned.
 \item If the sentence was found by applying a modal operator, the
atom \type{modal} is returned.
 \item Under other conditions, the instantiated version of the
proposition that resulted in the match is returned.
 \end{enumerate}
 \end{enumerate}

As already mentioned, \type{lookups} returns two lists of values.

Here are summaries of \type{lookup} and \type{lookups}:

 \begin{description}
 \item [\itype{lookup(q)}] Searches the database for a sentence
matching the query $q$.  Allowed keywords: \type{cutoffs}, \type{inv},
\apa, \somo.  Returns two values: an answer and the source of the
match.  Returns \NIL\ if the query cannot be found in the database.
 \item [\itype{lookups(q)}] Similar to \type{lookup}, but looks
for {\em all\/} sentences matching the query.  Returns two lists of
values.
 \end{description}

\subsection{Searching the entire database}
\label{s.contents}

In addition to searching for a particular sentence, it is also
possible to display large portions of the entire \mvl\ database.  The
basic functions for doing this are \type{contents}, which displays
the entire database, and \type{prfacts}, which accepts a \lisp\ atom
as an argument and displays all of those facts concerning that atom.

Each of these two functions also accepts a variety of keywords:
 \begin{description}
 \item [\itype{cutoffs}] As for \type{lookup}, it is possible to
supply a test that must be passed by the truth value of any database
sentence before that sentence is displayed.
 \end{description}
 In addition, \type{contents} accepts a keyword \type{props}, which
defaults to \type{(props)} (see Section \ref{s.database}) and is a
list of the sentences to be displayed.

Related to \type{prfacts} is \itype{unstash-facts-with-atom}, which
accepts an atom as an argument and removes from the database any
fact concerning that atom.  \type{Unstash-facts-with-atom} accepts
the same keywords as \type{unstash} does.

Here, then, are the functions described in this section:

\begin{description}
 \item [\itype{contents}] Displays the contents of the database.  Keywords:
\type{cutoffs}, \type{props}.
 \item [\itype{prfacts(atom)}] Displays the database facts concerning the
given atom.  Keyword: \type{cutoffs}.
 \item [\itype{unstash-facts-with-atom(atom)}] Unstashes from the
database any facts concerning the given atom.  Keywords:
\type{value}, \type{decrement-p}, \apa.
 \end{description}

\subsubsection{Command line interface}

\paragraph{Sun}
In addition to the above functions, on Sparcstations, the key
$\langle$meta$\rangle$-* may be typed to the command line processor
running from the editor.  The result is that any information in the \mvl\
database concerning the currently open predicate will be displayed to the
user.

\paragraph{Symbolics}
On Symbolics systems, the key $\langle$super-shift$\rangle$-A may be typed
to the command line processor.  The result is that any information in the
\mvl\ database concerning the currently open predicate will be displayed
to the user.  This feature is the \mvl\ analog to the
$\langle$ctrl-shift$\rangle$-A key, which displays information about the
arguments to the currently open function.

\section{The first-order prover}
\label{s.fo}

As discussed in the technical papers distributed with \mvl, inference
in \mvl\ proceeds using a conventional first-order theorem prover; in
this section, we will describe this prover and the interface to it.
Be warned that the information in this section is unlikely to be used
by other than the most sophisticated users of \mvl; most readers will
want to proceed directly to Section \ref{s.inf}.

In general, there will be a variety of invocations to the first-order
prover active at any particular time.  Each such invocation is
associated with a task that the multivalued prover needs to complete,
and each such task is a \lisp\ object that is an instance of the
\type{consider} structure.\footnote{In fact, each task is an instance
of a structure that {\em includes\/} the \type{consider} structure.
See the file \type{bc.lisp} for further details.}

The multiple invocations of the first-order prover are run in
parallel by \mvl.  To each such invocation is associated an instance
of the structure \type{proof} that is used to record temporary
information about the progress of the proof effort.  We will discuss
first the first-order prover itself, then the construction and
termination of the invocations of it, and finally the structure of
the multivalued proving tasks.

\subsection{Nature}
\label{s.simple}

The prover itself is a simple natural-deduction style complete
theorem prover.  Only inferences against the initial clauses of
database statements are allowed; the result of this is that after
applying \type{state} to a sentence of the form
 \[\type{(<= CONCLUSION PREMISE-1 ... PREMISE-N)}\]
 \type{PREMISE-1} will always be tried first when the prover attempts
to prove \type{CONCLUSION}.  In a similar way, stating the
forward-chaining rule
 \[\type{(=> PREMISE-1 ... PREMISE-N CONCLUSION)}\]
 will result in a situation where only \type{PREMISE-1} will trigger
forward chaining.

Control of the proof search is handled by a function that evaluates
the ``merit'' of each of the ongoing proof attempts, and then
attempts to select the best one.  The evaluation function is stored
in the value of the global constant \itype{*node-evaluation-fn*},
which defaults to \itype{best-first-fn}, a function that attempts to
determine the value of a proof by looking at the truth value being
accumulated.  An alternative to \type{best-first-fn} is the function
\itype{depth-fn}, which evaluates a node simply by its depth in the
proof tree.

Three global parameters are used to control the proof search:
 \begin{description}
 \item [\itype{*final-limit*}] is the limit point beyond which the
search is abandoned, so that any node that evaluates to a value worse
than this will be pruned.  The default value of this parameter is
\NIL, indicating that no nodes are to be pruned.
 \item [\itype{*initial-limit*}] is used to indicate the type of search.
If \NIL\ (the default), then the values assigned to the nodes are not
used to order the search, and depth-first search is used.
 \item [\itype{*limit-delta*}] is used if \type{*initial-limit*} is
non-\NIL.  In this case, the nodes are ordered by increasing
 \[\frac{\type{value(node)}-\type{*initial-limit*}}{\type{*limit-delta*}}\]
 A typical value for *initial-limit* is 0, and for *limit-delta* is
1.  The value of this parameter is irrelevant if
\type{*initial-limit*} is \NIL.
 \end{description}

These three parameters may be reset by the user:
 \begin{description}
 \item [\itype{set-search-parameters}] This function is the primitive
used to modify the parameters used to control the first-order theorem
prover's search.  It accepts as keywords \type{initial-limit},
\type{limit-delta}, and \type{final-limit}.
 \item [\itype{activate-breadth-first}] Activate breadth-first search
by setting \type{*node-evaluation-fn*} to \type{depth-fn}.  Allowed
keywords: \type{initial-limit}, \type{limit-delta}, and
\type{final-limit}.
 \item [\itype{activate-depth-first}] Activate depth-first search by
setting \type{*initial-limit*} to \NIL.  Allowed keyword:
\type{final-limit}.
 \item [\itype{activate-best-first}] Activate best-first search by
setting \type{*node-evaluation-fn*} to \type{best-first-fn}.  Allowed
keywords: \type{initial-limit}, \type{limit-delta}, and
\type{final-limit}.
 \end{description}

\subsection{Invoking the prover}

The prover is invoked using the function \type{init-prover}.  This
function accepts as arguments a proposition and a pointer to the task
responsible for this proof attempt.

When \type{init-prover} is invoked, an instance of \itype{proof}
is created and returned that contains auxiliary information about the
proof attempt.  Here are the fields in the structure:
 \begin{description}
 \item [\itype{active-nodes}] A list of the nodes that are still
under active consideration.
 \item [\itype{answers}] A list of answers that have been found.
 \item [\itype{indexp} and \itype{indexn}] These entries are used to
index the various subtasks generated by the prover.
 \item [\itype{initial-limit}, \itype{final-limit}, \itype{limit-delta}
and \itype{evaluation-fn}] These entries are the values of the search
control parameters used in this proof attempt.
 \item [\itype{task}] A pointer back to the task responsible
for this invocation of the first-order prover.
 \item [\itype{vars}] The variables that appear in the proposition
being proven.
 \end{description}

Each node is an instance of the structure \itype{goal}, which is
described in the file \type{first-order.lisp}.

To continue a proof process, the function \type{cont-prover} is used.
This function accepts as an argument an instance of \type{proof} that
contains the auxiliary information for the proof process being
resumed.  The value returned by \type{cont-prover} is a structure
describing the results of the proof process itself;
\type{cont-prover} returns as soon as the prover finds a single
solution to its query.  The value returned is an instance of the
\type{bcanswer} structure\index{\type{bcanswer} structure}; the
fields in this structure are as follows:
 \begin{description}
 \item[\type{answer}] This field contains a binding list and truth
value that the system has determined can be assigned to the given
instantiation of the original query.
 \item[\type{just}] This field contains a list of those database
sentences that were used in the proof.
 \end{description}

If the proof attempt fails, \NIL\ is returned.

Here is a summary of the functions described in this section:

\begin{description}
 \item [\itype{init-prover(p effort)}] This function
initializes and returns an instance of \type{proof} containing
auxiliary information about the proof process.
 \item [\itype{cont-prover(proof)}] Resume the task associated to the
auxiliary information in the given instance of the \type{proof}
structure.  This function returns an instance of \type{bcanswer}
corresponding to the result obtained by the prover, and returns as
soon as a solution to the original query is found.
 \end{description}

\subsection{The multivalued proof tasks}
\label{s.consider}

The individual proof attempts generated by a multivalued query are
instances of structures that include the structure
\type{consider}\index{consider structure}.  The exact details of this
and related structures are unlikely to be of much interest to a user
of the \mvl\ system, but we include them here in the interests of
completeness.  These are the fields in the
\type{consider} structure:
 \begin{description}
 \item [\itype{prop}] The proposition being investigated by this task.
 \item [\itype{parent}] The task for which the current task is a
child.
 \item [\itype{children}] Children of the current task.
 \item [\itype{truth-value}] The current best estimate of the truth value
that should be assigned to the proposition being considered.
 \item [\itype{relevant}] A list of answers that, if returned by this
task, are known to be relevant to the original query.
 \item [\itype{irrelevant}] A list of answers know {\em not\/} to be
relevant to the original query.
 \item [\itype{paths}] The \mvl\ system goes to a great deal of effort
to ensure that it does not consider proofs that will not affect the
eventual response to a user's query; intermediate information used in
this computation is stored in this field.  Further details can be found
in the file \type{bc.lisp} and the files to which it points.
 \end{description}

When the prover must select a task to work on next, it does so via the
following function:
 \begin{description}
 \item [\itype{choose-task}] Selects a task to work on next.  The
current implementation of \mvl\ simply selects any active task (see
Section \ref{s.bc}).
 \end{description}

\section{Inference}
\label{s.inf}

\subsection{Backward chaining}
\label{s.bc}

The information in Section \ref{s.fo} is more technical than most users
of \mvl\ are likely to need.  What is of much greater interest are those
database query tools that are like \type{lookup} and \type{lookups} but,
instead of simply looking for a fact in the database, attempt to derive
it as a consequence of other database information.  These functions use
those of the previous section as described in the accompanying technical
papers.

The functions accept three sorts of cutoff information, instead of the
simple information provided to \type{lookup} and \type{lookups}.  The
first, provided via the \type{cutoffs} keyword, is used to construct a
test to decide whether an answer computed by the prover is good enough
to be returned when the proof effort is complete.  The second,
provided via the \type{succeeds} keyword, indicates that an answer
should be returned as soon as it can be shown that the final answer
will pass the given test.  The final cutoff, supplied via the
\type{terminate} keyword, indicates that the answer should be returned
as soon as there is any indication that the final answer {\em might\/}
pass the given test.

As an example, suppose that we were trying to find a plan to achieve some
goal.  We might well be prepared to terminate our search as soon as we found
a plan that was {\em guaranteed\/} to achieve the goal; if we completed our
search {\em without\/} finding such a plan, we might then want to return
all plans that we {\em expected\/} would achieve the goal.  It is to
facilitate such possibilities that the backward-chaining functions expect
three sets of cutoff information.  The default value for \type{cutoffs}
is the standard one associated to the active bilattice (see Section
\ref{s.bilattice}); the default choice for \type{succeeds} corresponds
to complete information for the singular versions of the queries and
guarantees failure for the plural versions (since an answer might be
about to be discovered for a different instantiation of the query).
The default value for \type{terminate} always fails.

Here, then are the various backward-chaining functions.  They all
return two values: an answer (or list of answers for the plural
versions \type{bcs} and \type{consider}) and an instance of the
\type{analysis} structure describing tasks remaining to be considered
by the prover.  In addition, they all accept an optional
\itype{analysis} keyword that, if supplied, should be an instance of
the \type{analysis} structure that is to be continued instead of a
new one being created.  This is to allow the user to continue proof
efforts that have returned for some reason.
 \begin{description}
 \item [\itype{bc(q)}] Attempts to prove the query $q$ from information
in the database, looking only for a single answer.
 \item [\itype{bcs(q)}] Attempts to prove the query $q$ from information
in the database, looking for all possible answers.
 \end{description}

All of these functions work by maintaining an instance of the
\type{analysis} structure\index{\type{analysis} structure} that
includes a variety of information about the status of the current
proof attempt.  Here are the fields in this structure:
 \begin{description}
 \item [\type{prop}] The proposition being considered by the prover.
 \item [\type{cutoffs}] The cutoff information that must be satisfied
by any result before it is to be returned.
 \item [\type{succeeds}] The cutoff information that will allow the
prover to terminate early if an answer is guaranteed to pass the test.
 \item [\type{terminate}] The cutoff information that will allow the
prover to terminate early if an answer might pass the test.
 \item [\type{active-tasks}] A list of the proof tasks (see Section
\ref{s.consider}) that still need to be investigated.
 \item [\type{main-task}] The task that was created upon the 
original invocation of the prover.
 \item [\type{returned-answers}] A list of answers that have already
been returned to the user; this is to keep \mvl\ from compulsively
re-returning these answers if the \type{analysis} is reinvoked.
 \end{description}

The top-level functions using this structure are as follows:
 \begin{description}
 \item [\itype{init-analysis(prop)}] The \mvl\ primitive upon which
the proving functions are based, this function accepts a proposition
and creates and returns an instance of the \type{analysis} structure.
 \item [\itype{cont-analysis(analysis)}] This function continues the
investigation of a specific analysis, and can be used to resume a
proof effort that has been curtailed for some reason.  It returns two
values: an instance of the \type{answer} structure giving the results
of the proof attempts, and a modified version of the \type{analysis}
structure that can be used to continue the proof process if desired.
 \end{description}

The functions \type{init-analysis} and the user functions such as
\type{bc} all accept the following keywords:
 \begin{description}
 \item [\type{cutoffs}] Cutoff test as described above.
 \item [\type{succeeds}] Success test as described above.
 \item [\type{terminate}] Termination test as described above.
 \item [\type{initial-value}] The truth value assigned to the query in
the absence of database information.
 \item [\type{analysis}]  If supplied, use this analysis instead of
creating a new one.
 \end{description}

\subsection{Forward chaining}
\label{s.fc}

Forward chaining in \mvl\ is not as simple as it is in many other
logic programming languages because of the \nm\ nature of multivalued
inference.  It follows from this that inserting a new sentence into
the database may cause a previous conclusion to be retracted, and
this may in turn cause other sentences to become unsupported.

There are two functions that use this information to modify the
database.  The more common is called \itype{fc} and accepts as an
argument a proposition to be added to the database.  The proposition
is added and provided that it was not already in the database when the
forward chainer was invoked, its consequences are also added to the
database and forward chained upon.  The \type{fc} function accepts
the following keywords:
 \begin{description}
 \item [\itype{value}] This keyword is similar to that used by
\type{stash} and \type{state}.
 \item [\itype{increment-p}] This keyword is also similar to that used
by \type{stash} and \type{state}.  The default value is \NIL.
 \item [\itype{fc-increment-p}] When the information content of some
sentence in the database has {\em fallen}, should this drop be
propagated to other sentences that appear to depend on the sentence
in question?  The default value is \T.
 \item [\itype{cutoffs}] Only forward chain if the sentence being
inserted into the database has a truth value that passes the test
corresponding to \type{cutoffs}.  (See Section \ref{s.cutoffs}.)  The
default value used is the one associated with the active bilattice.
 \item [\itype{include-modal}] If \T\ (the default), then the forward
chainer also looks for rules that fire when the truth value of a
modal expression involving $p$ changes (where $p$ is the sentence
being forward chained upon).  See Section \ref{s.modal}.  The reason
this keyword is included is that searching for rules of this form is
computationally expensive.
 \end{description}

Slightly less common is \itype{erase}, the function used to remove a
proposition from the database.  This function accepts a database
sentence as an argument, and also the keywords \itype{value},
\itype{decrement-p} and \itype{cutoffs}, which are similar to the
corresponding keywords for \type{fc}.

Summarizing:
 \begin{description}
 \item [\itype{fc(p)}] Inserts the sentence $p$ into the database and
forward chains.  Allowed keywords: \type{value}, \type{increment-p},
\type{fc-increment-p}, \type{cutoffs} and \type{include-modal}.
 \item [\itype{erase(p)}] Removes $p$ from the database and forward
chains.  Allowed keywords: \type{value}, \type{decrement-p} and
\type{cutoffs}.
 \end{description}

\subsection{Debugging tools}

Debugging facilities are limited in \mvl.  The facilities that exist
describe the prover's activities during backward-chaining,
forward-chaining, or while the first-order theorem prover is looking
for a conventional proof of some proposition.

\subsubsection{Backward chaining}

Diagnostics are displayed during backward chaining using the two
global variables \itype{*bc-trace*} and \at\index{*anytime-trace*}.

Whenever a task is encountered that involves a proposition that is an
instance of some element of \type{*bc-trace*}, information concerning
the task will be displayed to the user.  The variable \at\ is
similar, but only reports when an {\em inference\/} task is
encountered.  An inference task is one that is trying to prove
something, as opposed to one analyzing a modal operator; this
distinction is described in greater detail in the file {\type
bc.lisp}.  In practice, the distinction between \type{*bc-trace*} and
\at\ is that the former should be used for debugging purposes, while
the latter should be used if you are interested in watching the
prover approximate and then refine its eventual answer to some
declarative query.

Although the user can manipulate \type{*bc-trace*} and \at\ directly,
there are four functions written to help do so:
 \begin{description}
 \item [\itype{bc-trace(pat)}] Include \type{pat} on
\type{*bc-trace*}, so that any instance of \type{pat} will be
traced.
 \item [\itype{bc-untrace(pat)}] Remove any instance of \type{pat}
from \type{*bc-trace*}.
 \item [\itype{anytime-trace(pat)}] Include \type{pat} on \at, so
that any instance of \type{pat} will be traced.
 \item [\itype{anytime-untrace(pat)}] Remove any instance of
\type{pat} from \at.
 \end{description}
 For all of these functions, the argument is optional.  If omitted,
\type{(?*)} (of which everything is an instance) is used.

\subsubsection{Forward chaining}

Tracing the activity of the forward chainer is essentially identical
to tracing the behavior of the backward chainer; the global variable
\type{*fc-trace*} is used and is manipulated using the functions:
 \begin{description}
 \item [\itype{fc-trace(pat)}] Include \type{pat} on
\type{*fc-trace*}, so that any instance of \type{pat} will be
traced.
 \item [\itype{fc-untrace(pat)}] Remove any instance of \type{pat}
from \type{*fc-trace*}.
 \end{description}
 Once again, the argument is optional.

\subsubsection{The first-order prover}

Tracing the activity of the first-order theorem prover itself is
somewhat different.  The facility for doing this is similar to the
Common Lisp \type{trace} facility, but instead of reporting when a
function is invoked, it reports when an attempt is made to prove an
instance of a particular predicate, and also when such an instance is
found successfully in the database.

The functions used to control this are \itype{bc-watch} and
\itype{bc-unwatch}.  They are similar to \type{trace} and
\type{untrace}:

 \begin{description}
 \item [\itype{bc-watch(preds)}] Accepts as arguments any number of
predicates that should be watched as the prover proceeds.  Each
predicate can also be a list of the form \type{(predicate depth)}
where \type{depth} is the number of successive ancestors to be
displayed when the prover encounters an instance of the predicate.
(This can be used to understand {\em why\/} a particular clause is
relevant to the proof being attempted.)  \type{bc-watch} returns a
list of the predicates currently being {\sc watch}ed.
 \item [\itype{bc-unwatch(preds)}] Removes the given predicates from
the list of {\sc watch}ed predicates.  If no arguments are supplied,
all predicates are un{\sc watch}ed.
 \end{description}

There are two additional features involved in tracing the activities
of the first-order prover.  The first is used whenever the global
parameter \itype{*save-nodes*} is \T\ (which it is by default).  When
the first-order prover displays information regarding a particular
proof node, the node is labelled with a number.  Subsequent nodes
will then use this label if they need to refer to the node in
question (e.g., as a parent).

If \type{*save-nodes*} is \T, then \type{*break-node*} can be set
to an integer; when a node is displayed that is labelled with this
integer, a \type{(break)} will be executed.  This is to allow the
user to interrupt the first-order prover if he needs to examine the
proof stack or some other \mvl\ internal.

 \begin{description}
 \item [\itype{*save-nodes*}] If \T\ (the default), nodes are
labelled with consecutive integers when they are displayed by the
first-order prover.
 \item [\itype{*break-node*}] The number of a node at which the
proof process should be interrupted.
 \end{description}

\section{Procedural attachment}
\label{s.attach}

In some cases, it may be desirable to store a fact in the database
{\em without\/} using the functions supplied in \mvl.  Suppose, for
example, that we have an $n+1$-ary relation $R$ that corresponds to
an $n$-ary function $f$, so that $R(x_1,\dots,x_n,y)$ holds just in
case $f(x_1,\dots,x_n)=y$.  Now, if we assert $R(c_1,\dots,c_n,d)$
into the database, where $d$ and the $c_i$ are constants, we should
probably also {\em remove\/} any sentences of the form
$R(c_1,\dots,c_n,d')$, where $d'\neq d$.  This sort of a situation
might arise if we were monitoring the temperature in a room, and
represented the observed temperature using a database sentence such
as
 \begin{equation}
\type{(TEMPERATURE ROOM 72)}
\label{e.temp}
\end{equation}

We handle this by informing \mvl\ that to stash a sentence of the
form \type{(TEMPERATURE ...)}, it should use not the usual function
\type{stash}, but some other user-defined function, say
\type{unique-stash}.  This new function would first \type{unstash}
any sentence that agreed with the given one except for the value of
the final term, and would then invoke \type{stash} on the original
sentence, but with \apa\ set to \NIL\ so that \type{unique-stash} was
not invoked a second time.

To modify the function used by \mvl\ to stash sentences such as this,
the user should simply place the desired function (in this case,
\type{unique-stash}) on the \type{stash} property of the relation in
question (in this case, \type{TEMPERATURE}).\footnote{As described in
Section \ref{s.invocation}, the \type{stash} property of a relation
can also be an association list, if multiple invocations of the
system are running in parallel.} The system will then automatically
invoke \type{unique-stash} instead of \type{stash} when a fact such
as (\ref{e.temp}) is asserted.

The functions within \mvl\ that can be modified in this fashion are
\type{stash}, \type{unstash}, \type{lookup} and \type{lookups}.  In
addition, the procedural attachments for \type{lookup} and
\type{lookups} are linked, so that if \type{lookup} has a procedural
attachment and \type{lookups} does not, a call to \type{lookups} will
invoke the procedural attachment for \type{lookup} and then pluralize
the result.  Similar remarks hold if \type{lookups} and not
\type{lookup} has a procedural attachment.

The procedural attachment is effected using the following functions:
 \begin{description}
 \item [\itype{invoke(fn p)}] This function invokes \type{fn} on the
sentence $p$.  In addition, any keyword arguments are passed to
\type{fn}, but with \type{\&allow-other-keys} so that the keywords
will be ignored if \type{fn} is not expecting them.
 \item [\itype{lookup-call(fn p)}] This function is similar to
\type{invoke}, but modifies the values returned to make them agree
with those expected to be returned by \type{lookup}.  Unless the
function \type{fn} supplies them, the truth value is assumed to be
\type{true}.  The second value is \type{attach}.
 \item [\itype{lookups-call(fn p)}] The plural version of
\type{lookup-call}.
 \end{description}

\subsection{Facilities to create procedural attachments}

Some facilities exist to create procedural attachments within the
\mvl\ system.  The following two macros are used for this purpose:
 \begin{description}
 \item [\itype{lisp-predicate(pred)}] This macro is used to indicate
that the given predicate can be evaluated by a \lisp\ call.  It will
also cause \mvl\ to signal an error if an attempt is made to stash
or unstash an instance of the function.
 \item [\itype{lisp-defined(fn)}] This macro is used to indicate that
the given function can be evaluated by a \lisp\ call.  It will also
cause \mvl\ to signal an error if an attempt is made to stash or
unstash an instance of the function.  \type{Lisp-defined} also
accepts an optional argument, which is an inverter function.  The
overall result of attempting to lookup an instance of the function
can be described as follows:
 \begin{enumerate}
 \item If none of the arguments to the function (i.e., all but the
last of the arguments of the relation) is a variable, then the function
is called and the result is unified with the final argument of the
relation.
 \item If the inverter function has been supplied and is equal to \T,
then the function is called and the result is unified with the final
argument of the relation even if some of the arguments are variables.
 \item If some of the arguments are variables but the result is not a
variable, then the inverter function is invoked if one has been
supplied.  The arguments passed to the inverter function are the
result (i.e., the last argument of the relation), and a list of all
of the arguments (i.e., all but the last argument of the relation).
If there is no inverter function, then \NIL\ is returned.  In
addition, if \wbli\ \index{*warn-on-bad-lisp-invocations*} is not
\NIL, the user is warned.  (The parameter \wbli\ defaults to \NIL.)
 \item If both some of the arguments and the result are variables, then
\NIL\ is returned.  In addition, if \wbli\ is not \NIL, the user is
warned.
 \end{enumerate}
 \end{description}

The current version of \mvl\ invokes \type{lisp-predicate} on the
following \lisp\ predicates: \type{numberp}, $<$, $>$, \type{equal},
\type{null}, \type{atom}, \type{listp}, \type{varp}, \type{varp?},
\type{varp$*$}, \type{groundp}, and $=$.  The procedural attachment
to $=$ attempts to unify its arguments.

The macro \type{lisp-defined} is invoked on the following:
\type{car}, \type{cdr}, \type{cadr}, \type{cons}, \type{list},
\type{list*}, \type{revappend}, \type{append}, \type{reverse},
\type{length}, \type{assoc}, \type{delete}, \type{+}, \type{*},
\type{-}, \type{/}, \type{floor}, \type{ceiling}, \type{truncate},
\type{round}, \type{float}, \type{mod}, \type{rem}, \type{max},
\type{min}, \type{sin}, \type{cos}, \type{atan}, \type{sqrt}, 
and \type{eval}.

\subsection{Other procedurally attached relations}

In addition to the \lisp\ functions and predicates described in the
previous section, the following relations have been defined and
are procedurally attached to suitable \lisp\ functions:
 \begin{description}
 \item [\itype{(bagof}\type{ ?x ?p ?b)}] This relation holds for $x$,
$p$ and $b$ if and only if $b$ is a list of all $x$ satisfying $p$.
 \item [\itype{(done exp-1 ... exp-n)}] This predicate, which is
useful for stashing only, causes \lisp\ to evaluate each of the subsequent
expressions.
 \item [\itype{(value symbol val)}] Stashing this predicate sets the
given symbol to the value \type{val}.  Unstashing it makes the symbol
unbound if its current value unifies with \type{val}.
 \item [\itype{(property symbol value indicator)}] Stashing this
predicate sets the symbol's \type{indicator} property to \type{value}
for the current invocation of \mvl.  Unstashing it removes the
\type{indicator} property for the current invocation of \mvl\ if the
previous value of this property unifies with \type{value}.  (See the
definitions of \type{get-mvl} and \type{rem-mvl} in Section
\ref{s.invocation}.)
 \item [\itype{(known p)}]  Stashing this predicate simply stashes $p$;
unstashing it causes an error.  Looking it up looks up $p$.
 \item [\itype{(unknown p)}]  Stashing this predicate unstashes $p$;
unstashing it causes an error.  Looking it up fails if $p$
can be found in the database, and succeeds otherwise.
 \item [\itype{(member item list)}] Tries to unify \type{item} with
the various elements in \type{list}.
 \end{description}

\section{Testing the MVL system}
\label{s.foln}

\subsection{As a diagnostic tool}

There are two functions available to help you to test both \mvl\ and
any applications you write using it.  The first, \type{test-mvl}, is
one you should already have run; if \mvl\ has been installed successfully,
it should return 0.

\type{Test-mvl} works by invoking the generic testing function
\type{mvl-test-file} on the file \type{mvl.test} in the \mvl\
directory.\footnote{\type{Test-mvl} works in its own invocation of
\mvl, so that the existing database is not modified.  This is not true
of \type{mvl-test-file}, however.} This file also contains
documentation regarding the details of the various tests being
executed.

In general, \type{mvl-test-file} accepts as input a file name, and
looks for a file with an extension \type{.TEST} on either the \mvl\
or the current directory.  It then works through the file by
evaluating the first s-expression it contains, and comparing the
value returned to the second s-expression.  If these two
s-expressions match, the next s-expression is evaluated, and so on.
If there is a mismatch, \type{mvl-test-file} reports the error;
eventually, it returns the total number of errors found in the entire
file.  When doing the comparison, \type{**} acts as a wildcard.
Thus, if the second s-expression is simply \type{**}, it will always
match; if it is \type{(** ** **)}, it will match whenever the first
expression returns a list of three values, and so on.

\begin{description}
 \item [\itype{test-mvl}] Tests the \mvl\ system by invoking the function
\type{mvl-test-file} on the file \type{MVL.TEST}.
 \item [\itype{mvl-test-file(string)}] Can be used to evaluate
various forms by the \mvl\ system and compare the results to
predefined values.  Both the forms and the expected answers are in
the \type{.TEST} file whose name is supplied to the function.
Returns the number of errors found.
 \end{description}

\subsection{As a tutorial}

The information used by the function \type{test-mvl} can also be used
as a tutorial in the \mvl\ system.  Take a look at the file
\type{mvl.test} that this function uses; you will see that the first
two lines are
 \begin{verbatim}
     (notime (mvl-test-file "bc" "backward chainer")
     0
\end{verbatim}
 This means that the function \type{mvl-test-file} will be invoked on
the file \type{bc.test}, and that this function should return 0
(i.e., 0 errors).  So let's have a look at \type{bc.test}:
 \begin{verbatim}
     (load-logic #\f)                    ;first-order logic
     **
\end{verbatim}
 Here, the system loads the bilattice corresponding to first-order
logic and is not interested in the value returned by the loading
operation.  Next, a declarative database corresponding to a full
adder is loaded (with the timer turned off).  The sentences in this
database can be found in the file \type{f1.mvl} and are also
described in Section 15 and Appendix B.2 of the paper ``Multivalued
logics: A uniform approach to inference in artificial intelligence.''

Next, we see in \type{mvl.test} the two lines
 \begin{verbatim}
     (get-bdg '?x (answer-binding (bc '(val (out 1 f1) 1 ?x))))
     OFF                   ;the first output (the sum line) should be off
\end{verbatim}
 We invoke the backward chainer on the sentence
 \[\type{(val (out 1 f1) 1 ?x))}\]
 to see if we can prove that the first output of the full adder takes
some value \type{?x}.  Then we get the binding of this variable to
see what the value actually is -- \type{OFF}, in this case.  The
remainder of the various test files can be analyzed similarly; many
of the examples in these files can also be found in the various
technical papers included in the \mvl\ distribution package.

\section{System support}
\label{s.support}

\subsection{Functions that compile and load the MVL system}

The following three functions can be used to manipulate the \mvl\
system as a whole:
 \begin{description}
 \item [\itype{load-mvl}] Loads the \mvl\ system.
 \item [\itype{compile-mvl}] Compiles and loads the \mvl\ system.  On
an Allegro system, the keyword \type{recompile} forces a recompile of
all of the files and the keyword \type{reload} is ignored (since the
files will all be reloaded anyway).  On a Symbolics system, the
keyword \type{reload} forces all of the files to be reloaded and the
keyword \type{recompile} forces them all to be recompiled.
 \item [\itype{release-mvl}] (Symbolics only)  Compiles and loads the
\mvl\ system, and declares the new version to be \type{released}.  The
keyword \type{no-recompile} can be used to skip the compilation step.
 \end{description}

\subsection{The MVL package}
\label{s.package}

In order to avoid name conflicts with existing code, the \mvl\ system
resides in its own package, \type{MVL}\index{MVL package}.  All
of the symbols and functions appearing in the appendix are exported
from this package and can therefore be accessed by invoking the
function \type{use-package(mvl)}.  When a new bilattice is loaded,
the modal operators it defines are also exported. 

To use the \mvl\ system, invoke the function \type{(use-package 'mvl)}
to make these symbols accessible.\footnote{As an alternative,
\type{(in-package 'mvl)} wil put you in the MVL package itself.}

\subsection{Running several copies of MVL in parallel}
\label{s.invocation}

Some users may wish to run several copies of the \mvl\ system
simultaneously, perhaps associating them with distinct {\sc lisp}
processes or with individual agents.  (The use of multiple copies of
\mvl\ to represent distinct components of a single agent's knowledge
goes against the grain of the multivalued approach and is not
recommended.)  In order to support this facility, the following
functions and macros have been defined:
 \begin{description}
 \item [\itype{existing-mvl-invocation}] This function returns a
single value that is an instance of the structure \mi.  It is useful
primarily for saving the current copy of \mvl\ so that it can be
restored later.
 \item [\itype{create-mvl-invocation}] This function sets up a new
copy of \mvl.  It accepts an optional argument that is an instance of
\mi; this can be used to reset the state of \mvl\ to a previous one.
If the optional argument is omitted, default values are used.  The
value returned is the \mvl\ invocation that was active before the
function was called, so that another call to
\type{create-mvl-invocation} can be used to restore the original
environment.
 \item [\itype{with-mvl-invocation}] This is a macro similar to the
Common Lisp macro \type{with-open-file}.  It accepts an optional
argument that is an instance of \mi\ and is followed by a body of
forms (presumably including a \type{load-logic}) that are evaluated
in a locally-bound copy of \mvl.
 \end{description}

These functions work by locally or globally rebinding the values of
the various parameters used by the \mvl\ system
(\type{*active-bilattice*}, \type{true} and so on).  The only
difficulty with this is that procedural attachments for a relation
$R$ use the property list of $R$, and this cannot be rebound.

In order to cater to this difficulty, there is an additional global
variable \itype{*mvl-invocation*} that labels the currently active
copy of \mvl.  To determine what procedure is attached to the
\type{stash} operation for $R$, the \type{stash} property of $R$ is
examined.  There are now the following possibilities:
 \begin{enumerate}
 \item If the value is \NIL, \type{stash} is not procedurally attached,
so the usual \mvl\ function is called.
 \item If the value is an atom, then it is assumed that this atom is
the name of a function that is the procedural attachment for \type{stash}
independent of the copy of \mvl\ that is currently active.
 \item If the value is a list, it is treated as an association list
and a cons whose car is \type{*mvl-invocation*} is searched
for.
 \begin{enumerate}
 \item If such a cons is found, its cdr is used as the procedural
attachment to \type{stash}. 
 \item If no such cons is found, but there is a cons whose car is \T,
then the cdr of that cons is used.
 \item If no cons has either \type{*mvl-invocation*} or \T\ as its
car, then \type{stash} is assumed to be unattached.
 \end{enumerate}
 \end{enumerate}
 The other functions described in Section \ref{s.attach} are handled
similarly.

The association lists described above are manipulated using the following
functions:
 \begin{description}
 \item [\itype{get-mvl (symbol indicator)}]  Like the Common Lisp
\type{get}, but if the value is a list, treats it as an a-list and
searches for an item whose car matches \type{*mvl-invocation*}.
The value taken can be changed with \type{setf}.
 \item [\itype{rem-mvl (symbol indicator)}] Like the Common Lisp
\type{remprop}, but if the value is a list, treats it as an a-list
and removes the item whose car matches \type{*mvl-invocation*}.
 \end{description}

\section{Selecting a bilattice}
\label{s.bilattice}

\subsection{Bilattice structure}

Of course, most of the power of the \mvl\ system comes not from the
features that were described in Sections \ref{s.fol1}--\ref{s.support}
of this manual, but from the fact that the user can change
representation language conveniently within the system.  He does this
by changing the {\em bilattice\/} that he is using.

A bilattice is the mathematical structure from which the truth values
used by the system are selected, and the formal ideas underlying this
approach are discussed at length in the accompanying technical
papers.  In the \mvl\ system, a bilattice is described as an instance
of the \type{bilattice} structure.  This structure has 27 fields; of
these, 17 are used to provide the formal information required by the
system, and the others are used to simplify the user interface.  Here
are 13 of the fields used to describe the bilattice itself:

\begin{description}
 \item [\itype{true}] A \lisp\ object that is the truth value assigned
to a sentence that is unconditionally true, such as the tautology
\type{(OR (P) (NOT (P)))}.
 \item [\itype{false}] A \lisp\ object that is the truth value assigned
to a sentence that is unconditionally false, such as the negation of a
tautology.
 \item [\itype{unknown}] A \lisp\ object that is the truth value assigned
to a sentence about which nothing is known.
 \item [\itype{bottom}] A \lisp\ object that is the truth value assigned
to a sentence that is both true and false.  If any sentence is assigned
this truth value (or can be shown to have it), it indicates the presence
of a contradiction in the database.
 \item [\itype{eq}] A \lisp\ function that takes two truth values as
arguments and returns \T\ if they are equal and \NIL\ if they are
not.  This function is needed because it may be possible to represent
a single element of the bilattice as a \lisp\ object in a variety of
fashions.
 \item [\itype{and}] A \lisp\ function that takes two truth values as
arguments and returns the truth value corresponding to their
conjunction.  If the two sentences $p$ and $q$ are labelled with
truth values $v_1$ and $v_2$ respectively, then the conjunction of
the truth values is the truth value that can be assigned to $p\wedge q$.
 \item [\itype{or}] A \lisp\ function that takes two truth values as
arguments and returns the truth value corresponding to their
disjunction.
 \item [\itype{not}] A \lisp\ function that takes a truth value as its
argument and returns the truth value corresponding to its negation.
 \item [\itype{plus}] A \lisp\ function that takes two truth values as
arguments and returns the truth value corresponding to their ``sum.''
If there is reason to believe that a sentence $p$ should be labelled
with truth value $v_1$, and independent reason to believe that it
should be labelled with truth value $v_2$, then the actual label is
given by the sum $v_1+v_2$.  This operation therefore corresponds to
combination of evidence.
 \item [\itype{dot}] A \lisp\ function that is the dual of the
\type{plus} operation.  It takes two truth values $v_1$ and $v_2$ as
arguments and returns the greatest truth value that could
subsequently be strengthened to a conclusion of either $v_1$ or
$v_2$.  This is a ``consensus'' operation.
 \item [\itype{kle}] A \lisp\ function that takes two truth values $x$
and $y$ and determines whether $x+y=y$.  Although this can also
be done using the bilattice $+$ and $=$ operations, this field can be
provided to give a more efficient implementation of this operation.
This field is optional.
 \item [\itype{tle}] A \lisp\ function that takes two truth values $x$
and $y$ and determines whether $x\vee y=y$.  This field is also
optional.
 \item [\itype{modal-ops}] A list of the modal operators associated
with the given bilattice.  See Section \ref{s.modal}.
 \end{description}

Three of the remaining formal slots are concerned with the incorporation
of binding information into truth values.  In some cases, when a
proposition is retrieved from the database, it is only an instance of
the database proposition that is used, and the truth value may have to
be modified to reflect that:
 \begin{description}
 \item [\itype{plug}]  This \lisp\ function accepts as arguments
a truth value and a binding list, and returns a potentially modified
truth value.  The default value is a function that leaves the truth
value unchanged.
 \end{description}

There is also a function that accepts a truth value and returns the
variables appearing within it:
 \begin{description}
 \item [\itype{vars}]  This \lisp\ function accepts a truth value
as its argument and returns a list of the variables appearing in it.
The default value is a function that returns \NIL\ independent of its
argument.
 \end{description}

The function \itype{absorb} is needed for subtle technical reasons
when working with bilattices with non-\NIL\ plugging functions;
further information on this can be found in the file \type{load.lisp}.
All three of these functions are \NIL\ if there is no interaction
between the bilattice truth values and the variables that they
contain.

The final formal field has to do with search control:
 \begin{description}
 \item [\itype{simplicity}] This \lisp\ function accepts a truth
value (which will frequently be the value obtained by conjoining the
premises in a particular proof attempt) and returns a figure
indicating the ``merit'' that should be assigned to the proof attempt
in question.  The value returned should be a nonempty list of
integers and is viewed lexicographically, with high values being
preferred to low ones.  The default value is a function that returns
the list \type{(0)} independent of the truth value supplied.
 \end{description}

The remaining ten fields in the \type{bilattice} structure are as
follows:
 \begin{description}
 \item [\itype{dws}] This \lisp\ function, ``dot with star,'' accepts
as arguments two truth values, $v_1$ and $v_2$.  The truth value returned
is $v_1$, modified so as to remove any contribution that could be
attributed to a sentence with truth value $v_2$.
 \item [\itype{stash-val}] A function that, given a sentence $p$,
returns the default truth value to be passed to \type{stash}.
 \item [\itype{unstash-val}] A function that, given a sentence $p$,
returns the default truth value to be passed to \type{unstash}.
 \item [\itype{cutoff-val}] The default cutoff point that must be
exceeded by any sentence if it is to be returned by \type{lookup}.
See Section \ref{s.cutoffs}.
 \item [\itype{print-fn}] This function is used to pretty-print
truth values in the associated bilattice.  Its three arguments
are a truth value, the bilattice in question, and a stream.
 \item [\itype{long-desc}] A long description of the bilattice, used
for input and output only.
 \item [\itype{short-desc}] A short description of the bilattice, also
used for input and output only.
 \item [\itype{char}] A one-character identifier for the bilattice.
This character can be used to change the bilattice using the function
\type{load-logic} described in Section \ref{s.change}.
 \item [\itype{type}] An overall description of the mathematical form
of the bilattice.  Perhaps it is a product of two other bilattices,
or formed from them in some other fashion.  This defaults to \type{basic}
if no other value is supplied; various methods for constructing new
bilattices from old ones are described in Section \ref{s.construct}. 
 \item [\itype{components}]  A list of the components used to construct
a bilattice whose \type{type} is not \type{basic}.
 \end{description}

\subsection{The active bilattice}

At any particular point in time, there will be a single bilattice
that has been selected by the user to be used in the application he
is developing.  This bilattice is stored as the value of
\type{*active-bilattice*}; in addition, many of the slots of the
active bilattice are directly available to the user.  Here they
are:

\begin{description}
 \item [\itype{*active-bilattice*}] The currently active bilattice.
 \item [\itype{true}] The truth value \type{true} in the active bilattice.
 \item [\itype{false}] The truth value \type{false} in the active bilattice.
 \item [\itype{unknown}] The truth value \type{unknown} in the active
bilattice.
 \item [\itype{bottom}] The truth value \type{bottom} in the active bilattice.
 \item [\itype{std-cutoffs}] Cutoff information constructed from the
value of \type{cutoff-val} in the active bilattice.  See Section
\ref{s.cutoffs}.
 \item [\itype{stash-value}] The function \type{stash-val} in the
active bilattice.
 \item [\itype{unstash-value}] The function \type{unstash-val} in the
active bilattice.
 \item [\itype{mvl-plus}] The function \type{plus} in the active bilattice.
 \item [\itype{mvl-dot}] The function \type{dot} in the active bilattice.
 \item [\itype{mvl-and}] The function \type{and} in the active bilattice.
 \item [\itype{mvl-or}] The function \type{or} in the active bilattice.
 \item [\itype{mvl-not}] The function \type{not} in the active bilattice.
 \item [\itype{mvl-eq}] The function \type{eq} in the active bilattice.
 \item [\itype{mvl-simplicity}] The function \type{simplicity} in the active
bilattice.
 \item [\itype{dot-with-star}] The function \type{dws} in the active 
bilattice.
 \item [\itype{mvl-vars}] The function \type{vars} in the active bilattice.
 \item [\itype{mvl-plug}] The function \type{plug} in the active bilattice.
 \item [\itype{mvl-print}] This function accepts a single argument, a truth
value, and modifies it in a way that will cause it to be pretty-printed.
The value returned by mvl-print should not be used for any purpose other
than printing.
 \end{description}

All of the functions above accept an optional argument that can be
used to evaluate their result using some bilattice other than the
active one.

\subsubsection{Functions derived from the active bilattice}

There are a variety of functions associated with the functions that
are defined when a new bilattice is made active.  In order to
simplify the description of these functions, we will denote
\type{mvl-and} by $\wedge$, \type{mvl-or} by $\vee$, \type{mvl-not}
by $\neg$, \type{mvl-plus} by $+$, and \type{mvl-dot} by $\cdot$.
The first auxiliary function defined is \type{dot-with-not}:
 \begin{description}
 \item [\itype{dot-with-not(x)}] Computes $x\cdot \neg x$.
 \end{description}

The other functions are constructed from the partial orders defined
using the bilattice operations.  We will write $x\leq_t y$ just in
case $x\wedge y = x$, for example, and $x\leq_k y$ just in case
$x\cdot y = x$.  (All of this notation follows that in the paper,
``Multivalued logics: A uniform approach to inference in artificial
intelligence''.)

Given these definitions, the following \lisp\ functions are defined:
 \begin{description}
 \item [\itype{k-le(x y)}] Returns \T\ if $x \leq_k y$, and \NIL\ otherwise.
 \item [\itype{k-ge(x y)}] Returns \T\ if $x \geq_k y$, and \NIL\ otherwise.
 \item [\itype{t-le(x y)}] Returns \T\ if $x \leq_t y$, and \NIL\ otherwise.
 \item [\itype{t-ge(x y)}] Returns \T\ if $x \geq_t y$, and \NIL\ otherwise.
 \item [\itype{k-lt(x y)}] Returns \T\ if $x <_k y$, and \NIL\ otherwise.
 \item [\itype{k-gt(x y)}] Returns \T\ if $x >_k y$, and \NIL\ otherwise.
 \item [\itype{t-lt(x y)}] Returns \T\ if $x <_t y$, and \NIL\ otherwise.
 \item [\itype{t-gt(x y)}] Returns \T\ if $x >_t y$, and \NIL\ otherwise.
 \item [\itype{k-not-le(x y)}] Returns \T\ if $x \not\leq_k y$, and
\NIL\ otherwise.
 \item [\itype{k-not-ge(x y)}] Returns \T\ if $x \not\geq_k y$, and
\NIL\ otherwise.
 \item [\itype{t-not-le(x y)}] Returns \T\ if $x \not\leq_t y$, and
\NIL\ otherwise.
 \item [\itype{t-not-ge(x y)}] Returns \T\ if $x \not\geq_t y$, and
\NIL\ otherwise.
 \item [\itype{k-not-lt(x y)}] Returns \T\ if $x \not<_k y$, and \NIL\
otherwise.
 \item [\itype{k-not-gt(x y)}] Returns \T\ if $x \not>_k y$, and \NIL\
otherwise.
 \item [\itype{t-not-lt(x y)}] Returns \T\ if $x \not<_t y$, and \NIL\
otherwise.
 \item [\itype{t-not-gt(x y)}] Returns \T\ if $x \not>_t y$, and \NIL\
otherwise.
 \end{description}

All of these functions accept an optional argument that is the
bilattice to be used in the comparison.  This argument defaults to
the currently active bilattice.

\subsubsection{Cutoff information}
\label{s.cutoffs}

We will refer to each of the 16 functions defined at the end of the
previous section (e.g., excluding \type{dot-with-not}) as a {\em tag}
\index{tag}.  Thus \type{k-le} is a tag, and so on.

We will now define a {\em simple test} to be a dotted pair
 \[\type{(v . tag)}\]
 where \type{v} is a truth value and \type{tag} is a tag.  We will
say that a truth value \type{x} {\em satisfies\/} the simple test
if and only if the function \type{tag}, when applied to the arguments
$x$ and $v$ (in that order), returns \T.

As an example, consider the simple test
 \[\type{(true . t-ge)}\]
 This test will be satisfied by any truth value that is $\geq_t
\type{true}$.  Since \type{true} itself is the only truth value
satisfying this test, the test will be satisfied for \type{true} and
no other truth values.  This is the test that is frequently used to
decide to terminate a proof effort.\footnote{Tests in the \mvl\ system
are in fact \lisp\ structures with \type{value} and \type{tag} fields;
we are using dotted lists in this manual only to make the presentation
clearer.}

Alternatively, the test used to determine whether or not to return
an answer after the theorem prover runs to completion is often given
by
 \[\type{(unknown . k-gt)}\]
 This test will succeed for any truth value about which more is known
that \type{unknown} -- in other words, for any truth value other than
\type{unknown} itself.

In general, the simple tests are combined into {\em cutoffs}.  A cutoff
is simply a list of lists of simple tests, and is interpreted as being
in conjunctive normal form.  Thus the cutoff corresponding to the simple
test $s$ is given simply by
 \[\type{((s))}\]

A cutoff corresponding to the conjunction of tests $s_1,\dots,s_n$ is given
by
 \[\type{((s-1) ... (s-n))}\]
 while that corresponding to their {\em disjunction\/} is given by
 \[\type{((s-1 ... s-n))}\]

In general, we can build up conjunctions of disjunctions of simple
tests in as complex a manner as we wish.  The cutoff \NIL\
corresponds to the empty conjunction and therefore always succeeds;
the cutoff \type{(NIL)} contains an empty disjunction and therefore
always fails.  For convenience, these two cutoffs are stored in the
global constants \itype{*always-succeeds*} and
\itype{*never-succeeds*} respectively.

Given a truth value $x$ and a cutoff $c$, the expression \type{(test x c)}
evaluates to \T\ if $x$ passes the test, and to \NIL\ if it fails it.
There is also a function that negates cutoffs:
 \begin{description}
 \item [\itype{test(value cutoffs \&optional bilattice)}] Returns \T\
if the truth value passes the test associated to \type{cutoffs}, and
\NIL\ otherwise.  The test is performed for the supplied bilattice,
or for the currently active one if the optional argument is omitted.
 \item [\itype{cutoffs-not(cutoffs \&optional bilattice)}] Returns a
new cutoffs $c$ such that
 \[\type{(test x c)} = \type{(test (mvl-not x) cutoffs)}\]
 for all truth values $x$.  In other words, $x$ passes $c$ if and
only if $\neg x$ passed the cutoffs supplied originally.
 \end{description}

\subsection{Selecting a new bilattice}
\label{s.change}

To change the bilattice being used, simply invoke the function
\type{load-logic}.  There are three ways to specify the new bilattice:
 \begin{enumerate}
 \item Do nothing.  You will then be asked to select the bilattice
from the list of bilattices known to the system.
 \item Supply the one-character identifier of the new bilattice as an
argument to \type{load-logic}.  The new bilattice must already be
known to the system.
 \item Supply the new bilattice itself as an argument to
\type{load-logic}.
 \end{enumerate}
 \type{Load-logic} empties the database before the bilattice is
changed.

The bilattices known to the system are kept in the global variable
\type{*bilattices*}.  To add a new bilattice to \type{*bilattices*},
invoke the function \type{bilattice} with the bilattice as an
argument.  This function will warn you if you are overwriting another
bilattice that has the same one-character designator.

Summarizing:
 \begin{description}
 \item [\itype{*bilattices*}] A list of the bilattices known to the system.
 \item [\itype{bilattice(b)}] Add the bilattice $b$ to \type{*bilattices*}.
 \item [\itype{load-logic}] Empty the database and make a new bilattice
active.  Accepts as an optional argument either a character designating
a bilattice, or the new bilattice itself.  If no argument is supplied,
the user is prompted for the new bilattice.
 \end{description}

\section{Bilattices available in the MVL system}

Here are the bilattices available in the \mvl\ system as distributed.
Each is discussed more fully in a subsequent section:
 \begin{description}
 \item [\itype{*first-order-bilattice*}] The bilattice corresponding to
first-order logic.
 \item [\itype{*atms-bilattice*}] The bilattice corresponding to an \atms,
where the justification information is maintained in disjunctive normal
form as suggested by de Kleer.
 \item [\itype{*cnf-atms-bilattice*}] An alternative version of the \atms\
bilattice where the justification information is maintained in conjunctive
normal form.
 \item [\itype{*default-bilattice*}] A bilattice that can be used for simple
default reasoning.
 \item [\itype{*hierarchy-bilattice*}] A bilattice similar to the
default bilattice but where the defaults are labelled with
integer-valued strengths.
 \item [\itype{*first-order-atms-bilattice*}] This bilattice corresponds
to a fully first-order \atms, as opposed to the more conventional one
dealing only with propositional logic.
 \item [\itype{*circumscription-bilattice*}] The bilattice used to compute
the results of the circumscription axiom.
 \item [\itype{*probability-bilattice*}] An extended version of the \atms\
bilattice that maintains probabilistic information as well.
 \item [\itype{*time-bilattice*}] A compound bilattice that can be used
for reasoning about events occurring against a temporal background.
 \end{description}

\subsection{First-order logic} \index{first-order logic}

The simplest bilattice is that corresponding to first-order logic, and
contains only the four truth values \type{true} (displayed as \type
T), \type{false} (displayed as \type F), \type{unknown} (\type U) and
\type{bottom} (\type{T/F}).  This bilattice is identified using the
character {\bf f}, and is active when \mvl\ is first loaded.

\subsection{Assumption-based truth maintenance}
\index{assumption-based truth maintenance}

Except for the \nm\ and probabilistic \atms\ bilattices, all of the
bilattices corresponding to \atms's have truth values that are dotted
pairs \type{(j . k)}, where $j$ is a justification for the sentence
being labelled, and $k$ is a justification for its negation.  The
bilattices differ in that the justifications are represented
differently.

\subsubsection{The ATMS bilattice}

For the simplest \atms\ bilattice, the justifications are simply
disjunctive normal form expressions.  Thus the justification
 \begin{equation}
\type{((P23 P37) (P1))}
\label{e.just}
\end{equation}
 is interpreted to mean that the sentence in question is a
consequence either of \type{P1}, or of \type{P23} and \type{P37} in
conjunction.

This bilattice is identified by the character {\bf a}.

\subsubsection{ATMS's in conjunctive normal form}

Alternatively, we can represent the justifications in conjunctive
normal form, as opposed to disjunctive normal form.  This representation
is useful, for example, in diagnosing multiple faults, since we want to
be able to explain an observed failure {\em either\/} in terms of the
collective failures of one set of components, or the collective failures
of another set, and so on.

This bilattice is identified by the character {\bf b}.

\subsubsection{First-order ATMS's}
\label{s.fo-atms}

It is also possible to extend these ideas slightly to obtain a fully
first-order \atms.  In order to do this, we need simply modify
justifications such as that appearing in (\ref{e.just}) so that the
primitive elements of the justification are not atoms referring to
propositions, but pairs
 \[\type{(atom . binding-list)}\]
 where the atom refers to a proposition and the binding list
indicates precisely what instantiation of that proposition is
actually needed to justify the sentence in question.

As an example, suppose that an assumption in our database is \type{P1},
given by
 \begin{equation}
\type{(FRIENDLY ?X)}
\label{e.assume}
\end{equation}
 and indicating that everything is friendly.  If we also have the rule
of inference
 \[\type{(<= (HAPPY ?X) (FRIENDLY ?X))}\]
 indicating that anything friendly is happy, and this rule is unconditionally
true (i.e., not an assumption), then the truth value assigned to the
query \type{(HAPPY HARRY)} will be
 \[\type{((((P1 . ((?X . HARRY))))))}\]
 indicating that Harry is happy by virtue of (\ref{e.assume}) applied
to $x=\mbox{Harry}$, and that there is no justification for the
conclusion that Harry is not happy.  (Recall that all of the \atms\
bilattices combine justification for and against each sentence into a
single truth value.)

This bilattice is identified by the character {\bf o}.

\subsubsection{Probabilistic ATMS's}
\label{s.prob}

Truth values in the probabilistic \atms\ bilattice consist of dotted
pairs, where the \type{car} of a truth value is simply an element of
the \atms\ bilattice and the \type{cdr} is the probability that the
given assumption is valid.

When entering an assumption, the user will be queried for the
probability in question; the probabilities of compound sentences are
then computed automatically by the system.  The algorithm used is
that described in Grnarov, Kleinrock and Gerla in the 1979
Proceedings of the Computer Networking Symposium.

The probabilistic \atms\ bilattice is identified by the character
{\bf p}.
  
\subsection{Default reasoning} \index{default reasoning}
\label{s.j}

The default bilattice consists of seven truth values.  In addition to
the usual elements \type{true}, \type{false}, \type{unknown} and
\type{bottom}, there are elements \type{dt}, \type{df} and \type{*}.

The truth value \type{dt} labels a conclusion that is true by default
(such as Tweety flying in the canonical example), while \type{df} labels
a conclusion that is {\em false\/} by default.  The final value,
\type{*}, labels conclusions that are both true {\em and\/} false
by default, as in the well-known Nixon-Republican-Quaker example.

This bilattice is identified by the character {\bf d}.

\subsubsection{Prioritized default reasoning}

Another version of the default bilattice allows for priorities among
the defaults.  In this bilattice, truth values are of the form
 \[\type{(n . x)}\]
 where $n$ is a nonnegative integer and $x$ is an element of the
first-order bilattice.  Truth values with $n=0$ are displayed as for
the first-order bilattice; with $n=1$ as for the default bilattice.
For $n>1$, we might see \type{D(3)T} or something like it.  This
particular truth value is the same as
 \[\type{(3 . true)}\]
 Truth values with small $n$ override those with larger $n$.
\type{Unknown} in this bilattice is given by the truth value \NIL.
The one-character identifier for the hierarchical default bilattice is
{\bf h}.

\subsection{Circumscription} \index{circumscription}
\label{s.circ}

The default bilattices are not able to deal with all of the
situations in which multiple extensions arise because single truth
values such as \type{*} label sentences that are true in some
extensions and false in others without providing any information
about precisely which extensions are which.  In order to handle this
difficulty, we need to work with a more complex bilattice.

One way to view the default bilattice is as a copy of the first-order
bilattice in which the point \type{unknown} is interpreted as, ``not
known with certainty,'' and is in fact replaced with another copy of
the first-order bilattice that is used to record uncertain
information (i.e., default truth, default falsity, or the presence of
a default contradiction).

The circumscriptive bilattice is formed similarly, except that
instead of simply recording default truth, falsity, and so on, we
record a {\em justification\/} for default conclusions.  This is done
by replacing the point \type{unknown} of the first-order bilattice
not with another copy of the first-order bilattice, but with a copy
of the first-order \atms\ bilattice.  (This idea is in fact a general
way to construct new bilattices from old ones, and is called a {\em
surgery}.  It is discussed in Section \ref{s.surgery}.)  The
designator for the circumscriptive bilattice is {\bf c}.
\type{Circum} and \type{circums} use this bilattice to accept a
sentence as an argument and attempt to prove it from the
circumscription axiom; they resemble \type{bc} and \type{bcs}.

The truth values in the circumscription bilattice corresponding to
certain information are of the form \type{(T . x)}, where $x$ is an
element of the first-order bilattice.  Those corresponding to uncertain
information are of the form \type{(NIL . y)}, where $y$ is an element
of the first-order \atms\ bilattice.

Finally, the circumscriptive theorem prover will output diagnostics
concerning its progress if the global variable \type{*circum-trace*}
has a non-\NIL\ value.

Here is a summary of the circumscriptive facilities:
 \begin{description} 
 \item [\itype{*circum-trace*}] If not \NIL, the circumscriptive proof
functions will output diagnostics describing their progress.
 \item [\itype{circum(p)}] Like \type{bc}, but looks for instances of
$p$ that follow from the circumscription axiom.  Allowed keywords:
\type{cutoffs} and \type{succeeds}.
 \item [\itype{circums(p)}] Like \type{bcs}, but looks for instances of
$p$ that follow from the circumscription axiom.  Allowed keywords:
\type{cutoffs} and \type{succeeds}.
 \end{description}

\subsection{Temporal reasoning}\index{temporal reasoning}

\type{*Time-bilattice*} is a bilattice of the truth values that might
be assigned to propositions that have a time-dependent nature, such as
whether or not a particular switch is on.  The truth values represent
the entire histories of the propositions, so that a single truth value
is a complete function from a set of time points (taken to be the
integers) into a set of ``base'' truth values (take to be elements of
the nonmonotonic {\sc atms} bilattice). This bilattice is discussed at
greater length in the paper, ``Computational considerations in
reasoning about action,'' and a sample use of this bilattice can be
found by examining the file \type{temporal.test} that is included with
\mvl.

\section{Facilities to create new bilattices}
\label{s.construct}

There are a variety of ways in which it is possible to construct
new bilattices from old ones, or from old ones in conjunction
with instances of some other structure.  In this section, we
describe the extent to which these facilities exist within the
\mvl\ system.

\subsection{Products}
\label{s.product}

Perhaps the simplest way in which a new bilattice can be constructed
is as the product of two existing bilattices.  If $A$ and $B$ are
bilattices, then a truth value in the product bilattice $A\times B$
will simply be a pair of truth values \type{(x . y)}, where $x$ is a
truth value in $A$ and $y$ is one in $B$.  All of the bilattice
operations on the product bilattice are computed pointwise using the
corresponding operations on the component bilattices $A$ and $B$.

The product construction is performed using the following function:
 \begin{description}
 \item [\itype{product-bilattice(b1 b2)}] Returns the bilattice
that is the product of the bilattices $b_1$ and $b_2$.
 \end{description}

\subsection{Functional attachments}

Another way in which to construct a bilattice is to start with a
bilattice $B$ and a function $f$ that maps $B$ into a fixed set $S$.
The elements of the new bilattice are simply pairs of the form
 \[\type{(x .~(f x))}\]
 where $x$ is an element of $B$.  The function $f$ here is simply
tagging along for the ride, and has no impact on the bilattice
operations.  It can be used to provide additional information about
the truth values, and is used to compute probabilistic information in
the bilattice {\bf p} described in Section \ref{s.prob}.

The relevant function is the following one:
 \begin{description}
 \item [\itype{append-bilattice(b f)}] Returns the bilattice obtained
by appending the result of applying the function $f$ to the points of
the original bilattice $b$.
 \end{description}

\subsection{Surgeries}
\label{s.surgery}

A somewhat more interesting construction is the one we described in
Section \ref{s.circ} as a {\em surgery}.  Here, we take one bilattice
$b_1$ and replace the point corresponding to \type{unknown} in this
bilattice with another bilattice $b_2$.  The truth values in the
resulting bilattice are either of the form \type{(T . x)} for a point
$x$ in the ``upper'' bilattice $b_1$, or \type{(NIL . y)} for a point
$y$ in the lower bilattice.  Truth values from $b_1$ take priority over
values from $b_2$.

To perform this construction generally, we have: 
 \begin{description}
 \item [\itype{splice-bilattice(upper lower)}] Returns the bilattice
obtained by replacing the point \type{unknown} in the bilattice
\type{upper} with a copy of the bilattice \type{lower}.
 \end{description}

As an example, the default bilattice {\bf d} of Section \ref{s.j} is
given by
 \[\type{(splice-bilattice *first-order-bilattice* *first-order-bilattice*)}\]
 The circumscriptive bilattice of Section \ref{s.circ} is given by
 \[\type{(splice-bilattice *first-order-bilattice*
*first-order-atms-bilattice*)}\]

\subsection{Hierarchies}

It is possible to iterate the surgery construction of the previous
section, producing a bilattice that is an infinite descending
``chain'' of copies of a basic bilattice $b$.  This is done by
the function \type{hierarchy-bilattice}:
 \begin{description}
 \item [\itype{hierarchy-bilattice(b)}] Returns the bilattice constructed
as an infinite descending chain of copies of $b$.
 \end{description}

Truth values in this bilattice have the form \type{(n . x)}, where
$x$ is an element of $b$, and $n$ is an integer giving the hierarchy
level.  The top hierarchy level of the bilattice has $n=0$; the
second level is $n=1$, and so on.  The point \type{unknown} is given
by \NIL.

This function is used to construct the hierarchical version of the
default bilattice.

\subsection{Lattices}

Bilattices can also be constructed from lattices.  A lattice is a set
equipped with a {\em single\/} partial order and greatest lower bound
and least upper bound operations, which we denote by $+$ and $\cdot$
respectively.

An element of the associated bilattice is a {\em pair\/} of elements
of the lattice, and the bilattice operations are given by:
 \begin{eqnarray*}
\neg(a,b) & = & (b,a) \\
(a,b) + (c,d) & = & (a+c,b+d) \\
(a,b) \cdot (c,d) & = & (a\cdot c,b\cdot d) \\
(a,b) \vee (c,d) & = & (a+c, b\cdot d) \\
(a,b) \wedge (c,d) & = & (a\cdot c,b+d)
 \end{eqnarray*}

Roughly speaking, the first component of the bilattice element labels
the ``truth'' of the sentence, and the second labels its falsity.  This
is how the \atms\ bilattices are constructed.

The function used is \type{lattice-to-bilattice}:

 \begin{description}
 \item [\itype{lattice-to-bilattice(l)}] Returns the bilattice constructed
from the lattice $l$.
 \end{description}

A lattice is a structure rather like a bilattice, but with fewer slots
(because it only has a single partial order).  The fields in an instance
of the \itype{lattice} structure are:

\begin{description}
 \item [\type{max}] A \lisp\ object that is the maximal element of
the lattice.
 \item [\type{min}] A \lisp\ object that is the minimal element of
the lattice.
 \item [\type{eq}] A \lisp\ function that takes two lattice elements
as arguments and returns \T\ if they are equal and \NIL\ if they are
not.
 \item [\type{and}] A \lisp\ function that takes two lattice elements
as arguments and returns their greatest lower bound.
 \item [\type{or}] A \lisp\ function that takes two lattice elements
as arguments and returns their least upper bound.
 \item [\type{le}] This optional field is a \lisp\ function that
accepts two arguments $x$ and $y$ and determines whether or not
$x\leq y$ under the partial order corresponding to the lattice.
 \item [\type{vars}] This \lisp\ function accepts a lattice element
as an argument and returns a list of the variables that it contains.
The default value returns \NIL\ independent of its argument.
 \item [\type{plug}] This \lisp\ function accepts as arguments a
lattice value and a binding list, and returns a potentially modified
lattice value.  The default value is a function that leaves the lattice
value unchanged.
 \item [\type{simplicity}] The lattice analog of the bilattice
\type{simplicity} slot.  When constructing a bilattice from a
lattice, the \type{simplicity} of a truth value $(a,b)$ is the result
of applying the simplicity function from the lattice to the lattice
element $a$.  The value $b$ is ignored in the simplicity computation.
 \item [\type{dws}] This \lisp\ function, ``dot with star,'' accepts
as arguments two lattice values, $v_1$ and $v_2$.  The lattice
element returned is $v_1$, modified so as to remove any contribution
that could be attributed to the lattice element $v_2$.
 \item [\type{stash-val}] The lattice element entering into
the bilattice function \type{stash-val}.
 \item [\type{long-desc}] A long description of the lattice, used
for input and output only.
 \item [\type{short-desc}] A short description of the lattice, also
used for input and output only.
 \item [\type{char}] A one-character identifier for the lattice.
 \end{description}

\subsection{Directed acyclic graphs}
\label{s.dag}

The final fashion in which \mvl\ supports the construction of new
bilattices is an extension of the product construction of Section
\ref{s.product}.  There, we formed the product of two bilattices
$A\times B$.  If the bilattices were identical, we would have formed
the product bilattice $B\times B$, or $B^2$.  Another way to view
this bilattice is as the collection of functions from the two point
set $\set{1,2}$ into the original bilattice $B$.  Since such a
function $f$ can be described as the pair $(f(1),f(2))$, there is no
difference between this description and the previous one.

Given this analogy, it is not hard to see that we can extend the
construction to the bilattice $B^S$ for an arbitrary set $S$, where
elements of this bilattice are functions from the set $S$ into $B$,
and all of the bilattice functions ($\wedge$, $\cdot$ and so on) are
once again computed pointwise.

The difficulty with this idea is that for large sets $S$, it is not
clear how we can represent a function from $S$ into $B$ in a
convenient way.  It turns out, however, that if $S$ has the structure
of a directed acyclic graph, or \dg, this problem can be overcome.

\subsubsection{Defining the DAG}

As a preliminary, we need some way to define the \dg\ itself.  In
\mvl, a \dg\ \index{dag} is an instance of the \type{dag} structure;
this structure has the following fields:
 \begin{description}
 \item [\itype{root}] The root of the \dg.
 \item [\itype{eq}] A function that accepts two elements of the \dg\
and returns \T\ if they are the same element and \NIL\ otherwise.
 \item [\itype{leq}] A function that accepts two elements of the \dg\
and returns \T\ if the first is farther from the root of the \dg\
than the second.  If not supplied by the user, this function can by
computed from the \type{dot} operation on the \dg.
 \item [\itype{dot}] A function that accepts two elements of the \dg\
and returns a list of their first common descendants, or \NIL\ if no
such descendant exists.
 \item [\itype{long}] A long description of the \dg.
 \item [\itype{short}] A short description of the \dg.
 \item [\itype{inherit}] A function describing the fashion in which
truth values are inherited within the \dg.  This function $f$ should take
four arguments: a bilattice $b$, two elements of the \dg\ $x$ and $y$,
and an element $z$ of the bilattice.  If $z$ is the element of the
bilattice associated to the \dg\ element $x$, then $f(b,x,y,z)$ will
be the value associated to the \dg\ element $y$ in the absence of
information to the contrary.  By default, this function ignores
its first three arguments and assigns $y$ the same value as $x$ has.
 \item [\type{vars}] This \lisp\ function accepts a \dg\ element
as an argument and returns a list of the variables that it contains.
The default value returns \NIL\ independent of its argument.
 \item [\type{plug}] This \lisp\ function accepts as arguments a
\dg\ value and a binding list, and returns a potentially modified
\dg\ value.  The default value is a function that leaves the \dg\
value unchanged.
 \item [\type{select}] In some applications (none of which is distributed
with \mvl), it appears that truth values should be inherited from only
{\em some\/} of the elements of the \dg\ which might potentially
impact the value of the function at some specific point.  The
selection function is given a point where the truth value is being
computed and a list of points that apparently contribute; it returns a
list of points that {\em actually\/} contribute.  By default, all
points contribute; most users of \mvl\ are very unlikely to need this
facility.
 \end{description}

\subsubsection{Constructing the bilattice}

Having selected a \dg\ $D$ and a bilattice $B$, we can now construct
the bilattice $B^D$.  Specific functions from $D$ to $B$ are described
as structures with the following slots:
 \begin{description}
 \item [\type{dag}] The domain of the function, $D$.
 \item [\type{bilattice}] The range of the function, $B$.
 \item [\type{list}] A list of elements of the form \type{(p . v)},
where \type p is a point where the value changes in a fashion {\em
other\/} than the one predicted by the inheritance function associated
with $D$ and \type v is the value taken there.
 \end{description}
 We will refer to these functions as {\em \dg-functions}.

There are a variety of functions available to manipulate \dg-functions,
and these are used to construct the bilattice $B^D$.  In some cases,
however, the user may want to access these functions directly.  Here
are some of them:
 \begin{description}
 \item [\itype{get-val(d f)}] Returns the value of the
function $f$ at the point $d$ of the \dg.
 \item [\itype{make-tv(dag bilattice d b)}] Returns a \dg-function that
takes the value $b$ at the point $d$.  If $d$ is not the root of the
\dg, the value at the root of the \dg\ is generally the point
\type{unknown} in the bilattice, but the user may supply the keyword
\type{root-fn} if he wishes, in which case the value at the root of
the \dg\ is the result of applying \type{root-fn} to the bilattice.
 \item [\itype{make-root(dag bilattice b)}] Returns a \dg-function that
takes the value $b$ at the root of the \dg.
 \item [\itype{dag-change-tv(f b-f)}] The argument
\type{b-f} should be a function from the bilattice to itself; a
modified version of the original \dg-function $f$ is returned in which
every point in the \dg-function has had the function \type{b-f}
applied to it.  There is an optional keyword \type{simplify}; if
supplied, the computed function is simplified before being returned.
(See \type{simplify}.)
 \item [\itype{dag-change-dag(f d-f)}] The argument
\type{d-f} should be a function from the \dg\ to itself; a modified
version of the original \dg-function $f$ is returned in which every
point at which the \dg-function takes a new value has had the function
\type{d-f} applied to it.  The bilattice argument is expected to be
the bilattice of functional truth values as opposed to the bilattice
that is the target of the functions.  There is an optional keyword
\type{simplify}; if supplied, the computed function is simplified
before being returned.  (See \type{simplify}.)
 \item [\itype{dag-list-of-answers(f)}] Given a \dg-function $f$, assumed
to be from the \dg\ of binding lists into the currently active bilattice, this
function constructs a list of answers corresponding to the points where
the \dg-function takes values.
 \item [\itype{dag-accumulate-fn(dag bilattice list)}] Given a \dg, a
bilattice, and a list of elements of the form
 \[\type{(dag-elt . bilattice-elt)}\]
a \dg-function is constructed that takes the given values at the
given elements of the \dg.

A keyword to \type{dag-accumulate-fn} is \type{modifier}.  This
function is used if the \type{d-list} contains the same \dg\ element
more than once, in which case the modifier function is called with
the new and old values (in that order) and the result is used as the
value of the \dg-function at that particular point.  The default
action is to simply overwrite the old value with the new one.  (In
other words, the default modifier function simply returns its first
argument while ignoring the second.)
 \item [\itype{dag-prune(f pred)}] Given a \dg-function
$f$ and a predicate on the bilattice, this function prunes from $f$
all points $d$ where $f(d)$ fails to satisfy the predicate.  There is
an optional keyword \type{modify-fn} that, if supplied, should be a
function $m$ from the bilattice to itself.  In this case, the points
where $m(f(d))$ fails to satisfy the predicate are pruned.
 \item [\itype{combine-fns(f1 f2 fn comp)}]
This function takes the two \dg-functions $f_1$ and $f_2$ and
combines them using the function \type{fn}, which should be a
function that takes two elements of the bilattice and returns a
single one.  The comparison function \type{comp} is used to determine
if one of the two functions can be returned directly; if, for example
\type{(comp f1 f2)} holds, than  we simply return $f_1$, and similarly
for $f_2$.  As an example, if the combining function is the bilattice
operation $\cdot$, the comparison function is \type{k-le}.
 \item [\itype{combine-1(f params f1 ... fn)}] This is an
$n$-ary version of \type{combine-fns}, where $f$ is an $n$-ary
bilattice function and there are $n$ \dg-functions supplied as well.
The additional argument \type{params} indicates which of the $n$ arguments
to this function are not necessarily \dg-functions, but are instead parameters
expected by the overall function $f$.  This additional argument is a list
of \T\ (if the $n$th argument is a parameter) or \NIL\ (if it isn't).
If the parameter list ends early, all of the remaining entries are
assumed to be \NIL.
 \item [\itype{add-dag(d b f)}] Given the \dg-function
$f$, add the value $b$ at the point $d$.  This function works by side
effect only.  There is a keyword \type{modifier} that, if supplied,
should be a function that takes two elements of the bilattice and
returns a single one.  Given such a function $m$, the \dg-function is
changed so that the new value at $d$ is $m(b,f(d))$, where $f(d)$ is
the old value.  (This is similar to the same keyword in
\type{dag-accumulate-fn}.)  The keyword \type{simplify} is also available.
 \item [\itype{simplify(f)}] In some cases,
manipulating the \dg-function $f$ may cause it to contain
redundancies, where the value at some point is presented explicitly
even though this value is in fact the one that would normally be
inherited from its parents.  This function returns a modification of
$f$ where these irrelevant entries have been removed.
 \item [\itype{fn-and(f1 f2)}] This function conjoins
the values of $f_1$ and $f_2$ using the bilattice $\wedge$ operation.
The new \dg-function is returned.
 \item [\itype{fn-or(f1 f2)}] This function disjoins
the values of $f_1$ and $f_2$ using the bilattice $\vee$ operation.
The new \dg-function is returned.
 \item [\itype{fn-plus(f1 f2)}] This function adds
the values of $f_1$ and $f_2$ using the bilattice $+$ operation.
The new \dg-function is returned.
 \item [\itype{fn-dot(f1 f2)}] This function combines
the values of $f_1$ and $f_2$ using the bilattice $\cdot$ operation.
The new \dg-function is returned.
 \item [\itype{fn-dws(f1 f2)}] This function combines
the values of $f_1$ and $f_2$ using the bilattice operation
\type{dot-with-star}.  The new \dg-function is returned.
 \item [\itype{fn-not(f1)}] This function negates the
values of $f_1$ using the bilattice $\neg$ operation.  The new
\dg-function is returned.
 \item [\itype{fn-eq(f1 f2)}] This function compares
the values of $f_1$ and $f_2$ using the bilattice comparison operation.
Either \T\ or \NIL\ is returned.
 \item [\itype{fn-comp(f1 f2 fn)}] This functions compares
the values of $f_1$ and $f_2$ using the bilattice comparison function
\type{fn}.  Either \T\ or \NIL\ is returned.
 \end{description}

Of course, the point of all of this is to enable us to define the
following function:
 \begin{description}
 \item [\itype{dag-bilattice(dag bilattice)}] Given a bilattice $B$
and a dag $D$, returns the bilattice $B^D$.
 \end{description}

The following functions are also available:
 \begin{description}
 \item [\itype{lattice-to-dag(lattice)}] Construct a \dg\ from the
given lattice.
 \item [\itype{lattice-to-dag-to-bilattice(lattice bilattice)}]  First
construct a \dg\ $D$ from the given lattice.  Then use the bilattice
$B$ to construct the bilattice $B^D$.
 \end{description}

\subsubsection{The DAG of theories}
\label{s.theory}

Two specific \dg s are provided for use in the \mvl\ system.  The first
uses a \dg\ corresponding to the existence of multiple theories within
a single database; the second uses a \dg\ to handle binding information.

The theory \dg\ consists of a hierarchy of theories.  In a single
database, for example, we might want to have rules that applied to
celestial bodies, rules applying to stars, rules applying only to red
giants or only to our sun, and so on.  The truth value of a rule
applicable to all celestial bodies should be inherited if we are
considering our sun (but it should be possible to overwrite it with
a different truth value if we wish), but a rule applicable to red
giants should not be inherited in this fashion.  So we see that our
\dg\ construction is well-suited for this sort of an application.

To activate the theory mechanism, simply invoke the \lisp\ function
\type{activate-theory-mechanism}.  This function replaces the active
bilattice, say $B$, with the bilattice constructed of functions from
the theory \dg\ into $B$.  The original bilattice can still be
accessed, since it is the second \type{component} of the newly formed
bilattice.  The function \type{theory-active-p} can be used to
determine whether or not the theory mechanism has been activated.

The \dg\ of theories is the value of the global variable
\type{theory-dag}, and has as its root the theory whose name is the
value of the constant \type{*global*}.  To indicate that a theory
\type{big} includes another theory \type{little} as a subtheory,
invoke the function 
 \[\type{(includes big little)}\]
 Thus, for example, one would introduce a new theory called
\type{theory} by writing
 \[\type{(includes *global* 'theory)}\]
 The effects of \type{includes} can be reversed by invoking the
function \type{unincludes}.

Note that since the theory \dg\ is used by the functions that
manipulate the bilattice $B^D$, it is important that this \dg\ be
maintained as new theories are introduced.  If \type{includes} is not
used to ``locate'' these new theories in the overall theory \dg\
itself, the theories will be inaccessible to many of the \mvl\
functions.

At any point in time, any new fact will be inserted into the database
as true in a particular theory; the theory into which it is inserted
is the value of \type{*active-theory*}.  The value of this variable
should not be changed directly, but should instead be modified by
invoking the function \type{activate} on the theory into which
subsequent facts should be written.  The variable
\type{true-in-theory} contains the truth value of a sentence that is
true in the currently active writing theory.

The theory \type{*active-theory*} is also the theory that is active
for reading at any point in time.  Of course, any parents of the
active theory are also active by virtue of the structure of the theory
\dg\ itself; if it is desired to read from multiple theories
simultaneously, a new theory should be created that is a child of all
of them.

Here is a summary of the functions described so far:
 \begin{description}
 \item [\itype{activate-theory-mechanism}] Activates the theory
mechanism.  Accepts as an optional argument either a one-character
bilattice descriptor, or a bilattice.
 \item [\itype{theory-active-p}] Returns \T\ or \NIL\ depending on
whether or not the theory mechanism has been activated.
 \item [\itype{theory-dag}] The \dg\ of theories.
 \item [\itype{*global*}] The root of the \dg\ of theories.
 \item [\itype{includes(big little)}] This function is used to state that
the theory \type{little} is a subtheory of the theory \type{big}, and
should inherit from it.
 \item [\itype{unincludes(big little)}] This function is used to state
that the theory \type{little} is no longer a subtheory of the theory
\type{big}, and should not inherit from it.
 \item [\itype{*active-theory*}] The theory that is currently active
for reading and writing.
 \item [\itype{activate(th)}] Activate the theory \type{th} for reading
and writing.
 \item [\itype{true-in-theory}]  The truth value of a sentence that is
true in the currently-active writing theory.
 \end{description}

Finally, five functions have been added to simplify the user
interface somewhat.  The function {\tt theory-\-contents} displays
the contents of a given theory, and \type{empty-theory} removes all
of the propositions in this theory.  To modify the contents of a
specific theory, the following functions can be used:
 \begin{description}
 \item [\itype{state-in-theory(p)}] Adds the proposition $p$ to the
currently active theory.  Allowed keywords: \type{theory} can
be used to specify a different theory to which $p$ should be added;
\type{value} and \type{increment-p} have the same meaning as for
\type{state}.  (See Section \ref{s.database}; if these keywords are
not used, the function \type{state} can be used equally well, since
the default stash value with the theory mechanism activated also
inserts the proposition into the currently active theory.)
 \item [\itype{unstate-from-theory(p)}] Removes the proposition $p$
from the currently active theory.  Allowed keywords:
\type{theory} can be used to specify a different theory from which
$p$ should be removed; \type{value} and \type{decrement-p} have the
same meaning as for \type{unstate}.  (See Section \ref{s.database}.)
 \item [\itype{erase-from-theory(p)}] Removes the proposition $p$
from the currently active theory, and forward chains on the result.
Allowed keywords: \type{theory} can be used to specify the theory from
which $p$ should be removed; \type{value}, \type{decrement-p} and
\type{cutoffs} have the same meaning as for \type{erase}.  (See
Section \ref{s.fc}.)
 \item [\itype{theory-contents()}] Displays the contents of the
currently active theory, which can be changed using the
keyword \type{theory}.  The keywords \type{cutoffs} and \type{props}
are allowed and are treated as for \type{contents}.  (See Section
\ref{s.contents}.)  There is an additional keyword,
\itype{this-theory-only}, which defaults to \T\ and, if supplied,
indicates that only sentences whose truth values {\em change\/} in the
given theory are to be displayed.  (This avoids displaying all
sentences in the global theory, for example.)
 \item [\itype{empty-theory()}] Empties the currently active
theory, although this theory may be changed using the keyword
\type{theory}.  The keyword \type{value} has a meaning similar to its
meaning as a keyword to \type{erase}.
 \end{description}

\subsubsection{The DAG of binding lists}
\label{s.binding-dag}

The other \dg\ supplied with the \mvl\ system is used for binding
lists since these, too, admit a natural inheritance structure.  The
root of the \dg\ is the empty binding list \NIL, and inheritance in
the \dg\ is determined by subsumption.  The bilattice constructed as
functions from this \dg\ into \type{*active-bilattice*} is used
extensively by the prover, and is available to the user as
\type{bdg-to-truth-val}.

The binding \dg\ itself is known as \type{binding-dag}, and is constructed
using the following functions:
 \begin{description}
 \item [\itype{equal-binding(b1 b2)}] Checks to see if two binding
lists are effectively the same.
 \item [\itype{append-binding-lists(b1 b2)}] Computes the least
specific binding list that is more specific than both $b_1$ and $b_2$
or returns \NIL\ if no such binding list exists ($b_1$ and $b_2$ may
bind the same variable to different constants, for example).
 \item [\itype{binding-le(b1 b2)}] Returns \NIL\ unless $b_1$ is a more
specific binding that $b_2$, in which case a list of relative binding
lists is returned.
 \end{description}

The following global variables are also used:
 \begin{description}
 \item [\itype{binding-dag}] The \dg\ of binding lists.
 \item [\itype{bdg-to-truth-val}] The bilattice of functions from
the \dg\ of binding lists into \type{*active-bilattice*}.
 \end{description}

\section{Modal operators}
\label{s.modal}

An additional feature that makes the \mvl\ system unique is its
ability to handle {\em modal operators}.  The approach taken by \mvl\
to modal operators is that they are functions from the bilattice of
truth values to itself.  Thus, for example, the modal operator that
takes a sentence $p$ into the sentence, ``I know that $p$,''
corresponds to the bilattice function (on the first-order bilattice)
that takes $t$ and $\bot$ into $t$ (since we know $p$ if we know it
to be true or know it to be true and false both), and takes $f$ and
$u$ into $f$ (since we do {\em not\/} know $p$ if we either know it
to be false or do not know if it is true or false).  Further details
can be found in the paper, ``Bilattices and modal operators.''

A modal expression in \mvl\ is any sentence of the form 
 \[\type{(op p-1 \dots\ p-n)}\]
 where \type{op} is recognizable as a modal operator because it has a
non-\NIL\ value for its \type{modal-op} property in the currently
active invocation of \mvl.  In this case, the value of the
\type{modal-op} property should be an instance of the structure
\itype{modal-operator}.  This structure has the following fields:
 \begin{description}
 \item [\itype{name}] The name of the modal operator.  This is used
to identify modal operators in different bilattices that have the
same commonsense meaning (and can therefore be usefully combined when
the bilattices are combined as in the previous section).
 \item [\itype{function}] A function that accepts as arguments the
truth values of the sentences under the scope of the modal operator
and returns the truth value of the modal expression itself.
 \item [\itype{parameters}] If the modal operator accepts $n$
arguments, this should be a list of $n$ entries indicating whether
each argument is a truth value or a parameter that should be passed
directly to the modal operator function.  (An entry of \T\ indicates
that the associated argument is parametric; an entry of \NIL\
indicates that it isn't.)  Suppose, for example, that we are working
with a bilattice where the truth values are functions describing the
behaviors of various fluents over the course of time, and have a modal
operator \type{delay}, where the truth value assigned to
 \[\type{(delay t p)}\]
 is expected to be the same as the truth value assigned to $p$ but
delayed by an amount of time $t$.  Here, the first argument ($t$) is
a parameter and the second argument ($p$) is not.

If there are fewer than $n$ entries on the parameter list, then all of
the missing entries are assumed to be \NIL\ (so that the associated
arguments to the modal function are not parameters).  Note: All
parametric arguments should be bound when attempting to derive a modal
expression.  This is similar to the requirement that arguments to a
procedurally-attached predicate be bound when the procedure is called.
 \end{description}

\subsection{Modal operators supplied with the MVL bilattices}

All bilattices in \mvl\ are equipped with a modal operator \itype{idn}
that does not change the truth value of its argument at all; this
modal operator can be used to semantically mark points where it may be
desirable to interrupt the inference process.  In addition, if the
prover is invoked on \type{(idn p)} instead of on \type p, attempts
will be made to prove both $p$ and its negation $\neg p$.

The bilattices supplied as part of the \mvl\ system are equipped with
the following modal operators:

 \begin{description}
 \item [\type{*first-order-bilattice*}] This bilattice has two
modal operators: $L$, which maps $p$ into, ``necessarily $p$,'' or,
``I know that $p$,'' and $M$, which maps $p$ into ``possibly $p$.''
 \item [\type{*atms-bilattice*}]  No modal operators are supplied with
this or any of the other \atms-type bilattices.
 \item [\type{*default-bilattice*}] This bilattice and the hierarchical
default bilattice both inherit the modal operators $L$ and $M$ from the
first-order bilattice.
 \item [\type{*circumscription-bilattice*}] This bilattice does not
come equipped with any modal operators.
 \item [\type{*time-bilattice*}] This bilattice inherits the modal
operators $L$ and $M$ from the bilattice of basic values.  The
following three additional modal operators are also defined:
 \begin{enumerate}
 \item \itype{delay(t p)} This is a modal operator that pushes the
truth value of $p$ into the future by an amount of time $t$.  Thus if
$p$ is true at time 2, $\type{delay}(p,2)$ will be true at time 4.
 \item \itype{propagate} is a modal operator that corresponds to the frame
axiom.  If $p$ is true at time 2, unknown between times 3 and 20, and false
at time 21, then $\type{propagate}(p)$ will be true from times 2 until 20
and false subsequently.
 \item \itype{at(t p)} is a modal operator that returns the value taken
by the truth value $p$ at the time $t$.  (In actuality, a function has
to be returned; the function is the one that takes value $p(t)$ at all
times instead of just at $t$.)
 \end{enumerate}
 \end{description}

\subsection{Quantification and modal operators}

In addition to \type{idn}, there are two other modal operators that
are available for use with any bilattice, \itype{exists} and
\itype{forall}.  Although not really modal operators on the bilattice
being used, these {\em are\/} legitimate modal operators on the bilattice
\type{bdg-to-truth-val} that was described in Section \ref{s.binding-dag}.

The functions associated to these modal operators mimic the usual behaviors
given to the quantifiers.  Thus an expression such as
 \begin{equation}
\type{(exists ?x (foo ?x))}
\label{e.exists}
\end{equation}
 is evaluated by first computing the truth value that could be assigned
to \type{(foo ?x)} as an element of \type{bdg-to-truth-val}, so that
a truth value is obtained for every binding of \type{?x}.  This truth
value is then manipulated so that the values for any particular binding
of \type{?x} are disjoined (the usual interpretation of the existential
quantifier), and the result is returned as the truth value assigned to
\ref{e.exists}.  The modal operator \type{forall} is handled similarly
and corresponds to universal quantification.  The first (parametric)
argument to these operators can also be a list of variables, in which
case all of the variables are quantified over.

\subsection{Modal operators created automatically}

Finally, many of the functions that construct new bilattices are able
to construct modal operators on these bilattice in terms of the modal
operators on the bilattices they are using in the construction.  More
specifically:

\begin{enumerate}
 \item A product bilattice $A \times B$ will inherit a modal operator
from $A$ and $B$ if these component bilattices have modal operators
with the same name.
 \item A bilattice constructed using \type{append-bilattice} will have
modal operators corresponding to whatever modal operators were available
on the bilattice used in the construction.
 \item A bilattice constructed by splicing two bilattices $A$ and $B$
together will inherit the modal operators from the ``upper'' bilattice.
 \item A bilattice constructed using \type{hierarchy-bilattice} will
have modal operators corresponding to whatever modal operators were
available on the bilattice used in the construction.
 \item A bilattice constructed from a lattice will not normally have
any modal operators associated to it.
 \item The bilattice $B^D$ will inherit whatever modal operators the
bilattice $B$ has.
 \end{enumerate}

\section*{Acknowledgement}

Support for the development of \mvl\ has been provided by the Air
Force Office of Scientific Research, the Defense Advance Research
Projects Agency, General Dynamics, the National Science Foundation,
Rockwell International, and Rome Labs.

\clearpage

\appendix

\section{Summary of MVL functions and constants}

\begin{list}{}{\setlength{\labelwidth}{1.5in}\setlength{\labelsep}{.25in}
\setlength{\leftmargin}{1.75in}}

% *A*

\item [\itype{activate(th)}] Activate the theory \type{th} for reading
and writing.

\item [\itype{activate-best-first}] Activate best-first search by
setting \type{*node-evaluation-fn*} to \type{best-first-fn}.  Allowed
keywords: \type{initial-limit}, \type{limit-delta}, and
\type{final-limit}.

\item [\itype{activate-breadth-first}] Activate breadth-first search
by setting \type{*node-evaluation-fn*} to \type{depth-fn}.  Allowed
keywords: \type{initial-limit}, \type{limit-delta}, and
\type{final-limit}.

\item [\itype{activate-depth-first}] Activate depth-first search by
setting \type{*initial-limit*} to \NIL.  Allowed keyword:
\type{final-limit}.

\item [\itype{activate-theory-mechanism}] Activates the theory
mechanism.  Accepts as an optional argument either a one-character
bilattice descriptor, or a bilattice.

\item [\itype{*active-bilattice*}] The currently active bilattice.

\item [\itype{*active-theory*}] The theory that is currently active
for reading and writing.

\item [\itype{add-dag(d b f)}] Given the \dg-function $f$, add the
value $b$ at the point $d$.  This function works by side effect only.
There is a keyword \type{modifier} that, if supplied, should be a
function that takes two elements of the bilattice and returns a single
one.  Given such a function $m$, the \dg-function is changed so that
the new value at $d$ is $m(b,f(d))$, where $f(d)$ is the old value.
(This is similar to the same keyword in \type{dag-accumulate-fn}.)
The keyword \type{simplify} is also available.

\item[\itype{answer-binding(a)}] Returns the binding field of the
answer \type a.

\item[\itype{answer-value(a)}] Returns the value field of the answer
\type a.

\item [\itype{anytime-trace(pat)}] Include \type{pat} on \at, so that
any instance of \type{pat} will be traced.  The argument is optional;
if omitted, \type{(?*)} (of which everything is an instance) will be
used.

\item [\itype{anytime-untrace(pat)}] Remove any instance of
\type{pat} from \at.  The argument is optional; if omitted,
\type{(?*)} (of which everything is an instance) will be used.

\item [\itype{append-bilattice(b f)}] Returns the bilattice obtained
by appending the result of applying the function $f$ to the points of
the original bilattice $b$.

\item [\itype{append-binding-lists(b1 b2)}] Computes the least
specific binding list that is more specific than both $b_1$ and $b_2$
or returns \NIL\ if no such binding list exists ($b_1$ and $b_2$ may
bind the same variable to different constants, for example).

\item [\itype{*atms-bilattice*}] The bilattice corresponding to an \atms,
where the justification information is maintained in disjunctive normal
form as suggested by de Kleer.

% *B*

\item [\itype{(bagof}\type{ ?x ?p ?b)}] This relation holds for
$x$, $p$ and $b$ if and only if $b$ is a list of all $x$ satisfying
$p$.

\item [\itype{bc(q)}] Attempts to prove the query $q$ from information
in the database, looking only for a single answer.

\item [\itype{bc-trace(pat)}] Include \type{pat} on
\type{*bc-trace*}, so that any instance of \type{pat} will be
traced.  The argument is optional; if omitted, \type{(?*)} (of which
everything is an instance) will be used.

\item [\itype{bc-untrace(pat)}] Remove any instance of \type{pat}
from \type{*bc-trace*}.  The argument is optional; if omitted,
\type{(?*)} (of which everything is an instance) will be used.

\item [\itype{bc-unwatch(preds)}] Removes the given predicates from
the list of {\sc WATCH}ed predicates.  If no arguments are supplied,
all predicates are un{\sc watch}ed.

\item [\itype{bc-watch(preds)}] Accepts as arguments any number of
predicates that should be watched as the prover proceeds.  Each
predicate can also be a list of the form \type{(predicate depth)}
where \type{depth} is the number of successive ancestors to be
displayed when the prover encounters an instance of the predicate.
Returns a list of the predicates currently being {\sc WATCH}ed.

\item [\itype{bcs(q)}] Attempts to prove the query $q$ from information
in the database, looking for all possible answers.

\item [\itype{bdg-to-truth-val}] The bilattice of functions from
the \dg\ of binding lists into \type{*active-bilattice*}.

\item [\itype{bilattice(b)}] Add the bilattice $b$ to \type{*bilattices*}.

\item [\itype{*bilattices*}] A list of the bilattices known to the system.

\item [\itype{binding-dag}] The \dg\ of binding lists.

\item [\itype{binding-le(b1 b2)}] Returns \NIL\ unless $b_1$ is a more
specific binding that $b_2$, in which case a list of relative binding
lists is returned.

\item [\itype{bottom}] The truth value \type{bottom} in the active bilattice.

\item [\itype{*break-node*}] The number of a node at which the
proof process should be interrupted.

% *C*

\item [\itype{choose-task}]  Selects a task to work on next.  The 
current implementation of \mvl\ simply selects any task that might
have some impact on the final answer being computed.

\item [\itype{circum(p)}] Like \type{bc}, but looks for instances of
$p$ that follow from the circumscription axiom.  Allowed keywords:
\type{cutoffs} and \type{succeeds}.

\item [\itype{*circum-trace*}] If not \NIL, the circumscriptive proof
functions will output diagnostics describing their progress.

\item [\itype{circums(p)}] Like \type{bcs}, but looks for instances of
$p$ that follow from the circumscription axiom.  Allowed keywords:
\type{cutoffs} and \type{succeeds}.

\item [\itype{*circumscription-bilattice*}] The bilattice used to compute
the results of the circumscription axiom.

\item [\itype{cnf(p)}] Returns $p$ in conjunctive normal form, as a list
of conjuncts, where each conjunct is a list of disjuncts.  Thus
 \[\type{(CNF '(AND A B))}\]
returns
 \[\type{((A) (B))}\]

\item [\itype{*cnf-atms-bilattice*}] An alternative version of the \atms\
bilattice where the justification information is maintained in conjunctive
normal form.

\item [\itype{combine-1(f params f1 ... fn)}] This is an $n$-ary
version of \type{combine-fns}, where $f$ is an $n$-ary bilattice
function and there are $n$ \dg-functions supplied as well.  The
additional argument \type{params} indicates which of the $n$ arguments
to this function are not necessarily \dg-functions, but are instead
parameters expected by the overall function $f$.  This additional
argument is a list of \T\ (if the $n$th argument is a parameter) or
\NIL\ (if it isn't).  If the parameter list ends early, all of the
remaining entries are assumed to be \NIL.

\item [\itype{combine-fns(f1 f2 fn comp)}] This function takes the two
\dg-functions $f_1$ and $f_2$ and combines them using the function
\type{fn}, which should be a function that takes two elements of the
bilattice and returns a single one.  The comparison function
\type{comp} is used to determine if one of the two functions can be
returned directly; if, for example \type{(comp f1 f2)} holds, than we
simply return $f_1$, and similarly for $f_2$.  As an example, if the
combining function is the bilattice operation $\cdot$, the comparison
function is \type{k-le}.

\item [\itype{compile-mvl}] Compiles and loads the \mvl\ system.
On a Symbolics system, the keyword \type{reload} forces all of the
files to be reloaded and the keyword \type{recompile} forces them all
to be recompiled.

\item [\itype{*completed-database*}] If \T, find all exceptions
to default proofs.  If \NIL\ (the default), accept defaults proofs
with possible exceptions as complete.

\item [\itype{cont-analysis(analysis)}] This function continues the
investigation of a specific analysis, and can be used to resume a
proof effort that has been curtailed for some reason.  It returns two
values: an instance of the \type{answer} structure giving the results
of the proof attempts, and a modified version of the \type{analysis}
structure that can be used to continue the proof process if desired.

\item [\itype{cont-prover(proof)}] Resume the task associated to the
auxiliary information in the given instance of the \type{proof}
structure.  This function returns an instance of \type{bcanswer}
corresponding to the result obtained by the prover, and returns as
soon as a solution to the original query is found.

\item [\itype{contents}] Displays the contents of the database.  Keywords:
\type{cutoffs}, \type{props}.

\item [\itype{create-mvl-invocation}] This function sets up a new
copy of \mvl.  It accepts an optional argument that is an instance of
\mi; this can be used to reset the state of \mvl\ to a previous one.
If the optional argument is omitted, default values are used.  The
value returned is the \mvl\ invocation that was active before the
function was called, so that another call to
\type{create-mvl-invocation} can be used to restore the original
environment.

\item [\itype{cutoffs-not(cutoffs \&optional bilattice)}] Returns a
new cutoffs $c$ such that
 \[\type{(test x c)} = \type{(test (mvl-not x) cutoffs)}\]
 for all truth values $x$.  In other words, $x$ passes $c$ if and
only if $\neg x$ passed the cutoffs supplied originally.

% *D*

\item [\itype{dag-accumulate-fn(dag bilattice list)}] Given a \dg, a
bilattice, and a list of elements of the form
 \[\type{(dag-elt . bilattice-elt)}\]
a \dg-function is constructed that takes the given values at the
given elements of the \dg.

A keyword to \type{dag-accumulate-fn} is \type{modifier}.  This
function is used if the \type{d-list} contains the same \dg\ element
more than once, in which case the modifier function is called with
the new and old values (in that order) and the result is used as the
value of the \dg-function at that particular point.  The default
action is to simply overwrite the old value with the new one.  (In
other words, the default modifier function simply returns its first
argument while ignoring the second.)

\item [\itype{dag-bilattice(dag bilattice)}] Given a bilattice $B$
and a dag $D$, returns the bilattice $B^D$.

\item [\itype{dag-change-dag(f d-f)}] The argument \type{d-f} should
be a function from the \dg\ to itself; a modified version of the
original \dg-function $f$ is returned in which every point at which
the \dg-function takes a new value has had the function \type{d-f}
applied to it.  The bilattice argument is expected to be the bilattice
of functional truth values as opposed to the bilattice that is the
target of the functions.  There is an optional keyword
\type{simplify}; if supplied, the computed function is simplified
before being returned.  (See \type{simplify}.)

\item [\itype{dag-change-tv(f b-f)}] The argument \type{b-f} should be
a function from the bilattice to itself; a modified version of the
original \dg-function $f$ is returned in which every point in the
\dg-function has had the function \type{b-f} applied to it.  There is
an optional keyword \type{simplify}; if supplied, the computed
function is simplified before being returned.  (See \type{simplify}.)

\item [\itype{dag-list-of-answers(f)}] Given a \dg-function $f$, assumed
to be from the \dg\ of binding lists into the currently active bilattice, this
function constructs a list of answers corresponding to the points where
the \dg-function takes values.

\item [\itype{dag-prune(f pred)}] Given a \dg-function $f$ and a
predicate on the bilattice, this function prunes from $f$ all points
$d$ where $f(d)$ fails to satisfy the predicate.  There is an optional
keyword \type{modify-fn} that, if supplied, should be a function $m$
from the bilattice to itself.  In this case, the points where
$m(f(d))$ fails to satisfy the predicate are pruned.

\item [\itype{*default-bilattice*}] A bilattice that can be used for simple
default reasoning.

\item [\itype{delete-bdg(var bdg-list)}] Destructively modifies
\type{bdg-list}, removing any binding for the variable \type{var}.

\item [\itype{denotes(a)}] Returns the proposition associated with a
particular \lisp\ atom.

\item [\itype{dnf(p)}] Returns $p$ in disjunctive normal form, as a list
of disjuncts, where each disjunct is a list of conjuncts.  Thus
 \[\type{(DNF '(AND A B))}\]
returns
 \[\type{((A B))}\]

\item [\itype{(done exp-1 ... exp-n)}] This predicate, which is
useful for stashing only, causes \lisp\ to evaluate each of the subsequent
expressions.

\item [\itype{dot-with-not(x)}] Computes $x\cdot \neg x$.

\item [\itype{dot-with-star}] The function \type{dws} in the active bilattice.

% *E*

\item [\itype{empty}] Unstashes every sentence in the database.  Keyword:
\type{value}.

\item [\itype{empty-theory()}] Empties the currently active theory,
although this theory may be changed using the optional theory
argument.  The keyword \type{value} has a meaning similar to its
meaning as a keyword to \type{erase}.

\item[\itype{equal-answer(a1 a2)}] Returns \T\ if the two answers
\type{a1} and \type{a2} are equal, and \NIL\ otherwise.  The
functions used are \type{binding-equal} and \type{mvl-eq}, the
functions normally used to test binding lists and truth values for
equality.

\item [\itype{equal-binding(b1 b2)}] Checks to see if two binding lists are
effectively the same.

\item [\itype{*equivalence-translation*}] Controls the fashion in which
\type{normal-form} deals with the connective \type{<=>}.

\item [\itype{erase(p)}] Removes $p$ from the database and forward
chains.  Allowed keywords: \type{value}, \type{decrement-p} and
\type{cutoffs}.

\item [\itype{erase-from-theory(p)}] Removes the proposition $p$ from
the currently active theory, and forward chains on the result.
Allowed keywords: \type{theory} can be used to specify the theory from
which $p$ should be removed; \type{value}, \type{decrement-p} and
\type{cutoffs} have the same meaning as for \type{erase}.  (See
Section \ref{s.fc}.)

\item [\itype{existing-mvl-invocation}] This function returns a
single value that is an instance of the structure \mi.  It is useful
primarily for saving the current copy of \mvl\ so that it can be
restored later.

% *F*

\item [\itype{false}] The truth value \type{false} in the active bilattice.

\item [\itype{fc(p)}] Inserts the sentence $p$ into the database and
forward chains.  Allowed keywords: \type{value}, \type{increment-p},
\type{fc-increment-p}, \type{cutoffs} and \type{include-modal}.

\item [\itype{fc-trace(pat)}] Include \type{pat} on
\type{*fc-trace*}, so that any instance of \type{pat} will be
traced.  The argument is optional; if omitted, \type{(?*)} (of which
everything is an instance) will be used.

\item [\itype{fc-untrace(pat)}] Remove any instance of \type{pat}
from \type{*fc-trace*}.  The argument is optional; if omitted,
\type{(?*)} (of which everything is an instance) will be used.

\item [\itype{*final-limit*}] is the limit point beyond which the
search is abandoned, so that any node that evaluates to a value worse
than this will be pruned.  The default value of this parameter is
\NIL, indicating that no nodes are to be pruned.

\item [\itype{*first-order-atms-bilattice*}] This bilattice corresponds
to a fully first-order \atms, as opposed to the more conventional one
dealing only with propositional logic.

\item [\itype{*first-order-bilattice*}] The bilattice corresponding to
first-order logic.

\item [\itype{fn-and(f1 f2)}] This function conjoins the values of
$f_1$ and $f_2$ using the bilattice $\wedge$ operation.  The new
\dg-function is returned.

\item [\itype{fn-comp(f1 f2 fn)}] This functions compares the values
of $f_1$ and $f_2$ using the bilattice comparison function \type{fn}.
Either \T\ or \NIL\ is returned.

\item [\itype{fn-dot(f1 f2)}] This function combines the values of
$f_1$ and $f_2$ using the bilattice $\cdot$ operation.  The new
\dg-function is returned.

\item [\itype{fn-dws(f1 f2)}] This function combines the values of
$f_1$ and $f_2$ using the bilattice operation \type{dot-with-star}.
The new \dg-function is returned.

\item [\itype{fn-eq(f1 f2)}] This function compares the values of
$f_1$ and $f_2$ using the bilattice comparison operation.  Either \T\
or \NIL\ is returned.

\item [\itype{fn-not(f1)}] This function negates the values of $f_1$
using the bilattice $\neg$ operation.  The new \dg-function is
returned.

\item [\itype{fn-or(f1 f2)}] This function disjoins the values of
$f_1$ and $f_2$ using the bilattice $\vee$ operation.  The new
\dg-function is returned.

\item [\itype{fn-plus(f1 f2)}] This function adds the values of $f_1$
and $f_2$ using the bilattice $+$ operation.  The new \dg-function is
returned.

% *G*

\item [\itype{get-bdg(var bdg-list)}] Returns the value assigned to
the variable \type{var} by the binding list \type{bdg-list}.  If the
variable does not appear in the binding list, \NIL\ is returned.

\item [\itype{get-mvl (symbol indicator)}]  Like the Common Lisp
\type{get}, but if the value is a list, treats it as an a-list and
searches for an item whose car matches \type{*mvl-invocation*}.
The value taken can be changed with \type{setf}.

\item [\itype{get-val(d f)}] Returns the value of the function $f$ at
the point $d$ of the \dg.

\item [\itype{*global*}] The root of the \dg\ of theories.

\item [\itype{groundp(p)}] Returns \T\ if the proposition $p$ is
ground (i.e., contains no variables) and \NIL\ otherwise.

% *H*

\item [\itype{hierarchy-bilattice(b)}] Returns the bilattice constructed
as an infinite descending chain of copies of $b$.

\item [\itype{*hierarchy-bilattice*}] A bilattice similar to the
default bilattice but where the defaults are labelled with
integer-valued strengths.

% *I*

\item [\itype{*if-translation*}] Controls the fashion in which
\type{normal-form} deals with the connective \type{if}.

\item [\itype{includes(big little)}] This function is used to state that
the theory \type{little} is a subtheory of the theory \type{big}, and
should inherit from it.

\item [\itype{index(p)}] Updates the discrimination net as needed
to insert the proposition $p$ into the database.

\item [\itype{init-analysis(prop)}] The \mvl\ primitive upon which
the proving functions are based, this function accepts a proposition
and creates and returns an instance of the \type{analysis} structure.

\item [\itype{init-prover(p effort)}] This function initializes and
returns an instance of \type{proof} containing auxiliary information
about the proof process.

\item [\itype{*initial-limit*}] is used to indicate the type of search.
If \NIL\ (the default), then the values assigned to the nodes are not
used to order the search, and depth-first search is used.  If not
\NIL, \type{*initial-limit*} is typically 0.

\item [\itype{instp(p1 p2)}] If $p_2$ is an instance of $p_1$,
return a list of their most general unifiers; only variables in $p_1$
are bound.  Returns \NIL\ if $p_2$ is not an instance of $p_1$.

\item [\itype{invoke(fn p)}] This function invokes \type{fn} on the
sentence $p$.  In addition, any keyword arguments are passed to
\type{fn}, but with \type{\&allow-other-keys} so that the keywords
will be ignored if \type{fn} is not expecting them.

% *J*

% *K*

\item [\itype{k-ge(x y)}] Returns \T\ if $x \geq_k y$, and \NIL\ otherwise.

\item [\itype{k-gt(x y)}] Returns \T\ if $x >_k y$, and \NIL\ otherwise.

\item [\itype{k-le(x y)}] Returns \T\ if $x \leq_k y$, and \NIL\ otherwise.

\item [\itype{k-lt(x y)}] Returns \T\ if $x <_k y$, and \NIL\ otherwise.

\item [\itype{k-not-ge(x y)}] Returns \T\ if $x \not\geq_k y$, and
\NIL\ otherwise.

\item [\itype{k-not-gt(x y)}] Returns \T\ if $x \not>_k y$, and \NIL\
otherwise.

\item [\itype{k-not-le(x y)}] Returns \T\ if $x \not\leq_k y$, and
\NIL\ otherwise.

\item [\itype{k-not-lt(x y)}] Returns \T\ if $x \not<_k y$, and \NIL\
otherwise.

\item [\itype{(known p)}] Stashing this predicate simply stashes $p$;
unstashing it causes an error.  Looking it up looks up $p$.

% *L*

\item [\itype{lattice-to-bilattice(l)}] Returns the bilattice constructed
from the lattice $l$.

\item [\itype{lattice-to-dag(lattice)}] Construct a \dg\ from the
given lattice.

\item [\itype{lattice-to-dag-to-bilattice(lattice bilattice)}]  First
construct a \dg\ $D$ from the given lattice.  Then use the bilattice
$B$ to construct the bilattice $B^D$.

\item [\itype{*limit-delta*}] is used if \type{*initial-limit*} is
non-\NIL.  In this case, the nodes are ordered by increasing
 \[\frac{\type{value(node)}-\type{*initial-limit*}}{\type{*limit-delta*}}\]
 A typical value for *initial-limit* is 0, and for *limit-delta* is
1.  The value of this parameter is irrelevant if
\type{*initial-limit*} is \NIL.

\item [\itype{lisp-defined(fn)}] This macro is used to indicate that
the given function can be evaluated by a \lisp\ call.  It is also
used to cause \mvl\ to signal an error if an attempt is made to stash
or unstash an instance of the function.  \type{Lisp-defined} also
accepts an optional argument, which is an inverter function.  The
overall result of attempting to lookup an instance of the function
can be described as follows:
 \begin{enumerate}
 \item If none of the arguments to the function (i.e., all but the
last of the arguments of the relation) is a variable, then the function
is called and the result is unified with the final argument of the
relation.
 \item If the inverter function has been supplied and is equal to \T,
then the function is called and the result is unified with the final
argument of the relation even if some of the arguments are variables.
 \item If some of the arguments are variables but the result is not a
variable, then the inverter function is invoked if one has been
supplied.  The arguments passed to the inverter function are the
result (i.e., the last argument of the relation), and a list of all
of the arguments (i.e., all but the last argument of the relation).
If there is no inverter function, then \NIL\ is returned.  In
addition, if \wbli\ \index{*warn-on-bad-lisp-invocations*} is not
\NIL, the user is warned.
 \item If both some of the arguments and the result are variables, then
\NIL\ is returned.  In addition, if \wbli\ is not \NIL, the user is
warned.
 \end{enumerate}

\item [\itype{lisp-predicate(pred)}] This macro is used to indicate
that the given predicate can be evaluated by a \lisp\ call.  It is
also used to cause \mvl\ to signal an error if an attempt is made to
stash or unstash an instance of the function.

\item [\itype{load-logic}] Empty the database and make a new bilattice
active.  Accepts as an optional argument either a character designating
a bilattice, or the new bilattice itself.  If no argument is supplied,
the user is prompted for the new bilattice.

\item [\itype{load-mvl}] Loads the \mvl\ system.

\item [\itype{lookup(q)}] Searches the database for a sentence
matching the query $q$.  Allowed keywords: \type{cutoffs}, \type{inv},
\apa, \somo.  Returns two values: an answer and the source of the
match.  Returns \NIL\ if the query cannot be found in the database.

\item [\itype{lookup-call(fn p)}] This function is similar to
\type{invoke}, but modifies the values returned to make them agree
with those expected to be returned by \type{lookup}.  Unless the
function \type{fn} supplies them, the truth value is assumed to be
\type{true}.  The second value is \type{attach}.

\item [\itype{lookups(q)}] Similar to \type{lookup}, but looks
for {\em all\/} sentences matching the query.  Returns two lists of
values.

\item [\itype{lookups-call(fn p)}] The plural version of
\type{lookup-call}.

% *M*

\item [\itype{make-root(dag bilattice b)}] Returns a \dg-function that
takes the value $b$ at the root of the \dg.

\item [\itype{make-tv(dag bilattice d b)}] Returns a \dg-function that
takes the value $b$ at the point $d$.  If $d$ is not the root of the
\dg, the value at the root of the \dg\ is generally the point
\type{unknown} in the bilattice, but the user may supply the keyword
\type{root-fn} if he wishes, in which case the value at the root of
the \dg\ is the result of applying \type{root-fn} to the bilattice.

\item [\itype{matchp(pattern atom)}] Returns a list of elements of
the form \type{(b . d)} where $b$ is a most general binding list such
that \type{(plug pattern b)} is an instance of the proposition
labelled with \type{atom} and $d$ is a binding list for the variables
in the database proposition.  \NIL\ is returned if no such binding
list exists.  Thus if \type{P1} is \type{(TALL \#:?2)}, and \type{P2}
is \type{(TALL TOM)}, then
 \[\type{(matchp (TALL TOM) P1)}\]
returns a list consisting of the single entry 
 \[\type{(NIL . ((\#:?2 . TOM)))}\]
 The \type{car} of this entry is \NIL, indicating that no variables
in the query need to be bound to match \type{P1}; the \type{cdr} is
\type{((\#:?2 . TOM))} since the variable \type{\#:?2} in \type{P1}
needs to be bound to \type{TOM} to match the query.
 \[\type{(matchp (TALL ?X) P2)}\]
returns \type{((((?X . TOM)) . NIL))}.

\item [\itype{mvl-and}] The function \type{and} in the active bilattice.

\item [\itype{mvl-dot}] The function \type{dot} in the active bilattice.

\item [\itype{mvl-eq}] The function \type{eq} in the active bilattice.

\item [\itype{mvl-file(string)}] Returns the path for the \mvl\ file
associated with the given name.  The file name may include an
extension, one may be supplied as an optional argument, or
\type{.MVL} will be used by default.  There is an optional keyword
\type{no-error} indicating that even if the file cannot be found, the
path should be returned.

\item [\itype{*mvl-invocation*}] A global variable containing a unique
identifier for the currently active copy of \mvl.

\item [\itype{mvl-load(file)}] Loads the sentences in the given file into
the \mvl\ database.

\item [\itype{mvl-not}] The function \type{not} in the active bilattice.

\item [\itype{mvl-or}] The function \type{or} in the active bilattice.

\item [\itype{mvl-plug}] The function \type{plug} in the active bilattice.

\item [\itype{mvl-plus}] The function \type{plus} in the active bilattice.

\item [\itype{mvl-print}] This function accepts a single argument, a
truth value, and modifies it in a way that will cause it to be
pretty-printed.  The value returned by mvl-print should not be used
for any purpose other than printing.

\item [\itype{mvl-save(file)}] Writes the \mvl\ database out to the
given file.  This function also accepts the keywords \type{cutoffs}
and \type{props}, which are treated as for \type{contents} (see
Section \ref{s.contents}).  Finally, the keyword \type{if-exists} is
passed to the \type{with-open-file} macro.

\item [\itype{mvl-simplicity}] The function \type{simplicity} in the active
bilattice.

\item [\itype{mvl-test-file(string)}] Can be used to evaluate
various forms by the \mvl\ system and compare the results to
predefined values.  Both the forms and the expected answers are in
the \type{.TEST} file whose name is supplied to the function.
Returns the number of errors found.

\item [\itype{mvl-unload(file)}] Removes the sentences in the given file from
the \mvl\ database.

\item [\itype{mvl-vars}] The function \type{vars} in the active bilattice.

% *N*

\item [\itype{negate(p)}] Returns the negation of $p$.

\item [\type{new-$*$var()}]\index{new-star-var} Returns a new
temporary sequence variable.

\item [\type{new-?var()}]\index{new-query-var} Returns a new
temporary variable.

\item [\itype{new-val(a v)}] Sets the truth value of the sentence
associated with the atom $a$ to $v$.  If $v$ is the truth value
\type{unknown}, the sentence is removed from the database.

\item [\itype{normal-form(p)}] Returns an equivalent form of $p$
suitable for insertion into the database (i.e., with the connectives
\type{IFF}, \type{IF} and \type{<=>} reduced to \type{<=} and \type{=>}).

% *O*

% *P*

\item [\itype{plug(p bdg-list)}] Returns the result of instantiating
the proposition $p$ using the given binding list.  This function is
nondestructive.

\item [\itype{prfacts(atom)}] Displays the database facts concerning the
given atom.  Keyword: \type{cutoffs}.

\item [\itype{*probability-bilattice*}] An extended version of the \atms\
bilattice that maintains probabilistic information as well.

\item [\itype{product-bilattice(b1 b2)}] Returns the bilattice
that is the product of the bilattices $b_1$ and $b_2$.

\item [\itype{(property symbol value indicator)}] Stashing this
predicate sets the symbol's \type{indicator} property to \type{value}
for the current invocation of \mvl.  Unstashing it removes the
\type{indicator} property for the current invocation of \mvl\ if the
previous value of this property unifies with \type{value}.  (See the
definitions of \type{get-mvl} and \type{rem-mvl}.)

\item [\itype{props()}] Returns a list of all propositions about
which any information has been recorded.

% *Q*

% *R*

\item [\itype{release-mvl}] (Symbolics only)  Compiles and loads the
\mvl\ system, and declares the new version to be \type{released}.  The
keyword \type{no-recompile} can be used to skip the compilation step.

\item [\itype{rem-mvl (symbol indicator)}] Like the Common Lisp
\type{remprop}, but if the value is a list, treats it as an a-list
and removes the item whose car matches \type{*mvl-invocation*}.

% *S*

\item [\itype{samep(p1 p2)}] Returns \T\ if the propositions $p_1$ and
$p_2$ are identical (up to variable renaming), and \NIL\ otherwise.
If the optional \type{return-answer?} argument is supplied and the
propositions are the same, a binding list is returned instead of \T.

\item [\itype{*save-nodes*}] If \T\ (the default), nodes are
labelled with consecutive integers when they are displayed by the
first-order prover.

\item [\itype{set-search-parameters}] This function is the primitive
used to modify the parameters used to control the first-order theorem
prover's search.  It accepts as keywords \type{initial-limit},
\type{limit-delta}, and \type{final-limit}.

\item [\itype{simplify(f)}] In some cases, manipulating the
\dg-function $f$ may have cause it to contain redundancies, where the
value at some point is presented explicitly even though this value is
in fact the one that would normally be inherited from its parents.
This function returns a modification of $f$ where these irrelevant
entries have been removed.

\item [\itype{splice-bilattice(upper lower)}] Returns the bilattice
obtained by replacing the point \type{unknown} in the bilattice
\type{upper} with a copy of the bilattice \type{lower}.

\item [\itype{standardize-operators(p)}] Returns an equivalent form of $p$
with the connectives \type{<=>}, \type{IF} and \type{IFF} removed.
This function does not use the control variables
\type{*if-translation*} and \et, since it is used when proving
sentences, not when inserting them into the database.

\item [\itype{stash(p)}] Inserts $p$ into the database.  The allowed
keywords are \type{value}, \type{increment-p} and \apa.

\item [\itype{stash-value}] The function \type{stash-val} in the
active bilattice.

\item [\itype{state(p)}] Inserts $p$ into the database after
converting it to normal form.  Keywords: \type{value},
\type{increment-p}.

\item [\itype{state-in-theory(p)}] Adds the proposition $p$ to
the currently active theory.  Allowed keywords: \type{theory} can be
used to specify a different theory to which $p$ should be added;
\type{value} and \type{increment-p} have the same meaning as for
\type{state}.  (See Section \ref{s.database}; if these keywords are
not used, the function \type{state} can be used equally well, since
the default stash value with the theory mechanism activated also
inserts the proposition into the currently active theory.)

\item [\itype{std-cutoffs}] Cutoff information constructed from the
value of \type{cutoff-val} in the active bilattice.  See Section
\ref{s.cutoffs}.

% *T*

\item [\itype{t-ge(x y)}] Returns \T\ if $x \geq_t y$, and \NIL\ otherwise.

\item [\itype{t-gt(x y)}] Returns \T\ if $x >_t y$, and \NIL\ otherwise.

\item [\itype{t-le(x y)}] Returns \T\ if $x \leq_t y$, and \NIL\ otherwise.

\item [\itype{t-lt(x y)}] Returns \T\ if $x <_t y$, and \NIL\ otherwise.

\item [\itype{t-not-ge(x y)}] Returns \T\ if $x \not\geq_t y$, and
\NIL\ otherwise.

\item [\itype{t-not-gt(x y)}] Returns \T\ if $x \not>_t y$, and \NIL\
otherwise.

\item [\itype{t-not-le(x y)}] Returns \T\ if $x \not\leq_t y$, and
\NIL\ otherwise.

\item [\itype{t-not-lt(x y)}] Returns \T\ if $x \not<_t y$, and \NIL\
otherwise.

\item [\itype{test(value cutoffs \&optional bilattice)}] Returns \T\
if the truth value passes the test associated to \type{cutoffs}, and
\NIL\ otherwise.  The test is performed for the supplied bilattice,
or for the currently active one is the optional argument is omitted.

\item [\itype{test-mvl}] Tests the \mvl\ system by invoking the function
\type{mvl-test-file} on the file \type{MVL.TEST}.

\item [\itype{theory-active-p}] Returns \T\ or \NIL\ depending on
whether or not the theory mechanism has been activated.

\item [\itype{theory-contents()}] Displays the contents of the
currently active theory, which can be changed using the
keyword \type{theory}.  The keywords \type{cutoffs} and \type{props}
are allowed and are treated as for \type{contents}.  (See Section
\ref{s.contents}.)  There is an additional keyword,
\itype{this-theory-only}, which defaults to \T\ and, if supplied,
indicates that only sentences whose truth values {\em change\/} in the
given theory are to be displayed.

\item [\itype{theory-dag}] The \dg\ of theories.

\item [\itype{*time-bilattice*}] A compound bilattice that can be used
for reasoning about events occurring against a temporal background.

\item [\itype{true}] The truth value \type{true} in the active bilattice.

\item [\itype{true-in-theory}]  The truth value of a sentence that is
true in the currently-active theory.

\item [\itype{truth-value(a)}] Returns the truth value associated with
a particular \lisp\ atom.  Thus, if the result of asserting
\type{(not (short tom))} into the database was the atom \type{P38},
\type{(denotes 'p38)} would evaluate to \type{(short tom)}, and
\type{(truth-value 'p38)} would evaluate to \type{false}. 

% *U*

\item [\itype{unifyp(p1 p2)}] Returns a list of the most general
unifiers of $p_1$ and $p_2$, or \NIL\ if the propositions fail to
unify.

\item [\itype{unincludes(big little)}] This function is used to state
that the theory \type{little} is no longer a subtheory of the theory
\type{big}, and should not inherit from it.

\item [\itype{unindex(p)}] Removes from the discrimination net 
information relating to the proposition $p$.

\item [\itype{unknown}] The truth value \type{unknown} in the active
bilattice.

\item [\itype{(unknown p)}]  Stashing this predicate unstashes $p$;
unstashing it causes an error.  Looking it up fails if $p$
can be found in the database, and succeeds otherwise.

\item [\itype{unstash(p)}] Removes $p$ from the database.  Allowed
keywords: \type{value}, \type{decrement-p} and \apa.

\item [\itype{unstash-facts-with-atom(atom)}] Unstashes from the
database any facts concerning the given atom.  Keywords:
\type{value}, \type{decrement-p}, \apa.

\item [\itype{unstash-value}] The function \type{unstash-val} in the
active bilattice.

\item [\itype{unstate(p)}] Removes $p$ from the database after
converting it to normal form.  Keywords: \type{value},
\type{decrement-p}.

\item [\itype{unstate-from-theory(p)}] Removes the proposition $p$
from the currently active theory.  Allowed keywords:
\type{theory} can be used to specify a different theory from which
$p$ should be removed; \type{value} and \type{decrement-p} have the
same meaning as for \type{unstate}.  (See Section \ref{s.database}.)

% *V*

\item [\itype{(value symbol val)}] Stashing this predicate sets the
given symbol to the value \type{val}.  Unstashing it makes the symbol
unbound if its current value unifies with \type{val}.

\item [\itype{varp(x)}] Returns \type{T} if $x$ is a variable, and \NIL\
otherwise.

\item [\type{varp$*$(x)}]\index{varp-star} Returns \T\ if $x$ is a
sequence variable, and \NIL\ otherwise.

\item [\type{varp?(x)}]\index{varp-query} Returns \T\ if $x$ is a
standard (i.e., non-sequence) variable, and \NIL\ otherwise.

\item [\itype{vars-in(p)}] Returns a list of the variables appearing in $p$.

\item [\itype{vartype(x)}] Returns ? if $x$ is a standard variable,
$*$ if $x$ is a sequence variable, and \NIL\ is $x$ is not a
variable.

% *W*

\item [\itype{with-mvl-invocation}] This is a macro similar to the
Common Lisp macro \type{with-open-file}.  It accepts an optional
argument that is an instance of \mi\ and is followed by a body of
forms (presumably including a \type{load-logic}) that are evaluated
in a locally-bound copy of \mvl.

% *X*

% *Y*

% *Z*

\end{list}

\input{manual.ind}

\end{document}


\end{document}


\end{document}


\end{document}
